% આ ભાગનો પૂર્ણ સંપાદનીય પાઠ છે. શૈલી, ફોન્ટ, ચિત્રો અને પ્રકરણસંદર્ભ માટે સંપૂર્ણ સ્રોત ZIP તથા reader/full-reader.tex જુઓ. આ એકલી સ્વતંત્ર કમ્પાઇલ થતી મુખ્ય ફાઇલ નથી.

% BEGIN OLP-0685
\olpart{pt}{સાબિતી-સિદ્ધાંત}
\phantomsection\label{guunit:OLP-0685}


\begin{editorial}
  સાબિતી-સિદ્ધાંતની સામગ્રી અધૂરી અને પ્રાયોગિક છે,
  અને હજી મુખ્ય PDF સંસ્કરણમાં સમાવાઈ નથી.
\end{editorial}

\begingroup
\makeatletter\def\ol@sectioncs{section}\makeatother

% BEGIN OLP-0712
\olchapter{pt}{seq}{સિક્વન્ટ કલન}
\olchapterid{pt}{seq}
\phantomsection\label{guunit:OLP-0712}


\begingroup
\makeatletter\def\ol@sectioncs{section}\makeatother

% BEGIN OLP-0700
\olfileid{pt}{seq}{int}

\olsection{પરિચય}
\phantomsection\label{guunit:OLP-0700}


સિક્વન્ટ કલનો સાબિતી તંત્રો છે,
જેને ગેરહાર્ડ ગેન્ટ્ઝેને ૧૯૩૦ના દાયકામાં રજૂ કર્યાં.
તેમનો હેતુ વસ્તુઓને સાબિત કરવાની વ્યવહારુ રીત આપવાનો ઓછો હતો.
ગેન્ટ્ઝેને તેમનો ઉપયોગ સાબિતીઓની શક્ય સંરચના
અને ગુણધર્મો વિશે અમુક પરિણામો સાબિત કરી શકવા માટે કર્યો.
સાબિતી-સિદ્ધાંત બરાબર આવા પ્રશ્નોનો વિચાર કરે છે.

નામ સૂચવે છે તેમ સિક્વન્ટ કલનો \emph{સિક્વન્ટો} પર કામ કરે છે,
સૂત્રો પર નહીં.
સિક્વન્ટ એ સૂત્રોના બે ક્રમો કે બહુગણોની જોડી છે.
ગેન્ટ્ઝેનનાં મૂળ તંત્રો (\Log{LK} અને \Log{LJ})માં ક્રમો હતા;
હવે સાબિતી-સિદ્ધાંતના અભ્યાસીઓ બહુગણો પસંદ કરે છે.
બહુગણ સમૂહ જેવો છે, પરંતુ તેમાં ઘટક
એક કરતાં વધુ વાર આવી શકે છે.
સમૂહની જેમ તેમાં પણ ઘટકોનો ક્રમ મહત્ત્વનો નથી.
એટલે $\{!A, !A, !B\}$ અને $\{!A, !B, !A\}$
એ જ બહુગણ દર્શાવે છે,
પરંતુ તે બહુગણ $\{!A, !B\}$થી અને $\{!A, !A, !B, !B\}$થી જુદો છે,
કારણ કે છેલ્લાં બેમાં $!A$ અને $!B$ની સંખ્યા પહેલા બહુગણથી જુદી છે.

સાબિતી-સિદ્ધાંતમાં પરંપરાગત રીતે સૂત્રોના બહુગણો
મોટા ગ્રીક અક્ષરોથી લખાય છે.
એટલે $\Gamma$, $\Delta$, $\Pi$, $\Lambda$ કે $\Theta$ લખીએ,
તો આપણે જે તર્કનો વિચાર કરતાં હોઈએ તેનાં સૂત્રોના બહુગણો સમજીએ.
બીજી પરંપરા મુજબ $\tuple{\Gamma, \Delta}$ લખવાને બદલે
સાબિતી-સિદ્ધાંતના અભ્યાસીઓ બંને ઘટકોને સિક્વન્ટ-ચિહ્નથી જુદા પાડે છે;
અહીં~$\Sequent$ વાપરીએ છીએ.

\begin{defn}[સિક્વન્ટ]
\emph{સિક્વન્ટ} આ સ્વરૂપનો વ્યંજક છે:
\[
\Gamma \Sequent \Delta
\]
અહીં $\Gamma$ અને $\Delta$ એ સૂત્રોના સીમિત
(સંભવતઃ ખાલી) બહુગણો છે.
$\Gamma$ને \emph{પૂર્વાંગ} અને $\Delta$ને \emph{ઉત્તરાંગ} કહે છે.
\end{defn}

$\Gamma$નાં સૂત્રોનું સંયોજન $!G$~લઈએ
(ખાલી સંયોજન~$\ltrue$ ગણીએ)
અને $\Delta$નાં સૂત્રોનો વિકલ્પ $!D$~લઈએ
(ખાલી વિકલ્પ~$\lfalse$ ગણીએ).
ત્યારે સિક્વન્ટ $\Gamma \Sequent \Delta$
કોઈ સંરચના~$\Struct{M}$માં સત્ય છે
જો અને ફક્ત જો $!G \lif !D$ સત્ય હોય.
શાસ્ત્રીય તર્કમાં આ એ જ છે કે
કાં તો $\Gamma$નાં સૂત્રોમાંથી કોઈ અસત્ય હોય
અથવા $\Delta$નાં સૂત્રોમાંથી કોઈ સત્ય હોય.

$\Gamma$ એ સૂત્રોનો બહુગણ હોય,
તો $\Gamma$માં $!A$ ઉમેરવાનું પરિણામ
$\Gamma, !A$ (અથવા $!A, \Gamma$)થી લખીએ.
$\Delta$ પણ સૂત્રોનો બહુગણ હોય,
તો $\Gamma, \Delta$ બંને બહુગણોનો બહુગણ-સંઘ છે:
$\Gamma$માં $!A$ની \emph{આવૃત્તિઓની સંખ્યા}~n હોય
(એટલે તેમાં~$!A$ની $n$ નકલ હોય)
અને $\Delta$માં આવૃત્તિઓની સંખ્યા~$m$ હોય,
તો $\Gamma, \Delta$માં $!A$ની આવૃત્તિઓની સંખ્યા~$n+m$ છે.
સિક્વન્ટની એક બાજુનો બહુગણ ખાલી હોય,
તો સામાન્ય રીતે જગ્યા ખાલી રાખીએ.
દાખલા તરીકે $\Sequent !A$ એ ખાલી પૂર્વાંગવાળો સિક્વન્ટ છે,
જેના ઉત્તરાંગમાં $!A$ની આવૃત્તિઓની સંખ્યા~$1$ છે.
\sourcecorrection{OLMD-307}{મૂળ અહીં બે ક્રમોના બહુગણ-સંઘ કહે છે. આ પરિચ્છેદમાં બંને વસ્તુઓ બહુગણો છે અને આવૃત્તિઓની સંખ્યા ઉમેરાય છે; તેથી બંને બહુગણોનો સંઘ કહ્યું છે.}

સિક્વન્ટ કલનમાં સાબિતી એ સિક્વન્ટોનું વૃક્ષ છે.
સૌથી ઉપરના (અથવા આરંભિક) સિક્વન્ટો
વિચારેલા તંત્રના સ્વયંસિદ્ધો છે.
કોઈ સિક્વન્ટ બીજા એક કે બે સિક્વન્ટોની નીચે હોય,
તો તંત્રના કોઈ અનુમાન-નિયમથી તેમનામાંથી યોગ્ય રીતે અનુસરવો જોઈએ.
પહેલા શાસ્ત્રીય તર્ક માટેનાં બે જુદાં સિક્વન્ટ તંત્રો,
$\Log{G1c}$ અને~$\Log{G3c}$, વિચારીએ.
તેમના સ્વયંસિદ્ધો અને નિયમો
\cref{pt:seq:int:tab:G1c,pt:seq:int:tab:G3c}માં છે.
 (ગેન્ટ્ઝેનનું મૂળ તંત્ર \Log{LK}
\olref[fol][seq][]{chap}માં વર્ણવ્યું છે.)

\begingroup
\makeatletter\def\ol@sectioncs{section}\makeatother

% BEGIN OLP-0704
\phantomsection\label{guunit:OLP-0704}


\begin{table}
\begin{tabular}{@{}c@{}c@{}}
\multicolumn{2}{@{}c@{}}{સ્વયંસિદ્ધો}\\[2ex]
$!A \Sequent !A$ & $\lfalse \Sequent \quad$
\\[2ex]
\multicolumn{2}{@{}c@{}}{બંધારણીય નિયમો}\\[2ex]
\Axiom$ \Gamma \fCenter \Delta $
\RightLabel{\LeftR{\Weakening}}
\UnaryInf$ !A, \Gamma \fCenter \Delta$
\DisplayProof
&
\Axiom$ \Gamma \fCenter \Delta$
\RightLabel{\RightR{\Weakening}}
\UnaryInf$ \Gamma \fCenter \Delta, !A$
\DisplayProof
\\[4ex]
\Axiom$ !A, !A, \Gamma \fCenter \Delta $
\RightLabel{\LeftR{\Contraction}}
\UnaryInf$ !A, \Gamma \fCenter \Delta$
\DisplayProof
&
\Axiom$ \Gamma \fCenter \Delta, !A, !A$
\RightLabel{\RightR{\Contraction}}
\UnaryInf$ \Gamma \fCenter \Delta, !A$
\DisplayProof
\\[4ex]
\multicolumn{2}{@{}c@{}}{તાર્કિક નિયમો}\\[2ex]
\Axiom$ \Gamma \fCenter \Delta, !A $
\RightLabel{\LeftR{\lnot}}
\UnaryInf$ \lnot !A, \Gamma \fCenter \Delta$
\DisplayProof
&
\Axiom$!A, \Gamma \fCenter \Delta$
\RightLabel{\RightR{\lnot}}
\UnaryInf$ \Gamma \fCenter \Delta, \lnot !A $
\DisplayProof
\\[4ex]\begin{tabular}[t]{@{}l@{}}
\Axiom$ !A, \Gamma \fCenter \Delta$
\RightLabel{\LeftR{\land}}
\UnaryInf$ !A \land !B, \Gamma \fCenter \Delta$
\DisplayProof
\\[3ex]
\Axiom$!B, \Gamma \fCenter \Delta$
\RightLabel{\LeftR{\land}}
\UnaryInf$!A \land !B, \Gamma \fCenter \Delta$
\DisplayProof\\[3ex]
\end{tabular}
&
\Axiom$\Gamma \fCenter \Delta, !A$
\Axiom$ \Gamma \fCenter \Delta, !B$
\RightLabel{\RightR{\land}}
\BinaryInf$ \Gamma \fCenter \Delta, !A \land !B $
\DisplayProof
\\[4ex]\Axiom$!A, \Gamma \fCenter \Delta$
\Axiom$!B, \Gamma \fCenter \Delta$
\RightLabel{\LeftR{\lor}}
\BinaryInf$!A \lor !B, \Gamma \fCenter \Delta$
\DisplayProof
&
\begin{tabular}[t]{@{}r@{}}
\Axiom$\Gamma \fCenter \Delta, !A$
\RightLabel{\RightR{\lor}}
\UnaryInf$ \Gamma \fCenter \Delta, !A \lor !B$
\DisplayProof
\\[3ex]
\Axiom$ \Gamma \fCenter \Delta, !B$
\RightLabel{\RightR{\lor}}
\UnaryInf$ \Gamma \fCenter \Delta, !A \lor !B$
\DisplayProof\\[3ex]
\end{tabular}
\\[4ex]
\Axiom$ \Gamma \fCenter \Delta, !A$
\Axiom$ !B, \Gamma \fCenter \Delta$
\RightLabel{\LeftR{\lif}}
\BinaryInf$ !A \lif !B, \Gamma \fCenter \Delta$
\DisplayProof
&
\Axiom$ !A, \Gamma \fCenter \Delta, !B$
\RightLabel{\RightR{\lif}}
\UnaryInf$ \Gamma \fCenter \Delta, !A \lif !B $
\DisplayProof
\\[4ex]
\Axiom$ !A(t), \Gamma \fCenter \Delta$
\RightLabel{\LeftR{\lforall}}
\UnaryInf$ \lforall[x][!A(x)],\Gamma \fCenter \Delta$
\DisplayProof
&
\Axiom$ \Gamma \fCenter \Delta, !A(a) $
\RightLabel{\RightR{\lforall}}
\UnaryInf$ \Gamma \fCenter \Delta, \lforall[x][!A(x)]$
\DisplayProof
\\[4ex]
\Axiom$ !A(a), \Gamma \fCenter \Delta $
\RightLabel{\LeftR{\lexists}}
\UnaryInf$ \lexists[x][!A(x)], \Gamma \fCenter \Delta$
\DisplayProof
&
\Axiom$ \Gamma \fCenter \Delta, !A(t) $
\RightLabel{\RightR{\lexists}}
\UnaryInf$ \Gamma \fCenter \Delta, \lexists[x][!A(x)]$
\DisplayProof
\end{tabular}
\caption{\Log{G1c}ના નિયમો. $\RightR{\forall}$ અને
$\LeftR{\exists}$માં અચળ~$a$ ફલિત-સિક્વન્ટમાં આવવો નહીં.
$\Gamma$ અને $\Delta$ એ સૂત્રના બહુગણ છે.}
\ollabel{tab:G1c}
\end{table}
% END OLP-0704
\endgroup

\begingroup
\makeatletter\def\ol@sectioncs{section}\makeatother

% BEGIN OLP-0707
\phantomsection\label{guunit:OLP-0707}


\begin{table}
\begin{tabular}{@{}c@{}c@{}}
\multicolumn{2}{@{}c@{}}{સ્વયંસિદ્ધો}\\[2ex]
$!C, \Gamma \Sequent \Delta, !C$ & $\lfalse, \Gamma \Sequent \Delta$
\\[2ex]
\multicolumn{2}{@{}c@{}}{તાર્કિક નિયમો}\\[2ex]
\Axiom$ \Gamma \fCenter \Delta, !A $
\RightLabel{\LeftR{\lnot}}
\UnaryInf$ \lnot !A, \Gamma \fCenter \Delta$
\DisplayProof
&
\Axiom$!A, \Gamma \fCenter \Delta$
\RightLabel{\RightR{\lnot}}
\UnaryInf$ \Gamma \fCenter \Delta, \lnot !A $
\DisplayProof
\\[4ex]
\Axiom$ !A, !B, \Gamma \fCenter \Delta$
\RightLabel{\LeftR{\land}}
\UnaryInf$ !A \land !B, \Gamma \fCenter \Delta$
\DisplayProof
&
\Axiom$\Gamma \fCenter \Delta, !A$
\Axiom$ \Gamma \fCenter \Delta, !B$
\RightLabel{\RightR{\land}}
\BinaryInf$ \Gamma \fCenter \Delta, !A \land !B $
\DisplayProof
\\[4ex]\Axiom$!A, \Gamma \fCenter \Delta$
\Axiom$!B, \Gamma \fCenter \Delta$
\RightLabel{\LeftR{\lor}}
\BinaryInf$!A \lor !B, \Gamma \fCenter \Delta$
\DisplayProof
&
\Axiom$\Gamma \fCenter \Delta, !A, !B$
\RightLabel{\RightR{\lor}}
\UnaryInf$ \Gamma \fCenter \Delta, !A \lor !B$
\DisplayProof
\\[4ex]
\Axiom$ \Gamma \fCenter \Delta, !A$
\Axiom$ !B, \Gamma \fCenter \Delta$
\RightLabel{\LeftR{\lif}}
\BinaryInf$ !A \lif !B, \Gamma \fCenter \Delta$
\DisplayProof
&
\Axiom$ !A, \Gamma \fCenter \Delta, !B$
\RightLabel{\RightR{\lif}}
\UnaryInf$ \Gamma \fCenter \Delta, !A \lif !B $
\DisplayProof
\\[4ex]
\Axiom$ !A(t), \lforall[x][!A(x)], \Gamma \fCenter \Delta$
\RightLabel{\LeftR{\lforall}}
\UnaryInf$ \lforall[x][!A(x)],\Gamma \fCenter \Delta$
\DisplayProof
&
\Axiom$ \Gamma \fCenter \Delta, !A(a) $
\RightLabel{\RightR{\lforall}}
\UnaryInf$ \Gamma \fCenter \Delta, \lforall[x][!A(x)]$
\DisplayProof
\\[4ex]
\Axiom$ !A(a), \Gamma \fCenter \Delta $
\RightLabel{\LeftR{\lexists}}
\UnaryInf$ \lexists[x][!A(x)], \Gamma \fCenter \Delta$
\DisplayProof
&
\Axiom$ \Gamma \fCenter \Delta, \lexists[x][!A(x)], !A(t) $
\RightLabel{\RightR{\lexists}}
\UnaryInf$ \Gamma \fCenter \Delta, \lexists[x][!A(x)]$
\DisplayProof
\end{tabular}
\caption{\Log{G3c}ના નિયમો. $\RightR{\forall}$ અને
$\LeftR{\exists}$માં અચળ~$a$ ફલિત-સિક્વન્ટમાં આવવો નહીં.
$\Gamma$ અને $\Delta$ એ સૂત્રના બહુગણ છે.
સ્વયંસિદ્ધોમાં સૂત્ર~$!C$ પરમાણ્વીય હોવું જોઈએ.}
\ollabel{tab:G3c}
\end{table}
\sourcecorrection{OLMD-333}{મૂળ સર્વપરિમાણકના ડાબા નિયમના આધારમાં જાળવેલા મુખ્ય સૂત્ર અને સંદર્ભ વચ્ચે અલ્પવિરામ નથી. બંને પૂર્વાંગની અલગ નોંધો હોવાથી તે અલ્પવિરામ ઉમેર્યો છે.}
% END OLP-0707
\endgroup


તંત્રોમાં સૂક્ષ્મ તફાવતો છે, અને તેથી તેમના
સાબિતી-સૈદ્ધાંતિક ગુણધર્મો જુદા છે.
\Log{G1c}માં $!A \Sequent !A$ સ્વયંસિદ્ધો છે;
બીજાં સૂત્રો આવવા દેવાતાં નથી.
\Log{G3c}માં $!C, \Gamma \Sequent \Delta, !C$
સ્વરૂપનાં સ્વયંસિદ્ધો છે, જેમાં $\Gamma$ અને~$\Delta$માં
મનસ્વી ``બાજુનાં'' સૂત્રો હોય,
પરંતુ બંને બાજુ આવતું સૂત્ર~$!C$ \emph{આણ્વિક} હોવું જોઈએ.
\Log{G1c}માં $\LeftR{\land}$ તથા $\RightR{\lor}$ના
બબ્બે સ્વરૂપ છે, જ્યારે \Log{G3c}માં એકેક છે.
વળી \Log{G1c}માં વધારાના ``બંધારણીય નિયમો'' છે,
જેને ``સંકોચન'' (\LeftR{\Contraction}, \RightR{\Contraction})
અને ``શિથિલીકરણ'' (\LeftR{\Weakening}, \RightR{\Weakening}) કહે છે.

\Log{G1c} અને \Log{G3c} તંત્રો શાસ્ત્રીય તર્ક માટે યથાર્થ અને પૂર્ણ છે.
બીજા શબ્દોમાં, આ તંત્રોમાં સાબિતીક્ષમ દરેક સિક્વન્ટ
દરેક સંરચનામાં સત્ય છે,
અને દરેક સંરચનામાં સત્ય હોય તે દરેક સિક્વન્ટની સાબિતી છે.
એટલે કોઈ સિક્વન્ટની સાબિતી છે \emph{કે નહીં} તે પ્રશ્નનો
જવાબ તર્કની બીજી રીતોથી પણ આપી શકાય
(દાખલા તરીકે નિદર્શ-સિદ્ધાંતની રીતોથી).
બંને તંત્રો દરેક સંરચનામાં સત્ય હોય
બરાબર તે જ સિક્વન્ટોને સાબિત કરે છે;
તેથી ખાસ કરીને તેઓ એ જ સિક્વન્ટોને સાબિત કરે છે.

પરંતુ સાબિતી-સિદ્ધાંત ફક્ત કંઈક સાબિત થઈ શકે છે
\emph{કે નહીં} તેનો જ નહીં, \emph{કેવી રીતે} તેનો પણ વિચાર કરે છે.
તેના અભ્યાસીઓ આવા પ્રશ્નો વિચારે છે:
``કોઈ ચોક્કસ નિયમ હંમેશાં ટાળી શકાય?'',
``તંત્રમાં આ નિયમ ઉમેરવાથી નવા સિક્વન્ટો સાબિત થઈ ન જાય તે શક્ય છે?'',
અથવા ``એક તંત્રની સાબિતીઓને બીજાની સાબિતીઓમાં ફેરવી શકાય?''
જવાબ ``હા'' હોય, તો એ હકીકતની સાબિતી
ઘણી વાર સાબિતીઓની જટિલતા પર આગમનથી થાય,
અથવા સમકક્ષ રીતે સાબિતીઓની સંરચના પર આગમનથી થાય.
આથી ફક્ત હકીકતની સાબિતી જ નહીં
(દાખલા તરીકે \Log{G1c} સિક્વન્ટ સાબિત કરે તો \Log{G3c} પણ એ જ સાબિત કરે),
પણ પ્રક્રિયા પણ મળે
(દાખલા તરીકે~\Log{G1c}ની સાબિતીઓને~\Log{G3c}ની સાબિતીઓમાં ફેરવવાની).

સાબિતી-સિદ્ધાંતનું કેન્દ્રીય પરિણામ \emph{કટ-લોપનું પ્રમેય} છે.
તે બતાવે છે કે સિક્વન્ટ તંત્રોમાં કટ-નિયમ
\[
\Axiom$\Gamma \fCenter \Delta, !A$
\Axiom$!A, \Pi \fCenter \Lambda$
\RightLabel{\Cut}
\BinaryInf$\Gamma, \Pi \fCenter \Delta, \Lambda$
\DisplayProof
\]
ઉમેરી શકાય અને કયા સિક્વન્ટો સાબિતીક્ષમ છે તે બદલાતું નથી.
વળી તેની સાબિતી એવી રીત આપે છે
જે \Log{G1c} (અથવા \Log{G3c})ના નિયમો સાથે
\Cut{} વાપરતી કોઈ પણ સાબિતીને,
ફક્ત \Log{G1c} (અથવા \Log{G3c})ના નિયમો વાપરતી સાબિતીમાં ફેરવે છે.
બીજા શબ્દોમાં આ પ્રક્રિયા \Cut{} વાપરતી સાબિતીઓમાંથી
એ નિયમનો ``લોપ'' કરે છે; તેથી આ નામ છે.
% END OLP-0700
\endgroup

\begingroup
\makeatletter\def\ol@sectioncs{section}\makeatother

% BEGIN OLP-0711
\olfileid{pt}{seq}{rp}

\olsection{નિયમો અને સાબિતીઓ}
\phantomsection\label{guunit:OLP-0711}


શરૂઆતમાં કેટલીક વ્યાખ્યાઓ અને પરિભાષાઓ નક્કી કરીએ.

દાખલા તરીકે, \Log{G3c}ના એવા બે નિયમો લો જે
$!A$ અને $!B$ ધરાવતા સિક્વન્ટોમાંથી $!A \land !B$
ધરાવતા સિક્વન્ટ સાબિત કરવાની છૂટ આપે છે:
\[
\Axiom$ !A, !B, \Gamma \fCenter \Delta$
\RightLabel{\LeftR{\land}}
\UnaryInf$ !A \land !B, \Gamma \fCenter \Delta$
\DisplayProof
\qquad
\Axiom$\Gamma \fCenter \Delta, !A$
\Axiom$ \Gamma \fCenter \Delta, !B$
\RightLabel{\RightR{\land}}
\BinaryInf$ \Gamma \fCenter \Delta, !A \land !B $
\DisplayProof
\]
રેખાની ઉપરના સિક્વન્ટોને \emph{આધાર-સિક્વન્ટો} અને
રેખાની નીચેના સિક્વન્ટને \emph{ફલિત-સિક્વન્ટ} કહે છે.
દરેક નિયમમાં $\Gamma$ અને $\Delta$નાં સૂત્રને
\emph{સંદર્ભ} અથવા \emph{ગૌણ સૂત્ર} કહે છે.
ફલિતમાં સંદર્ભનો ભાગ ન હોય તે સૂત્રને
\emph{મુખ્ય} સૂત્ર કહે છે, અને આધારોમાં સંદર્ભનો ભાગ
ન હોય તે સૂત્રને \emph{સક્રિય} (અથવા \emph{સહાયક})
સૂત્ર કહે છે.

પરિમાણકના નિયમો થોડા વધુ જટિલ છે. દાખલા તરીકે,
\Log{G1c}માં $\lforall$ માટેના નિયમો આ છે:
\[\Axiom$ !A(t), \Gamma \fCenter \Delta$
\RightLabel{\LeftR{\lforall}}
\UnaryInf$ \lforall[x][!A(x)],\Gamma \fCenter \Delta$
\DisplayProof
\qquad
\Axiom$ \Gamma \fCenter \Delta, !A(a) $
\RightLabel{\RightR{\lforall}}
\UnaryInf$ \Gamma \fCenter \Delta, \lforall[x][!A(x)]$
\DisplayProof
\]
\LeftR{\lforall}ના નિયમમાં સક્રિય સૂત્ર એ $!A(t)$ છે,
જ્યાં $t$ બંધ પદ છે, એટલે કે ચલ વિનાનું પદ.
આવું અનુમાન યોગ્ય છે, એટલે કે \LeftR{\lforall}ના ઢાંચારૂપ
નિયમ સાથે મેળ ખાય છે, જો કોઈ સૂત્ર~$!A$, કોઈ ચલ~$x$
અને કોઈ બંધ પદ~$t$ એવાં હોય કે ફલિતમાં મુખ્ય
સૂત્ર એ $\lforall[x][!A]$ હોય અને આધારમાં સક્રિય
સૂત્ર એ $\Subst{!A}{t}{x}$ હોય. \RightR{\lforall}નો
નિયમ એ જ રીતે કામ કરે છે, પરંતુ કોઈ પણ પદ~$t$ને બદલે
આધારમાં સક્રિય સૂત્ર એ $\Subst{!A}{a}{x}$ હોવું જોઈએ,
જ્યાં $a$ એ અચળ છે. આ અચળને
\RightR{\lforall}ના અનુમાનનો \emph{આઇગનચલ} કહે છે.
અનુમાન ત્યારે જ યોગ્ય છે જ્યારે નીચેના સિક્વન્ટમાં $a$
ન હોય, એટલે કે $a$ એ $\Gamma$, $\Delta$ કે $!A$માં ન આવે.
આને \emph{આઇગનચલ-શરત} કહે છે.

\Log{G1c} અને \Log{G3c} જેવાં સિક્વન્ટ-તંત્રોમાં
આ રીતે ઢાંચારૂપે આપેલા નિયમોનો સમૂહ તથા સ્વયંસિદ્ધ ગણાતા
કેટલાક સિક્વન્ટો હોય છે; તેઓ પણ ઢાંચારૂપે અપાય છે.
આવા તંત્રની સાબિતી એ સ્વયંસિદ્ધોમાંથી અનુમાનના નિયમો
વડે આગમનથી ઉત્પન્ન કરેલાં સિક્વન્ટોનાં સાંત વૃક્ષો છે.
દરેક પર્ણ સ્વયંસિદ્ધ છે, અને બીજા દરેક સિક્વન્ટને તેની તરત
ઉપરના સિક્વન્ટોમાંથી તંત્રના કોઈ અનુમાન-નિયમ વડે મેળવાય છે.

\begin{defn}\ollabel{defn:proof}
સિક્વન્ટ-તંત્રની \emph{સાબિતી} એ નીચે પ્રમાણે
આગમનથી રચાતાં સિક્વન્ટોનાં સાંત વૃક્ષો છે:
\begin{enumerate}
  \item\ollabel{defn:prf:axiom} દરેક સ્વયંસિદ્ધ સાબિતી છે.
  \item\ollabel{defn:prf:one-prem} જો $\pi_1$ એ
  સિક્વન્ટ~$S_1$માં અંત પામતી સાબિતી હોય, અને $S$ એ
  એક આધારવાળા નિયમ~$R$ વડે $S_1$માંથી મળતો સિક્વન્ટ હોય, તો
  \[
  \AxiomC{}
  \RightLabel{$\pi_1$}
  \DeduceC{$S_1$}
  \RightLabel{$R$}
  \UnaryInfC{$S$}
  \DisplayProof
  \]
  એ સાબિતી છે.
  \item\ollabel{defn:prf:two-prem} જો $\pi_1$ અને $\pi_2$
  અનુક્રમે સિક્વન્ટ~$S_1$ અને $S_2$માં અંત પામતી સાબિતી
  હોય, અને $S$ એ બે આધારવાળા નિયમ~$R$ વડે $S_1$ અને $S_2$માંથી
  મળતો સિક્વન્ટ હોય, તો
  \[
  \AxiomC{}
  \RightLabel{$\pi_1$}
  \DeduceC{$S_1$}
  \AxiomC{}
  \RightLabel{$\pi_2$}
  \DeduceC{$S_2$}
  \RightLabel{$R$}
  \BinaryInfC{$S$}
  \DisplayProof
  \]
  એ સાબિતી છે.
\end{enumerate}
\end{defn}

\Olref{defn:proof}માં સાબિતીને સિક્વન્ટોનાં વૃક્ષો
તરીકે વ્યાખ્યાયિત કરી છે. આ વૃક્ષો ``ઉપર તરફ''
વધે છે તેમ વિચારીએ. વૃક્ષની સૌથી ઉપરની પર્ણ-ગાંઠો
સ્વયંસિદ્ધો છે અને તેમને \emph{આરંભિક સિક્વન્ટો} કહે છે.
સૌથી નીચેના સિક્વન્ટને \emph{અંતિમ સિક્વન્ટ} કહે છે.
જો $\pi$ એ અંતિમ સિક્વન્ટ~$S$ ધરાવતી સાબિતી હોય તો
$\pi \Proves S$ પણ લખીએ. સિક્વન્ટ~$S$ની સાબિતી
હોય તો $\Proves S$ લખીએ; જરૂર પડે તો $\Proves$ ચિહ્નની
ડાબે તે સાબિતીનું તંત્ર બતાવીએ, દાખલા તરીકે
$\Log{G1c} \Proves !A, !B \Sequent !A \land
!B$. સાબિતીમાં આરંભિક સિક્વન્ટો સિવાયનો દરેક સિક્વન્ટ
કોઈ નિયમ~$R$ના અનુમાનનું \emph{ફલિત} છે.
સાબિતીમાં સિક્વન્ટની તરત ઉપર આવેલા એક કે બે સિક્વન્ટોને,
એટલે કે તેની ગાંઠની પિતૃ-ગાંઠોને, અનુમાનના \emph{આધારો} કહે છે.

સાબિતીની આગમનથી થતી રચનામાં વપરાતી
સાબિતીને તેની \emph{ઉપ-સાબિતી} કહે છે.
અનુમાનના આધારોમાં અંત પામતી ઉપ-સાબિતીને
તે અનુમાનની \emph{તાત્કાલિક ઉપ-સાબિતી} કહે છે.

સાબિતીની જટિલતા પ્રાકૃતિક સંખ્યા વડે માપવી ઘણી વાર
જરૂરી બને છે. સાબિતીમાં સિક્વન્ટોની કુલ સંખ્યા ગણી શકાય,
પણ ઘણી વાર વધુ ઉપયોગી માપ સાબિતીની ઊંચાઈ છે:
અંતિમ સિક્વન્ટથી આરંભિક સિક્વન્ટ સુધીની શાખામાં અનુમાન-પગલાંની
મહત્તમ સંખ્યા. તેની આગમનથી થતી વ્યાખ્યા પણ આપી શકીએ.

\begin{defn}
સાબિતી~$\pi$ની \emph{ઊંચાઈ}~$\pheight{\pi}$
નીચે પ્રમાણે આગમનથી વ્યાખ્યાયિત છે.
\begin{enumerate}
  \item જો $\pi$ ફક્ત એક સ્વયંસિદ્ધની બનેલી હોય, તો
  $\pheight{\pi} = 0$.
  \item જો $\pi$નો અંત એક આધારવાળા અનુમાનમાં થતો હોય અને તેની
  તાત્કાલિક ઉપ-સાબિતી એ $\pi_1$ હોય, તો
  $\pheight{\pi} = \pheight{\pi_1} + 1$.
  \item જો $\pi$નો અંત બે આધારવાળા અનુમાનમાં થતો હોય અને તેની
  તાત્કાલિક ઉપ-સાબિતી એ $\pi_1$ અને $\pi_2$ હોય, તો
  $\pheight{\pi} =
  \max(\pheight{\pi_1}, \pheight{\pi_2}) + 1$.
\end{enumerate}
જો સિક્વન્ટ~$S$ને ઊંચાઈ $\le n$ની સાબિતી
હોય, તો $\Proves[n] S$ લખીએ.
\end{defn}

\begin{ex}
\Log{G3c}માં $!C, \Gamma \Sequent \Delta, !C$ સ્વરૂપનો
દરેક સિક્વન્ટ સ્વયંસિદ્ધ છે, જ્યાં $!C$ પરમાણ્વીય હોવું જોઈએ.
ધારો કે $!D$ અને $!E$ પરમાણ્વીય વાક્યો છે.
તો $!D, !E \Sequent !D$ અને $!D, !E \Sequent !E$ એ
સ્વયંસિદ્ધ $!C, \Gamma \Sequent \Delta, !C$ના દાખલા છે.
દાખલા તરીકે, $!C, \Gamma \Sequent \Delta, !C$માં
$!C \ident !E$, $\Gamma = \{!D\}$ અને
$\Delta = \emptyset$ લઈએ તો બરાબર
$!D, !E \Sequent !E$ સિક્વન્ટ મળે છે.
(યાદ રાખો કે અહીં બહુ\emph{ગણ} છે, એટલે
$!D, !E \Sequent !E$ અને $!E, !D \Sequent !E$ એક જ સિક્વન્ટ છે.)
\sourcecorrection{OLMD-339}{મૂળના છેલ્લા તુલનાત્મક સિક્વન્ટના ઉત્તરાંગમાં પહેલાથી જુદું સૂત્ર હતું. બહુગણનો ક્રમ બદલવાથી પૂર્વાંગમાં ફેર નથી પડતો, પરંતુ ઉત્તરાંગનું સૂત્ર બદલવું તે સમાન સિક્વન્ટ આપતું નથી. તેથી બંને ઉત્તરાંગમાં એક જ સૂત્ર રાખ્યું છે.}

$!D, !E \Sequent !D$ એ \Log{G3c}ના સ્વયંસિદ્ધનો દાખલો
હોવાથી તે એકલો પણ \Log{G3c}માં સાબિતી ગણાય છે;
$!D, !E \Sequent !E$ માટે પણ
\olref{defn:proof}\olref{defn:prf:axiom} પ્રમાણે એમ જ છે.
તેમને $\pi_1$ અને $\pi_2$ કહીએ.
\olref{defn:prf:two-prem} પ્રમાણે આ બે સાબિતી લઈને
બે આધારવાળા નિયમ~$\RightR{\land}$ વડે નવી
સાબિતી~$\pi_3$ ઉત્પન્ન કરી શકીએ:
\[
\Axiom$!D, !E \fCenter !D$
\Axiom$!D, !E \fCenter !E$
\RightLabel{\RightR{\land}}
\BinaryInf$!D, !E  \fCenter !D \land !E $
\DisplayProof
\]
પછી $\pi_3$માંથી સાબિતી~$\pi_4$ મળે છે:
\[
\Axiom$!D, !E \fCenter !D$
\Axiom$!D, !E \fCenter !E$
\RightLabel{\RightR{\land}}
\insertBetweenHyps{\hspace{5ex}}
\BinaryInf$!D, !E  \fCenter !D \land !E $
\LeftSubproofLabel{$\pi_3$}
\RightLabel{\LeftR{\land}}
\UnaryInf$!D \land !E  \fCenter !D \land !E$
\RightSubproofLabel{$\pi_4$}
\DisplayProof
\]
અહીં એક આધારવાળો નિયમ~$\LeftR{\land}$ વાપર્યો છે
(\olref{defn:proof}\olref{defn:prf:one-prem} પ્રમાણે).
$\pi_4$ના છેલ્લા અનુમાનની તાત્કાલિક ઉપ-સાબિતી એ $\pi_3$ છે.
$\RightR{\land}$ના અનુમાનની તાત્કાલિક ઉપ-સાબિતી એ
$\pi_1$ અને $\pi_2$ છે. $\pi_4$ની ઉપ-સાબિતીમાં
$\pi_1$, $\pi_2$, $\pi_3$ અને $\pi_4$ પોતે આવે છે.
અહીં $\pheight{\pi_1} =
\pheight{\pi_2} = 0$, $\pheight{\pi_3} = 1$ અને
$\pheight{\pi_4} = 2$ છે. કોઈ પણ પરમાણ્વીય $!D$ અને
$!E$ માટે $\Log{G3c} \Proves[4] !D \land !E \fCenter !D
\land !E$ બતાવ્યું છે.
\end{ex}
% END OLP-0711
\endgroup

\begingroup
\makeatletter\def\ol@sectioncs{section}\makeatother

% BEGIN OLP-0702
\olfileid{pt}{seq}{ex}

\olsection{સાબિતીઓનાં ઉદાહરણો}
\phantomsection\label{guunit:OLP-0702}


સિક્વન્ટોની સાબિતીઓ આપવા
સામાન્ય રીતે અંતિમ સિક્વન્ટથી ઉપર તરફ કામ કરવું અનુકૂળ છે:
તેમાંના સૂત્રને મુખ્ય સૂત્ર તરીકે પસંદ કરી,
અનુરૂપ નિયમના આધારો સિક્વન્ટની ઉપર લખીએ;
પછી સ્વયંસિદ્ધો સુધી પહોંચીએ ત્યાં સુધી પ્રક્રિયા દોહરાવીએ.
\sourcecorrection{OLMD-315}{મૂળ પાછી તરફ સાબિતી-શોધમાં નિષ્કર્ષમાંથી પસંદ કરેલા સૂત્રને સક્રિય કહે છે. નિયમની વ્યાખ્યા મુજબ નિષ્કર્ષનું પસંદ કરેલું સૂત્ર મુખ્ય છે; આધારોમાં તેનાથી મળતાં સૂત્રો સક્રિય છે. અહીં શરૂઆત અને બીજા અનુસરણના પગલામાં મુખ્ય સૂત્ર કહ્યું છે.}
દાખલા તરીકે આ સિક્વન્ટ સાબિત કરવો છે એમ ધારો:
\[(!C \land !D) \lif !E \Sequent !C \lif (!D \lif !E)\]
$!C \lif (!D \lif !E)$ વિચારીએ:
તે $!A \lif !B$ સ્વરૂપનું છે
અને સિક્વન્ટની જમણી બાજુ છે.
આખા સિક્વન્ટને \Log{G1c}ના $\RightR{\lif}$ નિયમ સાથે સરખાવીએ,
\[\Axiom$ !A, \Gamma \fCenter \Delta, !B$
\RightLabel{\RightR{\lif}}
\UnaryInf$ \Gamma \fCenter \Delta, !A \lif !B $
\DisplayProof\]
તો $\Gamma = \{(!C \land !D) \lif !E\}$,
$\Delta = \emptyset$, $!A \ident !C$
અને $!B \ident !D \lif !E$ લેવા જોઈએ તે સૂચવે છે.
\sourcecorrection{OLMD-316}{મૂળ આ પગલામાં સંદર્ભને નિષ્કર્ષના જમણા સૂત્રથી ભરે છે. દર્શાવેલા સિક્વન્ટનો પૂર્વાંગ તો સંયોજનમાંથી અનુસરણ છે; Gammaમાં એ જ પૂર્વાંગ રાખ્યો છે, જેમ તરત પછીનું વૃક્ષ કરે છે.}
એટલે સાબિતીનું છેલ્લું અનુમાન આ હોઈ શકે:
\[\Axiom$ !C, (!C \land !D) \lif !E \fCenter !D \lif !E$
\RightLabel{\RightR{\lif}}
\UnaryInf$(!C \land !D) \lif !E \fCenter !C \lif (!D \lif !E)$
\DisplayProof\]
$!D \lif !E$ને મુખ્ય સૂત્ર ગણીને વિચાર દોહરાવીએ,
તો આ મળે:
\[\Axiom$ !D, !C, (!C \land !D) \lif !E \fCenter !E$
\RightLabel{\RightR{\lif}}
\UnaryInf$!C, (!C \land !D) \lif !E \fCenter !D \lif !E$
\RightLabel{\RightR{\lif}}
\UnaryInf$(!C \land !D) \lif !E \fCenter !C \lif (!D \lif !E)$
\DisplayProof\]
હવે $(!C \land !D) \lif !E$ પર ધ્યાન આપવું
અને \Log{G1c}નો $\LeftR{\lif}$ નિયમ વાપરવો જોઈએ:
\[\Axiom$ \Gamma \fCenter \Delta, !A$
\Axiom$ !B, \Gamma \fCenter \Delta$
\RightLabel{\LeftR{\lif}}
\BinaryInf$ !A \lif !B, \Gamma \fCenter \Delta$
\DisplayProof\]
તેથી આ મળે:
\[\Axiom$!D, !C \fCenter !E, !C \land !D$
\Axiom$!D, !C, !E \fCenter !E$
\RightLabel{\LeftR{\lif}}
\BinaryInf$ !D, !C, (!C \land !D) \lif !E \fCenter !E$
\RightLabel{\RightR{\lif}}
\UnaryInf$!C, (!C \land !D) \lif !E \fCenter !D \lif !E$
\RightLabel{\RightR{\lif}}
\UnaryInf$(!C \land !D) \lif !E \fCenter !C \lif (!D \lif !E)$
\DisplayProof\]
સૌથી ઉપરના ડાબા સિક્વન્ટમાં સૂત્ર $!C \land !D$ને
મુખ્ય ગણી $\RightR{\land}$ લગાવીએ, તો આ મળે:
\[
\Axiom$!D, !C \fCenter !E, !C$
\Axiom$!D, !C \fCenter !E, !D$
\RightLabel{\RightR{\land}}
\BinaryInf$!D, !C \fCenter !E, !C \land !D$
\Axiom$!D, !C, !E \fCenter !E$
\RightLabel{\LeftR{\lif}}
\BinaryInf$ !D, !C, (!C \land !D) \lif !E \fCenter !E$
\RightLabel{\RightR{\lif}}
\UnaryInf$!C, (!C \land !D) \lif !E \fCenter !D \lif !E$
\RightLabel{\RightR{\lif}}
\UnaryInf$(!C \land !D) \lif !E \fCenter !C \lif (!D \lif !E)$
\DisplayProof
\]
\sourcecorrection{OLMD-317}{મૂળના આ વૃક્ષમાં બે આધારવાળા અનુસરણ-ડાબા પગલાનું લેબલ અનુસરણ-જમણું છે. બંને આધાર અને મુખ્ય પૂર્વાંગ પ્રમાણે તેને ડાબું કર્યું છે.}
અત્યાર સુધી બધું યોગ્ય છે.
અત્યાર સુધી વાપરેલા નિયમો \Log{G1c} અને~\Log{G3c}માં સમાન હોવાથી
આ બધું \Log{G3c}માં પણ ચાલે છે.
આપણે સાબિતીના એવા અંશ સુધી પહોંચ્યા છીએ
જેના દરેક સૌથી ઉપરના સિક્વન્ટમાં ડાબે અને જમણે એ જ સૂત્ર છે.
પરંતુ તેઓ હજી સ્વયંસિદ્ધો નથી.

\Log{G1c}માં $!A \Sequent !A$ સ્વરૂપના સ્વયંસિદ્ધોમાંથી
સૌથી ઉપરના સિક્વન્ટોની સાબિતીઓ આપવી પડે.
શિથિલીકરણના નિયમોથી આ સરળતાથી થઈ શકે.
દાખલા તરીકે ડાબે સૌથી ઉપરનો સિક્વન્ટ
$!D, !C \Sequent !E, !C$ આ રીતે સાબિત થઈ શકે:
\[
\Axiom$!C \fCenter !C$
\RightLabel{\LeftR{\Weakening}}
\UnaryInf$!D, !C \fCenter !C$
\RightLabel{\RightR{\Weakening}}
\UnaryInf$!D, !C \fCenter !E, !C$
\DisplayProof
\]
\sourcecorrection{OLMD-318}{મૂળ આ નાનાં શિથિલીકરણ-વૃક્ષ પછી તેને દર્શાવતો DisplayProof આદેશ છોડે છે. માત્ર તે આદેશ ઉમેર્યો છે; વૃક્ષનાં સૂત્રો અને નિયમો અચળ છે.}
આખરે \Log{G1c}માં આ સ્વરૂપની સાબિતી મળે:
\[
\Axiom$!C \fCenter !C$
\RightLabel{\LeftR{\Weakening}}
\UnaryInf$!D, !C \fCenter !C$
\RightLabel{\RightR{\Weakening}}
\UnaryInf$!D, !C \fCenter !E, !C$
\Axiom$!D \fCenter !D$
\RightLabel{\LeftR{\Weakening}}
\UnaryInf$!D, !C \fCenter !D$
\RightLabel{\RightR{\Weakening}}
\UnaryInf$!D, !C \fCenter !E, !D$
\RightLabel{\RightR{\land}}
\BinaryInf$!D, !C \fCenter !E, !C \land !D$
\Axiom$!E \fCenter !E$
\RightLabel{\LeftR{\Weakening}}
\UnaryInf$!C, !E \fCenter !E$
\RightLabel{\LeftR{\Weakening}}
\UnaryInf$!D, !C, !E \fCenter !E$
\RightLabel{\LeftR{\lif}}
\BinaryInf$ !D, !C, (!C \land !D) \lif !E \fCenter !E$
\RightLabel{\RightR{\lif}}
\UnaryInf$!C, (!C \land !D) \lif !E \fCenter !D \lif !E$
\RightLabel{\RightR{\lif}}
\UnaryInf$(!C \land !D) \lif !E \fCenter !C \lif (!D \lif !E)$
\DisplayProof
\]
\sourcecorrection{OLMD-319}{મૂળના જમણા સ્વયંસિદ્ધમાંથી પહેલા શિથિલીકરણના પગલામાં ઉત્તરાંગ બદલીને બીજું સૂત્ર કરે છે. અહીં ઉત્તરાંગનું એ જ $!E$ રાખ્યું છે; શિથિલીકરણ ફક્ત ડાબે $!C$ ઉમેરે છે.}
આ સાબિતીની સંરચના (ખાસ કરીને ઊંચાઈ)
સૂત્રો $!C$, $!D$ અને~$!E$ કયા છે તેનાથી સ્વતંત્ર છે.
ઉપરની સાબિતીમાં આ ચિહ્નોની જગ્યાએ
કોઈ પણ સૂત્રો મૂકી શકાય,
અને પરિણામ હજી~\Log{G1c}ની સાબિતી હોય.
\Log{G1c}ની સાબિતી \emph{ઢાંચારૂપ} છે.

\Log{G3c}માં સ્વયંસિદ્ધો
$!A, \Gamma \Sequent \Delta, !A$ સ્વરૂપના સિક્વન્ટો છે,
જ્યાં $!A$ આણ્વિક હોવું જોઈએ.
એટલે $!D, !C \Sequent !E, !C$ એ \Log{G3c}ના
સ્વયંસિદ્ધ જેવો લાગે છે
($\Gamma = \{!D\}$; $\Delta = \{!E\}$),
પરંતુ $!C$ ખરેખર આણ્વિક હોય ત્યારે
તે આ મેળ ખાતી આવૃત્તિઓથી \emph{ખરેખર} સ્વયંસિદ્ધ છે.
$!C$, $!D$ અને $!E$ આણ્વિક સૂત્રો હોય,
તો આપણો સાબિતી-અંશ \Log{G3c}માં સંપૂર્ણ સાબિતી છે.
\sourcecorrection{OLMD-320}{મૂળ આણ્વિકતાને આ સિક્વન્ટો સ્વયંસિદ્ધ હોવાની આવશ્યક શરત કહે છે. અહીં બતાવેલી મેળ ખાતી આવૃત્તિઓ આણ્વિક હોય તે પૂરતી શરત છે; બીજા આણ્વિક સૂત્રનો બંને બાજુ મેળ અથવા ડાબે અસત્ય હોવાથી પણ સિક્વન્ટ સ્વયંસિદ્ધ હોઈ શકે. તેથી પૂરતી શરત તરીકે કહ્યું છે.}

સખત રીતે જોતાં આપણે
\[\Log{G3c} \Proves (!C \land !D) \lif !E \fCenter !C \lif (!D \lif
!E),\]
બતાવ્યું નથી, છતાં વ્યવહારમાં આપણે આટલું જ કરવાની
જરૂર છે (અને આટલું જ કરી શકીએ).
કારણ કે $!A$ પોતે આણ્વિક ન હોય તો પણ
$!A, \Gamma \Sequent \Delta, !A$ સ્વરૂપનો દરેક સિક્વન્ટ
\Log{G3c}માં સાબિત થઈ શકે છે.
$!A$ની ઊંડાઈ~$\depth{!A}$ પર આગમનથી આ બતાવી શકીએ.
\sourcecorrection{OLMD-321}{મૂળ આ વાક્યના પ્રદર્શનને G1c કહે છે, પરંતુ એ તંત્રની સંપૂર્ણ ઢાંચારૂપ સાબિતી ઉપર આપી ચૂક્યું છે. અહીં હજી પૂર્ણ ન થયેલી વાત G3cમાં અણઆણ્વિક સ્વયંસિદ્ધોના સ્થાને સાબિતીઓ આપવાની છે; તેથી તંત્રનું નામ G3c કર્યું છે.}

\begin{defn}\ollabel{defn:depth}
સૂત્રની \emph{ઊંડાઈ}~$\depth{!A}$ આ રીતે વ્યાખ્યાયિત છે:
\begin{enumerate}
  \item $!A$ આણ્વિક હોય, તો $\depth{!A} = 0$.
  \item $!A \ident \lnot !B$, $!A \ident \lforall[x][!B]$
  અથવા $!A \ident \lexists[x][!B]$ હોય,
  તો $\depth{!A} = \depth{!B} + 1$.
  \item $!A \ident (!B \land !C)$, $(!B \lor !C)$
  અથવા $(!B \lif !C)$ હોય,
  તો $\depth{!A} = \max(\depth{!B},\depth{!C}) + 1$.
\end{enumerate}
\sourcecorrection{OLMD-322}{મૂળ પરિમાણકનાં બંને કિસ્સામાં અંદરનું સૂત્ર $A$ લખે છે, પરંતુ ઊંડાઈની જમણી બાજુ અસંબંધિત $!B$ છે. બંને અંદરના સૂત્રોને એ જ $!B$ કર્યા છે, જેમ નિષેધના કિસ્સામાં છે.}
\end{defn}

કોઈ પણ પદ~$t$ માટે $\depth{A(x)} = \depth{A(t)}$
સાબિત કરવું મુશ્કેલ નથી.
આપણા માટે મહત્ત્વની વાત એ છે કે વિઘટનના તાર્કિક નિયમોમાં
મુખ્ય સૂત્રની ઊંડાઈ તેમાંથી મળતાં સક્રિય
ઉપ-સૂત્રોની ઊંડાઈ કરતાં મોટી છે.
\sourcecorrection{OLMD-323}{મૂળ આ ઊંડાઈની કડક અસમાનતા દરેક નિયમ માટે કહે છે. સંકોચન જેવા બંધારણીય નિયમમાં એક જ સૂત્રની ઊંડાઈ બદલાતી નથી; મુખ્ય સૂત્રની નકલ જાળવતા પરિમાણક-નિયમમાં પણ એ નકલની ઊંડાઈ સમાન છે. અહીં ઊંડાઈ ઘટતી હોય તે સાચા વિઘટિત ઉપસૂત્રોની વાત સ્પષ્ટ કરી છે.}
દાખલા તરીકે:
\begin{align*}
\depth{P(x)} = \depth{P(a)} & = 0,\\
\depth{\lforall[x][P(x)]} & = 1,\\
\depth{(\lforall[x][P(x)] \lif P(a))} & = \max(1,0) + 1 = 2.
\end{align*}
આથી~$\depth{!A}$ પર આગમનથી પરિણામો સાબિત કરી શકીએ,
જેમ કે આ પરિણામ.

\begin{prop}\ollabel{prop:G1c-axioms}
કોઈ પણ~$!A$ માટે $!A \Sequent !A$ની $ \Log{G1c}$માં
એવી સાબિતી છે જેના બધા સ્વયંસિદ્ધો આણ્વિક છે.
\end{prop}

\begin{proof}
  $\depth{!A}$ પર આગમનથી સાબિતી કરીએ.
  $\depth{!A} = 0$ હોય, તો $!A$ આણ્વિક છે
  અને $!A \Sequent !A$ પોતે~\Log{G1c}ની
  જરૂરી સ્વરૂપની સાબિતી છે.
  $\depth{!A} > 0$ હોય, તો $!A$ ઓછી ઊંડાઈનાં
  સૂત્રોમાંથી બનેલું છે,
  દાખલા તરીકે $!A \ident \lnot !B$,
  $!A \ident (!B \land !C)$
  અથવા $!A \ident \lforall[x][A(x)]$.
  આગમન-ધારણા એ છે કે
  $\depth{!D} < \depth{!A}$ હોય તે બધાં $!D$ માટે
  આણ્વિક સ્વયંસિદ્ધોમાંથી $\Log{G1c} \Proves !D \Sequent !D$ છે.

  $!A \ident \lnot !B$ હોય, તો $\depth{!B} < \depth{!A}$ છે.
  આગમન-ધારણાથી \Log{G1c}માં આણ્વિક સ્વયંસિદ્ધોમાંથી
  $!B \Sequent !B$ની સાબિતી~$\pi_1$ છે.
  તેને આ રીતે $ \lnot !B \Sequent \lnot !B$ની સાબિતીમાં વિસ્તારી શકીએ:
\[
  \AxiomC{}
  \RightLabel{$\pi_1$}
  \Deduce$!B \fCenter !B$
  \RightLabel{\LeftR{\lnot}}
  \UnaryInf$\lnot !B, !B \fCenter $
  \RightLabel{\RightR{\lnot}}
  \UnaryInf$\lnot !B \fCenter \lnot !B$
  \DisplayProof
\]

  $!A \ident (!B \land !C)$ હોય,
  તો $\depth{!B} < \depth{!A}$ અને $\depth{!C} < \depth{!A}$ છે.
  આગમન-ધારણાથી \Log{G1c}માં આણ્વિક સ્વયંસિદ્ધોમાંથી
  $!B \Sequent !B$ અને $!C \Sequent !C$ની
  સાબિતીઓ~$\pi_1$ અને~$\pi_2$ છે.
  તેમને $!B \land !C \Sequent !B \land !C$ની
  સાબિતીમાં આ રીતે વિસ્તારી શકીએ:
\[
  \AxiomC{}
  \RightLabel{$\pi_1$}
  \Deduce$!B, \fCenter !B$
  \RightLabel{\LeftR{\Weakening}}
  \UnaryInf$!B, !C, \fCenter !B$
  \RightLabel{\LeftR{\land}}
  \UnaryInf$!B \land !C, !C\fCenter !B$
  \RightLabel{\LeftR{\land}}
  \UnaryInf$!B \land !C, !B \land !C \fCenter !B$
  \AxiomC{}
  \RightLabel{$\pi_2$}
  \Deduce$!C, \fCenter !C$
  \RightLabel{\LeftR{\Weakening}}
  \UnaryInf$!B, !C, \fCenter !C$
  \RightLabel{\LeftR{\land}}
  \UnaryInf$!B \land !C, !C \fCenter !C$
  \RightLabel{\LeftR{\land}}
  \UnaryInf$!B \land !C, !B \land !C \fCenter !C$
  \RightLabel{\RightR{\land}}
  \BinaryInf$!B \land !C, !B \land !C \fCenter !B \land !C$
  \RightLabel{\LeftR{\Contraction}}
  \UnaryInf$!B \land !C \fCenter !B \land !C$
  \DisplayProof
  \]
  \sourcecorrection{OLMD-324}{મૂળના છેલ્લાં બે પગલામાં સંયોજન-જમણાનો નિષ્કર્ષ એક ડાબી આવૃત્તિ ખોટી રીતે ગુમાવે છે અને પછી ડાબું સંકોચન જમણું સંયોજન બદલી નાખે છે. પહેલા પગલામાં ડાબે બે સંયોજન રાખ્યાં છે; સંકોચન પછી ડાબે એક સંયોજન અને જમણે એ જ સંયોજન રહે છે.}

  હવે $!A \ident \lforall[x][!B(x)]$ છે એમ ધારો.
  $\lforall[x][!B(x)]$માં ન આવતો અચળ~$c$ લો.
  ત્યારે $\depth{!B(c)} = \depth{!B(x)} < \depth{!A}$ છે.
  આગમન-ધારણાથી \Log{G1c}માં આણ્વિક સ્વયંસિદ્ધોમાંથી
  $!B(c) \Sequent !B(c)$ની સાબિતી~$\pi_1$ છે.
  તેને $\lforall[x][!B(x)] \Sequent \lforall[x][!B(x)]$ની
  સાબિતીમાં આ રીતે વિસ્તારી શકીએ:
\[
  \AxiomC{}
  \RightLabel{$\pi_1$}
  \Deduce$!B(c) \fCenter !B(c)$
  \RightLabel{\LeftR{\lforall}}
  \UnaryInf$\lforall[x][!B(x)] \fCenter !B(c)$
  \RightLabel{\RightR{\lforall}}
  \UnaryInf$\lforall[x][!B(x)] \fCenter \lforall[x][!B(x)]$
  \DisplayProof
\]
  \sourcecorrection{OLMD-325}{મૂળ આ પરિચ્છેદના કેટલાક સૂત્રોમાં સામાન્ય $B$ અને કેટલાકમાં સૂત્ર-સ્થાનધારક $!B$ વાપરે છે; વૃક્ષમાં $!B$ છે. ગદ્યની સર્વપરિમાણક-ઓળખ અને દાખલની ઊંડાઈ પણ એ જ $!B$ સાથે સુસંગત કરી છે.}

  બાકીના કિસ્સા સ્વાધ્યાય તરીકે છોડ્યા છે.
\end{proof}

\begin{prob}
  $!A \ident (!B \lor !C)$, $!A \ident (!B \lif !C)$
  અને $!A \ident \lexists[x][B(x)]$ હોય તે કિસ્સા કરીને
  \olref{prop:G1c-axioms}ની સાબિતી પૂર્ણ કરો.
\end{prob}

આપણા હેતુ માટે $!A \ident \lnot !B$ના કિસ્સામાં
આ સાબિતી પણ વાપરી શક્યા હોત:
\[
\AxiomC{}
\RightLabel{$\pi_1$}
\Deduce$!B \fCenter !B$
\RightLabel{\RightR{\lnot}}
\UnaryInf$\fCenter !B, \lnot !B$
\RightLabel{\LeftR{\lnot}}
\UnaryInf$\lnot !B \fCenter \lnot !B$
\DisplayProof
\]
પરંતુ કેટલાંક સિક્વન્ટ તંત્રોમાં આ ન ચાલ્યું હોત.
અંતઃપ્રજ્ઞાવાદી સિક્વન્ટ કલન \Log{G1i} એ \Log{G1c}ને અનુરૂપ છે;
તેમાં દરેક સિક્વન્ટની જમણી બાજુ વધુમાં વધુ એક સૂત્રની મર્યાદા છે.
$\Delta = \emptyset$ હોય, તો પહેલી સાબિતી
\Log{G1i}માં પણ સાબિતી હોય,
પરંતુ બીજી ન હોય, કારણ કે એક સિક્વન્ટની જમણી બાજુ
બંને $!B$ અને~$\lnot !B$ છે.
\sourcecorrection{OLMD-326}{મૂળ અહીં ફરી G1iનો તફાવત માત્ર જમણી બાજુની સંખ્યાની મર્યાદા તરીકે કહે છે. અગાઉના અર્થઘટન-વિભાગમાં નોંધ્યું છે તેમ કેટલાક નિયમોના આધારનું સ્વરૂપ પણ બદલાય છે. અહીં આ મર્યાદા બંને નિષેધ-વૃક્ષો વચ્ચેનો બતાવેલો તફાવત સમજાવે છે.}

\Log{G1c}માં પણ $!A \ident \lforall[x][!B(x)]$
હોય ત્યારે નિયમોનો બીજો ક્રમ વાપરી ન શક્યા હોત તે નોંધો.
આ સાબિતી \emph{નથી}:
\[
  \AxiomC{}
  \RightLabel{$\pi_1$}
  \Deduce$!B(c) \fCenter !B(c)$
  \RightLabel{\RightR{\lforall}}
  \UnaryInf$!B(c) \fCenter \lforall[x][!B(x)]$
  \RightLabel{\LeftR{\lforall}}
  \UnaryInf$\lforall[x][!B(x)] \fCenter \lforall[x][!B(x)]$
  \DisplayProof
\]
કારણ કે~$\RightR{\lforall}$ની આઇગનચલ-શરત તૂટે છે.

\begin{prop}\ollabel{prop:G3c-axioms}
$!A, \Gamma \Sequent \Delta, !A$ સ્વરૂપનો દરેક સિક્વન્ટ
($!A$ આણ્વિક હોવું \emph{જરૂરી નથી})
\Log{G3c}માં સાબિત થઈ શકે છે.
\end{prop}

\begin{proof}
સાબિતી \olref{prop:G1c-axioms}ની સાબિતી જેવી છે,
પરંતુ આગમન-ધારણા વિશે સાવચેતી લેવી જોઈએ.
$n = \depth{!A} = 0$નો કિસ્સો ફરી સરળ છે.
$n=0$ હોય, તો $!A$ આણ્વિક છે,
અને $!A, \Gamma \Sequent \Delta, !A$ એ~\Log{G3c}નું સ્વયંસિદ્ધ છે.

હવે $n>0$ ધારો. ફરી કિસ્સા જુદા પાડીએ.
આગમન-ધારણા આ છે:
$\depth{!D} < n$ હોય તે બધાં $!D$ માટે,
અને બધા $\Pi$ તથા $\Lambda$ માટે,
$\Log{G3c} \Proves !D, \Pi \Sequent \Lambda, !D$ છે.
\sourcecorrection{OLMD-327}{મૂળ આ સ્થાનધારકને $D!$ લખે છે; બાકીના સૂત્રો પ્રમાણે તેને $!D$ કર્યો છે.}
$!A \ident (!B \land !C)$નો કિસ્સો લઈએ.
આગમન-ધારણા બે વાર વાપરીએ:
એક વાર $!D = !B$, $\Pi = !C, \Gamma$
અને $\Lambda = \Delta$ લઈને
$!B, !C, \Gamma \Sequent \Delta, !B$ની સાબિતી~$\pi_1$ મેળવીએ;
બીજી વાર $!D = !C$, $\Pi = !B, \Gamma$
અને $\Lambda = \Delta$ લઈને
$!B, !C, \Gamma \Sequent \Delta, !C$ની સાબિતી~$\pi_2$ મેળવીએ.
\sourcecorrection{OLMD-328}{મૂળ બીજી આગમન-અરજીમાં ડાબા સંદર્ભ માટે જમણા સંદર્ભનું ચિહ્ન Lambda લખે છે. દર્શાવેલા આધાર અને વૃક્ષ મુજબ અહીં એ જ ડાબો Gamma સંદર્ભ રાખ્યો છે.}
$!B \land !C, \Gamma \Sequent \Delta, !B \land !C$ની
સાબિતી આ છે:
\[
\AxiomC{}
\RightLabel{$\pi_1$}
\Deduce$!B, !C, \Gamma \fCenter \Delta, !B$
\RightLabel{\LeftR{\land}}
\UnaryInf$!B \land !C, \Gamma \fCenter \Delta, !B$
\AxiomC{}
\RightLabel{$\pi_2$}
\Deduce$!B, !C, \Gamma \fCenter \Delta, !C$
\RightLabel{\LeftR{\land}}
\UnaryInf$!B \land !C, \Gamma \fCenter \Delta, !C$
\RightLabel{\RightR{\land}}
\BinaryInf$!B \land !C, \Gamma \fCenter \Delta, !B \land !C$
\DisplayProof
\]

$!A \ident \lforall[x][!B(x)]$ માટે
$\lforall[x][!B(x)]$, $\Gamma$ કે~$\Delta$માં
ન આવતો અચળ પસંદ કરો.
$!D = !B(c)$, $\Pi = \lforall[x][!B(x)], \Gamma$
અને $\Lambda = \Delta$ માટે આગમન-ધારણા લગાવો.
તેથી $!B(c), \lforall[x][!B(x)], \Gamma \Sequent \Delta, !B(c)$ની
સાબિતી~$\pi_1$ મળે છે, જેમાંથી આ મેળવીએ:
\[
\AxiomC{}
\RightLabel{$\pi_1$}
\Deduce$!B(c), \lforall[x][!B(x)], \Gamma \fCenter \Delta, !B(c)$
\RightLabel{\LeftR{\lforall}}
\UnaryInf$\lforall[x][!B(x)], \Gamma \fCenter \Delta, !B(c)$
\RightLabel{\RightR{\lforall}}
\UnaryInf$\lforall[x][!B(x)], \Gamma \fCenter \Delta, \lforall[x][!B(x)]$
\DisplayProof
\]
$c$ની પસંદગીથી $\RightR{\lforall}$ની આઇગનચલ-શરત સંતોષાય છે.
\end{proof}

\begin{prob}
  $!A \ident (!B \lif !C)$ અને $!A \ident \lexists[x][B(x)]$
  હોય તે કિસ્સા કરીને \olref{prop:G3c-axioms}ની સાબિતી આગળ વધારો.
\end{prob}
% END OLP-0702
\endgroup

\begingroup
\makeatletter\def\ol@sectioncs{section}\makeatother

% BEGIN OLP-0699
\olfileid{pt}{seq}{irl}

\olsection{નિયમોનું અર્થઘટન}
\phantomsection\label{guunit:OLP-0699}


સિક્વન્ટનું અર્થઘટન યાદ કરો:
$\Gamma \Sequent \Delta$ કોઈ સંરચના~$\Struct{M}$માં સત્ય છે,
જો ઓછામાં ઓછું એક $!A \in \Gamma$ અસત્ય હોય
અથવા ઓછામાં ઓછું એક~$!A \in \Delta$ સત્ય હોય.
\Log{G3c}ના તાર્કિક નિયમોને તેમના મુખ્ય સૂત્રોની
સત્યતાની શરતો રજૂ કરતા નિયમો તરીકે વાંચી શકીએ.
\RightR{\lif} વિચારીએ. તેનું મુખ્ય સૂત્ર $!A \lif !B$ છે.
$!A$ અસત્ય હોય અથવા~$!B$ સત્ય હોય, તો તે સત્ય છે.
એટલે $\ \Sequent !A \lif !B$ સત્ય છે
જો અને ફક્ત જો $!A \Sequent !B$ સત્ય હોય.
આ સિક્વન્ટોના પૂર્વાંગ અને ઉત્તરાંગમાં સંદર્ભનાં
સૂત્રો $\Gamma$, $\Delta$ ઉમેરીએ,
તો બરાબર \RightR{\lif} નિયમનો નિષ્કર્ષ અને આધાર મળે છે.
\LeftR{\lif} માટે પણ સ્થિતિ સમાન છે.
$!A$~સત્ય અને $!B$~અસત્ય હોય, તો $!A \lif !B$ અસત્ય છે.
એટલે $!A \lif !B \Sequent \ $ સત્ય છે
જો અને ફક્ત જો બંને $\ \Sequent !A$ અને $!B \Sequent \ $ સત્ય હોય.
\LeftR{\lif}નો નિષ્કર્ષ તેના આધારોના સંયોજનને સમકક્ષ છે.
આ સ્વરૂપ~\Log{G3c}ના નિયમોની યથાર્થતા સુનિશ્ચિત કરે છે
(બધા આધાર સત્ય હોય, તો નિષ્કર્ષ સત્ય છે).
તે પૂર્ણતા પણ સુનિશ્ચિત કરે છે.

કોઈ પણ સત્ય-વિધેયાત્મક સંયોજકની સત્યતાની શરતો
આ રીતે સિક્વન્ટોના સમૂહથી વર્ણવી શકાય.
આ એ હકીકતમાંથી અનુસરે છે કે શાસ્ત્રીય તર્કમાં
દરેક સત્ય-વિધેયને સંયોજક સામાન્ય સ્વરૂપના સૂત્રથી વર્ણવી શકાય:
આણ્વિક અને નિષેધિત આણ્વિક સૂત્રોના વિકલ્પોનું સંયોજન.
દાખલા તરીકે $!A \oplus !B$ને સત્ય ગણીએ
જો અને ફક્ત જો $!A$ અને $!B$માંથી બરાબર એક સત્ય હોય,
એટલે ``અપવર્જક વિકલ્પ.''
ત્યારે $!A \oplus !B$ સત્ય છે જો અને ફક્ત જો
$(!A \lor !B) \land (\lnot !A \lor \lnot !B)$ સત્ય હોય,
અને $(\lnot !A \lor !B) \land (!A \lor \lnot !B)$ સત્ય હોય
તો તે અસત્ય છે. એટલે $\oplus$ માટે
સિક્વન્ટ કલનના નિયમો આ રીતે રજૂ કરી શકીએ:
\begin{align*}
&
\Axiom$\Gamma \fCenter \Delta, !A, !B$
\Axiom$!A, !B, \Gamma \fCenter \Delta$
\RightLabel{\RightR{\oplus}}
\BinaryInf$\Gamma \fCenter \Delta, !A \oplus !B$
\DisplayProof
\\
& \Axiom$!A, \Gamma \fCenter \Delta, !B$
\Axiom$!B, \Gamma \fCenter \Delta, !A$
\RightLabel{\LeftR{\oplus}}
\BinaryInf$!A \oplus !B, \Gamma \fCenter \Delta$
\DisplayProof
\end{align*}
આ નિયમો પણ યથાર્થ અને પૂર્ણ હોય.
તેમનાં નિષ્કર્ષ સત્ય છે જો અને ફક્ત જો બધા આધાર સત્ય હોય.
\sourcecorrection{OLMD-304}{મૂળના બીજા અપવર્જક-વિકલ્પ નિયમમાં મુખ્ય સૂત્ર ડાબે છે પરંતુ નિયમ-લેબલ જમણો કહે છે. તેને ડાબો નિયમ કર્યો છે; આધાર અને નિષ્કર્ષ અચળ છે.}

\begin{prob}
$!A \downarrow !B$ સત્ય છે જો અને ફક્ત જો
$!A$ કે~$!B$માંથી કોઈ સત્ય ન હોય (``NOR'') એમ ધારો.
તેના \Log{G3c} નિયમો કયા હોય?
(એવા બે સિક્વન્ટો શોધો જે એકસાથે સત્ય હોય
જો અને ફક્ત જો $!A \downarrow !B$ સત્ય હોય.
તેથી \RightR{\downarrow} નિયમ મળે.
પછી એવો એક સિક્વન્ટ શોધો જે સત્ય હોય
જો અને ફક્ત જો $!A \downarrow !B$ અસત્ય હોય.
તેથી \LeftR{\downarrow} નિયમ મળે.)
\end{prob}

\Log{G3c}માં દરેક સંયોજક અને પરિમાણકના બરાબર બે નિયમ છે:
એક ડાબો અને એક જમણો.
પરંતુ \Log{G1c}માં બે \LeftR{\land} અને બે \RightR{\lor} નિયમ છે.
કારણ એ છે કે \Log{G1c} જેનું અનુસરણ કરે છે તે
ગેન્ટ્ઝેનના મૂળ તંત્ર~\Log{LK}નું,
મહત્ત્વના અશાસ્ત્રીય તર્ક એટલે અંતઃપ્રજ્ઞાવાદી તર્ક માટે,
\Log{LJ} નામનું સ્વરૂપ હતું.
અંતઃપ્રજ્ઞાવાદી તર્ક માટે યથાર્થ અને પૂર્ણ તંત્ર મેળવવા
\Log{LJ}ના દરેક સિક્વન્ટના ઉત્તરાંગમાં
વધુમાં વધુ એક સૂત્ર રાખવાની મર્યાદા છે.
એટલે~\Log{G3c}નો \RightR{\lor} નિયમ વાપરી ન શકાય,
કારણ કે તેના આધારના ઉત્તરાંગમાં બંને $!A$ અને~$!B$ છે.
તેથી ગેન્ટ્ઝેને \Log{LJ} માટે \RightR{\lor}ના
બે નિયમો વાપર્યા અને \Log{LK} માટે પણ એ જ વાપર્યા.
અંતઃપ્રજ્ઞાવાદી સ્વરૂપ~\Log{G1i} એ~\Log{G1c}ને અનુરૂપ છે,
અને તેમાં બધા ઉત્તરાંગો વધુમાં વધુ એક સૂત્ર ધરાવે છે.
\sourcecorrection{OLMD-305}{મૂળ G1iને ફક્ત ઉત્તરાંગની મર્યાદા ધરાવતું G1c કહે છે. પરંતુ સ્થિર G1i કોષ્ટકના અનુસરણ-ડાબા નિયમના પહેલા આધારમાં મુખ્ય પૂર્વાંગ સિવાયનો ઉત્તરાંગ નથી; બીજો આધાર પોતાની મર્યાદિત ઉત્તરાંગ-સામગ્રી રાખે છે. એટલે ફક્ત દરેક સિક્વન્ટને મર્યાદિત કરવાથી બધા અંતઃપ્રજ્ઞાવાદી નિયમો આપમેળે મળતા નથી; સંબંધિત નિયમનું સ્વરૂપ પણ બદલાય છે.}
(\Log{G3c}નું પણ અંતઃપ્રજ્ઞાવાદી સ્વરૂપ છે,
પરંતુ તેના કેટલાક નિયમો~\Log{G3c}ના અનુરૂપ નિયમોથી જુદા છે.
દાખલા તરીકે તેમાં પણ \RightR{\lor}ના બે નિયમો છે.)

\Log{G1c}ના બંધારણીય નિયમો યથાર્થ છે તે સરળતાથી દેખાય છે.
દાખલા તરીકે $\Gamma \Sequent \Delta$માં
જમણે સત્ય સૂત્ર અથવા ડાબે અસત્ય સૂત્ર હોય,
તો $\Gamma \Sequent \Delta, !A$માં પણ એ જ સૂત્ર છે;
$!A$ સત્ય છે કે અસત્ય તે અપ્રસ્તુત છે.
એટલે~\RightR{\Weakening}નો આધાર સત્ય હોય, તો નિષ્કર્ષ પણ સત્ય છે.
અહીં ઊલટું સાચું નથી તે નોંધો.
$!A$ સત્ય હોવાથી નિષ્કર્ષ સત્ય હોઈ શકે;
ત્યારે આધાર સત્ય હોવાની ખાતરી નથી.
\sourcecorrection{OLMD-306}{મૂળ આ વાક્યમાં જમણે અસત્ય અને ડાબે સત્ય સૂત્ર કહે છે. તે સિક્વન્ટની સત્યતાની શરતથી ઊલટું છે. જમણે સત્ય અથવા ડાબે અસત્ય કર્યું છે, જેમ વિભાગની શરૂઆતની વ્યાખ્યામાં છે.}

બંધારણીય નિયમો શા માટે જોઈએ?
દાખલા તરીકે $\ \Sequent !A \lor \lnot !A$ને
સાબિત કરવા સંકોચન (\RightR{\Contraction}, \LeftR{\Contraction}) જોઈએ.
$!A \Sequent !A$થી શરૂ કરીએ,
તો~\RightR{\lnot}થી પહેલા $\Sequent !A, \lnot !A$ મળે.
પછી \RightR{\lor}નાં બંને સ્વરૂપ એક પછી એક વાપરી શકીએ:
એક વાર $\lnot !A$માંથી $!A \lor \lnot !A$ મેળવવા
અને એક વાર~$!A$માંથી $!A \lor \lnot !A$ મેળવવા.
$\ \Sequent !A \lor \lnot !A, !A \lor \lnot !A$ મળે છે.
એટલે~\Log{G1c}માં $\Sequent !A \lor \lnot !A$ મેળવવા
સાબિતીમાં એ જ સૂત્રની બે આવૃત્તિનું
``સંકોચન'' કરી શકવું જોઈએ:
\[
\Axiom$!A \fCenter !A$
\RightLabel{\RightR{\lnot}}
\UnaryInf$\fCenter !A, \lnot !A$
\RightLabel{\RightR{\lor}}
\UnaryInf$\fCenter !A, !A \lor \lnot !A$
\RightLabel{\RightR{\lor}}
\UnaryInf$\fCenter !A \lor \lnot !A, !A \lor \lnot !A$
\RightLabel{\RightR{\Contraction}}
\UnaryInf$\fCenter !A \lor \lnot !A$
\DisplayProof
\]
ખરેખર કેટલાક માન્ય સિક્વન્ટો~\Log{G1c}માં
શિથિલીકરણ કે સંકોચન વિના સાબિત કરી ન શકાય.
દાખલા તરીકે~\Log{G1c}માં $!A \Sequent !B \lif !A$
સિક્વન્ટની સાબિતીનો અંત \RightR{\lif}થી જ થાય
(કારણ કે ફક્ત $!B \lif !A$ મુખ્ય સૂત્ર હોઈ શકે).
તે માટે આધાર તરીકે $!B, !A \Sequent !A$ જોઈએ.
તેને સ્વયંસિદ્ધ $!A \Sequent !A$માંથી મેળવવા
\LeftR{\Weakening} જોઈએ:
\[
\Axiom$!A \fCenter !A$
\RightLabel{\LeftR{\Weakening}}
\UnaryInf$!B, !A \fCenter !A$
\RightLabel{\RightR{\lif}}
\UnaryInf$!A \fCenter !B \lif !A$
\DisplayProof
\]

\Log{G3c}માં સંકોચન અને શિથિલીકરણના નિયમો જરાય જરૂરી નથી.
શિથિલીકરણ સ્વયંસિદ્ધોમાં જ સમાયેલું છે.
બે આધાર ધરાવતા નિયમોમાં સંદર્ભ $\Gamma$, $\Delta$ની
બે નકલને એકમાં ભેળવીને મોટેભાગે સંકોચન ટાળી શકાય.
\Log{G1c} અને \Log{G3c} બંનેમાં આવા
``સંદર્ભ વહેંચતા'' નિયમો છે.
એક આધારવાળા નિયમોમાં ફક્ત \RightR{\lor}, \LeftR{\land},
\RightR{\lexists} અને~\LeftR{\lforall} માટે સંકોચન જોઈએ.
\Log{G3c}માં આ નિયમોનાં સ્વરૂપ સંકોચનને સંપૂર્ણપણે બિનજરૂરી બનાવે છે.
એક \Log{G2c} તંત્ર પણ છે (\olref{tab:G2c} જુઓ),
જેમાં બે આધારવાળા નિયમોના સંદર્ભો સ્વતંત્ર હોય છે.
\Log{G2c}માં સંકોચન~\Log{G1c} કરતાં પણ વધુ મહત્ત્વનું છે.

\begingroup
\makeatletter\def\ol@sectioncs{section}\makeatother

% BEGIN OLP-0706
\phantomsection\label{guunit:OLP-0706}


\begin{table}
\begin{tabular}{@{}c@{\hskip-5ex}c@{}}
\multicolumn{2}{@{}c@{}}{સ્વયંસિદ્ધો}\\[2ex]
$!A \Sequent !A$ & $\lfalse \Sequent \quad$
\\[2ex]
\multicolumn{2}{@{}c@{}}{બંધારણીય નિયમો}\\[2ex]
\Axiom$ \Gamma \fCenter \Delta $
\RightLabel{\LeftR{\Weakening}}
\UnaryInf$ !A, \Gamma \fCenter \Delta$
\DisplayProof
&
\Axiom$ \Gamma \fCenter \Delta$
\RightLabel{\RightR{\Weakening}}
\UnaryInf$ \Gamma \fCenter \Delta, !A$
\DisplayProof
\\[4ex]
\Axiom$ !A, !A, \Gamma \fCenter \Delta $
\RightLabel{\LeftR{\Contraction}}
\UnaryInf$ !A, \Gamma \fCenter \Delta$
\DisplayProof
&
\Axiom$ \Gamma \fCenter \Delta, !A, !A$
\RightLabel{\RightR{\Contraction}}
\UnaryInf$ \Gamma \fCenter \Delta, !A$
\DisplayProof
\\[4ex]
\multicolumn{2}{@{}c@{}}{તાર્કિક નિયમો}\\[2ex]
\Axiom$ \Gamma \fCenter \Delta, !A $
\RightLabel{\LeftR{\lnot}}
\UnaryInf$ \lnot !A, \Gamma \fCenter \Delta$
\DisplayProof
&
\Axiom$!A, \Gamma \fCenter \Delta$
\RightLabel{\RightR{\lnot}}
\UnaryInf$ \Gamma \fCenter \Delta, \lnot !A $
\DisplayProof
\\[4ex]\begin{tabular}[t]{@{}l@{}}
\Axiom$ !A, \Gamma \fCenter \Delta$
\RightLabel{\LeftR{\land}}
\UnaryInf$ !A \land !B, \Gamma \fCenter \Delta$
\DisplayProof
\\[3ex]
\Axiom$!B, \Gamma \fCenter \Delta$
\RightLabel{\LeftR{\land}}
\UnaryInf$!A \land !B, \Gamma \fCenter \Delta$
\DisplayProof\\[3ex]
\end{tabular}
&
\Axiom$\Gamma_1 \fCenter \Delta_1, !A$
\Axiom$ \Gamma_2 \fCenter \Delta_2, !B$
\RightLabel{\RightR{\land}}
\BinaryInf$ \Gamma_1, \Gamma_2 \fCenter \Delta_1, \Delta_2, !A \land !B $
\DisplayProof
\\[4ex]\Axiom$!A, \Gamma_1 \fCenter \Delta_1$
\Axiom$!B, \Gamma_2 \fCenter \Delta_2$
\RightLabel{\LeftR{\lor}}
\BinaryInf$!A \lor !B, \Gamma_1, \Gamma_2 \fCenter \Delta_1, \Delta_2$
\DisplayProof
&
\begin{tabular}[t]{@{}r@{}}
\Axiom$\Gamma \fCenter \Delta, !A$
\RightLabel{\RightR{\lor}}
\UnaryInf$ \Gamma \fCenter \Delta, !A \lor !B$
\DisplayProof
\\[3ex]
\Axiom$ \Gamma \fCenter \Delta, !B$
\RightLabel{\RightR{\lor}}
\UnaryInf$ \Gamma \fCenter \Delta, !A \lor !B$
\DisplayProof\\[3ex]
\end{tabular}
\\[4ex]
\Axiom$ \Gamma_1 \fCenter \Delta_1, !A$
\Axiom$ !B, \Gamma_2 \fCenter \Delta_2$
\RightLabel{\LeftR{\lif}}
\BinaryInf$ !A \lif !B, \Gamma_1, \Gamma_2 \fCenter \Delta_1, \Delta_2$
\DisplayProof
&
\Axiom$ !A, \Gamma \fCenter \Delta, !B$
\RightLabel{\RightR{\lif}}
\UnaryInf$ \Gamma \fCenter \Delta, !A \lif !B $
\DisplayProof
\\[4ex]
\Axiom$ !A(t), \Gamma \fCenter \Delta$
\RightLabel{\LeftR{\lforall}}
\UnaryInf$ \lforall[x][!A(x)],\Gamma \fCenter \Delta$
\DisplayProof
&
\Axiom$ \Gamma \fCenter \Delta, !A(a) $
\RightLabel{\RightR{\lforall}}
\UnaryInf$ \Gamma \fCenter \Delta, \lforall[x][!A(x)]$
\DisplayProof
\\[4ex]
\Axiom$ !A(a), \Gamma \fCenter \Delta $
\RightLabel{\LeftR{\lexists}}
\UnaryInf$ \lexists[x][!A(x)], \Gamma \fCenter \Delta$
\DisplayProof
&
\Axiom$ \Gamma \fCenter \Delta, !A(t) $
\RightLabel{\RightR{\lexists}}
\UnaryInf$ \Gamma \fCenter \Delta, \lexists[x][!A(x)]$
\DisplayProof
\end{tabular}
\caption{\Log{G2c}ના નિયમો. $\RightR{\forall}$ અને
$\LeftR{\exists}$માં અચળ~$a$ ફલિત-સિક્વન્ટમાં આવવો નહીં.
$\Gamma$ અને $\Delta$ એ સૂત્રના બહુગણ છે.}
\ollabel{tab:G2c}
\end{table}
% END OLP-0706
\endgroup
% END OLP-0699
\endgroup

\begingroup
\makeatletter\def\ol@sectioncs{section}\makeatother

% BEGIN OLP-0703
\olfileid{pt}{seq}{qua}

\olsection{નિયમિત સાબિતીઓ અને પ્રતિસ્થાપન}
\phantomsection\label{guunit:OLP-0703}


પરિમાણકના નિયમોમાં કેટલીક જટિલતાઓ ઊભી થાય છે,
કારણ કે \LeftR{\lexists} અને~\RightR{\lforall} પર
``આઇગનચલ-શરતો'' લાગુ પડે છે.
આવા અનુમાનોમાં અચળ~$a$નો ક્યારેય પુનઃઉપયોગ ન થાય
તે સુનિશ્ચિત કરીને મોટાભાગની જટિલતાઓ ઉકેલી શકાય.
આ વ્યાખ્યા રજૂ કરીએ.

\begin{defn}
\Log{G1c} કે \Log{G3c}માં સાબિતી~$\pi$ \emph{નિયમિત} છે, જો
\begin{enumerate}
  \item કોઈ આઇગનચલ~$a$ એક કરતાં વધુ
  \LeftR{\lexists} કે~\RightR{\lforall} અનુમાનનું આઇગનચલ ન હોય, અને
  \item $a$ એ \LeftR{\lexists} કે~\RightR{\lforall}નું આઇગનચલ હોય,
  તો તે $\pi$માં ફક્ત એ અનુમાનની ઉપર જ આવે.
\end{enumerate}
\end{defn}

\begin{lem}\ollabel{lem:subst-regular}
$\pi$ નિયમિત હોય અને તેનો અંત $\Gamma \Sequent \Delta$ હોય,
$c$ એ~$\pi$માં આઇગનચલ તરીકે ન વપરાતો અચળ હોય,
અને $t$~એ~$\pi$નું કોઈ આઇગનચલ ન ધરાવતું બંધ પદ હોય,
તો~$\pi$માં $c$ની દરેક આવૃત્તિને બદલે~$t$ મૂકવાથી મળતું
$\Subst{\pi}{t}{c}$ નિયમિત સાબિતી છે.
$\Subst{\pi}{t}{c}$નો અંતિમ સિક્વન્ટ
$\Subst{\Gamma}{t}{c} \Sequent \Subst{\Delta}{t}{c}$ છે,
જ્યાં $\Subst{\Gamma}{t}{c} =
\Setabs{!A(t)}{!A(c) \in \Gamma}$.
સંદર્ભો બહુગણો હોવાથી આ રૂપાંતરમાં દરેક આવૃત્તિ જળવાય છે.
\sourcecorrection{OLMD-329}{મૂળ પ્રતિસ્થાપિત પદ માટે ફક્ત આઇગનચલો ન ધરાવવાની શરત આપે છે. આ અધ્યાયના નિયમ-વિભાગમાં પરિમાણકનો દાખલ બંધ પદ હોવાનો છે; ખુલ્લું પદ મૂકવાથી એ શરત અથવા પરિમાણકની અંદર ચલ-બંધન બદલાઈ શકે. તેથી અહીં બંધ પદની શરત સ્પષ્ટ કરી છે.}
\end{lem}

\begin{proof}
  $\pheight{\pi}$ પર આગમનથી સાબિતી કરીએ.
  $\pheight{\pi} = 0$ હોય, તો ફક્ત સ્વયંસિદ્ધ છે.
  ત્યારે $\Subst{\pi}{t}{c}$ પણ સ્વયંસિદ્ધ છે,
  કારણ કે તેનું કોઈ સૂત્ર~$!A(c)$ એ $!A(t)$ બને છે.

  $\pheight{\pi} > 0$ હોય,
  તો~$\pi$ના છેલ્લા અનુમાન પ્રમાણે કિસ્સા જુદા પાડીએ.
  બંધારણીય નિયમોના (\Log{G1c}માં) અને વિધાનતર્કીય સંયોજકોના
  તાર્કિક નિયમોના કિસ્સા સરળ છે અને સ્વાધ્યાય તરીકે છોડીએ.
  $\pi$નો અંત પરિમાણકના અનુમાનથી થાય તે કિસ્સા વિચારીએ.
  \begin{enumerate}
    \item છેલ્લું અનુમાન $\RightR{\lforall}$ હોય,
    તો $\pi$નું સ્વરૂપ આ છે:
\[
    \AxiomC{}
    \RightLabel{$\pi'$}
    \Deduce$\Gamma(c) \fCenter \Delta'(c), !A(a,c)$
    \RightLabel{\RightR{\lforall}}
    \UnaryInf$\Gamma(c) \fCenter \Delta'(c), \lforall[x][!A(x,c)]$
    \DisplayProof
    \]
    આગમન-ધારણાથી $\Subst{\pi'}{t}{c}$ એ
    $\Gamma(t) \fCenter \Delta'(t), !A(a,t)$ની
    નિયમિત સાબિતી છે.
    $a$ એ~$\pi$નું આઇગનચલ હોવાથી $a$~એ~$t$માં આવતું નથી.
    એટલે $\Subst{\pi}{t}{c}$નું છેલ્લું અનુમાન,
\[
    \AxiomC{}
    \RightLabel{$\Subst{\pi'}{t}{c}$}
    \Deduce$\Gamma(t) \fCenter \Delta'(t), !A(a,t)$
    \RightLabel{\RightR{\lforall}}
    \UnaryInf$\Gamma(t) \fCenter \Delta'(t), \lforall[x][!A(x,t)]$
    \DisplayProof
    \]
    યોગ્ય છે: નિષ્કર્ષમાં~$a$ આવતું નથી.
    \item \Log{G3c}માં છેલ્લું અનુમાન $\LeftR{\lforall}$ હોય,
    તો $\pi$નું સ્વરૂપ આ છે:
\[
    \AxiomC{}
    \RightLabel{$\pi'$}
    \Deduce$!A(s(c),c), \lforall[x][!A(x,c)], \Gamma(c) \fCenter \Delta'(c)$
    \RightLabel{\LeftR{\lforall}}
    \UnaryInf$\lforall[x][!A(x,c)], \Gamma(c) \fCenter \Delta'(c)$
    \DisplayProof
    \]
    આગમન-ધારણાથી $\Subst{\pi'}{t}{c}$ એ
    $!A(s(t),t), \lforall[x][!A(x,t)], \Gamma(t) \fCenter \Delta'(t)$ની
    નિયમિત સાબિતી છે.
    $\Subst{\pi}{t}{c}$નું છેલ્લું અનુમાન,
\[
    \AxiomC{}
    \RightLabel{$\Subst{\pi'}{t}{c}$}
    \Deduce$!A(s(t),t), \lforall[x][!A(x,t)], \Gamma(t) \fCenter \Delta'(t)$
    \RightLabel{\LeftR{\lforall}}
    \UnaryInf$\lforall[x][!A(x,t)], \Gamma(t) \fCenter \Delta'(t)$
    \DisplayProof
    \]
    યોગ્ય છે. \Log{G1c}માં~\LeftR{\lforall}નો કિસ્સો સમાન,
    પરંતુ થોડો સરળ છે.
    \sourcecorrection{OLMD-330}{મૂળના આ સર્વપરિમાણક-ડાબા કિસ્સાનાં બંને વૃક્ષોમાં જમણા નિયમનું લેબલ છે. મુખ્ય સૂત્ર તથા આધારમાં જાળવેલી તેની આવૃત્તિ મુજબ બંને લેબલ ડાબા કર્યા છે.}
    \item છેલ્લું અનુમાન \RightR{\lexists} કે \LeftR{\lexists} હોય: સ્વાધ્યાય.
  \end{enumerate}
  દરેક કિસ્સામાં નિયમિત સાબિતીના અંતે સિક્વન્ટ ઉમેર્યો છે
  (બે આધારવાળા નિયમોમાં બે નિયમિત સાબિતીઓના અંતે).
  નવું સિક્વન્ટ $\pi$ના અંતિમ સિક્વન્ટથી ફક્ત એટલું જુદું છે
  કે $c$ને બદલે પદ~$t$ મૂક્યું છે.
  વળી $\pi$ અને $\Subst{\pi}{t}{c}$ના આઇગનચલો એ જ છે.
  $t$માં~$\pi$નું કોઈ આઇગનચલ નથી એમ ધાર્યું હોવાથી,
  નવા અંતિમ સિક્વન્ટમાં $\Subst{\pi}{t}{c}$નું કોઈ આઇગનચલ નથી.
  એટલે $\Subst{\pi}{t}{c}$ નિયમિત છે.
\end{proof}

\begin{prop}
$\pi$ એ \Log{G1c} કે \Log{G3c}ની સાબિતી હોય,
તો એ જ સંરચના (ખાસ કરીને એ જ ઊંચાઈ) ધરાવતી
નિયમિત સાબિતી~$\pi'$ છે,
જે $\pi$થી ફક્ત આઇગનચલોનાં પ્રતિસ્થાપનથી જુદી પડે છે.
\end{prop}

\begin{proof}
ફક્ત આ સાબિતી પૂરતું, $\pi$ના આઇગનચલ-અનુમાનને \emph{સ્વચ્છ} કહીએ,
જો તેનું આઇગનચલ બીજા અનુમાનનું આઇગનચલ ન હોય
અને~$\pi$માં તે અનુમાનની તરત પહેલાંની ઉપ-સાબિતી
સિવાય ક્યાંય ન આવે; નહીં તો \emph{અસ્વચ્છ} કહીએ.
$\pi$માં અસ્વચ્છ આઇગનચલ-અનુમાનોની સંખ્યા $n$ લો.
$n$ પર આગમનથી બતાવીએ કે $\pi$~ને
જરૂરી નિયમિત સાબિતી~$\pi'$માં ફેરવી શકાય.
$n=0$ હોય, તો~$\pi$નાં બધાં આઇગનચલ-અનુમાન સ્વચ્છ છે;
એટલે $\pi$ નિયમિત છે અને કશું સાબિત કરવાનું નથી.
નહિતર સૌથી ઉપરનું અસ્વચ્છ આઇગનચલ-અનુમાન પસંદ કરો,
દાખલા તરીકે \RightR{\lforall}.
તે અનુમાનથી પૂરી થતી $\pi$ની ઉપ-સાબિતી~$\pi_1$નું સ્વરૂપ આ છે:
\[
    \AxiomC{}
    \RightLabel{$\pi_2$}
    \Deduce$\Gamma \fCenter \Delta, !A(c)$
    \RightLabel{\RightR{\lforall}}
    \UnaryInf$\Gamma \fCenter \Delta, \lforall[x][!A(x)]$
    \DisplayProof
    \]
$c$ આઇગનચલ હોવાથી~$\Gamma$, $\Delta$
કે $\lforall[x][!A(x)]$માં આવતું નથી તે નોંધો.
સાબિતી $\pi_2$ નિયમિત છે:
પસંદ કરેલા અસ્વચ્છ અનુમાનની ઉપરના આઇગનચલ-અનુમાન સ્વચ્છ છે.
પસંદ કરેલું આઇગનચલ તેમાં બીજા અનુમાનનું આઇગનચલ નથી.
\sourcecorrection{OLMD-331}{મૂળ સૌથી ઉપરનું આઇગનચલ-અનુમાન પસંદ કરે છે, પરંતુ તે સ્વચ્છ પણ હોઈ શકે, એટલે અસ્વચ્છ અનુમાનોની સંખ્યા ઘટવાની ખાતરી નથી. સૌથી ઉપરનું અસ્વચ્છ અનુમાન પસંદ કર્યું છે. તેની ઉપર સ્વચ્છ અનુમાન હોઈ શકે; તેથી ઉપસાબિતી નિયમિત છે તે કહ્યું છે, એકેય આઇગનચલ-અનુમાન નથી એવો મૂળનો દાવો કર્યો નથી.}
$\pi$માં ન આવતો અચળ $a$ લો.
\olref{lem:subst-regular} મુજબ $\Subst{\pi_2}{a}{c}$
એ $\Gamma \fCenter \Delta, !A(a)$ની નિયમિત સાબિતી છે;
તેના પર~$\RightR{\lforall}$ લગાવી શકીએ.
$\pi$માં $\pi_1$ને બદલે સાબિતી $\pi_1'$,
\[
    \AxiomC{}
    \RightLabel{$\Subst{\pi_2}{a}{c}$}
    \Deduce$\Gamma \fCenter \Delta, !A(a)$
    \RightLabel{\RightR{\lforall}}
    \UnaryInf$\Gamma \fCenter \Delta, \lforall[x][!A(x)]$
    \DisplayProof
    \]
મૂકીએ, તો અસ્વચ્છ આઇગનચલ-અનુમાનને બદલે સ્વચ્છ અનુમાન મૂક્યું છે;
ફક્ત પસંદ કરેલા આઇગનચલ $c$ને બદલે~$a$ મૂકવાનો તફાવત છે.
મળેલી સાબિતીમાં વધુમાં વધુ $n-1$ અસ્વચ્છ આઇગનચલ-અનુમાનો છે,
અને પરિણામ આગમન-ધારણાથી અનુસરે છે.
\sourcecorrection{OLMD-332}{મૂળ નવી સંખ્યા બરાબર એકથી ઘટે છે કહે છે. બીજા અનુમાનોની એક જ આઇગનચલના પુનઃઉપયોગની ખામી પણ આ બદલાવથી દૂર થઈ શકે; તેથી ઓછામાં ઓછા એક અસ્વચ્છ અનુમાનનો ઘટાડો, એટલે વધુમાં વધુ જૂની સંખ્યાથી એક ઓછાં, એવો પૂરતો દાવો કર્યો છે.}
\end{proof}

હમણાં સાબિત કરેલા પ્રતિજ્ઞાપનને આગળ સ્પષ્ટ રીતે ટાંકશું નહીં,
પરંતુ વિચારેલી બધી સાબિતીઓ નિયમિત છે એમ માનીશું:
તે ન હોય, તો આઇગનચલો બદલીને હંમેશાં નિયમિત બનાવી શકીએ.
% END OLP-0703
\endgroup

\begingroup
\makeatletter\def\ol@sectioncs{section}\makeatother

% BEGIN OLP-0698
\olfileid{pt}{seq}{adm}

\olsection{સ્વીકાર્ય અને નિષ્પન્નક્ષમ નિયમો}
\phantomsection\label{guunit:OLP-0698}


સાબિતી-સિદ્ધાંતનાં ઘણાં પરિણામો આ પ્રશ્નને લગતાં
પરિણામો વિશે છે અથવા તેમનો ઉપયોગ કરે છે:
કોઈ ચોક્કસ નિયમથી જે બધું સાબિત થઈ શકે
તે એ નિયમ વિના પણ સાબિત થઈ શકે?
કટ-લોપનું પ્રમેય આનું સૌથી જાણીતું ઉદાહરણ છે.
તે બતાવે છે કે~\Cut{} સાથેના સિક્વન્ટ તંત્રના નિયમોથી
જે સિક્વન્ટ સાબિત થઈ શકે તે~\Cut વિના પણ
સાબિત થઈ શકે છે. આ પ્રકારનો ગુણધર્મ
મહત્ત્વનો હોવાથી તેને નામ આપ્યું છે.

\begin{defn}
સિક્વન્ટ તંત્રમાં નિયમ~$R$ \emph{સ્વીકાર્ય} છે,
જો તંત્ર સાથે નિયમ~$R$ વાપરીને સાબિત થઈ શકે તે દરેક
સિક્વન્ટ~$R$ વિના પણ સાબિત થઈ શકે.
\end{defn}

\begin{ex}\ollabel{ex:land-deriv}
\Log{G1c} અને \Log{G3c}માં $\LeftR{\land}$ નિયમનાં
સ્વરૂપ જુદાં છે. \Log{G1c}માં નિયમનાં બે સ્વરૂપ છે:
એકમાં મુખ્ય સૂત્ર~$!A \land !B$નું
ડાબું સંયોજ્ય $!A$ આધારમાં છે અને
બીજામાં જમણું સંયોજ્ય~$!B$ છે:
\[
\Axiom$ !A, \Gamma \fCenter \Delta$
\RightLabel{$\LeftR{\land}_1$}
\UnaryInf$ !A \land !B, \Gamma \fCenter \Delta$
\DisplayProof
\qquad
\Axiom$!B, \Gamma \fCenter \Delta$
\RightLabel{$\LeftR{\land}_2$}
\UnaryInf$!A \land !B, \Gamma \fCenter \Delta$
\DisplayProof
\]
\Log{G3c}માં ફક્ત એક નિયમ છે:
\[
\Axiom$ !A, !B, \Gamma \fCenter \Delta$
\RightLabel{$\LeftR{\land}_3$}
\UnaryInf$ !A \land !B, \Gamma \fCenter \Delta$
\DisplayProof
\]
હવે પૂછી શકીએ: ``\Log{G3c}નો $\LeftR{\land}_3$ નિયમ
\Log{G1c}માં સ્વીકાર્ય છે?'' હા છે:
$\LeftR{\land}_3$થી થતાં દરેક અનુમાનનું
ફક્ત~\Log{G1c}ના નિયમોથી સીધું અનુકરણ કરી શકીએ.
$\LeftR{\land}_1$ અને $\LeftR{\land}_2$ ઉપરાંત
બંધારણીય નિયમ~$\LeftR{\Contraction}$ પણ જોઈએ:
\[
\Axiom$!A, !B, \Gamma \fCenter \Delta$
\RightLabel{$\LeftR{\land}_1$}
\UnaryInf$!A \land !B, !B, \Gamma \fCenter \Delta$
\RightLabel{$\LeftR{\land}_2$}
\UnaryInf$!A \land !B, !A \land !B, \Gamma \fCenter \Delta$
\RightLabel{\LeftR{\Contraction}}
\UnaryInf$!A \land !B, \Gamma \fCenter \Delta$
\DisplayProof
\]
\sourcecorrection{OLMD-300}{મૂળના અંતિમ સંકોચન-પગલામાં મુખ્ય સંયોજન અને બાજુના સંદર્ભ વચ્ચેનો અલ્પવિરામ ખૂટે છે; ઉમેર્યો છે.}
\end{ex}

એક નિયમથી થતા અનુમાનનું બીજા નિયમોથી અનુકરણ કરવું
એ નિયમ સ્વીકાર્ય છે તે બતાવવાની એક રીત છે.
પરંતુ એ એકમાત્ર રીત નથી; કેટલાક નિયમો સ્વીકાર્ય હોય
પરંતુ આ રીતે તેમનું અનુકરણ ન થઈ શકે.
જેમનું થઈ શકે તેમને \emph{નિષ્પન્નક્ષમ} કહે છે.

\begin{defn}
નિયમ~$R$,
\[
\AxiomC{$S_1$}
\AxiomC{$\dots$}
\AxiomC{$S_n$}
\RightLabel{$R$}
\TrinaryInfC{$S$}
\DisplayProof
\]
તંત્રમાં \emph{નિષ્પન્નક્ષમ} કહેવાય, જો તંત્રના નિયમો
અને સ્વયંસિદ્ધો સાથે~$S_1$, \dots,~$S_n$નાં સ્વરૂપના
વધારાના આરંભિક સિક્વન્ટો વાપરીને
સિક્વન્ટ~$S$ની ઢાંચારૂપ સાબિતી મળે.
\end{defn}

\olref{ex:land-deriv}માં આપેલી સાબિતી બતાવે છે
કે $\LeftR{\land}_3$ એ~\Log{G1c}માં નિષ્પન્નક્ષમ નિયમ છે:
તે સાબિતી, $\LeftR{\land}_3$ના આધારને
આરંભિક સિક્વન્ટો તરીકે લઈ, ફક્ત \Log{G1c}ના નિયમોથી
$\LeftR{\land}_3$ના નિષ્કર્ષની ઢાંચારૂપ સાબિતી આપે છે.
(તેમાં \Log{G1c}ના કોઈ સ્વયંસિદ્ધ વપરાયા નથી.)

\begin{prop}\ollabel{prop:lor-G3-adm}
\Log{G3c}ના $\LeftR{\land}$ અને $\RightR{\lor}$
નિયમો~\Log{G1c}માં નિષ્પન્નક્ષમ છે.
\end{prop}

\begin{proof}
  $\LeftR{\land}$ માટેની સાબિતી \olref{ex:land-deriv}માં આપી છે.
  $\RightR{\lor}$ માટેની સાબિતી સ્વાધ્યાય તરીકે છોડી છે.
\end{proof}

\begin{prob}
\Log{G3c}નો $\RightR{\lor}$ નિયમ~\Log{G1c}માં
નિષ્પન્નક્ષમ છે તે બતાવો.
\end{prob}

\begin{prob}
$\liff$ વ્યાખ્યાયિત કારક છે એમ ધારો:
$!A \liff !B$ એ $(!A \lif !B) \land (!B \lif !A)$નું
ટૂંકું રૂપ છે. બતાવો કે આ નિયમો
\begin{enumerate}
\item \Axiom$\Gamma, !A, !B \fCenter \Delta$
\Axiom$\Gamma \fCenter \Delta, !A, !B$
\RightLabel{\LeftR{\liff}}
\BinaryInf$!A \liff !B, \Gamma \fCenter \Delta$
\DisplayProof
\item
\Axiom$!A, \Gamma \fCenter \Delta, !B$
\Axiom$!B, \Gamma \fCenter \Delta, !A$
\RightLabel{\RightR{\liff}}
\BinaryInf$\Gamma \fCenter \Delta, !A \liff !B$
\DisplayProof
\end{enumerate}
(a)~\Log{G1c}માં અને (b)~\Log{G3c}માં નિષ્પન્નક્ષમ છે.
\sourcecorrection{OLMD-301}{મૂળમાં ડાબા દ્વિઅનુસરણ-નિયમનું મુખ્ય સૂત્ર પણ જમણી બાજુ છે. તેને નિષ્કર્ષની ડાબી બાજુ મૂક્યું છે: બંને આધાર અનુક્રમે બંને સૂત્રો સત્ય હોય અને બંને અસત્ય હોય તે કિસ્સા આવરે છે.}
\end{prob}


\begin{ex}\ollabel{ex:weak-adm}
\Log{G1c}નો $\RightR{\Weakening}$ નિયમ,
\[
\Axiom$ \Gamma \fCenter \Delta$
\RightLabel{\RightR{\Weakening}}
\UnaryInf$ \Gamma \fCenter \Delta, !A$
\DisplayProof
\]
\Log{G3c}માં સ્વીકાર્ય છે.
સાબિતીનો વિચાર સરળ છે:~\Log{G3c}માં
આધાર $\Gamma \Sequent \Delta$ની સાબિતી~$\pi$ છે એમ ધારો.
$\pi$ના દરેક સિક્વન્ટના ઉત્તરાંગમાં સૂત્ર~$!A$ ઉમેરો.
સ્વયંસિદ્ધો સ્વયંસિદ્ધ જ રહે છે;
યોગ્ય અનુમાનો પણ યોગ્ય અનુમાન જ રહે છે.
નવી સાબિતીનો અંતિમ સિક્વન્ટ $\Gamma \Sequent
\Delta, !A$ છે.
\sourcecorrection{OLMD-302}{મૂળ નવી સાબિતીના આ અંતને સૂત્ર કહે છે, પરંતુ દર્શાવેલું પદ આખું સિક્વન્ટ છે. તે પ્રમાણે અહીં અંતિમ સિક્વન્ટ કહ્યું છે.}

પરંતુ $\RightR{\Weakening}$ નિયમ~\Log{G3c}માં નિષ્પન્નક્ષમ નથી.
તે હોય એમ ધારો: ફક્ત~\Log{G3c}ના સ્વયંસિદ્ધો અને નિયમો,
તથા જરૂર હોય તો વધારાનો આરંભિક સિક્વન્ટ~$\Gamma \Sequent \Delta$,
વાપરીને $\Gamma \Sequent \Delta, !A$ની સાબિતી મળે એમ ધારો.
$\RightR{\Weakening}$ નિષ્પન્નક્ષમ હોવા માટે જરૂરી છે તેમ
આ સાબિતી સંપૂર્ણ ઢાંચારૂપ હોય,
તો $\Gamma$ અને~$\Delta$ કાઢીને તેને~$\Sequent !A$ની
સાબિતીમાં ફેરવી શકીએ.
એથી વધારાના આરંભિક સિક્વન્ટો~$\Gamma \Sequent \Delta$
ખાલી સિક્વન્ટ~$\quad \Sequent \quad$ બની જાય.
ખાલી સિક્વન્ટમાં કોઈ સૂત્રો નથી,
પરંતુ~\Log{G3c}ના દરેક નિયમના આધારમાં ઓછામાં ઓછું
એક સક્રિય સૂત્ર જરૂરી છે.
એટલે ખાલી સિક્વન્ટ આ ઢાંચારૂપ સાબિતીના
કોઈ અનુમાનનો આધાર ન બની શકે.
તેથી સાબિતી ઢાંચારૂપ ત્યારે જ હોઈ શકે
જ્યારે વધારાનો સિક્વન્ટ $\Gamma \Sequent \Delta$
તેમાં ખરેખર આરંભિક સિક્વન્ટ તરીકે વપરાયો ન હોય.
બીજા શબ્દોમાં, તે ફક્ત~\Log{G3c}ના સ્વયંસિદ્ધો અને નિયમોથી
$\Gamma \Sequent \Delta, !A$ની ઢાંચારૂપ નિષ્પત્તિ હોય,
અને આ સ્વરૂપના દરેક સિક્વન્ટની~\Log{G3c}માં
સાબિતી છે તે બતાવ્યું હોય.
પરંતુ આ અશક્ય છે, કારણ કે \Log{G3c} યથાર્થ છે;
એટલે દરેક સંરચનામાં સત્ય હોય તેવા સિક્વન્ટો જ
સાબિત કરી શકે. $\Gamma = \Delta = \emptyset$
અને $!A \ident \lfalse$ લઈએ, તો \Log{G3c}એ
દાખલા તરીકે $\quad \Sequent \lfalse$ સિક્વન્ટ સાબિત
કરવો પડે, જે તે કરી શકતું નથી.
\end{ex}

હવે શિથિલીકરણ~\Log{G3c}માં સ્વીકાર્ય નિયમ છે
તે પરિણામની આગમનથી સાબિતી આપીએ.
ખરેખર સહજ દલીલ બતાવે છે તેમ થોડું વધુ સબળ પરિણામ સાચું છે:
$\pi$ના દરેક સિક્વન્ટની જમણી બાજુ $!A$ ઉમેરવાથી
સાબિતીની સંરચના બદલાતી નથી;
ખાસ કરીને ઊંચાઈ વધતી નથી.
આ મહત્ત્વના ગુણધર્મની અલગ વ્યાખ્યા આપીએ.

\begin{defn}
નિયમ~$R$,
\[
\AxiomC{$S_1$}
\AxiomC{$\dots$}
\AxiomC{$S_n$}
\RightLabel{$R$}
\TrinaryInfC{$S$}
\DisplayProof
\]
તંત્રમાં \emph{ઊંચાઈ-સંરક્ષક સ્વીકાર્ય}
(અથવા \emph{ઊંચાઈ-સંરક્ષક-સ્વીકાર્ય}) છે, જો
$\Proves[h_i] S_i$ એ $i = 1$, \dots,~$n$ માટે હોય,
ત્યારે $\Proves_m S$ હોય, જ્યાં $m = \max(h_1, \dots,
h_n)$.
\end{defn}

નિયમ ઊંચાઈ-સંરક્ષક સ્વીકાર્ય થવા માટે આ પૂરતું છે:
આધારો~$S_i$ની સાબિતીઓ~$\pi_i$ને
ઊંચાઈ~$h_i = \pheight{\pi_i}$ હોય,
ત્યારે નિષ્કર્ષની એવી સાબિતી~$\pi$ મળે
જેની ઊંચાઈ~$\pheight{\pi}$ એ~$h_i$ના મહત્તમ કરતાં મોટી ન હોય.
ખાસ કરીને એક જ આધાર~$S_1$ હોય,
તો $\pheight{\pi} \le \pheight{\pi_1}$ પૂરતું છે.
$\RightR{\Weakening}$ અને $\LeftR{\Weakening}$ના કિસ્સામાં
ખરેખર $\pheight{\pi} = \pheight{\pi_1}$ છે,
પરંતુ સમાનતા જરૂરી નથી.

\begin{prop}\ollabel{prop:weak-G3c-adm}
\Log{G1c}ના $\RightR{\Weakening}$ અને $\LeftR{\Weakening}$
નિયમો~\Log{G3c}માં ઊંચાઈ-સંરક્ષક સ્વીકાર્ય છે.
\end{prop}

\begin{proof}
  બતાવવાનું છે કે $\Log{G3c} \Proves \Gamma \Sequent \Delta$
  માટે ઊંચાઈ~$\pheight{\pi} = n$ની સાબિતી~$\pi$ હોય,
  તો $\pheight{\pi'} = n$ સાથે
  $\Gamma \Sequent \Delta, !A$ની \Log{G3c}-સાબિતી $\pi'$ મળે.
  $n$ પર આગમનથી સાબિતી કરીએ.
  $n=0$ હોય, તો $\pi$ ફક્ત સ્વયંસિદ્ધ $\Gamma \Sequent \Delta$ છે
  (સામાન્ય ઓળખ-સ્વયંસિદ્ધના કિસ્સામાં કોઈ આણ્વિક $!C$
  બંને $\Gamma$ અને~$\Delta$માં આવે છે);
  $\Gamma \Sequent \Delta, !A$ પણ સ્વયંસિદ્ધ છે.
  \sourcecorrection{OLMD-303}{મૂળના આધાર-કિસ્સાના કૌંસમાં ફક્ત આણ્વિક ઓળખ-સ્વયંસિદ્ધનો ઉલ્લેખ છે. G3cના કોષ્ટકમાં અસત્ય ડાબે હોય તે સ્વયંસિદ્ધ પણ છે; તેની જમણી બાજુ સૂત્ર ઉમેરવાથી તે પણ સ્વયંસિદ્ધ રહે છે. આ બીજો આધાર-કિસ્સો અહીં સ્પષ્ટ નોંધ્યો છે.}
  $n = \pheight{\pi} >0$ હોય,
  તો સાબિતી~$\pi$નું છેલ્લું અનુમાન હોય.
  વપરાયેલા નિયમ પ્રમાણે કિસ્સા જુદા પાડીએ.
  ફક્ત એક ઉદાહરણ આપીએ:
  છેલ્લું અનુમાન $\RightR{\land}$ છે એમ ધારો.
  ત્યારે $\Delta = \Delta', !B \land !C$ છે
  અને છેલ્લાં અનુમાનની તરત પહેલાંની ઉપ-સાબિતીઓ
  $\pi_1$ અને~$\pi_2$નાં અંતિમ સિક્વન્ટો અનુક્રમે
  $\Gamma \Sequent \Delta', !B$ અને
  $\Gamma \Sequent \Delta', !C$ છે.
  બીજા શબ્દોમાં $\pi$નું સ્વરૂપ આ છે:
  \[
  \AxiomC{}
  \RightLabel{$\pi_1$}
  \Deduce$\Gamma \fCenter \Delta', !B$
  \AxiomC{}
  \RightLabel{$\pi_2$}
  \Deduce$\Gamma \fCenter \Delta', !C$
  \RightLabel{\RightR{\land}}
  \BinaryInf$\Gamma \fCenter \Delta', !B \land !C$
  \DisplayProof
  \]
  વળી $n = \max(\pheight{\pi_1}, \pheight{\pi_2}) + 1$ છે.
  એટલે $\pheight{\pi_1} < n$ અને $\pheight{\pi_2} < n$ છે,
  તેથી આગમન-ધારણા લાગુ પડે છે:
  $\Gamma \Sequent \Delta', !B, !A$ તથા
  $\Gamma \Sequent \Delta', !C, !A$ની
  સાબિતીઓ $\pi_1'$ અને~$\pi_2'$ મળે છે.
  $\RightR{\land}$ નિયમથી
  $\Gamma \Sequent \Delta', !B \land !C, !A$ની
  સાબિતી~$\pi'$ મળે છે:
  \[
  \AxiomC{}
  \RightLabel{$\pi_1'$}
  \Deduce$\Gamma \fCenter \Delta', !B, !A$
  \AxiomC{}
  \RightLabel{$\pi_2'$}
  \Deduce$\Gamma \fCenter \Delta', !C, !A$
  \RightLabel{\RightR{\land}}
  \BinaryInf$\Gamma \fCenter \Delta', !B \land !C, !A$
  \DisplayProof
  \]
  અહીં $\pheight{\pi'} =
  \max(\pheight{\pi_1'}, \pheight{\pi_2'}) + 1 = \max(\pheight{\pi_1},
  \pheight{\pi_2}) + 1 = \pheight{\pi}$ છે.
\end{proof}

\olref{prop:weak-G3c-adm} ઘણું ઉપયોગી છે.
તેને વારંવાર લગાવવાની નોંધ આપીએ:
$\pi$ એ~$\Gamma \Sequent \Delta$ની સાબિતી હોય,
તો $\pi[\Pi \Sequent \Lambda]$ એ
$\Pi$નાં બધાં સૂત્રો માટે ડાબે શિથિલીકરણની સ્વીકાર્યતા
અને~$\Lambda$નાં બધાં સૂત્રો માટે
જમણે શિથિલીકરણની સ્વીકાર્યતા લગાવવાનું પરિણામ છે.
તે $\Gamma, \Pi \Sequent \Delta, \Lambda$ની સાબિતી છે.
% END OLP-0698
\endgroup

\begingroup
\makeatletter\def\ol@sectioncs{section}\makeatother

% BEGIN OLP-0701
\olfileid{pt}{seq}{inv}

\olsection{નિયમોની પ્રતિલોમક્ષમતા}
\phantomsection\label{guunit:OLP-0701}


\begin{defn}
નિયમ~$R$,
\[
\AxiomC{$S_1$}
\AxiomC{$\dots$}
\AxiomC{$S_k$}
\RightLabel{$R$}
\TrinaryInfC{$S$}
\DisplayProof
\]
તંત્રમાં \emph{પ્રતિલોમક્ષમ} છે,
જો $\Proves S$ હોય ત્યારે
$i = 1$, \dots,~$k$ માટે $\Proves S_i$ હોય.
તે તંત્રમાં \emph{ઊંચાઈ-સંરક્ષક પ્રતિલોમક્ષમ}
(અથવા \emph{ઊંચાઈ-સંરક્ષક-પ્રતિલોમક્ષમ}) છે,
જો $\Proves[n] S$ હોય ત્યારે $i = 1$, \dots,~$k$ માટે
$\Proves[n] S_i$ હોય.
\sourcecorrection{OLMD-308}{મૂળ એક જ $n$ને આધારોની સંખ્યા અને ઊંચાઈની સીમા બંને માટે વાપરે છે; આધારોની સંખ્યા $k$ કરી છે. સામાન્ય પ્રતિલોમક્ષમતામાં અનિર્ધારિત ઊંચાઈ-સૂચક $h_i$ દૂર કર્યો છે અને ઊંચાઈની નોંધ બાકીના અધ્યાય પ્રમાણે ચોરસ કૌંસમાં કરી છે. અર્થ એ છે કે નિષ્કર્ષની સાબિતીથી દરેક આધારની સાબિતી મળે, અને ઊંચાઈ-સંરક્ષક કિસ્સામાં તેની ઊંચાઈ ન વધે.}
\end{defn}

\begin{lem}\ollabel{lem:G3c-invert}
\Log{G3c}ના $\lnot$, $\land$, $\lor$ અને $\lif$ માટેના નિયમો
ઊંચાઈ-સંરક્ષક પ્રતિલોમક્ષમ છે, એટલે:
\begin{enumerate}
  \item \RightR{\lnot}: જો $\Log{G3c} \Proves[n] \Gamma \Sequent
  \Delta, \lnot !A$, તો $\Log{G3c} \Proves[n] !A, \Gamma \Sequent
  \Delta$.
  \item \LeftR{\lnot}: જો $\Log{G3c} \Proves[n] \Gamma, \lnot !A
  \Sequent \Delta$, તો $\Log{G3c} \Proves[n] \Gamma \Sequent \Delta,
  !A$.
  \item \RightR{\land}: જો $\Log{G3c} \Proves[n] \Gamma \Sequent
  \Delta, !A \land !B$, તો $\Log{G3c} \Proves[n] \Gamma \Sequent
  \Delta, !A$ અને $\Log{G3c} \Proves[n] \Gamma \Sequent
  \Delta, !B$.
  \item \LeftR{\land}: જો $\Log{G3c} \Proves[n] \Gamma, !A \land !B
  \Sequent \Delta$, તો $\Log{G3c} \Proves[n] !A, !B, \Gamma \Sequent
  \Delta$.
  \item \RightR{\lor}: જો $\Log{G3c} \Proves[n] \Gamma \Sequent
  \Delta, !A \lor !B$, તો $\Log{G3c} \Proves[n] \Gamma \Sequent
  \Delta, !A, !B$.
  \item \LeftR{\lor}: જો $\Log{G3c} \Proves[n] !A \lor !B, \Gamma \Sequent
  \Delta$, તો $\Log{G3c} \Proves[n] !A, \Gamma \Sequent
  \Delta$ અને $\Log{G3c} \Proves[n] !B, \Gamma \Sequent
  \Delta$.
  \item \RightR{\lif}: જો $\Log{G3c} \Proves[n] \Gamma \Sequent
  \Delta, !A \lif !B$, તો $\Log{G3c} \Proves[n] !A, \Gamma \Sequent
  \Delta, !B$.
  \item \LeftR{\lif}: જો $\Log{G3c} \Proves[n] !A \lif !B, \Gamma
  \Sequent \Delta$, તો $\Log{G3c} \Proves[n] \Gamma \Sequent \Delta,
  !A$ અને $\Log{G3c} \Proves[n] !B, \Gamma \Sequent \Delta$.
\end{enumerate}
\end{lem}

\begin{proof}
બતાવવાનું છે કે~\Log{G3c}ના કોઈ નિયમ~$R$ માટે,
$R$ના નિષ્કર્ષ~$S$ની સાબિતી~$\pi$ હોય,
તો~$R$ના આધારો $S_i$ની સાબિતીઓ~$\pi_i$ પણ હોય,
અને $\pheight{\pi_i} \le \pheight{\pi}$ હોય.
\LeftR{\land} નિયમ વિચારીએ:
$!A \land !B, \Gamma \Sequent \Delta$ સ્વરૂપના
અંતિમ સિક્વન્ટની સાબિતી~$\pi$ છે એમ ધારો.
બતાવવાનું છે કે $!A, !B, \Gamma \Sequent \Delta$ની
સાબિતી $\pi'$ મળે અને
$\pheight{\pi'} \le \pheight{\pi}$ હોય.
આ $n = \pheight{\pi}$ પર આગમનથી કરીએ.

$n = 0$ હોય, તો $\pi$માં ફક્ત સ્વયંસિદ્ધ
$!A \land !B, \Gamma \Sequent \Delta$ છે.
\Log{G3c}ના ઓળખ-સ્વયંસિદ્ધોનું સ્વરૂપ
$!C, \Gamma \Sequent \Delta, !C$ છે, જ્યાં $!C$ આણ્વિક છે.
એટલે $!C$ એ $!A \land !B$ ન હોઈ શકે.
$!A \land !B, \Gamma \Sequent \Delta$ ઓળખ-સ્વયંસિદ્ધ હોય,
તો કોઈ આણ્વિક $!C$ બંને $\Gamma$ અને~$\Delta$નું ઘટક છે.
ત્યારે $!A, !B, \Gamma \Sequent \Delta$ પણ સ્વયંસિદ્ધ છે,
અને આપણી સાબિતી~$\pi'$ મળી ગઈ
(તેની પણ ઊંચાઈ~$0$ છે).
\sourcecorrection{OLMD-309}{મૂળ આ આધાર-કિસ્સામાં G3cના સ્વયંસિદ્ધોને ફક્ત આણ્વિક ઓળખ-સ્વયંસિદ્ધ ગણે છે. કોષ્ટકમાં અસત્ય ડાબે હોય તે સ્વયંસિદ્ધ પણ છે. અહીં મુખ્ય સંયોજનને બે સંયોજ્યોથી બદલતાં સંદર્ભનું અસત્ય સૂત્ર રહે છે; તેથી એ કિસ્સામાં પણ પરિણામ શૂન્ય ઊંચાઈનું સ્વયંસિદ્ધ છે.}

$n>0$ હોય, તો છેલ્લે અનુમાન આવે છે.
$n$ માટે દાવો સાબિત કરવો છે, પરંતુ
ઊંચાઈ~$<n$ની કોઈ પણ સાબિતી માટે તે સાચો છે એમ માની શકીએ,
સંદર્ભ જુદો હોય તો પણ.
એટલે કોઈ પણ $k<n$ અને $\Pi$, $\Lambda$ માટે,
$!A \land !B, \Pi \Sequent \Lambda$ની
ઊંચાઈ~$k$ની સાબિતી હોય,
તો $!A, !B, \Pi \Sequent \Lambda$ની
ઊંચાઈ~$\le k$ની સાબિતી મળે.

બે કિસ્સા છે: ડાબેનું $!A \land !B$ મુખ્ય સૂત્ર છે અથવા નથી.
મુખ્ય હોય, તો છેલ્લું અનુમાન \LeftR{\land} છે
અને તેની તરત પહેલાંની ઉપ-સાબિતી~$\pi_1$નો અંત
$!A, !B, \Gamma \Sequent \Delta$ છે.
$\pheight{\pi_1} < \pheight{\pi}$ હોવાથી પરિણામ આ કિસ્સામાં સાચું છે.

$!A \land !B$ મુખ્ય ન હોય,
તો બીજું કોઈ સૂત્ર મુખ્ય છે અને $!A \land !B$
છેલ્લાં અનુમાનના સંદર્ભમાં છે.
ત્યારે છેલ્લા નિયમ~$R$ની તરત પહેલાંની ઉપ-સાબિતીઓ પર
આગમન-ધારણા લાગુ પડે છે.
તેમાંથી એક સાબિતી~$\pi_1'$ અથવા બે
સાબિતીઓ~$\pi_1'$ અને~$\pi_2'$ મળે છે,
જેમની ઊંચાઈ~$\pi$ની તરત પહેલાંની ઉપ-સાબિતીઓની
ઊંચાઈ કરતાં મોટી નથી.
$\pi_1'$ અને $\pi_2'$ના અંતિમ સિક્વન્ટોના ડાબા સંદર્ભમાં
$!A \land !B$ને બદલે $!A, !B$ હોય છે.
$R$ લગાવતાં સાબિતી~$\pi'$ મળે.
દાખલા તરીકે $\pi$નો અંત~\LeftR{\lif}થી થાય એમ ધારો:
\[
\AxiomC{}
\RightLabel{$\pi_1$}
\Deduce$!A \land !B, \Gamma' \fCenter \Delta, !C$
\AxiomC{}
\RightLabel{$\pi_2$}
\Deduce$!A \land !B, \Gamma', !D \fCenter \Delta$
\RightLabel{\LeftR{\lif}}
\BinaryInf$!A \land !B, \Gamma', !C \lif !D \fCenter \Delta$
\DisplayProof
\]
અહીં ડાબેનું $!C \lif !D$ મુખ્ય છે
અને ડાબો સંદર્ભ $!A \land !B, \Gamma'$ છે
(એટલે $\Gamma = \Gamma', !C \lif !D$).
આગમન-ધારણાથી અનુક્રમે
$!A, !B, \Gamma' \Sequent \Delta, !C$
અને $!A, !B, \Gamma', !D \Sequent \Delta$ની
સાબિતીઓ $\pi_1'$ અને $\pi_2'$ મળે છે.
આ બે સિક્વન્ટોને ફરી \LeftR{\lif}ના આધાર તરીકે
લઈને~$\pi'$ મળે:
\[
\AxiomC{}
\RightLabel{$\pi_1'$}
\Deduce$!A, !B, \Gamma' \fCenter \Delta, !C$
\AxiomC{}
\RightLabel{$\pi_2'$}
\Deduce$!A, !B, \Gamma', !D \fCenter \Delta$
\RightLabel{\LeftR{\lif}}
\BinaryInf$!A, !B, \Gamma', !C \lif !D \fCenter \Delta$
\DisplayProof
\]
\sourcecorrection{OLMD-310}{મૂળ આ આગમન-પગલાના ગદ્યમાં મળેલા બંને આધારોમાં હજી $!A \land !B$ લખે છે, જ્યારે તરત પછીનું વૃક્ષ યોગ્ય રીતે $!A, !B$ ધરાવે છે. ગદ્યના બંને આધાર એ જ વિસ્તૃત સંયોજન પ્રમાણે કર્યા છે.}
આપણી પાસે
\begin{align*}
\pheight{\pi} & = \max(\pheight{\pi_1}, \pheight{\pi_2}) + 1\\
\pheight{\pi'} & = \max(\pheight{\pi_1'}, \pheight{\pi_2'}) + 1
\intertext{વ્યાખ્યા મુજબ, અને}
\pheight{\pi_1'} & \le \pheight{\pi_1}\\
\pheight{\pi_2'} & \le \pheight{\pi_2}
\intertext{આગમન-ધારણા મુજબ. એટલે,}
\pheight{\pi'} & \le \pheight{\pi}.
\end{align*}
\end{proof}

\begin{prob}
\RightR{\lif} માટે \olref{lem:G3c-invert}ની સાબિતી આપો.
\end{prob}

\begin{lem}\ollabel{lem:invert-quant}
\RightR{\lforall} અને \LeftR{\lexists} નિયમો
આ અર્થમાં ઊંચાઈ-સંરક્ષક પ્રતિલોમક્ષમ છે:
\begin{enumerate}
  \item $\Log{G3c} \Proves[n] \Gamma \Sequent \Delta,
  \lforall[x][!A(x)]$ હોય, તો
  $\Log{G3c} \Proves[n] \Gamma \Sequent \Delta, !A(c)$ છે, અને
  \item $\Log{G3c} \Proves[n] \lexists[x][!A(x)], \Gamma \Sequent
  \Delta$ હોય, તો
  $\Log{G3c} \Proves[n] !A(c), \Gamma \Sequent \Delta$ છે.
\end{enumerate}
આ કોઈ પણ એવા $c$ માટે છે
જે $\Gamma$, $\Delta$ અને અનુક્રમે
$\lforall[x][!A(x)]$ અથવા $\lexists[x][!A(x)]$માં ન આવે.
\end{lem}

\begin{proof}
  (1) બતાવીએ અને (2) સ્વાધ્યાય તરીકે છોડીએ.
  દલીલ \olref{lem:G3c-invert}ની સાબિતી જેવી છે.
  આગમનના પગલામાં ફરી બે કિસ્સા છે.
  $\lforall[x][!A(x)]$ મુખ્ય હોય તે કિસ્સો સરળ છે:
  છેલ્લા $\RightR{\lforall}$ અનુમાનનો આધાર
  જરૂરી સ્વરૂપ $\Gamma \Sequent \Delta, !A(a)$નો છે,
  અને તેનાથી પૂરી થતી તરત પહેલાંની ઉપ-સાબિતી~$\pi_1$ માટે
  $\pheight{\pi_1} < \pheight{\pi}$ છે.
  વળી આઇગનચલ-શરતથી $a$ એ $\Gamma$, $\Delta$
  કે $\lforall[x][!A(x)]$માં આવતું નથી.
  $\pi_1$ નિયમિત છે એમ માની શકીએ,
  તેથી \olref[qua]{lem:subst-regular} મુજબ
  $\Subst{\pi_1}{c}{a}$ એ
  $\Gamma \Sequent \Delta, !A(c)$ની
  એ જ ઊંચાઈની સાબિતી છે.

બીજા કિસ્સામાં $\lforall[x][!A(x)]$
છેલ્લાં અનુમાનમાં મુખ્ય નથી; બીજું સૂત્ર મુખ્ય છે.
ફરી એક કિસ્સાનું ઉદાહરણ લઈએ:
છેલ્લું અનુમાન~\RightR{\land} છે એમ ધારો.
$\Delta$નું કોઈ સૂત્ર~$!B \land !C$ સ્વરૂપનું અને મુખ્ય છે.
સાબિતી~$\pi$નો અંત આ છે:
\[
\AxiomC{}
\RightLabel{$\pi_1$}
\Deduce$\Gamma \fCenter \Delta', \lforall[x][!A(x)], !B$
\AxiomC{}
\RightLabel{$\pi_2$}
\Deduce$\Gamma \fCenter \Delta', \lforall[x][!A(x)], !C$
\RightLabel{\RightR{\land}}
\BinaryInf$\Gamma \fCenter \Delta', \lforall[x][!A(x)], !B \land !C$
\DisplayProof
\]
અહીં $\Delta = \Delta', !B \land !C$ છે.
ઉપપાદની બધી તાજગી-શરતો સંતોષતો
અચળ $c$ લઈએ, જે $\Delta$માં ન આવે.
ત્યારે તે $\Delta', !B$ કે $\Delta', !C$માં પણ આવતો નથી.
\sourcecorrection{OLMD-311}{મૂળ ઉદાહરણમાં ફક્ત ઉત્તરાંગમાં ન આવતો અચળ પસંદ કરવાનું કહે છે. ઉપરના ઉપપાદની શરત મુજબ તે પૂર્વાંગ અને સંબંધિત પરિમાણિત સૂત્રમાં પણ ન આવવો જોઈએ; તે બધી શરતો જાળવી છે.}
એટલે આગમન-ધારણા $\pi_1$ અને~$\pi_2$ પર લાગુ પડે:
સાબિતીઓ~$\pi_1'$ અને $\pi_2'$ મળે છે,
જ્યાં $\pheight{\pi_1'} \le \pheight{\pi_1}$
અને $\pheight{\pi_2'} \le \pheight{\pi_2}$ છે.
$\Gamma \Sequent \Delta, !A(c)$ની
જરૂરી સાબિતી~$\pi'$ આ છે:
\[
\AxiomC{}
\RightLabel{$\pi_1'$}
\Deduce$\Gamma \fCenter \Delta', !A(c), !B$
\AxiomC{}
\RightLabel{$\pi_2'$}
\Deduce$\Gamma \fCenter \Delta', !A(c), !C$
\RightLabel{\RightR{\land}}
\BinaryInf$\Gamma \fCenter \Delta', !A(c), !B \land !C$
\DisplayProof
\]
અહીં $\pheight{\pi'} = \max(\pheight{\pi_1'}, \pheight{\pi_2'})+1 \le \pheight{\pi}$ છે.
\end{proof}

\Olref{lem:G3c-invert} કટ-લોપના પ્રમેયની સાબિતીમાં
મહત્ત્વનું છે. તેનો ઉપયોગ આ પરિણામ બતાવવા પણ કરી શકાય:

\begin{prop}\ollabel{prop:G3c-cont-adm}
\Log{G1c}ના સંકોચન-નિયમો \LeftR{\Contraction} અને \RightR{\Contraction}
\Log{G3c}માં ઊંચાઈ-સંરક્ષક સ્વીકાર્ય છે.
\end{prop}

\begin{proof}
\RightR{\Contraction} અને \LeftR{\Contraction} માટે એકસાથે દલીલ આપીએ.
\Log{G3c}માં $\Gamma \Sequent \Delta, !A, !A$
અથવા $!A, !A, \Gamma \Sequent \Delta$ની સાબિતી $\pi$ છે એમ ધારો.
બતાવવાનું છે કે અનુક્રમે $\Gamma \Sequent \Delta, !A$
અથવા $!A, \Gamma \Sequent \Delta$ની
\Log{G3c}-સાબિતી~$\pi'$ મળે,
અને $\pheight{\pi'} \le \pheight{\pi}$ હોય.
આ $n = \pheight{\pi}$ પર આગમનથી કરીએ.

$n=0$ હોય, તો $\pi$ ફક્ત સ્વયંસિદ્ધ છે.
$\Gamma \Sequent \Delta, !A, !A$ એ~\Log{G3c}નું
ઓળખ-સ્વયંસિદ્ધ હોય, તો $\Gamma \Sequent \Delta, !A$ પણ એવું છે:
કાં તો કોઈ આણ્વિક~$!C$ બંને $\Gamma$ અને~$\Delta$નું ઘટક છે,
અથવા $!A$ આણ્વિક છે અને~$\Gamma$નું ઘટક છે.
અંતિમ સિક્વન્ટ $!A, !A, \Gamma \Sequent \Delta$ હોય,
તો પણ એ જ દલીલ લાગુ પડે.
\sourcecorrection{OLMD-312}{મૂળ આ આધાર-કિસ્સામાં ફક્ત ઓળખ-સ્વયંસિદ્ધની શક્યતાઓ આપે છે. અસત્ય ડાબે હોય તે સ્વયંસિદ્ધ માટે પણ સંકોચન પછી અસત્યની ઓછામાં ઓછી એક આવૃત્તિ રહે છે; તેથી આ બીજા સ્વયંસિદ્ધનો કિસ્સો પણ શૂન્ય ઊંચાઈએ સાચો છે.}

$n>0$ હોય, તો $\pi$નો અંત કોઈ નિયમ~$R$થી થાય.
આ છેલ્લા અનુમાનમાં $!A$ મુખ્ય ન હોય,
તો \olref{lem:G3c-invert}ની સાબિતીની જેમ
$\pi$ની તરત પહેલાંની ઉપ-સાબિતીઓ પર આગમન-ધારણા લગાવીએ
અને મળેલી સાબિતીઓ પર એ જ નિયમ~$R$ લગાવીને
સાબિતી~$\pi'$ મેળવી શકીએ.
આગમન-ધારણા આ છે:
$\Pi \Sequent \Lambda, !D, !D$
અથવા $!D, !D, \Pi \Sequent \Lambda$ની
ઊંચાઈ~$k<n$ની સાબિતી હોય,
તો અનુક્રમે $\Pi \Sequent \Lambda, !D$
અથવા $!D, \Pi \Sequent \Lambda$ની
ઊંચાઈ~$\le k$ની સાબિતી મળે.
(સંદર્ભો કે સંકોચાતું સૂત્ર એ $\Gamma$, $\Delta$
કે~$!A$થી જુદાં હોય, તો પણ આગમન-ધારણા લાગુ પડે છે!)

દાખલા તરીકે $R$ એ $\RightR{\land}$ છે
અને $\pi$~નો અંત આ છે એમ ધારો:
\[
\AxiomC{}
\RightLabel{$\pi_1$}
\Deduce$\Gamma \fCenter \Delta', !B, !A, !A$
\AxiomC{}
\RightLabel{$\pi_2$}
\Deduce$\Gamma \fCenter \Delta', !C, !A, !A$
\RightLabel{\RightR{\land}}
\BinaryInf$\Gamma \fCenter \Delta', !B \land !C, !A, !A$
\DisplayProof
\]
અહીં $\Delta = \Delta', !B \land !C$ છે.
આગમન-ધારણાથી અનુક્રમે
$\Gamma \fCenter \Delta', !B, !A$
અને $\Gamma \fCenter \Delta', !C, !A$ની
સાબિતીઓ $\pi_1'$ અને $\pi_2'$ મળે છે,
જ્યાં $\pheight{\pi_1'} \le \pheight{\pi_1}$
અને $\pheight{\pi_2'} \le \pheight{\pi_2}$ છે.
સાબિતી~$\pi'$ આ છે:
\[
\AxiomC{}
\RightLabel{$\pi_1'$}
\Deduce$\Gamma \fCenter \Delta', !B, !A$
\AxiomC{}
\RightLabel{$\pi_2'$}
\Deduce$\Gamma \fCenter \Delta', !C, !A$
\RightLabel{\RightR{\land}}
\BinaryInf$\Gamma \fCenter \Delta', !B \land !C, !A$
\DisplayProof
\]

છેલ્લાં અનુમાનમાં $!A$ મુખ્ય હોય,
તો $\pi$ની તરત પહેલાંની ઉપ-સાબિતીઓ પર
પ્રતિલોમન-ઉપપાદ (\olref{lem:G3c-invert}) લગાવીએ,
પછી આગમન-ધારણા અને પછી ફરી નિયમ~$R$ લગાવીએ.
દાખલા તરીકે $\pi$ એ $\Gamma \Sequent \Delta, !A, !A$
સાબિત કરે છે, $!A \ident (!B \land !C)$ છે
અને છેલ્લા અનુમાનમાં મુખ્ય છે એમ ધારો.
ત્યારે $\pi$નું સ્વરૂપ આ છે:
\[
\AxiomC{}
\RightLabel{$\pi_1$}
\Deduce$\Gamma \fCenter \Delta, !B \land !C, !B$
\AxiomC{}
\RightLabel{$\pi_2$}
\Deduce$\Gamma \fCenter \Delta, !B \land !C, !C$
\RightLabel{\RightR{\land}}
\BinaryInf$\Gamma \fCenter \Delta, !B \land !C, !B \land !C$
\DisplayProof
\]
$\pi_1$ પર \olref{lem:G3c-invert} લગાવીને
$\Gamma \Sequent \Delta, !B, !B$ની સાબિતી~$\pi_1'$ મળે.
(તે $\Gamma \Sequent \Delta, !C, !B$ની
સાબિતી પણ આપે છે, પરંતુ તેની જરૂર નથી.)
એ જ રીતે~$\pi_2$ પર \olref{lem:G3c-invert} લગાવીને
$\Gamma \Sequent \Delta, !C, !C$ની સાબિતી~$\pi_2'$ મળે.
\RightR{\land}નું પ્રતિલોમન ઊંચાઈ-સંરક્ષક હોવાથી
$\pheight{\pi_1'} \le \pheight{\pi_1} < \pheight{\pi}$
અને $\pheight{\pi_2'} \le \pheight{\pi_2} < \pheight{\pi}$ છે.
એટલે આગમન-ધારણા $\pi_1'$ અને $\pi_2'$ પર લાગુ પડે:
અનુક્રમે $\Gamma \Sequent \Delta, !B$
અને $\Gamma \Sequent \Delta, !C$ની
સાબિતીઓ~$\pi_1''$ અને $\pi_2''$ મળે,
જ્યાં $\pheight{\pi_1''} \le \pheight{\pi_1'} \le \pheight{\pi_1}$
અને $\pheight{\pi_2''} \le \pheight{\pi_2'} \le \pheight{\pi_2}$ છે.
પછી સાબિતી~$\pi'$ આ છે:
\[
\AxiomC{}
\RightLabel{$\pi_1''$}
\Deduce$\Gamma \fCenter \Delta, !B$
\AxiomC{}
\RightLabel{$\pi_2''$}
\Deduce$\Gamma \fCenter \Delta, !C$
\RightLabel{\RightR{\land}}
\BinaryInf$\Gamma \fCenter \Delta, !B \land !C$
\DisplayProof
\]
અને $\pheight{\pi'} \le \pheight{\pi}$ છે.

પરિમાણક-નિયમોમાં સૂત્ર મુખ્ય હોય
તે કિસ્સામાં ખાસ સાવચેતી લેવી જોઈએ.

$!A \ident \lforall[x][!B(x)]$ જમણે મુખ્ય છે એમ ધારો.
ત્યારે $\pi$નો અંત આ છે:
\[
\AxiomC{}
\RightLabel{$\pi_1$}
\Deduce$\Gamma \fCenter \Delta, \lforall[x][!B(x)], !B(a)$
\RightLabel{\RightR{\lforall}}
\UnaryInf$\Gamma \fCenter \Delta, \lforall[x][!B(x)], \lforall[x][!B(x)]$
\DisplayProof
\]
\sourcecorrection{OLMD-313}{મૂળ આ સર્વપરિમાણકના વૃક્ષમાં અસ્તિત્વપરિમાણકના જમણા નિયમનું લેબલ આપે છે. આધાર અને નિષ્કર્ષ બંને મુજબ તેને સર્વપરિમાણકનો જમણો નિયમ કર્યો છે.}
\olref{lem:invert-quant} મુજબ કોઈ પણ એવા~$c$ માટે
જે $\Gamma \fCenter \Delta, \lforall[x][!B(x)], !B(a)$માં ન આવે,
$\Gamma \fCenter \Delta, !B(c), !B(a)$ની
ઊંચાઈ $\le \pheight{\pi_1}$ની સાબિતી~$\pi_1'$ છે.
$\pi_1'$ નિયમિત છે એમ માની શકીએ;
એટલે \olref[qua]{lem:subst-regular} મુજબ
સાબિતી $\Subst{\pi_1'}{a}{c}$ એ
$\Gamma \fCenter \Delta, !B(a), !B(a)$ની
ઊંચાઈ $\le \pheight{\pi_1}$ની સાબિતી છે.
હવે આગમન-ધારણા લાગુ પડે છે અને
$\Gamma \fCenter \Delta, !B(a)$ની સાબિતી~$\pi_1''$ આપે છે.
\RightR{\lforall} લગાવીને~$\pi'$ મેળવો:
\[
\AxiomC{}
\RightLabel{$\pi_1''$}
\Deduce$\Gamma \fCenter \Delta, !B(a)$
\RightLabel{\RightR{\lforall}}
\UnaryInf$\Gamma \fCenter \Delta, \lforall[x][!B(x)]$
\DisplayProof
\]
અહીં $\pheight{\pi'} \le \pheight{\pi}$ છે.

હવે $!A \ident \lexists[x][!B(x)]$ છે
અને જમણે મુખ્ય છે એમ ધારો. ત્યારે $\pi$નું સ્વરૂપ આ છે:
\[
\AxiomC{}
\RightLabel{$\pi_1$}
\Deduce$\Gamma \fCenter \Delta, \lexists[x][!B(x)], \lexists[x][!B(x)], !B(t)$
\RightLabel{\RightR{\lexists}}
\UnaryInf$\Gamma \fCenter \Delta, \lexists[x][!B(x)], \lexists[x][!B(x)]$
\DisplayProof
\]
આગમન-ધારણા~$\pi_1$ પર લાગુ પડે છે
અને $\Gamma \Sequent \Delta, \lexists[x][!B(x)], !B(t)$ની
સાબિતી~$\pi_1'$ મળે છે.
(અહીં પ્રતિલોમન-ઉપપાદની જરૂર નથી તે નોંધો!)
પછી સાબિતી~$\pi'$ આ છે:
\[
\AxiomC{}
\RightLabel{$\pi_1'$}
\Deduce$\Gamma \fCenter \Delta,  \lexists[x][!B(x)], !B(t)$
\RightLabel{\RightR{\lexists}}
\UnaryInf$\Gamma \fCenter \Delta, \lexists[x][!B(x)]$
\DisplayProof
\]
અહીં $\pheight{\pi'} \le \pheight{\pi}$ છે.
\sourcecorrection{OLMD-314}{મૂળ આગમનથી મળેલા આધારના ગદ્ય અને વૃક્ષમાં પરિમાણકનું દાખલ-સૂત્ર ફક્ત $!B$ લખે છે. બંને સ્થાનોમાં મૂળ પહેલાના આધારનું એ જ દાખલ $!B(t)$ રાખ્યું છે. એક મળેલી સાબિતી માટે મૂળનું ભૂલથી આવેલું બહુવચન પણ એકવચન કર્યું છે.}
\end{proof}

\begin{prob}
\LeftR{\Contraction} માટે \olref{prop:G3c-cont-adm}ની
સાબિતીની રૂપરેખા આપો.
$!A \ident (!B \land !C)$ અને $!A \ident (!B \lif !C)$
હોય તે કિસ્સા વર્ણવો.
\end{prob}

\begin{prob}
\LeftR{\Contraction} માટે \olref{prop:G3c-cont-adm}ની
સાબિતીના બાકી પરિમાણક-કિસ્સાની રૂપરેખા આપો:
$!A \ident \lforall[x][!B(x)]$
અને $!A \ident \lexists[x][!B(x)]$ હોય તે કિસ્સા.
\end{prob}
% END OLP-0701
\endgroup

\begingroup
\makeatletter\def\ol@sectioncs{section}\makeatother

% BEGIN OLP-0713
\olfileid{pt}{seq}{tr}

\olsection{\Log{G1c} અને \Log{G3c} વચ્ચેનું રૂપાંતરણ}
\phantomsection\label{guunit:OLP-0713}


અગાઉ કહ્યું હતું કે \Log{G1c} અને \Log{G3c} બંને
શાસ્ત્રીય તર્ક માટે યથાર્થ અને પૂર્ણ છે: દરેક સંરચનામાં
સાચા હોય તે અને ફક્ત તે જ સિક્વન્ટો બંને તંત્રમાં સાબિત થાય છે.
ખાસ કરીને બંને એક જ સિક્વન્ટો સાબિત કરે છે. હવે આ હકીકતને
યથાર્થતા અને પૂર્ણતાનાં પ્રમેયોમાંથી પસાર થયા વિના,
ફક્ત સાબિતી-સિદ્ધાંત વડે સાબિત કરી શકીએ.
આવી સાબિતી \Log{G1c}ની સાબિતીને \Log{G3c}ની
સાબિતીમાં અને ઊલટી દિશામાં \emph{રૂપાંતરણો} આપે છે.

\begin{prop}
જો $\Log{G1c} \Proves S$ હોય, તો $\Log{G3c} \Proves S$.
\end{prop}

\begin{proof}
  બતાવીએ કે જો $\pi$ એ \Log{G1c}માં સાબિતી હોય, તો
  \Log{G3c}માં $S$ની સાબિતી~$\pi'$ છે.
  આ માટે $n = \pheight{\pi}$ પર આગમન વાપરીએ.

  જો $n = 0$ હોય, તો $\pi$ સ્વયંસિદ્ધ છે.
  સ્વયંસિદ્ધ $\lfalse \Sequent
  \quad$ એ \Log{G3c}નું પણ સ્વયંસિદ્ધ છે:
  સંદર્ભ $\Gamma,
  \Delta$ ખાલી લો. \Log{G1c}માં કોઈ પણ $!A$ માટે,
  ફક્ત પરમાણ્વીય માટે નહીં, $!A \Sequent !A$ સ્વયંસિદ્ધ માન્ય છે.
  પરંતુ \olref[ex]{prop:G3c-axioms}માં સાબિત કર્યું છે કે
  આવા બધા સિક્વન્ટો \Log{G3c}માં સાબિતીક્ષમ છે.
  \sourcecorrection{OLMD-340}{મૂળ અહીં કોઈ પણ સૂત્રના ઓળખ-સિક્વન્ટને સ્વયંસિદ્ધ તરીકે મંજૂરી આપતું તંત્ર G3c કહે છે. આ G1cની શરત છે; G3cમાં પરમાણ્વીય ઓળખ-સિક્વન્ટો જ સ્વયંસિદ્ધ છે, જ્યારે સામાન્ય ઓળખ-સિક્વન્ટો નિષ્પન્ન થાય છે. તંત્રનું નામ તે મુજબ સુધાર્યું છે.}

  જો $n>0$ હોય, તો $\pi$નો અંત નિયમ~$R$માં થાય છે.
  આગમનની ધારણાથી તે નિયમના આધારોની સાબિતી
  \Log{G3c}માં મળે છે, એટલે કે તાત્કાલિક ઉપ-સાબિતીનાં
  \Log{G3c}માં રૂપાંતરણો મળે છે.
  જો નિયમ~$R$ એ \Log{G3c}નો પણ નિયમ હોય તો
  આધારોની રૂપાંતરિત સાબિતી પર તે નિયમ લાગુ પાડીએ.
  જે નિયમો \Log{G3c}માં એ જ સ્વરૂપે નથી, તેમને
  સ્વીકાર્ય નિયમો અને વધારાનાં શિથિલીકરણથી અનુકરિત કરીએ.
  જો નિયમ~$R$ એ $\LeftR{\land}$ અથવા $\RightR{\lor}$
  હોય, તો \olref[adm]{prop:weak-G3c-adm} મુજબ
  ચૂકી ગયેલું સક્રિય સૂત્ર ઉમેર્યા પછી સંબંધિત તાર્કિક નિયમ વાપરીએ.
  શિથિલીકરણના નિયમ માટે \olref[adm]{prop:weak-G3c-adm}
  અને સંકોચનના નિયમ માટે \olref[inv]{prop:G3c-cont-adm} વાપરીએ.
  \sourcecorrection{OLMD-341}{મૂળનો સંદર્ભ વિરુદ્ધ દિશાના નિયમ-નિષ્પન્નન અંગે છે. યોગ્ય દિશામાં, સંયોજનના ડાબા નિયમ માટે આધારના પૂર્વાંગમાં ચૂકી ગયેલું બીજું જોડક-સૂત્ર ઉમેરો અને પછી G3cનો સંયોજન-ડાબો નિયમ વાપરો. વિકલ્પના જમણા નિયમ માટે ઉત્તરાંગમાં ચૂકી ગયેલું બીજું વિકલ્પક ઉમેરો અને પછી G3cનો વિકલ્પ-જમણો નિયમ વાપરો. તેથી અહીં સ્વીકાર્ય શિથિલીકરણના પરિણામનો સંદર્ભ આપ્યો છે.}
  \sourcecorrection{OLMD-342}{મૂળ બંને દિશાની ચર્ચામાં બે પરિમાણક-નિયમોના તફાવતને ચૂકી જાય છે. G1cથી G3cમાં સર્વપરિમાણક-ડાબા આધારના પૂર્વાંગમાં મુખ્ય સર્વપરિમાણિત સૂત્ર શિથિલીકરણથી ઉમેરો; અસ્તિત્વપરિમાણક-જમણા આધારના ઉત્તરાંગમાં મુખ્ય અસ્તિત્વપરિમાણિત સૂત્ર ઉમેરો. પછી સંબંધિત G3c નિયમ લાગુ પડે છે. ઊલટી દિશામાં, G3cના આધારમાં મુખ્ય સૂત્ર પહેલેથી જ હોવાથી G1cનો સંબંધિત પરિમાણક-નિયમ લાગુ કરતાં મુખ્ય સૂત્રની બે આવૃત્તિ મળે છે; સર્વપરિમાણક માટે ડાબું અને અસ્તિત્વપરિમાણક માટે જમણું સંકોચન વાપરો. બાકીના સામાન્ય નિયમો સીધા લાગુ પડે છે. આ પગલાં બંને આગમનોના ચૂકી ગયેલા કિસ્સાઓ પૂર્ણ કરે છે.}
\end{proof}

\begin{prop}
જો $\Log{G3c} \Proves S$ હોય, તો $\Log{G1c} \Proves S$.
\end{prop}

\begin{proof}
  ફરી $S$ની સાબિતી~$\pi$ની ઊંચાઈ પર આગમન વાપરીએ.
  \Log{G3c}નું દરેક સ્વયંસિદ્ધ \Log{G1c}માં કોઈ સ્વયંસિદ્ધ
  તથા \LeftR{\Weakening} અને \RightR{\Weakening}ના
  વારંવાર ઉપયોગથી સાબિત કરી શકાય:
  \[
  \Axiom$!C \fCenter !C$
  \RightLabel{\LeftR{\Weakening}}
  \Deduce$!C, \Gamma \fCenter !C$
  \RightLabel{\RightR{\Weakening}}
  \Deduce$!C, \Gamma \fCenter \Delta, !C$
  \DisplayProof
\qquad
  \Axiom$\lfalse \fCenter$
  \RightLabel{\LeftR{\Weakening}}
  \Deduce$\lfalse, \Gamma \fCenter$
  \RightLabel{\RightR{\Weakening}}
  \Deduce$\lfalse, \Gamma \fCenter \Delta$
  \DisplayProof
  \]
  જો $\pi$નો અંત નિયમ~$R$માં થાય, તો આગમનની ધારણાથી
  આધારોની સાબિતી \Log{G1c}માં મળે છે.
  બંને તંત્રમાં એકસરખો હોય તે નિયમ~$R$ આ સાબિતી પર સીધો લાગુ પડે છે.
  \LeftR{\land} અને \RightR{\lor}ના સ્વરૂપો
  \Log{G1c} અને \Log{G3c}માં જુદાં છે;
  પરિમાણકના બે વધારાના કિસ્સા ઉપરની નોંધમાં સમજાવ્યા છે.
  \Log{G3c}ના આ બે જોડક-નિયમો \Log{G1c}માં નીચે પ્રમાણે નિષ્પન્ન થાય છે:
  \[
  \Axiom$!A, !B, \Gamma \fCenter \Delta$
  \RightLabel{\LeftR{\land}}
  \UnaryInf$!A \land !B, !B, \Gamma \fCenter \Delta$
  \RightLabel{\LeftR{\land}}
  \UnaryInf$!A \land !B, !A \land !B, \Gamma \fCenter \Delta$
  \RightLabel{\LeftR{\Contraction}}
  \UnaryInf$!A \land !B, \Gamma \fCenter \Delta$
    \DisplayProof
\qquad
  \Axiom$\Gamma \fCenter \Delta, !A, !B$
  \RightLabel{\RightR{\lor}}
  \UnaryInf$\Gamma \fCenter \Delta, !A \lor !B, !B$
  \RightLabel{\RightR{\lor}}
  \UnaryInf$\Gamma \fCenter \Delta, !A \lor !B, !A \lor !B$
  \RightLabel{\RightR{\Contraction}}
  \UnaryInf$\Gamma \fCenter \Delta, !A \lor !B$
    \DisplayProof
  \]
  \sourcecorrection{OLMD-343}{મૂળના વિકલ્પ-જમણા વૃક્ષની વચલી પંક્તિમાં બે લગોલગ અલ્પવિરામ છે. સંદર્ભ અને બે મુખ્ય સૂત્ર-આવૃત્તિઓ વચ્ચે ફક્ત જરૂરી અલ્પવિરામ જાળવ્યો છે.}
\end{proof}
% END OLP-0713
\endgroup


\OLEndChapterHook


\begin{editorial}
આગળનો સંબંધિત અંશ મૂળના સક્રિય આયાતક્રમમાં નથી. તેને અહીં સ્પષ્ટ વૈકલ્પિક સંદર્ભમાં જાળવ્યો છે.
\end{editorial}
\begingroup
\makeatletter\def\ol@sectioncs{section}\makeatother

% BEGIN OLP-0705
\phantomsection\label{guunit:OLP-0705}
\olfileid{pt}{seq}{alt0705}


\begin{table}
\begin{tabular}{@{}c@{}c@{}}
\multicolumn{2}{@{}c@{}}{સ્વયંસિદ્ધો}\\[2ex]
$!A \Sequent !A$ & $\lfalse \Sequent \quad$
\\[2ex]
\multicolumn{2}{@{}c@{}}{બંધારણીય નિયમો}\\[2ex]
\Axiom$ \Gamma \fCenter \Delta $
\RightLabel{\LeftR{\Weakening}}
\UnaryInf$ !A, \Gamma \fCenter \Delta$
\DisplayProof
&
\Axiom$ \Gamma \fCenter $
\RightLabel{\RightR{\Weakening}}
\UnaryInf$ \Gamma \fCenter !A$
\DisplayProof
\\[4ex]
\Axiom$ !A, !A, \Gamma \fCenter \Delta $
\RightLabel{\LeftR{\Contraction}}
\UnaryInf$ !A, \Gamma \fCenter \Delta$
\DisplayProof
&
\\[4ex]
\multicolumn{2}{@{}c@{}}{તાર્કિક નિયમો}\\[2ex]
\Axiom$ \Gamma \fCenter !A $
\RightLabel{\LeftR{\lnot}}
\UnaryInf$ \lnot !A, \Gamma \fCenter $
\DisplayProof
&
\Axiom$!A, \Gamma \fCenter $
\RightLabel{\RightR{\lnot}}
\UnaryInf$ \Gamma \fCenter \lnot !A $
\DisplayProof
\\[4ex]\begin{tabular}[t]{@{}l@{}}
\Axiom$ !A, \Gamma \fCenter \Delta$
\RightLabel{\LeftR{\land}}
\UnaryInf$ !A \land !B, \Gamma \fCenter \Delta$
\DisplayProof
\\[3ex]
\Axiom$!B, \Gamma \fCenter \Delta$
\RightLabel{\LeftR{\land}}
\UnaryInf$!A \land !B, \Gamma \fCenter \Delta$
\DisplayProof\\[3ex]
\end{tabular}
&
\Axiom$\Gamma \fCenter !A$
\Axiom$ \Gamma \fCenter !B$
\RightLabel{\RightR{\land}}
\BinaryInf$ \Gamma \fCenter !A \land !B $
\DisplayProof
\\[4ex]\Axiom$!A, \Gamma \fCenter \Delta$
\Axiom$!B, \Gamma \fCenter \Delta$
\RightLabel{\LeftR{\lor}}
\BinaryInf$!A \lor !B, \Gamma \fCenter \Delta$
\DisplayProof
&
\begin{tabular}[t]{@{}r@{}}
\Axiom$\Gamma \fCenter !A$
\RightLabel{\RightR{\lor}}
\UnaryInf$ \Gamma \fCenter !A \lor !B$
\DisplayProof
\\[3ex]
\Axiom$ \Gamma \fCenter !B$
\RightLabel{\RightR{\lor}}
\UnaryInf$ \Gamma \fCenter !A \lor !B$
\DisplayProof\\[3ex]
\end{tabular}
\\[4ex]
\Axiom$ \Gamma \fCenter !A$
\Axiom$ !B, \Gamma \fCenter \Delta$
\RightLabel{\LeftR{\lif}}
\BinaryInf$ !A \lif !B, \Gamma \fCenter \Delta$
\DisplayProof
&
\Axiom$ !A, \Gamma \fCenter !B$
\RightLabel{\RightR{\lif}}
\UnaryInf$ \Gamma \fCenter !A \lif !B $
\DisplayProof
\\[4ex]
\Axiom$ !A(t), \Gamma \fCenter \Delta$
\RightLabel{\LeftR{\lforall}}
\UnaryInf$ \lforall[x][!A(x)],\Gamma \fCenter \Delta$
\DisplayProof
&
\Axiom$ \Gamma \fCenter !A(a) $
\RightLabel{\RightR{\lforall}}
\UnaryInf$ \Gamma \fCenter \lforall[x][!A(x)]$
\DisplayProof
\\[4ex]
\Axiom$ !A(a), \Gamma \fCenter \Delta $
\RightLabel{\LeftR{\lexists}}
\UnaryInf$ \lexists[x][!A(x)], \Gamma \fCenter \Delta$
\DisplayProof
&
\Axiom$ \Gamma \fCenter !A(t) $
\RightLabel{\RightR{\lexists}}
\UnaryInf$ \Gamma \fCenter \lexists[x][!A(x)]$
\DisplayProof
\end{tabular}
\caption{અંતઃપ્રજ્ઞાવાદી તર્ક માટે \Log{G1i}ના નિયમો.
$\RightR{\forall}$ અને $\LeftR{\exists}$માં અચળ~$a$
ફલિત-સિક્વન્ટમાં આવવો નહીં. $\Gamma$ એ સૂત્રનો બહુગણ છે,
અને $\Delta$ ખાલી હોઈ શકે અથવા એક જ સૂત્ર ધરાવી શકે.
\RightR{\Contraction}નો નિયમ નથી. \Log{G1m} એ
\RightR{\Weakening} અને $\lfalse,
\Gamma \Sequent \Delta$ સ્વરૂપનાં સ્વયંસિદ્ધો વિનાનું \Log{G1i} છે.}
\ollabel{tab:G1c}
\end{table}
% END OLP-0705
\endgroup

\begin{editorial}
આગળનો સંબંધિત અંશ મૂળના સક્રિય આયાતક્રમમાં નથી. તેને અહીં સ્પષ્ટ વૈકલ્પિક સંદર્ભમાં જાળવ્યો છે.
\end{editorial}
\begingroup
\makeatletter\def\ol@sectioncs{section}\makeatother

% BEGIN OLP-0708
\phantomsection\label{guunit:OLP-0708}
\olfileid{pt}{seq}{alt0708}


\begin{table}
\begin{tabular}{@{}c@{}c@{}}
\multicolumn{2}{@{}c@{}}{સ્વયંસિદ્ધો}\\[2ex]
$!C, \Gamma \Sequent !C$ & $\lfalse, \Gamma \Sequent \Delta$
\\[2ex]
\multicolumn{2}{@{}c@{}}{તાર્કિક નિયમો}\\[2ex]
\Axiom$\lnot !A, \Gamma \fCenter !A $
\RightLabel{\LeftR{\lnot}}
\UnaryInf$\lnot !A, \Gamma \fCenter \Delta$
\DisplayProof
&
\Axiom$!A, \Gamma \fCenter$
\RightLabel{\RightR{\lnot}}
\UnaryInf$ \Gamma \fCenter \lnot !A $
\DisplayProof
\\[4ex]
\Axiom$ !A, !B, \Gamma \fCenter \Delta$
\RightLabel{\LeftR{\land}}
\UnaryInf$ !A \land !B, \Gamma \fCenter \Delta$
\DisplayProof
&
\Axiom$\Gamma \fCenter !A$
\Axiom$ \Gamma \fCenter !B$
\RightLabel{\RightR{\land}}
\BinaryInf$ \Gamma \fCenter !A \land !B $
\DisplayProof
\\[4ex]\Axiom$!A, \Gamma \fCenter \Delta$
\Axiom$!B, \Gamma \fCenter \Delta$
\RightLabel{\LeftR{\lor}}
\BinaryInf$!A \lor !B, \Gamma \fCenter \Delta$
\DisplayProof
&
\begin{tabular}[t]{@{}r@{}}
\Axiom$\Gamma \fCenter !A$
\RightLabel{\RightR{\lor}}
\UnaryInf$ \Gamma \fCenter !A \lor !B$
\DisplayProof
\\[3ex]
\Axiom$\Gamma \fCenter !B$
\RightLabel{\RightR{\lor}}
\UnaryInf$ \Gamma \fCenter !A \lor !B$
\DisplayProof\\[3ex]
\end{tabular}
\\[4ex]
\Axiom$!A \lif !B, \Gamma \fCenter !A$
\Axiom$!B, \Gamma \fCenter \Delta$
\RightLabel{\LeftR{\lif}}
\BinaryInf$!A \lif !B, \Gamma \fCenter \Delta$
\DisplayProof
&
\Axiom$!A, \Gamma \fCenter !B$
\RightLabel{\RightR{\lif}}
\UnaryInf$ \Gamma \fCenter !A \lif !B $
\DisplayProof
\\[4ex]
\Axiom$ !A(t), \lforall[x][!A(x)], \Gamma \fCenter \Delta$
\RightLabel{\LeftR{\lforall}}
\UnaryInf$ \lforall[x][!A(x)],\Gamma \fCenter \Delta$
\DisplayProof
&
\Axiom$ \Gamma \fCenter !A(a) $
\RightLabel{\RightR{\lforall}}
\UnaryInf$ \Gamma \fCenter \lforall[x][!A(x)]$
\DisplayProof
\\[4ex]
\Axiom$ !A(a), \Gamma \fCenter \Delta $
\RightLabel{\LeftR{\lexists}}
\UnaryInf$ \lexists[x][!A(x)], \Gamma \fCenter \Delta$
\DisplayProof
&
\Axiom$ \Gamma \fCenter !A(t) $
\RightLabel{\RightR{\lexists}}
\UnaryInf$ \Gamma \fCenter \lexists[x][!A(x)]$
\DisplayProof
\end{tabular}
\caption{અંતઃપ્રજ્ઞાવાદી તર્ક માટે \Log{G3i}ના નિયમો.
$\RightR{\forall}$ અને $\LeftR{\exists}$માં અચળ~$a$
ફલિત-સિક્વન્ટમાં આવવો નહીં. $\Gamma$ એ સૂત્રનો બહુગણ છે,
અને $\Delta$ ખાલી હોઈ શકે અથવા એક જ સૂત્ર ધરાવી શકે.
સ્વયંસિદ્ધોમાં સૂત્ર~$!C$ પરમાણ્વીય હોવું જોઈએ.
અલગ તંત્ર અંગે નોંધ: \Log{G1m} એ જમણા શિથિલીકરણ અને
$\lfalse, \Gamma \Sequent
\Delta$ સ્વરૂપનાં સ્વયંસિદ્ધો વિનાનું \Log{G1i} છે.}
\ollabel{tab:G3c}
\end{table}
\sourcecorrection{OLMD-334}{મૂળ સર્વપરિમાણકના ડાબા નિયમના આધારમાં મુખ્ય સૂત્ર અને સંદર્ભ વચ્ચેનો અલ્પવિરામ ઉમેર્યો છે.}
\sourcecorrection{OLMD-335}{મૂળ વિકલ્પના જમણા નિયમના આધારમાં બે સૂત્રો છે, જ્યારે આ જ કોષ્ટકની શરત ઉત્તરાંગમાં વધુમાં વધુ એક સૂત્રની છૂટ આપે છે. તેથી એક સૂત્રવાળા બે જુદા દાખલ-નિયમ દર્શાવ્યા છે: એકમાં પહેલું વિકલ્પક અને બીજામાં બીજું વિકલ્પક આધાર છે. આ સ્વરૂપ G1iના કોષ્ટક અને નિયમોના અર્થઘટન સાથે મેળ ખાય છે.}
\sourcecorrection{OLMD-336}{મૂળની છેલ્લી કૅપ્શન-નોંધ G1નાં તંત્રો અંગે છે; તેને આ G3i કોષ્ટકની વ્યાખ્યા ગણવી નહીં. G1iના અલગ કોષ્ટક અનુસાર G1mમાં જમણું શિથિલીકરણ પણ દૂર કરવું પડે છે. અહીં એ ચૂકી ગયેલી શરત ઉમેરેલી છે; ફક્ત અસત્યનાં સ્વયંસિદ્ધો દૂર કરવાથી મૂળ નકાર-નિયમો સાથે ઇચ્છિત ન્યૂનતમ તર્ક મળતું નથી.}
% END OLP-0708
\endgroup

\begin{editorial}
આગળનો સંબંધિત અંશ મૂળના સક્રિય આયાતક્રમમાં નથી. તેને અહીં સ્પષ્ટ વૈકલ્પિક સંદર્ભમાં જાળવ્યો છે.
\end{editorial}
\begingroup
\makeatletter\def\ol@sectioncs{section}\makeatother

% BEGIN OLP-0709
\phantomsection\label{guunit:OLP-0709}
\olfileid{pt}{seq}{alt0709}


\begin{table}
\begin{tabular}{@{}c@{}c@{}}
\multicolumn{2}{@{}c@{}}{સ્વયંસિદ્ધો}\\[2ex]
$!A \Sequent !A$ & $\lfalse \Sequent \quad$
\\[2ex]
\multicolumn{2}{@{}c@{}}{બંધારણીય નિયમો}\\[2ex]
\Axiom$ \Gamma \fCenter \Delta $
\RightLabel{\LeftR{\Weakening}}
\UnaryInf$ !A, \Gamma \fCenter \Delta$
\DisplayProof
&
\Axiom$ \Gamma \fCenter \Delta$
\RightLabel{\RightR{\Weakening}}
\UnaryInf$ \Gamma \fCenter \Delta, !A$
\DisplayProof
\\[4ex]
\Axiom$ !A, !A, \Gamma \fCenter \Delta $
\RightLabel{\LeftR{\Contraction}}
\UnaryInf$ !A, \Gamma \fCenter \Delta$
\DisplayProof
&
\Axiom$ \Gamma \fCenter \Delta, !A, !A$
\RightLabel{\RightR{\Contraction}}
\UnaryInf$ \Gamma \fCenter \Delta, !A$
\DisplayProof
\\[4ex]
\Axiom$ \Gamma, !A, !B, \Pi \fCenter \Delta $
\RightLabel{\LeftR{\Exchange}}
\UnaryInf$\Gamma, !B, !A, \Pi \fCenter \Delta$
\DisplayProof
&
\Axiom$ \Gamma \fCenter \Delta, !A, !B, \Lambda$
\RightLabel{\RightR{\Exchange}}
\UnaryInf$ \Gamma \fCenter \Delta, !B, !A, \Lambda$
\DisplayProof
\\[4ex]
\multicolumn{2}{@{}c@{}}{તાર્કિક નિયમો}\\[2ex]
\Axiom$ \Gamma \fCenter \Delta, !A $
\RightLabel{\LeftR{\lnot}}
\UnaryInf$ \lnot !A, \Gamma \fCenter \Delta$
\DisplayProof
&
\Axiom$!A, \Gamma \fCenter \Delta$
\RightLabel{\RightR{\lnot}}
\UnaryInf$ \Gamma \fCenter \Delta, \lnot !A $
\DisplayProof
\\[4ex]\begin{tabular}[t]{@{}l@{}}
\Axiom$ !A, \Gamma \fCenter \Delta$
\RightLabel{\LeftR{\land}}
\UnaryInf$ !A \land !B, \Gamma \fCenter \Delta$
\DisplayProof
\\[3ex]
\Axiom$!B, \Gamma \fCenter \Delta$
\RightLabel{\LeftR{\land}}
\UnaryInf$!A \land !B, \Gamma \fCenter \Delta$
\DisplayProof\\[3ex]
\end{tabular}
&
\Axiom$\Gamma \fCenter \Delta, !A$
\Axiom$ \Gamma \fCenter \Delta, !B$
\RightLabel{\RightR{\land}}
\BinaryInf$ \Gamma \fCenter \Delta, !A \land !B $
\DisplayProof
\\[4ex]\Axiom$!A, \Gamma \fCenter \Delta$
\Axiom$!B, \Gamma \fCenter \Delta$
\RightLabel{\LeftR{\lor}}
\BinaryInf$!A \lor !B, \Gamma \fCenter \Delta$
\DisplayProof
&
\begin{tabular}[t]{@{}r@{}}
\Axiom$\Gamma \fCenter \Delta, !A$
\RightLabel{\RightR{\lor}}
\UnaryInf$ \Gamma \fCenter \Delta, !A \lor !B$
\DisplayProof
\\[3ex]
\Axiom$ \Gamma \fCenter \Delta, !B$
\RightLabel{\RightR{\lor}}
\UnaryInf$ \Gamma \fCenter \Delta, !A \lor !B$
\DisplayProof\\[3ex]
\end{tabular}
\\[4ex]
\Axiom$ \Gamma \fCenter \Delta, !A$
\Axiom$ !B, \Gamma \fCenter \Delta$
\RightLabel{\LeftR{\lif}}
\BinaryInf$ !A \lif !B, \Gamma \fCenter \Delta$
\DisplayProof
&
\Axiom$ !A, \Gamma \fCenter \Delta, !B$
\RightLabel{\RightR{\lif}}
\UnaryInf$ \Gamma \fCenter \Delta, !A \lif !B $
\DisplayProof
\\[4ex]
\Axiom$ !A(t), \Gamma \fCenter \Delta$
\RightLabel{\LeftR{\lforall}}
\UnaryInf$ \lforall[x][!A(x)],\Gamma \fCenter \Delta$
\DisplayProof
&
\Axiom$ \Gamma \fCenter \Delta, !A(a) $
\RightLabel{\RightR{\lforall}}
\UnaryInf$ \Gamma \fCenter \Delta, \lforall[x][!A(x)]$
\DisplayProof
\\[4ex]
\Axiom$ !A(a), \Gamma \fCenter \Delta $
\RightLabel{\LeftR{\lexists}}
\UnaryInf$ \lexists[x][!A(x)], \Gamma \fCenter \Delta$
\DisplayProof
&
\Axiom$ \Gamma \fCenter \Delta, !A(t) $
\RightLabel{\RightR{\lexists}}
\UnaryInf$ \Gamma \fCenter \Delta, \lexists[x][!A(x)]$
\DisplayProof
\end{tabular}
\caption{\Log{LK}ના નિયમો. $\RightR{\forall}$ અને
$\LeftR{\exists}$માં અચળ~$a$ ફલિત-સિક્વન્ટમાં આવવો નહીં.
$\Gamma$, $\Delta$ અને $\Pi$ એ સૂત્રના ક્રમો છે.}
\ollabel{tab:G1c}
\end{table}
% END OLP-0709
\endgroup

\begin{editorial}
આગળનો સંબંધિત અંશ મૂળના સક્રિય આયાતક્રમમાં નથી. તેને અહીં સ્પષ્ટ વૈકલ્પિક સંદર્ભમાં જાળવ્યો છે.
\end{editorial}
\begingroup
\makeatletter\def\ol@sectioncs{section}\makeatother

% BEGIN OLP-0710
\phantomsection\label{guunit:OLP-0710}
\olfileid{pt}{seq}{alt0710}


\begin{table}
\begin{tabular}{@{}c@{}c@{}}
\multicolumn{2}{@{}c@{}}{સ્વયંસિદ્ધો}\\[2ex]
$!C, \Gamma \Sequent \Delta, !C$ & $\lfalse, \Gamma \Sequent \Delta$
\\[2ex]
\multicolumn{2}{@{}c@{}}{તાર્કિક નિયમો}\\[2ex]
\Axiom$\lnot !A, \Gamma \fCenter \Delta, !A $
\RightLabel{\LeftR{\lnot}}
\UnaryInf$\lnot !A, \Gamma \fCenter \Delta$
\DisplayProof
&
\Axiom$!A, \Gamma \fCenter$
\RightLabel{\RightR{\lnot}}
\UnaryInf$ \Gamma \fCenter \Delta, \lnot !A $
\DisplayProof
\\[4ex]
\Axiom$ !A, !B, \Gamma \fCenter \Delta$
\RightLabel{\LeftR{\land}}
\UnaryInf$ !A \land !B, \Gamma \fCenter \Delta$
\DisplayProof
&
\Axiom$\Gamma \fCenter \Delta, !A$
\Axiom$ \Gamma \fCenter \Delta, !B$
\RightLabel{\RightR{\land}}
\BinaryInf$ \Gamma \fCenter \Delta, !A \land !B $
\DisplayProof
\\[4ex]\Axiom$!A, \Gamma \fCenter \Delta$
\Axiom$!B, \Gamma \fCenter \Delta$
\RightLabel{\LeftR{\lor}}
\BinaryInf$!A \lor !B, \Gamma \fCenter \Delta$
\DisplayProof
&
\Axiom$\Gamma \fCenter \Delta, !A, !B$
\RightLabel{\RightR{\lor}}
\UnaryInf$ \Gamma \fCenter \Delta, !A \lor !B$
\DisplayProof
\\[4ex]
\Axiom$!A \lif !B, \Gamma \fCenter \Delta, !A$
\Axiom$ !B, \Gamma \fCenter \Delta$
\RightLabel{\LeftR{\lif}}
\BinaryInf$!A \lif !B, \Gamma \fCenter \Delta$
\DisplayProof
&
\Axiom$ !A, \Gamma \fCenter !B$
\RightLabel{\RightR{\lif}}
\UnaryInf$ \Gamma \fCenter \Delta, !A \lif !B $
\DisplayProof
\\[4ex]
\Axiom$ !A(t), \lforall[x][!A(x)], \Gamma \fCenter \Delta$
\RightLabel{\LeftR{\lforall}}
\UnaryInf$ \lforall[x][!A(x)],\Gamma \fCenter \Delta$
\DisplayProof
&
\Axiom$ \Gamma \fCenter !A(a) $
\RightLabel{\RightR{\lforall}}
\UnaryInf$ \Gamma \fCenter \lforall[x][!A(x)]$
\DisplayProof
\\[4ex]
\Axiom$ !A(a), \Gamma \fCenter \Delta $
\RightLabel{\LeftR{\lexists}}
\UnaryInf$ \lexists[x][!A(x)], \Gamma \fCenter \Delta$
\DisplayProof
&
\Axiom$ \Gamma \fCenter \Delta, \lexists[x][!A(x)], !A(t) $
\RightLabel{\RightR{\lexists}}
\UnaryInf$ \Gamma \fCenter \Delta, \lexists[x][!A(x)]$
\DisplayProof
\end{tabular}
\caption{અંતઃપ્રજ્ઞાવાદી તર્ક માટે અનેક ફલિત ધરાવતા \Log{mG3i}ના નિયમો.
$\RightR{\forall}$ અને $\LeftR{\exists}$માં અચળ~$a$
ફલિત-સિક્વન્ટમાં આવવો નહીં. $\Gamma$ અને $\Delta$ એ
સૂત્રના બહુગણ છે. સ્વયંસિદ્ધોમાં સૂત્ર~$!C$
પરમાણ્વીય હોવું જોઈએ. મૂળ કૅપ્શન \Log{mG1m}ને
$\lfalse,
\Gamma \Sequent \Delta$ સ્વરૂપનાં સ્વયંસિદ્ધો વિનાના \Log{mG1i}
તરીકે પણ વર્ણવે છે.}
\ollabel{tab:G3c}
\end{table}
\sourcecorrection{OLMD-337}{મૂળ સર્વપરિમાણકના ડાબા નિયમના આધારમાં મુખ્ય સૂત્ર અને સંદર્ભ વચ્ચેનો અલ્પવિરામ ઉમેર્યો છે.}
\sourcecorrection{OLMD-338}{મૂળ કૅપ્શનના અંતે mG1i અને mG1m આવે છે, પરંતુ આ અધ્યાય તેમના નિયમો આપતો નથી. તેથી આ વાક્યને મૂળનો ઉલ્લેખ માની જાળવ્યું છે. તેને mG3iમાંથી ન્યૂનતમ તંત્ર બનાવવાની વ્યાખ્યા ગણવી નહીં; એવા તંત્રના નિયમો માટે જરૂરી માહિતી અહીં મળતી નથી.}
% END OLP-0710
\endgroup
% END OLP-0712
\endgroup

\begingroup
\makeatletter\def\ol@sectioncs{section}\makeatother

% BEGIN OLP-0657
\olchapter{pt}{cut}{કટ-નિરાકરણ}
\olchapterid{pt}{cut}
\phantomsection\label{guunit:OLP-0657}


\begingroup
\makeatletter\def\ol@sectioncs{section}\makeatother

% BEGIN OLP-0659
\olfileid{pt}{cut}{int}

\olsection{પરિચય}
\phantomsection\label{guunit:OLP-0659}


\Cut{} નિયમ આ છે:
\[
\Axiom$\Gamma \fCenter \Delta, !A$
\Axiom$!A, \Pi \fCenter \Lambda$
\RightLabel{\Cut}
\BinaryInf$\Gamma, \Pi \fCenter \Delta, \Lambda$
\DisplayProof
\]
(સૂત્ર~$!A$ને
\emph{કટ-સૂત્ર} કહે છે.)

\Cut{} એ ગેન્ઝેનના મૂળ
સિક્વન્ટ કલન~\Log{LK}માં હતું.
ગેન્ઝેનનું \emph{Hauptsatz}
(``મુખ્ય પ્રમેય'') કહે છે કે
~\Log{LK}માં સાબિતીક્ષમ
દરેક સિક્વન્ટ \Cut{} નિયમ વિના
પણ સાબિતીક્ષમ છે;
એટલે~\Log{LK}-સાબિતીઓમાંથી
\Cut{} ``દૂર'' કરી શકાય.
આપણી પ્રણાલીઓ~\Log{G1c}
અને~\Log{G3c}માં \Cut{} નથી,
પણ તેમના માટે પણ પ્રશ્ન ઊભો થાય:
\Log{G1c}ના (અથવા~\Log{G3c}ના)
બીજા નિયમો સાથે~\Cut{} વાપરતી
સાબિતીઓમાંથી \Cut{} દૂર કરી શકાય?
પહેલેથી નક્કી કરેલી પરિભાષામાં
એ જ પ્રશ્ન છે: શું~\Cut{}
એ~\Log{G1c}નો
(અથવા~\Log{G3c}નો)
સ્વીકાર્ય નિયમ છે?

કટ-નિરાકરણ પ્રમેય સાબિતી સિદ્ધાંતનું
સંભવતઃ સૌથી મહત્ત્વનું અને જાણીતું
પરિણામ છે. પહેલાં તો તે બતાવે છે કે
~\Log{G1c} અને~\Log{G3c}ને
પ્રથમ નજરે અનિવાર્ય લાગે એવો નિયમ
જરૂરી નથી. વાસ્તવિક સાબિતીઓનું
ઔપચારિકીકરણ કરો ત્યારે ઉપપ્રમેયો
વાપરવાની રીત જોઈએ જ.
ગણિતજ્ઞો એક પરિણામ સાબિત કરે
અને પછી બીજી સાબિતીઓમાં તેનો
આધાર લે છે. તેથી પ્રથમ પરિણામ~$L$
(ઉપપ્રમેય)ને
$\Sequent L$ સિક્વન્ટની
સાબિતી તરીકે ઔપચારિક કરી શકાય.
પછી~$L$ વાપરતા બીજા
પરિણામ~$R$ની સાબિતીને
$L \Sequent R$ સિક્વન્ટની
સાબિતી તરીકે લખો.
માત્ર~$R$ની સાબિતી,
એટલે $\Sequent R$ની
સાબિતી, છે તે કેવી રીતે જાણવું?
બે સાબિતીઓ જોડવાની રીત
જોઈએ; \Cut{} નિયમ એ કરવા દે છે.

આપણે કહ્યું હતું કે~\Log{G1c}
અને~\Log{G3c} પૂર્ણ છે;
એટલે સાબિત થવા જોઈએ એવા
બધા પ્રામાણ્ય ધરાવતા સિક્વન્ટ
તેઓ સાબિત કરે છે. આ સીધું
પણ સાબિત કરી શકાય.
પરંતુ ગેન્ઝેન લખતા હતા ત્યારે
માત્ર સ્વયંસિદ્ધિમૂલક નિષ્પત્તિ
પ્રણાલીઓની પૂર્ણતા જાણીતી હતી.
~\Log{LK}ની પૂર્ણતા બતાવવા
તેમણે સ્વયંસિદ્ધિમૂલક પ્રણાલીની
નિષ્પત્તિઓને
\Log{LK}-સાબિતીઓમાં રૂપાંતરિત
કરવાની રીત આપી. તેમાં~\Cut{}
અનિવાર્ય રીતે વપરાતો હતો.
કટ-નિરાકરણ માત્ર શાસ્ત્રીય તર્ક
માટે મહત્ત્વનું નથી: બીજા અનેક
તર્કો માટે પણ સિક્વન્ટ કલન છે.
કેટલાક તર્કોના સિક્વન્ટ કલનની
પૂર્ણતા સીધી સાબિત કરવી મુશ્કેલ
અથવા અશક્ય હોય છે; તેથી
બીજી પ્રણાલીઓમાંથી~\Cut{}
વાપરતા રૂપાંતરો જ પૂર્ણતા
બતાવવાનો માર્ગ હોય છે.
પછી~\Cut{} બિનજરૂરી છે અને
તે વિનાની પ્રણાલી પૂર્ણ છે
એ સ્થાપવા સાબિતી-સૈદ્ધાંતિક
રીતે કટ-નિરાકરણ સાબિત કરવું પડે.

કટ-નિરાકરણ વડે સ્વતંત્ર રસનાં
બીજાં પરિણામો પણ મેળવી શકાય
એથી પણ તે મહત્ત્વનું છે.
આપણે ક્રેગના અંતર્વેશન ઉપપ્રમેય
અને હર્બ્રાન્ડના પ્રમેયને જોઈશું.

કટ-નિરાકરણની સાબિતીમાં સામાન્ય રીતે
~\Cut{} વાપરતી સાબિતીને
તે વિનાની સાબિતીમાં કેવી રીતે
ફેરવવી તે બતાવાય છે.
આ રૂપાંતરો મોટા ભાગે
``સ્થાનિક'' હોય છે:
સાબિતીનો એક ભાગ
(સામાન્ય રીતે~\Cut{} અનુમાનમાં
પૂરી થતી ઉપ-સાબિતી) લઈને
તેને એ જ સિક્વન્ટની બીજી
ઉપ-સાબિતીથી બદલીયે,
જેમાં~\Cut{} અનુમાન દૂર થયું
હોય અથવા બીજા અનુમાનોથી
બદલાયું હોય.
\sourcecorrection{OLMD-213}{મૂળમાં glossary ચિહ્ન સાબિતી બે પંક્તિ વચ્ચે તૂટેલું છે; અહીં અખંડ ચિહ્ન પુનઃસ્થાપિત કર્યું છે.}
આવી દલીલો
\Cut{} ધરાવતી સાબિતીઓને
તે વિનાની સાબિતીઓમાં ફેરવવાના
પગલું-દર-પગલું કલનવિધિ આપે છે.
ઉપયોગોમાં કટ-નિરાકરણ પ્રક્રિયાઓ
લઈએ ત્યારે આ વધુ માહિતી આપે છે.
દાખલા તરીકે, અંતર્વેશન ઉપપ્રમેયની
નિદર્શ-સૈદ્ધાંતિક સાબિતી કહી શકે
કે જો~$!A \lif !B$ પ્રામાણ્ય ધરાવે
તો અંતર્વેશક છે
(હમણાં અંતર્વેશક શું છે તે છોડીએ).
કટ-નિરાકરણ વાપરતી
સાબિતી-સૈદ્ધાંતિક દલીલ
$!A \lif !B$ની સાબિતી
લઈ અંતર્વેશક આપતો કલનવિધિ આપે છે.
નિદર્શ-સૈદ્ધાંતિક દલીલની તુલનામાં
આ દલીલ \emph{રચનાત્મક} છે.

કટ-નિરાકરણ પ્રમેય સાબિત કરવાની
બે પ્રચલિત રીતો છે.
પહેલી (ગેન્ઝેનથી શરૂ થયેલી)
સાબિતીમાં સૌથી ઉપરના કટ દૂર કરી
શકાય છે એ બતાવે છે, અને
પછી એક પછી એક એવા કટ
દૂર કરે છે. બીજી
(મૂળે શ્યુટે અને ટાઇટની)
સાબિતીમાં સૌથી જટિલ કટ
પર ધ્યાન આપે છે અને એવા
મહત્તમ કટ ક્રમશઃ દૂર કરે છે
જેનાથી વધુ ઊંડાઈનું કટ
સૂત્ર બીજા કટમાં ન હોય.
દરેક રીતના લાભ અને મર્યાદાઓ છે.
કેટલીક પ્રણાલીઓમાં એક રીત ચાલે,
જ્યારે બીજી અમલમાં મૂકવી
મુશ્કેલ કે અશક્ય હોય.
તેથી બંને રીતો વર્ણવીશું.
કટ-નિરાકરણ~\Log{G3c} માટે
સાબિત કરીશું. વાસ્તવમાં
~\Cut{}ને નહીં, પણ
\emph{સંદર્ભ-સહભાગી કટ}
~\CutCS{}ને દૂર કરી શકાય
એ બતાવીશું:
\[
\Axiom$\Gamma \fCenter \Delta, !A$
\Axiom$!A, \Gamma \fCenter \Delta$
\RightLabel{\CutCS}
\BinaryInf$\Gamma \fCenter \Delta$
\DisplayProof
\]
\sourcecorrection{OLMD-214}{મૂળ બીજા વૃક્ષનું વર્ણન $\CutCS$નું છે, પણ વૃક્ષનું નિયમ-લેબલ $\Cut$ છે; સમાન સંદર્ભવાળા કટનું યોગ્ય લેબલ અહીં પુનઃસ્થાપિત કર્યું છે.}
જો સિક્વન્ટ કલન શિથિલીકરણ અને
સંકોચનને મંજૂરી આપે
(~\Log{G1c}ની જેમ સ્પષ્ટ
નિયમો તરીકે કે~\Log{G3c}ની
જેમ સ્વીકાર્ય નિયમો તરીકે),
તો~\Cut{} વડે~\CutCS{}નું
અને તેનાથી વિરુદ્ધ પણ
અનુકરણ કરી શકાય.
અમે~\Cut{}ને બદલે~\CutCS{}
પસંદ કરીએ છીએ, કારણ કે પહેલી
કટ-નિરાકરણ રીત~\CutCS{} માટે
વર્ણવવી સહેલી છે અને બીજી
રીત માત્ર~\CutCS{} માટે ચાલે છે.
% END OLP-0659
\endgroup

\begingroup
\makeatletter\def\ol@sectioncs{section}\makeatother

% BEGIN OLP-0656
\olfileid{pt}{cut}{top}

\olsection{સૌથી ઉપરના કટ દૂર કરવાં}
\phantomsection\label{guunit:OLP-0656}


\begin{defn}
\CutCS{} નિયમથી પૂરી થતી સાબિતી~$\pi$ લો,
\[
\AxiomC{}
\RightLabel{$\pi_1$}
\Deduce$\Gamma \fCenter \Delta, !A$
\AxiomC{}
\RightLabel{$\pi_2$}
\Deduce$!A, \Gamma \fCenter \Delta$
\RightLabel{\CutCS}
\BinaryInf$\Gamma \fCenter \Delta$
\DisplayProof
\]
અને તેમાં બીજું કોઈ \CutCS{} અનુમાન ન હોય.
~$\pi$ની \emph{કટ-ઊંચાઈ}
$\cheight{\pi} = \pheight{\pi_1} + \pheight{\pi_2}$ છે.
~$\pi$નો \emph{કટ-ક્રમ}
$\cutrank{\pi} = \depth{!A}$ છે.
\end{defn}

\begin{lem}\ollabel{lem:cut-adm-G3c}
\CutCS{} નિયમ~\Log{G3c}માં સ્વીકાર્ય છે.
એટલે, જો~$\pi$ એ માત્ર એક~\CutCS{}થી
પૂરી થતી અને અન્યથા માત્ર~$\Log{G3c}$ના નિયમો
વાપરતી સાબિતી હોય, તો એ જ અંતિમ સિક્વન્ટની
સાબિતી~$\pi'$ એ~\Log{G3c}માં છે.
\end{lem}

\begin{proof}
સાબિતી કટની ઊંચાઈ અને ક્રમ પર
દ્વિ-આગમનથી છે. સાબિતી~$\pi$નો અંત આ રીતે થાય છે:
\[
\AxiomC{}
\RightLabel{$\pi_1$}
\Deduce$\Gamma \fCenter \Delta, !A$
\AxiomC{}
\RightLabel{$\pi_2$}
\Deduce$!A, \Gamma \fCenter \Delta$
\RightLabel{\CutCS}
\BinaryInf$\Gamma \fCenter \Delta$
\DisplayProof
\]

આગમનનો આધાર: $\cheight{\pi} = 0$ અને
$\cutrank{\pi} = 0$. કારણ કે
$\cheight{\pi} = \pheight{\pi_1} + \pheight{\pi_2} = 0$,
$\Gamma \Sequent \Delta, !A$ અને
$!A, \Gamma \Sequent \Delta$ બંને સ્વયંસિદ્ધ છે.
બે કિસ્સા છે: (a)~તેમાંથી કોઈ એકમાં~$!A$
મુખ્ય નથી, અથવા (b)~બંનેમાં~$!A$ મુખ્ય છે.
નીચે બંને કિસ્સામાં
$\Gamma \Sequent \Delta$ સ્વયંસિદ્ધ છે
એ બતાવીશું (અને તે માટે આગમન-પૂર્વધારણા
જરૂરી નથી).

હવે~$\pi_1$ અને $\pi_2$ સ્વયંસિદ્ધ છે કે
અનુમાનથી પૂરી થાય છે, અને એ અનુમાનમાં
$!A$ મુખ્ય છે કે નહીં, તેના આધારે
કિસ્સા જુદા પાડીએ. કિસ્સા~A અને~B
આગમનનો આધાર આવરી લે છે.
આગમન-પગલામાં
$\cheight{\pi} > 0$ અથવા
$\cutrank{\pi} > 0$ (અથવા બંને)
ધારીએ છીએ. આગમન-પૂર્વધારણા મુજબ,
માત્ર એક છેલ્લા~\CutCS{}થી પૂરી થતી
એવી કોઈ પણ સાબિતીમાંથી કટ દૂર કરી શકાય,
જેનો કટ-ક્રમ $= \cutrank{\pi}$ અને
કટ-ઊંચાઈ $< \cheight{\pi}$ હોય,
અથવા કટ-ક્રમ~$< \cutrank{\pi}$ હોય.
જ્યારે $\cheight{\pi} = 0$ હોય
(બંને આધારવિધાન સ્વયંસિદ્ધ હોય), ત્યારે
કિસ્સા~A અને~B લાગુ પડે છે.
$\cheight{\pi} > 0$ના કિસ્સા C--Dમાં આવે છે.

\paragraph{કિસ્સો A:}
કોઈ~$!B$ એ~$\Gamma$ અને~$\Delta$ બંનેમાં છે.
તો~$\Gamma \Sequent \Delta$ સ્વયંસિદ્ધ છે
(જો~$!B$ પરમાણ્વીય હોય), અથવા
\olref[seq][ex]{prop:G3c-axioms} મુજબ
\Log{G3c}માં તેની \CutCS{} વિનાની
સાબિતી~$\pi'$ છે. આ આગમનના આધારનો
કિસ્સો~(a) આવરી લે છે: જો
$\Gamma \Sequent \Delta, !A$ કે
$!A, \Gamma \Sequent \Delta$ સ્વયંસિદ્ધ હોય,
પણ તેમાં~$!A$ મુખ્ય ન હોય,
તો કોઈ બીજું પરમાણ્વીય~$!B$
એ~$\Gamma$ અને~$\Delta$ બંનેમાં હોવું જોઈએ.
જો કોઈ એક આધારવિધાન એવું સ્વયંસિદ્ધ હોય
જેમાં~$!A$ મુખ્ય નથી,
તો બીજું આધારવિધાન નિયમનો નિષ્કર્ષ હોય
ત્યારે પણ આ આગમન-પગલું આવરી લે છે.

\paragraph{કિસ્સો B:}
બંને આધારવિધાન સ્વયંસિદ્ધ છે અને
$!A$ મુખ્ય છે, તેથી પરમાણ્વીય છે.
ડાબું આધારવિધાન
$\Gamma \Sequent \Delta, !A$ લો.
$!A$ મુખ્ય હોવાથી તે~$\Gamma$માં હોવું જોઈએ
(નહીં તો~$\Gamma \Sequent \Delta, !A$
સ્વયંસિદ્ધ ન હોત). એ જ રીતે
$!A$ એ~$\Delta$માં હોવું જોઈએ,
નહીં તો~$!A, \Gamma \Sequent \Delta$
સ્વયંસિદ્ધ ન હોત. તેથી પરમાણ્વીય~$!A$
એ~$\Gamma$ અને~$\Delta$ બંનેમાં છે,
એટલે~$\Gamma \Sequent \Delta$
સ્વયંસિદ્ધ છે. આ આગમનના આધારનો
કિસ્સો~(b) આવરી લે છે.

\begin{center}
કિસ્સા A અને Bની વૈકલ્પિક રજૂઆત
\end{center}

\paragraph{કિસ્સો A:}
કટનાં બે આધારવિધાનમાંથી એક,
પરમાણ્વીય~$!C$ સાથે
$!C, \Gamma' \Sequent \Delta', !C$
સ્વરૂપનું સ્વયંસિદ્ધ છે.
જો~$!C$ એ કટ સૂત્ર~$!A$ ન હોય,
તો~$!C$ એ~$\Gamma$ અને~$\Delta$ બંનેમાં
હોવું જોઈએ. ત્યારે
$\Gamma \Sequent \Delta$ સ્વયંસિદ્ધ છે
અને તેને~$\pi'$ લઈ શકાય.
જો~$!C \ident !A$ એ કટ સૂત્ર હોય,
તો બે શક્યતાઓ છે: ડાબું આધારવિધાન સ્વયંસિદ્ધ
હોય તો $!A \in \Gamma$ અને
$\Gamma = \Gamma', !A$;
જમણું સ્વયંસિદ્ધ હોય તો
$!A \in \Delta$ અને
$\Delta = \Delta', !A$.
પહેલા કિસ્સામાં~$\pi_2$ એ
$!A, !A, \Gamma' \Sequent \Delta$ની
સાબિતી છે; \LeftR{\Contraction}ની
સ્વીકાર્યતા
(\olref[seq][inv]{prop:G3c-cont-adm}) વડે
$\Gamma \Sequent \Delta$ની સાબિતી મળે છે.
બીજા કિસ્સામાં~$\pi_1$ એ
$\Gamma \Sequent \Delta', !A, !A$ની
સાબિતી છે; \RightR{\Contraction}ની
સ્વીકાર્યતા વડે
$\Gamma \Sequent \Delta$ની સાબિતી મળે છે.
\sourcecorrection{OLMD-200}{મૂળમાં બીજા કિસ્સામાં $\pi_2$ લખ્યું છે; બેવડા ઉત્તરાંગવાળી સાબિતી ડાબા આધારવિધાનની $\pi_1$ છે.}

\paragraph{કિસ્સો B:}
આધારવિધાનોમાંથી એક
$\lfalse, \Gamma' \Sequent \Delta'$
સ્વરૂપનું સ્વયંસિદ્ધ છે.
જો ડાબું આધારવિધાન આવું હોય,
તો $\lfalse \in \Gamma$ અને
$\Gamma \Sequent \Delta$ પણ સ્વયંસિદ્ધ છે.
જો જમણું આવું હોય, તો
$\lfalse \in \Gamma$ હોઈ શકે
(ત્યારે પણ $\Gamma \Sequent \Delta$
સ્વયંસિદ્ધ છે), અથવા
$!A \ident \lfalse$ હોઈ શકે.
પછીના કિસ્સામાં~$\pi_1$ એ
$\Gamma \Sequent \Delta, \lfalse$ની
સાબિતી છે. ~$\pi_1$ના દરેક
સિક્વન્ટના ઉત્તરાંગમાંથી~$\lfalse$ની
એક નકલ કાઢી તેને
~$\Gamma \Sequent \Delta$ની
સાબિતીમાં ફેરવી શકાય.
(કસરત: $\pheight{\pi_1}$ પર
આગમનથી આ ચકાસો.)

હવે માનો કે કિસ્સો~A કે~B લાગુ પડતો નથી.
બાકી રહેલા કિસ્સાઓમાં
$\cheight{\pi} > 0$ છે, એટલે
ઓછામાં ઓછું એક આધારવિધાન
કોઈ અનુમાનનો નિષ્કર્ષ છે.

\paragraph{કિસ્સા C અને D: કટનું સ્થાનાંતર.}
કિસ્સા~C અને~Dની બધી શક્યતાઓમાં
એક સામાન્ય ઢબ છે. શરૂઆતમાં
~\CutCS{}થી પૂરી થતી સાબિતી લઈએ.
તેનું એક આધારવિધાન નિયમ~$R$થી મળે છે,
જેમાં કટ સૂત્ર~$!A$ મુખ્ય નથી
(બીજું કોઈ સૂત્ર~$!D$ મુખ્ય છે).
આ સાબિતીમાં \CutCS{} એ નિયમ~$R$
પછી આવે છે; યોજનાત્મક રીતે આ પ્રમાણે:
\[
\AxiomC{}
\RightLabel{$\pi_1$}
\Deduce$\Gamma \fCenter \Delta', !D, !A$
\AxiomC{}
\RightLabel{$\pi_3$}
\Deduce$!A, \Gamma, \Pi \fCenter \Delta', \Lambda$
\RightLabel{$R$}
\UnaryInf$!A, \Gamma \fCenter \Delta', !D$
\RightSubproofLabel{$\pi_2$}
\RightLabel{\CutCS}
\BinaryInf$\Gamma \fCenter \Delta', !D$
\DisplayProof
\]
આ માત્ર ત્યારેનું ચિત્ર છે જ્યારે~$R$
એક આધારવિધાનવાળો જમણો નિયમ હોય,
$A$ એ~$R$નું મુખ્ય સૂત્ર ન હોય,
અને જે સૂત્ર~$!D$ \emph{મુખ્ય}
છે તે~\CutCS{}ના જમણા આધારવિધાનના
ઉત્તરાંગમાં હોય. જો~$!D$ પૂર્વાંગમાં
હોય (એટલે~$R$ ડાબો નિયમ હોય)
અથવા~\CutCS{}ના ડાબા આધારવિધાનમાં
હોય, તો ચિત્ર અલગ હશે.
$\Pi$ અને~$\Lambda$નાં સૂત્રો
નિયમ~$R$નાં સક્રિય સૂત્રો દર્શાવે છે.

આ ક્રમ ઉલટાવી એવી સાબિતી લઈએ
જેમાં~$R$ એ~\CutCS{} પછી આવે છે:
\[
\AxiomC{}
\RightLabel{$\pi_1'$}
\Deduce$\Gamma, \Pi \fCenter \Delta', !D, \Lambda, !A$
\AxiomC{}
\RightLabel{$\pi_3'$}
\Deduce$!A, \Gamma, \Pi \fCenter \Delta', !D, \Lambda$
\RightLabel{\CutCS}
\BinaryInf$\Gamma, \Pi \fCenter \Delta', !D, \Lambda$
\RightLabel{$R$}
\UnaryInf$\Gamma \fCenter \Delta', !D, !D$
\DisplayProof
\]
\sourcecorrection{OLMD-201}{મૂળ યોજનાત્મક વૃક્ષના અંતિમ સિક્વન્ટમાં $\Gamma$ પછી વધારાનો અલ્પવિરામ છે; પૂર્વાંગ માત્ર $\Gamma$ છે.}
ઊંચાઈ જાળવતા શિથિલીકરણથી
અનુક્રમે $\pi_1$ અને $\pi_3$માંથી
$\pi_1'$ અને $\pi_3'$ મેળવીએ છીએ.

\CutCS{} અને~$R$નો ક્રમ બદલી દીધો હોવાથી
તેને ``કટને~$R$ની ઉપર ખસેડવો'' કહે છે.
એથી~\CutCS{}ના જમણા આધારવિધાનમાં
પૂરી થતી સાબિતીની ઊંચાઈ,
અને તેથી કટની ઊંચાઈ, ઘટે છે.
એટલે નવા~\CutCS{}માં પૂરી થતી
ઉપ-સાબિતીઓની ઊંચાઈ અનુક્રમે
$\pi_1$ અને $\pi_3$થી વધારે નથી;
જમણી બાજુનું એક અનુમાન કાઢ્યું હોવાથી
કટની ઊંચાઈ ઘટે છે.
હવે આગમન-પૂર્વધારણા લાગુ પડે છે
અને~\CutCS{} અનુમાન દૂર થાય છે.
છેલ્લે~$\RightR{\Contraction}$ની
સ્વીકાર્યતા વડે~$!D$ની વધારાની નકલ
દૂર કરીએ છીએ.

\paragraph{કિસ્સો C:}
$\pi_1$ અને $\pi_2$માંથી ઓછામાં ઓછી એક
એવા એક-આધારવિધાનવાળા નિયમ~$R$થી પૂરી થાય છે
જેમાં~$!A$ મુખ્ય નથી.
કુલ 18 કિસ્સા છે: \CutCS{}ના ડાબા
કે જમણા આધારવિધાનમાં~$!A$ મુખ્ય નથી
અને નવ એક-આધારવિધાનવાળા નિયમોમાંથી
(\LeftR{\lnot}, \RightR{\lnot},
\RightR{\lor}, \LeftR{\land},
\RightR{\lif} તથા ચાર પરિમાણક નિયમો)
કયો~$R$ છે, તેના પરથી.
માત્ર બે ઉદાહરણ જોઈએ:
~$\pi_2$ના છેલ્લા અનુમાનમાં
$!A$ મુખ્ય નથી અને તેનો અંત
$\RightR{\lif}$ કે
\LeftR{\lexists}થી થાય છે.
બીજા કિસ્સા સમાન છે અને
કસરત તરીકે છોડીએ છીએ.

માનો કે~$\pi$નો અંત આ રીતે થાય છે:
\[
\AxiomC{}
\RightLabel{$\pi_1$}
\Deduce$\Gamma \fCenter \Delta', !B \lif !C, !A$
\AxiomC{}
\RightLabel{$\pi_3$}
\Deduce$!A, !B, \Gamma \fCenter \Delta', !C$
\RightLabel{$\RightR{\lif}$}
\UnaryInf$!A, \Gamma \fCenter \Delta', !B \lif !C$
\RightSubproofLabel{$\pi_2$}
\RightLabel{\CutCS}
\BinaryInf$\Gamma \fCenter \Delta', !B \lif !C$
\DisplayProof
\]
અહીં $\Delta = \Delta', !B \lif !C$.
~$\pi_1$ પર ઊંચાઈ જાળવતું
શિથિલીકરણ
(\olref[seq][adm]{prop:weak-G3c-adm}) લગાવી
$!B, \Gamma \fCenter \Delta', !B \lif !C, !C, !A$ની
સાબિતી~$\pi_1'$ મેળવીએ
(ડાબે~$!B$ અને જમણે~$!C$ ઉમેરીએ).
એ જ રીતે~$\pi_3$ પર
\olref[seq][adm]{prop:weak-G3c-adm} લગાવી
$!A, !B, \Gamma \fCenter \Delta', !B \lif !C, !C$ની
સાબિતી~$\pi_3'$ મેળવીએ
(જમણે~$!B \lif !C$ ઉમેરીને).
આગમન-પૂર્વધારણા નીચેની સાબિતી પર લાગુ પડે છે:
\[
\AxiomC{}
\RightLabel{$\pi_1'$}
\Deduce$!B, \Gamma \fCenter \Delta', !B \lif !C, !C, !A$
\AxiomC{}
\RightLabel{$\pi_3'$}
\Deduce$!A, !B, \Gamma \fCenter \Delta', !B \lif !C, !C$
\RightLabel{\CutCS}
\BinaryInf$!B, \Gamma \fCenter \Delta', !B \lif !C, !C$
\DisplayProof
\]
કટ સૂત્ર~$!A$ છે. આ સાબિતીનો
કટ-ક્રમ~$\pi$ જેટલો જ છે,
પણ કટની ઊંચાઈ ઓછી છે.
આથી \Log{G3c}માં \CutCS{} વિના
$!B, \Gamma \fCenter \Delta', !B \lif !C, !C$ની
સાબિતી~$\pi_4$ મળે છે.
તેના પર~$\RightR{\lif}$ લગાવતાં
$\Gamma \fCenter \Delta', !B \lif !C, !B \lif !C$ની
સાબિતી મળે છે:
\[
\AxiomC{}
\RightLabel{$\pi_4$}
\Deduce$!B, \Gamma \fCenter \Delta', !B \lif !C, !C$
\RightLabel{$\RightR{\lif}$}
\UnaryInf$\Gamma \fCenter \Delta', !B \lif !C, !B \lif !C$
\DisplayProof
\]
$\Gamma \Sequent \Delta$,
એટલે~$\Gamma \fCenter \Delta', !B \lif !C$ની
સાબિતી~$\pi'$ મેળવવા
\olref[seq][inv]{prop:G3c-cont-adm}
(સંકોચનની સ્વીકાર્યતા) લગાવો.

હવે માનો કે~$\pi$નો અંત આ રીતે થાય છે:
\[
\AxiomC{}
\RightLabel{$\pi_1$}
\Deduce$\lexists[x][!B(x)], \Gamma' \fCenter \Delta, !A$
\AxiomC{}
\RightLabel{$\pi_3$}
\Deduce$!A, !B(c), \Gamma' \fCenter \Delta$
\RightLabel{$\LeftR{\lexists}$}
\UnaryInf$!A, \lexists[x][!B(x)], \Gamma' \fCenter \Delta$
\RightSubproofLabel{$\pi_2$}
\RightLabel{\CutCS}
\BinaryInf$\Gamma' \fCenter \Delta$
\DisplayProof
\]
ફરી આગમન-પૂર્વધારણા આના પર લગાવીએ:
\[
\AxiomC{}
\RightLabel{$\pi_1[!B(c) \Sequent {}]$}
\Deduce$\lexists[x][!B(x)], !B(c), \Gamma' \fCenter \Delta, !A$
\AxiomC{}
\LeftLabel{$\pi_3[\lexists[x][!B(x)] \Sequent {}]$}
\Deduce$!A, \lexists[x][!B(x)], !B(c), \Gamma' \fCenter \Delta$
\RightLabel{\CutCS}
\BinaryInf$\lexists[x][!B(x)], !B(c), \Gamma' \fCenter \Delta$
\RightLabel{\LeftR{\lexists}}
\UnaryInf$\lexists[x][!B(x)], \lexists[x][!B(x)], \Gamma' \fCenter \Delta$
\DisplayProof
\]
ધ્યાન આપો કે સ્વચલની શરત પૂરી થાય છે:
~$\pi_2$માં \LeftR{\lexists}ની સ્વચલ-શરત
ખાતરી આપે છે કે~$c$ એ
$\lexists[x][!B(x)]$, $\Gamma'$
કે~$\Delta$માં આવતો નથી.
અંતે \olref[seq][inv]{prop:G3c-cont-adm} લગાવો.


\paragraph{કિસ્સો D:}
$\pi_1$ અને $\pi_2$માંથી ઓછામાં ઓછી એક
એવા બે-આધારવિધાનવાળા નિયમ~$R$થી પૂરી થાય છે
જેમાં~$!A$ મુખ્ય નથી.
છ કિસ્સા છે. અહીં~$\pi_2$નું છેલ્લું
અનુમાન~$\RightR{\land}$ હોય
અને તેમાં~$!A$ મુખ્ય ન હોય તે કિસ્સો
જોઈએ; બીજા સમાન છે.

માનો કે~$\pi$નો અંત આ રીતે થાય છે:
\[
\AxiomC{}
\RightLabel{$\pi_1$}
\Deduce$\Gamma \fCenter \Delta', !B \land !C, !A$
\AxiomC{}
\RightLabel{$\pi_3$}
\Deduce$!A, \Gamma \fCenter \Delta', !B$
\AxiomC{}
\RightLabel{$\pi_4$}
\Deduce$!A, \Gamma \fCenter \Delta', !C$
\RightLabel{$\RightR{\land}$}
\BinaryInf$!A, \Gamma \fCenter \Delta', !B \land !C$
\RightSubproofLabel{$\pi_2$}
\RightLabel{\CutCS}
\BinaryInf$\Gamma \fCenter \Delta$
\DisplayProof
\]
આગમન-પૂર્વધારણા નીચેની સાબિતીઓ
પર લાગુ પડે છે:
\[
\AxiomC{}
\RightLabel{$\pi_1'$}
\Deduce$\Gamma \fCenter \Delta', !B \land !C, !B, !A$
\AxiomC{}
\RightLabel{$\pi_3'$}
\Deduce$!A, \Gamma \fCenter \Delta', !B \land !C, !B$
\RightLabel{\CutCS}
\BinaryInf$\Gamma \fCenter \Delta', !B \land !C, !B$
\DisplayProof
\]
અને આના પર પણ:
\[
\AxiomC{}
\RightLabel{$\pi_1''$}
\Deduce$\Gamma \fCenter \Delta', !B \land !C, !C, !A$
\AxiomC{}
\RightLabel{$\pi_4'$}
\Deduce$!A, \Gamma \fCenter \Delta', !B \land !C, !C$
\RightLabel{\CutCS}
\BinaryInf$\Gamma \fCenter \Delta', !B \land !C, !C$
\DisplayProof
\]
\sourcecorrection{OLMD-202}{મૂળ બીજા કટ-વૃક્ષમાં $\Delta$ પછી ફરી $!B\land !C$ ઉમેરે છે, જોકે $\Delta=\Delta',!B\land !C$; અહીં $\Delta'$ રાખવાથી બંને આધારવિધાન અને નિષ્કર્ષ મેળ ખાય છે.}
અહીં \Log{G3c}ની સાબિતીઓ
$\pi_1'$ અને $\pi_1''$ એ~$\pi_1$માંથી
ઊંચાઈ જાળવતા શિથિલીકરણથી
(\olref[seq][adm]{prop:weak-G3c-adm}) મળે છે
(એક વાર જમણે~$!B$ અને બીજી વાર~$!C$
ઉમેરીને). સાબિતીઓ
$\pi_3'$ અને~$\pi_4'$ પણ અનુક્રમે
$\pi_3$ અને~$\pi_4$માંથી
ઊંચાઈ જાળવતા શિથિલીકરણથી મળે છે
(જમણે~$!B \land !C$ ઉમેરીને).
આમ~$\pi_3$ અને~$\pi_4$ પર પણ
શિથિલીકરણ લાગુ પડે છે.

આગમન-પૂર્વધારણાથી \Log{G3c}માં
$\Gamma \fCenter \Delta', !B \land !C, !B$
અને $\Gamma \fCenter \Delta', !B \land !C, !C$ની
સાબિતીઓ~$\pi_5$ અને~$\pi_6$ મળે છે.
નિયમ~$\RightR{\land}$ વડે તેમને
$\Gamma \Sequent \Delta', !B \land !C, !B \land !C$ની
સાબિતીમાં જોડીએ:
\[
\AxiomC{}
\RightLabel{$\pi_5$}
\Deduce$\Gamma \fCenter \Delta', !B \land !C, !B$
\AxiomC{}
\RightLabel{$\pi_6$}
\Deduce$\Gamma \fCenter \Delta', !B \land !C, !C$
\RightLabel{\RightR{\land}}
\BinaryInf$\Gamma \fCenter \Delta', !B \land !C, !B \land !C$
\DisplayProof
\]
\sourcecorrection{OLMD-203}{મૂળ સંયુક્ત વૃક્ષના અંતિમ સિક્વન્ટમાં $\Delta$ પછી $!B\land !C$ની બે નકલો વધુ છે; $\Delta'$ સાથે બે નકલો મળે છે અને સંકોચનથી ઇચ્છિત $\Delta$ મળે છે.}
$\Gamma \Sequent \Delta$,
એટલે~$\Gamma \fCenter \Delta', !B \land !C$ની
સાબિતી માટે
\olref[seq][inv]{prop:G3c-cont-adm}
(સંકોચનની સ્વીકાર્યતા) લગાવો.

\paragraph{કિસ્સો E: કટ ઘટાડવા.}
છેલ્લો કિસ્સો ત્યારે છે જ્યારે
$\pi_1$ અને $\pi_2$ બંને એવા
અનુમાનોથી પૂરી થાય છે જેમાં~$!A$
મુખ્ય છે. ~ $!A$ના દરેક શક્ય
મુખ્ય કારક માટે એક,
એમ છ કિસ્સા છે.

દરેક કિસ્સામાં કટ સૂત્ર~$!A$
ધરાવતા \CutCS{}ને~$!A$ના
ઉપ-સૂત્રો પરના એક અથવા વધુ કટમાં
ફેરવીએ છીએ. આ એ અનુમાનોના સક્રિય
સૂત્રો છે જેમાં~$!A$ મુખ્ય
સૂત્ર છે. તેથી આ કિસ્સાઓને
કટ અનુમાનનું ``ઘટાડવું'' અથવા
``રૂપાંતર'' પણ કહેવાય છે.

\begin{enumerate}
\item જો $!A \ident \lnot !B$ હોય,
તો~$\pi$નો અંત આ રીતે થાય છે:
\[
\AxiomC{}
\RightLabel{$\pi_1'$}
\Deduce$!B, \Gamma \fCenter \Delta$
\RightLabel{\RightR{\lnot}}
\UnaryInf$\Gamma \fCenter \Delta, \lnot !B$
\LeftSubproofLabel{$\pi_1$}
\AxiomC{}
\RightLabel{$\pi_2'$}
\Deduce$\Gamma \fCenter \Delta, !B$
\RightLabel{\LeftR{\lnot}}
\UnaryInf$\lnot !B, \Gamma \fCenter \Delta$
\RightSubproofLabel{$\pi_2$}
\RightLabel{\CutCS}
\BinaryInf$\Gamma \fCenter \Delta$
\DisplayProof
\]
આગમન-પૂર્વધારણા વડે આમાંથી
\CutCS{} દૂર કરી શકાય છે:
\[
\AxiomC{}
\RightLabel{$\pi_2'$}
\Deduce$\Gamma \fCenter \Delta, !B$
\AxiomC{}
\RightLabel{$\pi_1'$}
\Deduce$!B, \Gamma \fCenter \Delta$
\RightLabel{\CutCS}
\BinaryInf$\Gamma \fCenter \Delta$
\DisplayProof
\]

\item
જો $!A \ident (!B \lor !C)$ હોય,
તો~$\pi$નો અંત આ રીતે થાય છે:
\[
\AxiomC{}
\RightLabel{$\pi_1'$}
\Deduce$\Gamma \fCenter \Delta, !B, !C$
\RightLabel{\RightR{\lor}}
\UnaryInf$\Gamma \fCenter \Delta, !B \lor !C$
\LeftSubproofLabel{$\pi_1$}
\AxiomC{}
\RightLabel{$\pi_2'$}
\Deduce$!B, \Gamma \fCenter \Delta$
\AxiomC{}
\RightLabel{$\pi_2''$}
\Deduce$!C, \Gamma \fCenter \Delta$
\RightLabel{\LeftR{\lor}}
\BinaryInf$!B \lor !C, \Gamma \fCenter \Delta$
\RightSubproofLabel{$\pi_2$}
\RightLabel{\CutCS}
\BinaryInf$\Gamma \fCenter \Delta$
\DisplayProof
\]
\CutCS{}થી પૂરી થતી આ સાબિતી લો:
\[
\AxiomC{}
\RightLabel{$\pi_1'$}
\Deduce$\Gamma \fCenter \Delta, !B, !C$
\AxiomC{}
\RightLabel{$\pi_2'''$}
\Deduce$!C, \Gamma \fCenter \Delta, !B$
\RightLabel{\CutCS}
\BinaryInf$\Gamma \fCenter \Delta, !B$
\DisplayProof
\]
અહીં~$\pi_2'''$ એ~$\pi_2''$માંથી
ઊંચાઈ જાળવતા શિથિલીકરણથી મળે છે.
તેનો કટ-ક્રમ $< \cutrank{\pi}$ છે,
કારણ કે $\depth{!C} < \depth{!B \lor !C}$.
\sourcecorrection{OLMD-204}{મૂળમાં અપરિભાષિત $\cutr{\pi}$ છે; આ વિભાગની વ્યાખ્યા પ્રમાણે યોગ્ય સંજ્ઞા $\cutrank{\pi}$ છે.}
તેથી આગમન-પૂર્વધારણાથી
\Log{G3c}માં
$\Gamma \fCenter \Delta, !B$ની
સાબિતી~$\pi_1''$ મળે છે.
હવે \CutCS{}થી પૂરી થતી આ સાબિતી લો:
\[
\AxiomC{}
\RightLabel{$\pi_1''$}
\Deduce$\Gamma \fCenter \Delta, !B$
\AxiomC{}
\RightLabel{$\pi_2'$}
\Deduce$!B, \Gamma \fCenter \Delta$
\RightLabel{\CutCS}
\BinaryInf$\Gamma \fCenter \Delta$
\DisplayProof
\]
તેનો કટ-ક્રમ
$= \depth{!B} < \depth{!B \lor !C}$ છે.
આથી આગમન-પૂર્વધારણા ફરી લાગુ પડે છે
અને $\Gamma \fCenter \Delta$ની
સાબિતી~$\pi'$ મળે છે.

\item $!A \ident (!B \land !C)$. કસરત.

\item જો $!A \ident (!B \lif !C)$ હોય,
તો~$\pi$નો અંત આ રીતે થાય છે:
\[
\AxiomC{}
\RightLabel{$\pi_1'$}
\Deduce$!B, \Gamma \fCenter \Delta, !C$
\RightLabel{\RightR{\lif}}
\UnaryInf$\Gamma \fCenter \Delta, !B \lif !C$
\LeftSubproofLabel{$\pi_1$}
\AxiomC{}
\RightLabel{$\pi_2'$}
\Deduce$\Gamma \fCenter \Delta, !B$
\AxiomC{}
\RightLabel{$\pi_2''$}
\Deduce$!C, \Gamma \fCenter \Delta$
\RightLabel{\LeftR{\lif}}
\BinaryInf$!B \lif !C, \Gamma \fCenter \Delta$
\RightSubproofLabel{$\pi_2$}
\RightLabel{\CutCS}
\BinaryInf$\Gamma \fCenter \Delta$
\DisplayProof
\]
\CutCS{}થી પૂરી થતી આ સાબિતી લો:
\[
\AxiomC{}
\RightLabel{$\pi_1'$}
\Deduce$!B, \Gamma \fCenter \Delta, !C$
\AxiomC{}
\RightLabel{$\pi_2''$}
\Deduce$!C, \Gamma \fCenter \Delta$
\RightLabel{\CutCS}
\BinaryInf$!B, \Gamma \fCenter \Delta$
\DisplayProof
\]
તેનો કટ-ક્રમ~$\pi$ કરતાં ઓછો છે.
આગમન-પૂર્વધારણા
$!B, \Gamma \fCenter \Delta$ની
સાબિતી~$\pi_1''$ આપે છે.
હવે \CutCS{}થી પૂરી થતી આ સાબિતી લો:
\[
\AxiomC{}
\RightLabel{$\pi_2'$}
\Deduce$\Gamma \fCenter \Delta, !B$
\AxiomC{}
\RightLabel{$\pi_1''$}
\Deduce$!B, \Gamma \fCenter \Delta$
\RightLabel{\CutCS}
\BinaryInf$\Gamma \fCenter \Delta$
\DisplayProof
\]
તેનો કટ-ક્રમ પણ~$\pi$ના કટ-ક્રમ કરતાં
ઓછો છે. તેથી આગમન-પૂર્વધારણા
$\Gamma \fCenter \Delta$ની
\Log{G3c}-સાબિતી~$\pi'$ આપે છે.
\item જો $!A \ident \lforall[x][!B(x)]$ હોય,
તો~$\pi$નો અંત આ રીતે થાય છે:
\[
\AxiomC{}
\RightLabel{$\pi_1'$}
\Deduce$\Gamma \fCenter \Delta, !B(c)$
\RightLabel{\RightR{\lforall}}
\UnaryInf$\Gamma \fCenter \Delta, \lforall[x][!B(x)]$
\LeftSubproofLabel{$\pi_1$}
\AxiomC{}
\RightLabel{$\pi_2'$}
\Deduce$!B(t), \lforall[x][!B(x)], \Gamma \fCenter \Delta$
\RightLabel{\LeftR{\lforall}}
\UnaryInf$\lforall[x][!B(x)], \Gamma \fCenter \Delta$
\RightSubproofLabel{$\pi_2$}
\RightLabel{\CutCS}
\BinaryInf$\Gamma \fCenter \Delta$
\DisplayProof
\]
આગમન-પૂર્વધારણા વડે આમાંથી
\CutCS{} દૂર કરી શકાય છે:
\[
\AxiomC{}
\RightLabel{$\pi_1''$}
\Deduce$!B(t), \Gamma \fCenter \Delta, \lforall[x][!B(x)]$
\AxiomC{}
\RightLabel{$\pi_2'$}
\Deduce$!B(t), \lforall[x][!B(x)], \Gamma \fCenter \Delta$
\RightLabel{\CutCS}
\BinaryInf$!B(t), \Gamma \fCenter \Delta$
\DisplayProof
\]
અહીં~$\pi_1''$ એ~$\pi_1$માંથી
(~$\pi_1'$માંથી નહીં!)
ઊંચાઈ જાળવતા શિથિલીકરણથી મળે છે.
(આ સાબિતીનો કટ-ક્રમ સમાન છે,
પણ કટની ઊંચાઈ ઓછી છે.)
તેથી $!B(t), \Gamma \fCenter \Delta$ની
સાબિતી~$\pi_2''$ મળે છે.

સાબિતી~$\pi_1'$નો અંતિમ સિક્વન્ટ
$\Gamma \Sequent \Delta, !B(c)$ છે,
જ્યાં~$c$ એ~$\RightR{\lforall}$
અનુમાનનો સ્વચલ છે.
તેથી~$c$ એ~$!B(x)$, $\Gamma$
કે~$\Delta$માં આવતો નથી.
સાબિતી~$\pi$ નિયમિત છે એમ ધારીએ,
એટલે~$c$ માત્ર ત્યાર પછી આવતા
$\RightR{\lforall}$ અનુમાનનો જ સ્વચલ છે
અને તેની ઉપર જ આવે છે;
અર્થાત્, તે માત્ર~$\pi_1'$માં આવે છે
અને~$\pi_1'$ની અંદરના કોઈ અનુમાનનો
સ્વચલ નથી. કારણ કે~$c$ એ~$\pi_2$માં
આવતો નથી, તે~$t$માં પણ આવતો નથી.
\olref[seq][qua]{lem:subst-regular} મુજબ
સાબિતી~$\Subst{\pi_1'}{t}{c}$ એ
$\Gamma \Sequent \Delta, !B(t)$ની
સાબિતી છે. હવે નીચેની
સાબિતી પર આગમન-પૂર્વધારણા લગાવીએ:
\[
\AxiomC{}
\RightLabel{$\Subst{\pi_1'}{t}{c}$}
\Deduce$\Gamma \fCenter \Delta, !B(t)$
\AxiomC{}
\RightLabel{$\pi_2''$}
\Deduce$!B(t), \Gamma \fCenter \Delta$
\RightLabel{\CutCS}
\BinaryInf$\Gamma \fCenter \Delta$
\DisplayProof
\]
(આગમન-પૂર્વધારણા લાગુ પડે છે,
કારણ કે કટ સૂત્ર~$!B(t)$ છે;
એટલે આ સાબિતીનો કટ-ક્રમ~$\pi$
કરતાં ઓછો છે.)
\item $!A \ident \lexists[x][!B(x)]$નો
કિસ્સો કસરત તરીકે છોડીએ છીએ.
\end{enumerate}
\end{proof}

\begin{prob}
\olref{lem:cut-adm-G3c}ની સાબિતીના કિસ્સા~Dનો
એ ઉપકિસ્સો ચકાસો જેમાં~$!A$ એ~$\pi_1$ના
છેલ્લા અનુમાનમાં મુખ્ય નથી અને તે અનુમાનનો
અંત~$\LeftR{\lif}$થી થાય છે.
(ધ્યાન આપો: આ કસરતનો એક ભાગ એ પણ છે કે
શું સાબિત કરવાનું છે તે સ્પષ્ટ કહો,
એટલે આ કિસ્સામાં~$\pi$નું સ્વરૂપ જણાવો.)
\end{prob}

\begin{prob}
\olref{lem:cut-adm-G3c}ની સાબિતીના કિસ્સા~Dનો
એ ઉપકિસ્સો ચકાસો જેમાં~$!A$ એ~$\pi_1$ના
છેલ્લા અનુમાનમાં મુખ્ય નથી અને તે અનુમાનનો
અંત~$\LeftR{\lforall}$થી થાય છે.
\end{prob}

\begin{prob}
\olref{lem:cut-adm-G3c}ની સાબિતીના કિસ્સા~Eનો
એ ઉપકિસ્સો ચકાસો જેમાં~$!A$ એ
$\pi_1$ અને $\pi_2$ બંનેના
છેલ્લા અનુમાનમાં મુખ્ય છે અને
$!A \ident (!B \land !C)$ છે.
\end{prob}

\begin{prob}
\olref{lem:cut-adm-G3c}ની સાબિતીના કિસ્સા~Eનો
એ ઉપકિસ્સો ચકાસો જેમાં~$!A$ એ
$\pi_1$ અને $\pi_2$ બંનેના
છેલ્લા અનુમાનમાં મુખ્ય છે અને
$!A \ident \lexists[x][!B(x)]$ છે.
\end{prob}

ધ્યાન આપો કે કિસ્સા~Eમાં આગમન-પૂર્વધારણા
જે~\CutCS{}થી પૂરી થતી સાબિતી પર
લગાડીએ તેની કટની ઊંચાઈ મૂળ
સાબિતીની કટની ઊંચાઈ કરતાં વધારે
હોઈ શકે. દાખલા તરીકે,
$!A \ident (!B \lif !C)$ હોય ત્યારે
~$\pi$ને બે \CutCS{} અનુમાનો ધરાવતી
સાબિતીમાં ફેરવ્યું:
\[
\AxiomC{}
\RightLabel{$\pi_2'$}
\Deduce$\Gamma \fCenter \Delta, !B$
\AxiomC{}
\RightLabel{$\pi_1'$}
\Deduce$!B, \Gamma \fCenter \Delta, !C$
\AxiomC{}
\RightLabel{$\pi_2''$}
\Deduce$!C, \Gamma \fCenter \Delta$
\RightLabel{\CutCS}
\BinaryInf$!B, \Gamma \fCenter \Delta$
\RightSubproofLabel{$\pi_1''$}
\RightLabel{\CutCS}
\BinaryInf$\Gamma \fCenter \Delta$
\DisplayProof
\]
\sourcecorrection{OLMD-205}{મૂળ અંતિમ વૃક્ષમાં બીજા કટ પછી પણ પૂર્વાંગમાં $!B$ રહી જાય છે; $!B$ પર કટ કરતાં અંતિમ સિક્વન્ટ $\Gamma\Sequent\Delta$ છે.}
ઉપરનો \CutCS{} એ~$\pi_1'$ અને $\pi_2''$
વચ્ચે છે; તેનો કટ-ક્રમ અને કટની
ઊંચાઈ બંને~$\pi$ કરતાં ઓછાં છે.
પણ આગમન-પૂર્વધારણાથી મળેલી
સાબિતી~$\pi_1''$ની કટ-ઊંચાઈ
હવે~$\pi$ કરતાં ઓછી હોવી જરૂરી નથી.
તેમ છતાં, નીચેના \CutCS{}નું
કટ સૂત્ર~$!B$ હોવાથી તેનો
કટ \emph{ક્રમ}~$\pi$ના કટ-ક્રમ કરતાં
ઓછો છે.

\olref{lem:cut-adm-G3c} બતાવે છે કે
સાબિતીમાંથી સૌથી ઉપરનું એક
\CutCS{} અનુમાન દૂર કરી શકાય.
સૌથી ઉપરના~\CutCS{}માં પૂરી થતી
ઉપ-સાબિતીઓ પર તેને વારંવાર
લગાવી કોઈ પણ સાબિતીમાંથી
ગમે તેટલાં \CutCS{} અનુમાનો
દૂર કરી શકાય તે બતાવીશું.

\begin{thm}\ollabel{thm:cut-elim-G3c-topmost}
$\Log{G3c} + \CutCS$ની કોઈ પણ
સાબિતી~$\pi$ને~\Log{G3c}ની
સાબિતીમાં ફેરવી શકાય છે.
\end{thm}

\begin{proof}
  ~$\pi$માંનાં \CutCS{} અનુમાનોની
  સંખ્યા~$n$ પર આગમન કરીએ.
  જો $n=0$, તો કશું કરવાનું નથી:
  $\pi$ પહેલેથી~\Log{G3c}ની
  સાબિતી છે. અન્યથા, સૌથી ઉપરના
  \CutCS{} અનુમાનમાં પૂરી થતી
  ઉપ-સાબિતી~$\pi'$ લો.
  તેમાં એક જ \CutCS{} નિયમ છે,
  અને તે છેલ્લું અનુમાન છે.
  \olref{lem:cut-adm-G3c} લગાવી
  તેને \CutCS{} વિનાની
  સાબિતી~$\pi''$માં ફેરવો,
  અને~$\pi'$ ઉપ-સાબિતીને
  $\pi''$થી બદલો.
  મળેલી સાબિતીમાં
  $n-1$ \CutCS{} અનુમાનો રહે છે,
  તેથી આગમન-પૂર્વધારણા લાગુ પડે છે.
\end{proof}
% END OLP-0656
\endgroup

\begingroup
\makeatletter\def\ol@sectioncs{section}\makeatother

% BEGIN OLP-0655
\olfileid{pt}{cut}{inv}

\olsection{સર્વોચ્ચ ક્રમના કટ દૂર કરવાં}
\phantomsection\label{guunit:OLP-0655}


\olref[seq][inv]{lem:G3c-invert}ની સાબિતી બતાવે છે કે જો
\Log{G3c} કોઈ તાર્કિક નિયમના નિષ્કર્ષને સાબિત કરે છે
(\RightR{\lexists} અને \LeftR{\lforall} સિવાય), તો તે
સાબિતીની ઊંચાઈ વધાર્યા વિના તે નિયમનાં દરેક આધારવિધાનને
પણ સાબિત કરે છે. આથી વધુ બતાવી શકાય છે: આ વાત
$\Log{G3c} + \CutCS$ માટે પણ સાચી છે.

અગાઉની જેમ, \CutCS{} અનુમાનનો \emph{કટ-ક્રમ} તેના
કટ સૂત્રની ઊંડાઈ ગણીએ છીએ.

\begin{lem}\ollabel{lem:inv-G3c-cut}
જો $R$ એ \RightR{\lexists} કે \LeftR{\lforall}
સિવાયનો તાર્કિક નિયમ હોય અને
$\Log{G3c} + \CutCS$ સાબિતી~$\pi$ વડે તેના
નિષ્કર્ષને સાબિત કરે છે, તો તે તેના દરેક આધારવિધાનને પણ
સાબિતી~$\pi'$ વડે સાબિત કરે છે (નિયમને બે
આધારવિધાનો હોય તો સાબિતીઓ~$\pi'$ અને~$\pi''$ વડે).
વધુમાં, (a)~$\pheight{\pi'}\le\pheight{\pi}$
(અને $\pheight{\pi''}\le\pheight{\pi}$) તથા
(b)~$\pi'$માં (અને $\pi''$માં) \CutCS{} અનુમાનોની
સંખ્યા તેમજ તેમના ક્રમો~$\pi$માં જેટલાં જ હોય છે.
\end{lem}

\begin{proof}
\cref{pt:seq:inv:lem:G3c-invert,pt:seq:inv:lem:invert-quant}ની
સાબિતીઓ તપાસવાથી આ મળે છે. ત્યાં
સાબિતી~$\pheight{\pi}$ પર આગમનથી હતી. આધારકિસ્સામાં
એક સ્વયંસિદ્ધને બીજા સ્વયંસિદ્ધથી બદલ્યું હતું.
આથી (a) અને~(b) સ્પષ્ટ રીતે સાચાં છે.
જો~$\pi$નો છેલ્લો નિયમ~$R$ હોય, તો તેની તત્કાળ
ઉપ-સાબિતીઓ જ જોઈતી સાબિતીઓ~$\pi_i$ છે,
અને (a) તથા~(b) બંને સાચાં છે. બીજા બધા
કિસ્સાઓમાં~$\pi$ની તત્કાળ ઉપ-સાબિતીઓ,
એટલે છેલ્લા અનુમાન~$R$નાં આધારવિધાનો તરફ દોરી જતી
ઉપસાબિતીઓ, પર આગમન-પૂર્વધારણા લગાવી
અને પરિણામો પર એ જ નિયમ~$R$ ફરી લગાવ્યો.
નવી સાબિતીની તત્કાળ ઉપ-સાબિતીઓ માટે (a) અને~(b)
સાચાં છે, તેથી સમગ્ર સાબિતી માટે પણ સાચાં છે.
છેલ્લું અનુમાન~\CutCS{} હોય તે કિસ્સો ચકાસવાનો રહે છે.
અહીં ફરી માત્ર એક ઉદાહરણ લઈએ: માનો કે
$\pi$ એ \LeftR{\land}ના નિષ્કર્ષ
$!B \land !C, \Gamma \Sequent \Delta$ની
સાબિતી છે અને તેનો અંત~\CutCS{}થી થાય છે:
\[
\AxiomC{}
\RightLabel{$\pi_1$}
\Deduce$!B \land !C, \Gamma \fCenter \Delta, !A$
\AxiomC{}
\RightLabel{$\pi_2$}
\Deduce$!A, !B \land !C, \Gamma \fCenter \Delta$
\RightLabel{\CutCS}
\BinaryInf$!B \land !C, \Gamma \fCenter \Delta$
\DisplayProof
\]
આગમન-પૂર્વધારણા $\pi_1$ અને~$\pi_2$ પર લાગુ પડે છે
(તેઓ પણ $\LeftR{\land}$ના નિષ્કર્ષનાં દૃષ્ટાંતોની
સાબિતીઓ છે). તે અનુક્રમે
$!B, !C, \Gamma \fCenter \Delta, !A$ અને
$!A, !B, !C, \Gamma \fCenter \Delta$ની
સાબિતીઓ~$\pi_1'$ અને~$\pi_2'$ આપે છે.
તેમને \CutCS{} વડે જોડતાં~$\pi'$ મળે છે:
\[
\AxiomC{}
\RightLabel{$\pi_1'$}
\Deduce$!B, !C, \Gamma \fCenter \Delta, !A$
\AxiomC{}
\RightLabel{$\pi_2'$}
\Deduce$!A, !B, !C, \Gamma \fCenter \Delta$
\RightLabel{\CutCS}
\BinaryInf$!B, !C, \Gamma \fCenter \Delta$
\DisplayProof
\]
અહીં (a) અને~(b) સ્પષ્ટ રીતે સચવાય છે.
\end{proof}

ધ્યાન આપો કે \CutCS{}ને બદલે \Cut{} વાપર્યું હોત
તો આ રીત ચાલત નહીં. ત્યારે છેલ્લી સાબિતીના
નિષ્કર્ષના પૂર્વાંગમાં $!B, !C, \Gamma$ની બે નકલો
અને ઉત્તરાંગમાં $\Delta$ની બે નકલો હોત.
સંકોચનની સ્વીકાર્યતા લેવી પડત, પરંતુ
\olref[seq][inv]{prop:G3c-cont-adm}ની સાબિતી
ઉલટાવી શકવાની ઉપપ્રમેયનો ઉપયોગ કરે છે.
એટલે દલીલ ચક્રીય બની જાય.

\cref{pt:cut:inv:lem:inv-G3c-cut} વડે કટ-નિરાકરણની
વૈકલ્પિક સાબિતી આપી શકાય. તેમાં \CutCS{}
અનુમાનોના સૌથી ઉપરના દૃષ્ટાંતોને એક પછી એક
દૂર કરવાને બદલે, સર્વોચ્ચ ક્રમનાં \CutCS{}
અનુમાનોને દૂર કરીએ છીએ. એટલે,
$\Log{G3c} + \CutCS$માં સાબિતી~$\pi$ લઈને
તેમાં ક્યાંય પણ રહેલો કટ દૂર કરીએ
(કદાચ તેના સ્થાને નીચા ક્રમના કટ આવે)
અને કટ ન રહે ત્યાં સુધી આ પ્રક્રિયા ચાલુ રાખીએ.
માત્ર એટલી શરત છે કે જે કટને લઈએ તેની ઉપર
સમાન કે વધુ ક્રમનો કોઈ કટ ન હોય.

\begin{lem}\ollabel{lem:max-cut-red-G3c}
જો $\pi$ એ $n$ ક્રમના~\CutCS{} અનુમાનથી પૂરી થતી
$\Log{G3c}+\CutCS$ની સાબિતી હોય, અને તે ઉપરાંત
માત્ર~$\Log{G3c}$ના નિયમો તથા $<n$ ક્રમનાં
\CutCS{} અનુમાનો વાપરતી હોય, તો એ જ અંતિમ
સિક્વન્ટની એવી સાબિતી~$\pi'$ પણ
$\Log{G3c} + \CutCS$માં છે જેમાં દરેક
\CutCS{} અનુમાનનો ક્રમ~$<n$ છે.
\end{lem}

\begin{proof}
સાબિતી~$\pi$નો અંત આ રીતે થાય છે:
\[
\AxiomC{}
\RightLabel{$\pi_1$}
\Deduce$\Gamma \fCenter \Delta, !A$
\AxiomC{}
\RightLabel{$\pi_2$}
\Deduce$!A, \Gamma \fCenter \Delta$
\RightLabel{\CutCS}
\BinaryInf$\Gamma \fCenter \Delta$
\DisplayProof
\]
હવે~$!A$ના સ્વરૂપ પ્રમાણે કિસ્સા જુદા પાડીએ.

માનો કે~$!A$ પરમાણ્વીય છે. શરત મુજબ
~$\pi_1$ અને $\pi_2$માંનાં બધાં \CutCS{} અનુમાનોના
ક્રમ $<\depth{!A} = 0$ હોવા જોઈએ;
આથી~$\pi_1$ કે~$\pi_2$માં \CutCS{} અનુમાન હોઈ જ શકે નહીં.
$\pheight{\pi_2}$ પર આગમનથી બતાવીએ કે
$\Gamma \Sequent \Delta$ની કટ-મુક્ત સાબિતી છે.
જો $\pheight{\pi_2} = 0$, તો
$!A,\Gamma \Sequent \Delta$ સ્વયંસિદ્ધ છે.
જો~$!A$ મુખ્ય ન હોય, તો કોઈ પરમાણ્વીય
$!C$ એ~$\Gamma$ તથા $\Delta$ બંનેનું
ઘટક છે, એટલે
$\Gamma \Sequent \Delta$ પોતે સ્વયંસિદ્ધ છે.
જો~$!A$ મુખ્ય હોય, તો
$\Delta = \Delta', !A$.
\sourcecorrection{OLMD-196}{મૂળમાં $\Delta = \Delta, !A$ છે; અહીં બાકીના ઉત્તરાંગ માટે $\Delta'$ જરૂરી છે.}
આ કિસ્સામાં~$\pi_1$નો અંત
$\Gamma \Sequent \Delta', !A, !A$થી થાય છે.
\olref[seq][inv]{prop:G3c-cont-adm}
(સંકોચનની સ્વીકાર્યતા) લગાવીને
$\Gamma \Sequent \Delta', !A$ની
કટ-મુક્ત સાબિતી મળે છે.
જો $\pheight{\pi_2} > 0$, તો તેનો અંત
કોઈ નિયમ~$R$થી થાય છે. તે નિયમનું એક
આધારવિધાન લો: તેનું સ્વરૂપ
$!A, \Gamma',\Pi \Sequent \Delta',\Lambda$
છે, જ્યાં $\Pi$ અને $\Lambda$ એ~$R$નાં
સક્રિય સૂત્રો છે.
\olref[seq][adm]{prop:weak-G3c-adm}ને
~$\pi_1$ અને~$\pi_2$નાં યોગ્ય તત્કાળ
ઉપ-સાબિતીઓ પર લગાવતાં
$\Gamma,\Gamma',\Pi \Sequent
\Delta,\Delta',\Lambda,!A$ અને
$!A,\Gamma,\Gamma',\Pi \Sequent
\Delta,\Delta',\Lambda$ની સાબિતીઓ મળે છે.
\sourcecorrection{OLMD-197}{મૂળ લખાણ $\pi_2$ પર સીધું શિથિલીકરણ કહે છે; નિયમના આધારવિધાન સુધી પહોંચવા તેની યોગ્ય તત્કાળ ઉપ-સાબિતી જોઈએ.}
તેના પર આગમન-પૂર્વધારણા લાગુ કરતાં
~$R$ના દરેક આધારવિધાન માટે
$\Gamma,\Gamma',\Pi \Sequent
\Delta,\Delta',\Lambda$ની સાબિતી મળે છે.
પછી~$R$ લગાવતાં
$\Gamma,\Gamma \Sequent \Delta,\Delta$ મળે છે;
\olref[seq][inv]{prop:G3c-cont-adm} વડે
$\Gamma \Sequent \Delta$ની કટ-મુક્ત
સાબિતી મળે છે.

જો~$!A$ પરમાણ્વીય ન હોય, તો તેના
મુખ્ય કારક પ્રમાણે~$!A$ના કિસ્સા જુદા પાડીએ.
દરેક કિસ્સામાં \olref{lem:inv-G3c-cut}
લગાવીને, $!A$ મુખ્ય સૂત્ર હોય તેવા
ડાબા અને જમણા નિયમોના આધારવિધાનોની
સાબિતીઓ મેળવીએ. પછી
\olref[top]{lem:cut-adm-G3c}ની સાબિતીના
કિસ્સા~E મુજબ આગળ વધીએ.
\begin{enumerate}
  \item $!A \ident \lnot !B$. સાબિતીઓ $\pi_1$ અને $\pi_2$નો
  અંત $\Gamma \Sequent \Delta, \lnot !B$ અને
  $\lnot !B, \Gamma \Sequent \Delta$થી થાય છે.
  \olref{lem:inv-G3c-cut} મુજબ
  $!B, \Gamma \Sequent \Delta$ અને
  $\Gamma \Sequent \Delta, !B$ની
  સાબિતીઓ $\pi_1'$ અને $\pi_2'$ છે.
  સાબિતી~$\pi'$ આ છે:
  \[
  \AxiomC{}
  \RightLabel{$\pi_2'$}
  \Deduce$\Gamma \fCenter \Delta, !B$
  \AxiomC{}
  \RightLabel{$\pi_1'$}
  \Deduce$!B, \Gamma \fCenter \Delta$
  \RightLabel{\CutCS}
  \BinaryInf$\Gamma \fCenter \Delta$
  \DisplayProof
  \]
  \item $!A \ident (!B \lor !C)$. સાબિતીઓ $\pi_1$ અને $\pi_2$નો
  અંત $\Gamma \Sequent \Delta, !B \lor !C$ અને
  $!B \lor !C, \Gamma \Sequent \Delta$થી થાય છે.
  \olref{lem:inv-G3c-cut} વડે
  ~$\pi_1$ પર \RightR{\lor}નું ઉલટાવવું
  લગાવતાં $\Gamma \Sequent \Delta, !B, !C$ની
  સાબિતી~$\pi_1'$ મળે છે.
  \olref{lem:inv-G3c-cut} વડે~$\pi_2$ પર
  \LeftR{\lor}નું ઉલટાવવું લગાવતાં
  $!B, \Gamma \Sequent \Delta$ની
  સાબિતી~$\pi_2'$ અને
  $!C, \Gamma \Sequent \Delta$ની
  સાબિતી~$\pi_2''$ મળે છે.
  $\pi_2''$ પર
  \olref[seq][adm]{prop:weak-G3c-adm}
  લગાવતાં $!C, \Gamma \Sequent \Delta, !B$ની
  સાબિતી~$\pi_2'''$ પણ મળે છે.
  સાબિતી~$\pi'$ આ છે:
  \[
  \AxiomC{}
  \RightLabel{$\pi_1'$}
  \Deduce$\Gamma \fCenter \Delta, !B, !C$
  \AxiomC{}
  \RightLabel{$\pi_2'''$}
  \Deduce$!C, \Gamma \fCenter \Delta, !B$
  \RightLabel{\CutCS}
  \BinaryInf$\Gamma \fCenter \Delta, !B$
  \AxiomC{}
  \RightLabel{$\pi_2'$}
  \Deduce$!B, \Gamma \fCenter \Delta$
  \RightLabel{\CutCS}
  \BinaryInf$\Gamma \fCenter \Delta$
  \DisplayProof
  \]
  \item $!A \ident (!B \land !C)$. કસરત.
  \item $!A \ident (!B \lif !C)$. સાબિતીઓ $\pi_1$ અને $\pi_2$નો
  અંત $\Gamma \Sequent \Delta, !B \lif !C$ અને
  $!B \lif !C, \Gamma \Sequent \Delta$થી થાય છે.
  \olref{lem:inv-G3c-cut} વડે
  ~$\pi_1$ પર \RightR{\lif}નું ઉલટાવવું
  લગાવતાં $!B, \Gamma \Sequent \Delta, !C$ની
  સાબિતી~$\pi_1'$ મળે છે.
  \olref{lem:inv-G3c-cut} વડે~$\pi_2$ પર
  \LeftR{\lif}નું ઉલટાવવું લગાવતાં
  $\Gamma \Sequent \Delta, !B$ની
  સાબિતી~$\pi_2'$ અને
  $!C, \Gamma \Sequent \Delta$ની
  સાબિતી~$\pi_2''$ મળે છે.
  $\pi_2''$ પર
  \olref[seq][adm]{prop:weak-G3c-adm}
  લગાવતાં $!C, !B, \Gamma \Sequent \Delta$ની
  સાબિતી~$\pi_2'''$ પણ મળે છે.
  સાબિતી~$\pi'$ આ છે:
  \[
  \AxiomC{}
  \RightLabel{$\pi_2'$}
  \Deduce$\Gamma \fCenter \Delta, !B$
  \AxiomC{}
  \RightLabel{$\pi_1'$}
  \Deduce$!B, \Gamma \fCenter \Delta, !C$
  \AxiomC{}
  \RightLabel{$\pi_2'''$}
  \Deduce$!C, !B, \Gamma \fCenter \Delta$
  \RightLabel{\CutCS}
  \BinaryInf$!B, \Gamma \fCenter \Delta$
  \RightLabel{\CutCS}
  \BinaryInf$\Gamma \fCenter \Delta$
  \DisplayProof
  \]
  \item $!A \ident \lforall[x][!B(x)]$. સાબિતીઓ $\pi_1$ અને
  $\pi_2$નો અંત $\Gamma \Sequent \Delta, \lforall[x][!B(x)]$
  અને $\lforall[x][!B(x)], \Gamma \Sequent \Delta$થી થાય છે.
  \olref{lem:inv-G3c-cut} મુજબ
  $\Gamma \Sequent \Delta, !B(c)$ની
  સાબિતી~$\pi_1'$ છે.

  હવે સાબિતી~$\pi_2$ લો. તેનો અંત
  $\lforall[x][!B(x)], \Gamma \Sequent \Delta$થી
  થાય છે. $\pi_2$ના અંતિમ સિક્વન્ટથી ઉપરની તરફ
  જતાં, સિક્વન્ટની ડાબી બાજુ રહેલા
  $\lforall[x][!B(x)]$ના એવા દરેક દૃષ્ટાંતને
  ચિહ્નિત કરો જે $\pi_2$ના અંતિમ સિક્વન્ટમાં
  એ જ $\lforall[x][!B(x)]$ના દૃષ્ટાંત સુધી પહોંચે છે. એટલે:
  (a)~$\pi_2$ના અંતિમ સિક્વન્ટમાં
  $\lforall[x][!B(x)]$ ચિહ્નિત કરો.
  (b)~નિયમના નિષ્કર્ષના સંદર્ભમાં
  $\lforall[x][!B(x)]$
  ચિહ્નિત હોય તો તેના આધારવિધાનમાં અનુરૂપ
  દૃષ્ટાંત ચિહ્નિત કરો. (c)~\LeftR{\lforall}
  નિયમના નિષ્કર્ષમાં $\lforall[x][!B(x)]$
  મુખ્ય અને ચિહ્નિત હોય, તો આધારવિધાનમાં
  $\lforall[x][!B(x)]$ તથા $!B(t)$ના
  અનુરૂપ સક્રિય દૃષ્ટાંતો ચિહ્નિત કરો.
  
  હવે $\lforall[x][!B(x)]$ના આ બધા
  ચિહ્નિત દૃષ્ટાંતો લો. \LeftR{\lforall}
  સિવાયના કોઈ નિયમમાં તેઓ સક્રિય નથી
  (ચિહ્નિત થનારા સક્રિય સૂત્રો
  માત્ર \LeftR{\lforall}નાં છે).
  $\pi_2$માંથી $\lforall[x][!B(x)]$ના
  બધા ચિહ્નિત દૃષ્ટાંતો કાઢી નાખો.
  જો મળેલું લખાણ સાચી સાબિતી હોય,
  તો $\lforall[x][!B(x)]$ના ચિહ્નિત
  દૃષ્ટાંતો કદી સક્રિય નહોતા;
  તેઓ બધા સીધા સ્વયંસિદ્ધોમાંથી આવ્યા હતા.
  સાબિતીનો અંત~$\Gamma \Sequent \Delta$થી
  થાય છે, તેથી~$\pi$ને તેનાથી બદલી શકાય.
  આમ સર્વોચ્ચ ક્રમનો એક કટ દૂર થયો.

  અન્યથા, મળેલું લખાણ સાચી સાબિતી નથી,
  કારણ કે કેટલીક \LeftR{\lforall} નિષ્પત્તિઓમાં
  સક્રિય રહેલાં $\lforall[x][!B(x)]$
  સૂત્રો કાઢી નાખ્યાં છે.
  તેથી તે નિષ્પત્તિઓ આ સ્વરૂપની
  ``નિષ્પત્તિઓ'' બની છે:
  \[
  \AxiomC{}
  \RightLabel{$\pi_t$}
  \Deduce$!B(t), \Pi \fCenter \Lambda$
  \RightLabel{\LeftR{\lforall}}
  \UnaryInf$\Pi \fCenter \Lambda$
  \DisplayProof
  \]
  આવા સૌથી ઉપરના દરેક ``અનુમાન''ને
  આથી બદલો:
  \[
  \AxiomC{}
  \RightLabel{$\Subst{\pi_1'}{t}{c}[\Pi \Sequent \Lambda]$}
  \Deduce$\Gamma, \Pi \fCenter \Delta, \Lambda, !B(t)$
  \AxiomC{}
  \RightLabel{$\pi_t[\Gamma \Sequent \Delta]$}
  \Deduce$!B(t), \Gamma, \Pi \fCenter \Delta, \Lambda$
  \RightLabel{\CutCS}
  \BinaryInf$\Gamma, \Pi \fCenter \Delta, \Lambda$
  \DisplayProof
  \]
  અને તેની નીચેના દરેક સિક્વન્ટની ડાબે
  $\Gamma$ તથા જમણે $\Delta$ ઉમેરો.
  આધારવિધાનમાં ચિહ્નિત $!B(t)$ ધરાવતાં
  બધા \LeftR{\lforall} ``અનુમાનો''ને
  કોઈ~$!B(t)$ પરના \CutCS{} અનુમાનથી
  બદલી ન દેવાય ત્યાં સુધી આ દોહરાવો.
  મળેલી સાબિતીનો અંતિમ સિક્વન્ટ
  (અને હવે તે સાચી સાબિતી છે)
  $\Gamma, \dots, \Gamma \Sequent
  \Delta, \dots, \Delta$ સ્વરૂપનો છે.
  સંકોચનની સ્વીકાર્યતા લગાવી તેને
  $\Gamma \Sequent \Delta$ની સાબિતીમાં
  ફેરવો અને~$\pi$ને તેનાથી બદલો.
  $\lforall[x][!B(x)]$ પરના એક કટને
  $!B(t)$ પરના (કદાચ ઘણા) કટથી બદલ્યો છે,
  પણ તે બધા નીચા ક્રમના છે.
  $\pi_1$ કે $\pi_2$માં પહેલેથી રહેલા
  બધા કટ પણ રહે છે, અને તે પણ
  નીચા ક્રમના જ છે.
\end{enumerate}
\end{proof}

\begin{prob}
\olref{lem:max-cut-red-G3c}ની સાબિતીમાં
$!A \ident (!B \land !C)$નો કિસ્સો ચકાસો.
\sourcecorrection{OLMD-198}{મૂળ કસરત ઉલટાવવાના ઉપપ્રમેયનો સંદર્ભ આપે છે; અહીં ચર્ચાતો કિસ્સો મહત્તમ-ક્રમના કટ ઘટાડવાના ઉપપ્રમેયનો છે.}
અર્થાત્, જો અંતમાં
\[
\AxiomC{}
\RightLabel{$\pi_1$}
\Deduce$\Gamma \fCenter \Delta, !B \land !C$
\AxiomC{}
\RightLabel{$\pi_2$}
\Deduce$!B \land !C, \Gamma \fCenter \Delta$
\RightLabel{\CutCS}
\BinaryInf$\Gamma \fCenter \Delta$
\DisplayProof
\]
ધરાવતી સાબિતી~$\pi$ હોય અને
$\pi_1$ તથા $\pi_2$માં
$<\depth{!B \land !C}$ ક્રમના કટ જ હોય,
તો $\Gamma \Sequent \Delta$ની એવી
સાબિતી~$\pi'$ છે જેમાં પણ
$<\depth{!B \land !C}$ ક્રમના કટ જ છે,
એ બતાવો.
\end{prob}



\olref{lem:max-cut-red-G3c} બતાવે છે કે
સાબિતીમાં સર્વોચ્ચ ક્રમના \CutCS{}
અનુમાનને નીચા ક્રમનાં \CutCS{} અનુમાનોથી
બદલી શકાય.
\sourcecorrection{OLMD-199}{મૂળ અનુચ્છેદ અગાઉના બીજા ઉપપ્રમેય તરફ નિર્દેશ કરે છે; અહીં જરૂરી પરિણામ આ વિભાગનું મહત્તમ-ક્રમના કટ ઘટાડવાનું ઉપપ્રમેય છે.}
આ ઉપપ્રમેય ફરી લગાવી,
સર્વોચ્ચ ક્રમના~\CutCS{} પર સમાપ્ત
થતી સૌથી ઉપરની ઉપ-સાબિતીઓ પર
ક્રમશઃ પ્રક્રિયા કરીને, કોઈ પણ
સાબિતીમાંથી ગમે તેટલા \CutCS{}
અનુમાનો દૂર કરી શકાય છે.

\begin{thm}\ollabel{thm:cut-elim-G3c-largest}
$\Log{G3c} + \CutCS$ની કોઈ પણ
સાબિતી~$\pi$ને~\Log{G3c}ની
સાબિતીમાં ફેરવી શકાય છે.
\end{thm}

\begin{proof}
  પહેલાં બતાવીએ કે~$\pi$માંના સર્વોચ્ચ
  ક્રમના બધા કટ દૂર કરી શકાય છે.
  માનો કે~$d$ એ~$\pi$માંના કટનો
  મહત્તમ ક્રમ છે અને~$n$ એ~$d$
  ક્રમના કટની સંખ્યા છે. સાબિતી
  ~$n$ પર આગમનથી છે.
  જો $n=0$, તો કશું કરવાનું નથી.
  જો $n>0$, તો~$d$ ક્રમનો
  સૌથી ઉપરનો કટ પસંદ કરો અને
  તેમાં પૂરી થતી ઉપ-સાબિતીને
  $\pi_1$ કહો.
  \olref{lem:max-cut-red-G3c} મુજબ
  એ જ અંતિમ સિક્વન્ટની એવી
  સાબિતી~$\pi_1'$ છે જેમાં
  બધા કટનો ક્રમ $<d$ છે.
  $\pi$માં $\pi_1$ને $\pi_1'$થી
  બદલો. મળેલી સાબિતીમાં
  $d$ ક્રમના $n-1$ કટ રહે છે,
  તેથી આગમન-પૂર્વધારણા લાગુ પડે છે.

  હવે~$d$ પર આગમનથી પરિણામ મળે છે.
  જો $d=0$, તો બધા કટ પરમાણ્વીય છે
  અને અગાઉના પરિણામથી તે બધા
  દૂર કરી શકાય છે. જો $d>0$,
  તો અગાઉનું પરિણામ એવી સાબિતી આપે
  છે જેમાં $d$ ક્રમના બધા કટ
  $<d$ ક્રમના કટથી બદલાયા છે.
  કારણ કે~$d$ એ~$\pi$માંના કટનો
  મહત્તમ ક્રમ હતો, મળેલી સાબિતીમાં
  કટનો મહત્તમ ક્રમ $<d$ છે,
  અને આગમન-પૂર્વધારણા લાગુ પડે છે.
\end{proof}
% END OLP-0655
\endgroup

\begingroup
\makeatletter\def\ol@sectioncs{section}\makeatother

% BEGIN OLP-0661
\olfileid{pt}{cut}{mid}

\olsection{મધ્ય-સિક્વન્ટ પ્રમેય અને હર્બ્રાન્ડનું પ્રમેય}
\phantomsection\label{guunit:OLP-0661}


કટ-નિરાકરણ પ્રમેયનું મહત્ત્વનું
પરિણામ મધ્ય-સિક્વન્ટ પ્રમેય છે.
તે કહે છે કે જો
$\Gamma \Sequent \Delta$ની
એવી સાબિતી~$\pi$ હોય
જેમાં દરેક સૂત્ર પરિમાણક-અગ્ર
સ્વરૂપનું હોય, તો એવી
સાબિતી~$\pi'$ છે જેમાં
$S \ident \Pi \Sequent \Lambda$
સિક્વન્ટ (\emph{મધ્ય-સિક્વન્ટ})
હોય અને~$\pi'$માં~$S$ની
ઉપરનાં બધાં અનુમાનો વિધાનાત્મક,
જ્યારે~$S$ની નીચેના બધાં અનુમાનો
પરિમાણકનાં હોય.
(સૂત્ર~$!A$ ત્યારે
\emph{પરિમાણક-અગ્ર સ્વરૂપનું}
હોય જ્યારે તેનું સ્વરૂપ
\[\mathsf{Q}_1x_1\dots\mathsf{Q}_nx_n\,!B(x_1,\dots,x_n)\]
હોય, જ્યાં દરેક
$\mathsf{Q}_i$ એ~$\lforall$
અથવા~$\lexists$ છે.)
મધ્ય-સિક્વન્ટ પ્રમેયનું બીજું
મહત્ત્વનું પરિણામ હર્બ્રાન્ડનું
પ્રમેય છે. તેનો સામાન્ય કિસ્સો
વર્ણવવો થોડો મુશ્કેલ છે,
પણ વ્યાપક રીતે તે કહે છે કે
પ્રથમ-ક્રમ તર્કના દરેક
સાબિતીક્ષમ વાક્ય~$!A$
માટે એવી વિધાનાત્મક
પુનરુક્તિ~$!A'$ છે જેમાંથી
$!A$ સાબિત કરી શકાય.
અહીં~$!B$ પરિમાણક-રહિત હોય ત્યારે
$\lexists[x][!B(x)]$ સ્વરૂપનાં
વાક્યો માટેનું રૂપ આપીએ:

\begin{thm}[હર્બ્રાન્ડનું પ્રમેય]
માનો કે~$!B(x)$ પરિમાણક-રહિત છે
અને $\Proves
\lexists[x][!B(x)]$.
તો એવાં પદો~$t_1$, \dots, $t_n$
છે કે
$\Proves !B(t_1) \lor \dots \lor !B(t_n)$.
\end{thm}

સૂત્ર
$!B(t_1) \lor \dots \lor !B(t_n)$ને
$\lexists[x][!B(x)]$નો
\emph{હર્બ્રાન્ડ વિકલ્પ} કહે છે.
કારણ કે~$!B$ પરિમાણક-રહિત છે,
આ વિકલ્પ પણ પરિમાણક-રહિત છે.
આથી તે સાબિતીક્ષમ હોય તો
કટ-મુક્ત સાબિતીથી સાબિતીક્ષમ છે,
અને તેથી માત્ર વિધાનાત્મક
અનુમાનો વાપરતી સાબિતી
પણ છે. એટલે તે પુનરુક્તિ છે.

હર્બ્રાન્ડના પ્રમેયનો મુશ્કેલ ભાગ
એ છે કે દરેક સાબિતીક્ષમ
સૂત્ર~$\lexists[x][!B(x)]$
માટે હર્બ્રાન્ડ વિકલ્પ છે
તે બતાવવું. વ્યસ્ત દિશા,
એટલે~$\lexists[x][!B(x)]$નો
હર્બ્રાન્ડ વિકલ્પ હોય તો
તે સાબિત થઈ શકે, સહેલી છે.
ખરેખર હર્બ્રાન્ડ વિકલ્પમાંથી
$\lexists[x][!B(x)]$ ફલિત થાય છે.
દાખલા તરીકે~$n=2$ માટે
~$\Log{G3c}$માં આ
સાબિતી છે:
\[
\Axiom$!B(t_1) \fCenter \lexists[x][!B(x)], !B(t_1), !B(t_2)$
\Axiom$!B(t_2) \fCenter \lexists[x][!B(x)], !B(t_1), !B(t_2)$
\RightLabel{\LeftR{\lor}}
\BinaryInf$!B(t_1) \lor !B(t_2) \fCenter \lexists[x][!B(x)], !B(t_1), !B(t_2)$
\RightLabel{\RightR{\lexists}}
\UnaryInf$!B(t_1) \lor !B(t_2) \fCenter \lexists[x][!B(x)], !B(t_2)$
\RightLabel{\RightR{\lexists}}
\UnaryInf$!B(t_1) \lor !B(t_2) \fCenter \lexists[x][!B(x)]$
\DisplayProof
\]

$\Sequent \lexists[x][!B(x)]$ની
સાબિતી~$\pi$માંથી
હર્બ્રાન્ડ વિકલ્પ કાઢવા,
પહેલાં એવો કિસ્સો લો જેમાં
કટ-મુક્ત સાબિતી~$\pi$નાં
બધાં $\RightR{\lexists}$
અનુમાનો અંતે આવે છે:
\[
\AxiomC{}
\RightLabel{$\pi_1$}
\Deduce$ \fCenter \lexists[x][!B(x)], !B(t_1), \dots, !B(t_n)$
\RightLabel{\RightR{\lexists}}
\UnaryInf$ \fCenter \lexists[x][!B(x)], !B(t_2), \dots, !B(t_n)$
\RightLabel{$(n-1) \times \RightR{\lexists}$}
\Deduce$ \fCenter \lexists[x][!B(x)]$
\DisplayProof
\]
જો~$\pi$માં
$\lexists[x][!B(x)]$
સક્રિય હોય એવાં
$\RightR{\lexists}$
અનુમાનો આ બધાં જ હોય,
તો~$\pi_1$માં બીજું કોઈ
પરિમાણક-અનુમાન નથી.
\sourcecorrection{OLMD-221}{મૂળમાં અહીં અપરિભાષિત $\pi_1'$ લખાયું છે; ઉપરના વૃક્ષની ઉપરની ઉપસાબિતી $\pi_1$ છે.}
કારણ કે~$!B(x)$માં બીજાં
પરિમાણકો નથી અને અંતિમ
સિક્વન્ટના પૂર્વાંગમાં
કોઈ સૂત્રો નથી.
અર્થાત્,
$ \fCenter \lexists[x][!B(x)],\allowbreak
!B(t_1),\allowbreak \dots,\allowbreak !B(t_n)$
એ~$\pi$નો મધ્ય-સિક્વન્ટ છે.
વધુમાં, તેની ઉપર પરિમાણક-અનુમાનો
ન હોવાથી ઉપરના બધા સિક્વન્ટમાં
$\lexists[x][!B(x)]$ હોય છે,
પણ તે સક્રિય કે મુખ્ય નથી.
તેથી તેને~$\pi_1$ સાબિતીમાંથી
કાઢી શકાય છે અને
$\fCenter !B(t_1), \dots, !B(t_n)$ની
સાબિતી મળે છે.
પછી~$\RightR{\lor}$ના
$n-1$ ઉપયોગોથી
હર્બ્રાન્ડ વિકલ્પ
$\Sequent !B(t_1) \lor \dots \lor !B(t_n)$ની
સાબિતી મળે છે.

મુશ્કેલી એ છે કે
$\Sequent \lexists[x][!B(x)]$ની
કટ-મુક્ત સાબિતીમાં પણ
બધાં~\RightR{\lexists}
અનુમાનો એક જ અંતિમ જૂથમાં
હોય એમ ધારવું ચાલે નહીં.
તેમની વચ્ચે એવા અનુમાનો
હોઈ શકે જેમાં કોઈ~$!B(t_i)$
મુખ્ય હોય. તેથી મધ્ય-સિક્વન્ટ
પ્રમેય જોઈએ.

\begin{thm}[મધ્ય-સિક્વન્ટ પ્રમેય]\ollabel{thm:midsequent}
જો $\Gamma \Sequent \Delta$ની
\Log{G3c}-સાબિતી~$\pi$ હોય
અને~$\Gamma$ તથા~$\Delta$નાં
બધાં સૂત્રો પરિમાણક-અગ્ર
સ્વરૂપનાં હોય, તો એવી
સાબિતી~$\pi'$ છે જેમાં
$\Pi \Sequent \Lambda$
સિક્વન્ટ છે અને તેની નીચેના
બધાં અનુમાનો પરિમાણકનાં,
ઉપરનાં બધાં વિધાનાત્મક છે.
\end{thm}

\begin{proof}
  ~$\pi$માંના પરિમાણક-અનુમાનનો
  \emph{ક્રમ} એટલે તેની નીચેના
  વિધાનાત્મક અનુમાનોની સંખ્યા,
  અને~$\pi$નો ક્રમ~$o(\pi)$ એટલે
  તેનાં બધાં પરિમાણક-અનુમાનોના
  ક્રમોનો સરવાળો, એમ વ્યાખ્યાયિત કરો.
  જો $o(\pi) = 0$, તો કોઈ
  પરિમાણક-અનુમાનની નીચે
  વિધાનાત્મક અનુમાન નથી,
  અને કામ પૂરું છે:
  પરિમાણક-અનુમાનો હોય તો
  બધાં એક-આધારવિધાનવાળાં હોવાથી
  સૌથી ઉપરનું એક પરિમાણક-અનુમાન
  છે; તેનું આધારવિધાન
  મધ્ય-સિક્વન્ટ છે.
  પરિમાણક-અનુમાન જ ન હોય તો
  અંતિમ સિક્વન્ટ મધ્ય-સિક્વન્ટ લઈએ.
  \sourcecorrection{OLMD-222}{મૂળ આધારકિસ્સો પરિમાણક-અનુમાન હોય ત્યારે જ સૌથી ઉપરના અનુમાનનું આધારવિધાન બતાવે છે; એકેય પરિમાણક-અનુમાન ન હોય તે કિસ્સામાં અંતિમ સિક્વન્ટ પૂરતું છે.}

  જો $o(\pi) > 0$, તો
  $>0$ ક્રમનું સૌથી નીચેનું
  પરિમાણક-અનુમાન પસંદ કરો.
  તેનો ક્રમ $>0$ હોવાથી
  તેની નીચે વિધાનાત્મક
  અનુમાન હોવું જોઈએ.
  તે સૌથી નીચેનું હોવાથી
  તેનો નિષ્કર્ષ બીજા
  પરિમાણક-અનુમાનનું આધારવિધાન
  હોઈ શકે નહીં; એવું હોય
  તો એ નીચેના પરિમાણક-અનુમાનની
  નીચે પણ વિધાનાત્મક અનુમાન
  હશે. કયું પરિમાણક-અનુમાન
  અને તેની નીચે કયું વિધાનાત્મક
  અનુમાન છે તેના (ઘણા)
  કિસ્સા જુદા પાડીએ.
  અંતિમ સિક્વન્ટમાં માત્ર
  પરિમાણક-અગ્ર સૂત્રો હોવાથી
  પરિમાણક-અનુમાનનું મુખ્ય
  સૂત્ર પછીના વિધાનાત્મક
  અનુમાનમાં સક્રિય હોઈ
  શકે નહીં. નહીં તો
  તે વિધાનાત્મક અનુમાનના
  મુખ્ય સૂત્રમાં વિધાનાત્મક
  સંયોજકના વ્યાપમાં
  પરિમાણક આવશે.
  ઉપ-સૂત્ર ગુણધર્મથી
  તે અંતિમ સિક્વન્ટના કોઈ
  સૂત્રનું
  ઉપ-સૂત્ર હોવું પડશે.
  તેથી પરિમાણક-અનુમાનનું
  મુખ્ય સૂત્ર નીચેના
  વિધાનાત્મક અનુમાનના સંદર્ભનું
  ઘટક છે.

  બે ઉદાહરણો જોઈએ:

પહેલાં માનો કે $\RightR{\lforall}$નું
નિષ્કર્ષ $\LeftR{\lif}$ના જમણા
આધારવિધાનમાં આવે છે. ત્યારે~$\pi$ની
સંબંધિત ઉપ-સાબિતીનો અંત
ઉપ-સાબિતી~$\pi_1$માં થાય છે.
\[
\AxiomC{}
\RightLabel{$\pi_2$}
\Deduce$\Gamma' \fCenter \Delta', \lforall[x][!A(x)], !B$
\AxiomC{}
\RightLabel{$\pi_3$}
\Deduce$!C, \Gamma' \fCenter \Delta', !A(c)$
\RightLabel{\RightR{\lforall}}
\UnaryInf$!C, \Gamma' \fCenter \Delta', \lforall[x][!A(x)]$
\RightLabel{\LeftR{\lif}}
\BinaryInf$!B \lif !C, \Gamma' \fCenter \Delta', \lforall[x][!A(x)]$
\DisplayProof
\]
આપણે ધારીએ કે~$\pi$ નિયમિત છે;
આથી સ્વનિર્ધારક ચલ~$c$,
$\pi_3$ની બહાર આવતો નથી.
ખાસ કરીને તે $!B
\lif !C$માં કે~$\pi_2$માં
આવતો નથી. હવે સાબિતી~$\pi_1'$
જુઓ:
\[
\AxiomC{}
\RightLabel{$\pi_2'$}
\Deduce$\Gamma' \fCenter \Delta', \lforall[x][!A(x)], !A(c), !B$
\AxiomC{}
\RightLabel{$\pi_3$}
\Deduce$!C, \Gamma' \fCenter \Delta', \lforall[x][!A(x)], !A(c)$
\RightLabel{\LeftR{\lif}}
\BinaryInf$!B \lif !C, \Gamma' \fCenter \Delta', \lforall[x][!A(x)], !A(c)$
\RightLabel{\RightR{\lforall}}
\UnaryInf$!B \lif !C, \Gamma' \fCenter \Delta', \lforall[x][!A(x)], \lforall[x][!A(x)]$
\DisplayProof
\]
અહીં~$\pi_2$ના જમણા ભાગમાં
$!A(c)$ ઉમેરી શિથિલીકરણ
કરવાથી~$\pi_2'$ મળે છે.
(આથી કોઈ અનુમાન, ખાસ કરીને
ક્રમને અસર કરે એવું પરિમાણક-અનુમાન,
ઉમેરાતું નથી. વળી~$\pi_2$ની
સ્વનિર્ધારક ચલની કોઈ શરત પણ
ભંગ થતી નથી, કારણ કે~$!A(c)$માં
આવતો કોઈ અચળ~$\pi_2$માં
સ્વનિર્ધારક ચલ તરીકે વપરાતો નથી.)
$c$ એ~$!B \lif !C$માં
આવતો ન હોવાથી
\RightR{\lforall}ની
સ્વનિર્ધારક ચલની શરત
હજી પણ પૂરી થાય છે.

હવે સંકોચનના સ્વીકાર્યપણાથી
$!B \lif !C, \Gamma' \fCenter \Delta',
\lforall[x][!A(x)]$ની
સાબિતી~$\pi_1''$ મળે છે.
$\lforall[x][!A(x)]$ માટે
$\RightR{\Contraction}$નું
સ્વીકાર્યપણું સાબિત કરતી
સાબિતી અને તેમાં વપરાયેલી
$\RightR{\lforall}$ની
ઉલટાવી શકાય તેવી હોવાની
સાબિતી કોઈ અનુમાનનો ક્રમ
બદલતી નથી: વધુમાં વધુ
અનુમાનો દૂર કરે છે.
આથી~$\pi_1''$માં
$\pi_1$ કરતાં વધુ
પરિમાણક-અનુમાનો નથી.
વળી~$\pi_1''$માં દરેક
પરિમાણક-અનુમાનની નીચે
તેનાં~$\pi_1$માંનાં
અનુરૂપ અનુમાન જેટલાં જ
વિધાનાત્મક અનુમાનો છે,
એક અપવાદ સિવાય:
છેલ્લા \RightR{\lforall}ની
નીચે~$\pi$માં \LeftR{\lif} હતું;
તેને~$\pi_1''$માં
ઉપર ખસેડ્યું છે.
હવે~$\pi$માં
$\pi_1$ને બદલે~$\pi_1''$
મૂકો. મળેલી સાબિતીનો
ક્રમ ઓછો છે, કારણ કે
ઓછામાં ઓછો
$\RightR{\lforall}$ અનુમાનનો
ક્રમ (ઓછામાં ઓછો)~$1$ ઘટ્યો છે.
હવે અનુગામી ધારણા લાગુ પડે છે.

પરિમાણક-અનુમાન
\RightR{\lexists} હોય તે
કિસ્સા વધુ સહેલા છે, કારણ કે
સંકોચન જરૂરી નથી અને
સ્વનિર્ધારક ચલની શરતની
ચિંતા પણ નથી. દાખલા તરીકે,
માનો કે~$\RightR{\lexists}$
પછી~$\RightR{\land}$ આવે છે
(આ વખતે તેના ડાબા
આધારવિધાન તરીકે).
સંબંધિત ઉપ-સાબિતી
$\pi_1$ આ છે:
\[
\AxiomC{}
\RightLabel{$\pi_2$}
\Deduce$\Gamma' \fCenter \Delta', \lexists[x][!A(x)], !A(t), !B$
\RightLabel{\RightR{\lexists}}
\UnaryInf$\Gamma' \fCenter \Delta', \lexists[x][!A(x)], !B$
\AxiomC{}
\RightLabel{$\pi_3$}
\Deduce$\Gamma' \fCenter \Delta', \lexists[x][!A(x)], !C$
\RightLabel{\RightR{\land}}
\BinaryInf$\Gamma' \fCenter \Delta', \lexists[x][!A(x)], !B \land !C$
\DisplayProof
\]
\sourcecorrection{OLMD-223}{મૂળ વૃક્ષમાં આ સિક્વન્ટના \(\Gamma'\) પહેલાં વધારાનું ઉદ્ગારચિહ્ન હતું; અનુમાનના આધારવિધાન અને નિષ્કર્ષ સાથે સુસંગત કરવા તે કાઢ્યું છે.}
હવે~$\pi$માં $\pi_1$ને
બદલે~$\pi_1'$ મૂકો:
\[
\AxiomC{}
\RightLabel{$\pi_2$}
\Deduce$\Gamma' \fCenter \Delta', \lexists[x][!A(x)], !A(t), !B$
\AxiomC{}
\RightLabel{$\pi_3'$}
\Deduce$\Gamma' \fCenter \Delta', \lexists[x][!A(x)], !A(t), !C$
\RightLabel{\RightR{\land}}
\BinaryInf$\Gamma' \fCenter \Delta', \lexists[x][!A(x)], !A(t), !B \land !C$
\RightLabel{\RightR{\lexists}}
\UnaryInf$\Gamma' \fCenter \Delta', \lexists[x][!A(x)], !B \land !C$
\DisplayProof
\]
અહીં~$\pi_3$ના જમણા ભાગમાં
$!A(t)$ ઉમેરી શિથિલીકરણ
કરવાથી~$\pi_3'$ મળે છે.
આ શિથિલીકરણમાં~$\pi_3$ના
દરેક સિક્વન્ટની જમણી બાજુ
માત્ર~$!A(t)$ ઉમેરાય છે,
તેથી કોઈ પરિમાણક-અનુમાનનો
ક્રમ બદલાતો નથી.
$\pi'$માં પરિમાણક-અનુમાનોના
ક્રમો~$\pi$ જેવા જ છે;
માત્ર~\RightR{\land} સાથે
અદલબદલ કરેલા
\RightR{\lexists} અનુમાનનો
ક્રમ~$1$ ઘટ્યો છે.
હવે અનુગામી ધારણા લાગુ પડે છે.
\end{proof}

\begin{prob}
\olref{thm:midsequent}ની
સાબિતીમાં એ કિસ્સા વર્ણવો
જ્યાં~\LeftR{\lexists}નું
નિષ્કર્ષ~\RightR{\lif}ના
આધારવિધાનમાં આવે છે,
અને જ્યાં~\LeftR{\lforall}નું
નિષ્કર્ષ~\LeftR{\lor}ના
ડાબા આધારવિધાનમાં આવે છે.
\end{prob}

\begin{proof}[હર્બ્રાન્ડના પ્રમેયની સાબિતી]
માનો કે~$\pi$નો અંત
$\Sequent \lexists[x][!A(x)]$માં
થાય છે.
\olref{thm:midsequent}થી
આપણે~$\pi$ને મધ્ય-સિક્વન્ટ~$S$
ધરાવતી સાબિતી~$\pi'$માં
ફેરવી શકીએ છીએ. આ
મધ્ય-સિક્વન્ટનું સ્વરૂપ
\[\Sequent
\lexists[x][!A(x)], !A(t_1), \dots, !A(t_n).\]
હોવું જોઈએ.
$S$ની ઉપરનાં બધાં
અનુમાનો વિધાનાત્મક હોવાથી
$\lexists[x][!A(x)]$
$S$ની ઉપરના કોઈ અનુમાનમાં
મુખ્ય હોઈ શકે નહીં.
તે પરમાણ્વીય નથી, તેથી
સ્વયંસિદ્ધમાં પણ મુખ્ય
હોઈ શકે નહીં.
વળી~$\lexists[x][!A(x)]$,
$!A(t_1)$નું યોગ્ય
ઉપસૂત્ર હોઈ શકે નહીં
(યાદ રાખો કે~$!A(x)$
પરિમાણક-રહિત છે), એટલે
$S$ની ઉપર સક્રિય પણ
હોઈ શકે નહીં.
તેથી~$S$માં પૂરી થતી
ઉપ-સાબિતીમાંથી
$\lexists[x][!A(x)]$
દૂર કરતાં
$\Sequent !A(t_1), \dots, !A(t_n)$ની
સાચી સાબિતી મળે છે.
તેમાંથી~$n-1$ વાર
$\RightR{\lor}$ લાગુ કરી
$\Sequent !A(t_1) \lor \dots \lor !A(t_n)$ની
સાબિતી મેળવી શકાય છે.
\end{proof}
% END OLP-0661
\endgroup

\begingroup
\makeatletter\def\ol@sectioncs{section}\makeatother

% BEGIN OLP-0658
\olfileid{pt}{cut}{itp}

\olsection{માએહારાનું ઉપપ્રમેય અને ક્રેગનું અંતર્વેશન પ્રમેય}
\phantomsection\label{guunit:OLP-0658}


વિલિયમ ક્રેગનું અંતર્વેશન પ્રમેય કહે છે:
જો $!A \Entails !B$, તો એવું
વાક્ય~$!C$ (\emph{અંતર્વેશક}) છે કે
$!A \Entails !C$ અને $!C \Entails !B$.
મહત્ત્વની વાત એ છે કે અંતર્વેશક~$!C$
એવો પસંદ કરી શકાય કે આ~$!C$માં આવતાં
બધાં અચળો અને વિધેયો
~$!A$ તથા~$!B$ બંનેમાં આવે.
(અહીં કોઈ ફલનો સામેલ નથી એમ ધારીએ છીએ.)
આ મર્યાદા ન હોય તો~$!A$ અને~$!B$
પોતે જ સરળતાથી અંતર્વેશક બની જાય.

અંતર્વેશન પ્રમેય શાસ્ત્રીય પ્રથમ-ક્રમ
તર્ક માટે સાચું છે. તેને શોધાયેલો
પ્રથમ-ક્રમ તર્કનો ``છેલ્લો મહત્ત્વનો ગુણધર્મ''
પણ કહેવામાં આવ્યો છે.
નિદર્શ સિદ્ધાંત અને સ્વયંસંચાલિત
તર્કવિચારમાં તેના મહત્ત્વના ઉપયોગો છે.
તેની મૂળ પ્રેરણા અને ઉપયોગ વિજ્ઞાનના
તત્ત્વજ્ઞાનમાં હતાં: કોઈ સિદ્ધાંતને
કયા મૂળભૂત ખ્યાલોથી સ્વયંસિદ્ધ કરી શકાય
કે ન શકાય, તે તપાસવામાં તે મદદ કરે છે.
તેમાંથી બેથનું વ્યાખ્યેયતા પ્રમેય પણ
ફલિત થાય છે: સિદ્ધાંત કોઈ ખ્યાલને
ગૂઢ રીતે વ્યાખ્યાયિત કરી શકે,
તો સ્પષ્ટ રીતે પણ કરી શકે છે.

ગણિતીય તર્કનાં અનેક પરિણામોની જેમ
અંતર્વેશન પ્રમેયને નિદર્શ સિદ્ધાંતની
અને સાબિતી સિદ્ધાંતની, બંને પદ્ધતિથી
સાબિત કરી શકાય છે.
સાબિતી-સૈદ્ધાંતિક પદ્ધતિનો લાભ એ છે કે
તે અંતર્વેશકનું અસ્તિત્વ જ નહીં,
પરંતુ $!A \Entails !B$ની
સાબિતીમાંથી તેને ગણતો કલનવિધિ
પણ આપે છે; દાખલા તરીકે,
~\Log{G3c}માં
$!A \Sequent !B$ની સાબિતીમાંથી.

$\Lang L(!A)$ એ~$!A$માંનાં
અચળો અને વિધેયોનો ગણ
અને
$\Lang L(\Gamma) = \bigcup \Setabs{\Lang L(!A)}{!A \in
\Gamma}$ માનો. આ સમગ્ર વિભાગમાં
ફલનો વિનાનાં સૂત્રો
લેવાં છે.

\begin{thm}[ક્રેગનું અંતર્વેશન પ્રમેય]\ollabel{thm:interpolation}
~$!A$ની ભાષાને~$\Lang L(!A)$
અને~$!B$ની ભાષાને~$\Lang L(!B)$ કહો.
જો $\Log{G3c} \Proves !A \Sequent !B$,
તો એવી સૂત્ર~$!C$ છે
જેની ભાષા માટે
$\Lang L(!C) \subseteq \Lang
L(!A) \cap \Lang L(!B)$ અને
$\Log{G3c} \Proves !A \Sequent !C$
તેમજ
$\Log{G3c} \Proves !C \Sequent !B$.
\end{thm}

~\Log{G3c}ની સાબિતીઓની રચનાનો
ઉપયોગ કરવા અંતર્વેશન પ્રમેયનું
સામાન્યીકરણ સાબિત કરીએ છીએ.
તે માટે સિક્વન્ટના પૂર્વાંગ અને
ઉત્તરાંગને બે બે ભાગમાં વહેંચેલા
વિચારો. $\maeh{\Gamma_1}{\Gamma_2} \Sequent
\maeh{\Delta_1}{\Delta_2}$ લખીએ ત્યારે
$\Gamma_1$ અને $\Delta_1$ પહેલા ભાગમાં,
અને $\Gamma_2$ અને $\Delta_2$
બીજા ભાગમાં છે. પછી નીચેના ઉપપ્રમેયથી
અંતર્વેશન પ્રમેય તરત ફલિત થાય છે.

\begin{lem}[માએહારાનું ઉપપ્રમેય]\ollabel{lem:maehara}
જો $\Log{G3c}
\Proves \maeh{\Gamma_1}{\Gamma_2} \Sequent \maeh{\Delta_1}{\Delta_2}$,
તો એવી સૂત્ર~$!C$ છે
જેની ભાષા માટે
$\Lang L(!C)
\subseteq \Lang L(\Gamma_1,
\Delta_1) \cap \Lang L(\Gamma_2, \Delta_2)$
અને
$\Log{G3c} \Proves \Gamma_1 \Sequent \Delta_1, !C$
તથા
$\Log{G3c} \Proves !C,
\Gamma_2 \Sequent \Delta_2$.
\end{lem}

\begin{proof}
  માનો કે $\pi \Proves \maeh{\Gamma_1}{\Gamma_2} \Sequent
  \maeh{\Delta_1}{\Delta_2}$.
  $n=\pheight{\pi}$ પર આગમનથી
  દાવો સાબિત કરીએ છીએ.

  જો $n=0$, તો
  $\maeh{\Gamma_1}{\Gamma_2} \Sequent
  \maeh{\Delta_1}{\Delta_2}$
  સ્વયંસિદ્ધ છે અને કોઈ પરમાણ્વીય
  સૂત્ર~$!A$ એ
  $\Gamma_1, \Gamma_2$ તથા
  $\Delta_1, \Delta_2$ બંનેમાં આવે છે.
  ડાબે~$!A$ પહેલા કે બીજા ભાગમાં,
  અને જમણે પહેલા કે બીજા ભાગમાં છે,
  તેના ચાર કિસ્સા છે.
  \begin{enumerate}
    \item $!A \in \Gamma_1$ અને $!A \in \Delta_1$.
    એવું~$!C$ જોઈએ કે
    $\Gamma_1 \Sequent \Delta_1, !C$ અને
    $!C, \Gamma_2 \Sequent \Delta_2$
    બંને સાબિતીક્ષમ હોય.
    પહેલું તો~$!C$ ગમે તે લઈએ
    તોય સ્વયંસિદ્ધ છે, કારણ કે~$!A$
    ડાબે અને જમણે બંને છે.
    $!C, \Gamma_2 \Sequent \Delta_2$
    સાબિતીક્ષમ રહે તેની ખાતરી માટે
    $!C \ident \lfalse$ લેવું પડે.
    \item $!A \in \Gamma_1$ અને
    $!A \in \Delta_2$. ત્યારે
    $\Gamma_1 \Sequent \Delta_1, !A$
    અને $!A, \Gamma_2 \Sequent \Delta_2$
    બંને સ્વયંસિદ્ધ છે; તેથી
    $!C \ident !A$ લઈ શકાય.
    \item $!A \in \Gamma_2$ અને
    $!A \in \Delta_1$. ત્યારે
    $!A, \Gamma_1 \Sequent \Delta_1$
    અને $\Gamma_2 \Sequent \Delta_2, !A$
    સ્વયંસિદ્ધ હોય, પરંતુ તેમાં~$!A$
    ખોટી બાજુએ છે. આથી~$!A$
    અંતર્વેશક નથી. તેમ છતાં
    $\RightR{\lnot}$ અને
    $\LeftR{\lnot}$ વડે
    $\Gamma_1 \Sequent \Delta_1, \lnot!A$
    અને $\lnot !A, \Gamma_2 \Sequent \Delta_2$
    સાબિત કરી શકાય.
    \item $!A \in \Gamma_2$ અને
    $!A \in \Delta_2$. કસરત.
  \end{enumerate}
  જો $\lfalse \in \Gamma_1$
  અથવા $\lfalse \in \Gamma_2$
  હોવાને કારણે
  $\maeh{\Gamma_1}{\Gamma_2} \Sequent
  \maeh{\Delta_1}{\Delta_2}$
  સ્વયંસિદ્ધ હોય, તો એ કિસ્સાઓ
  પણ કસરત તરીકે છોડીએ છીએ.

  જો $n>0$, તો~$\pi$નો અંત
  તાર્કિક અનુમાનથી થાય છે.
  કયો નિયમ વપરાયો અને મુખ્ય સૂત્ર
  કયા ભાગમાં છે, તેના કિસ્સા
  જુદા પાડવા પડે.
  દરેક કિસ્સામાં આગમન-પૂર્વધારણા
  વાપરી શકીએ: $\pheight{\pi_1} < n$
  ધરાવતી દરેક સાબિતી~$\pi_1$
  માટે, ~$\pi_1$ના અંતિમ સિક્વન્ટના
  બે ભાગ ગમે તે રીતે કર્યા હોય,
  ઉપપ્રમેય સાચું છે.
  હવે કેટલાક ચોક્કસ કિસ્સા જોઈએ.
  \begin{enumerate}
    \item $\pi$નો અંત~$\RightR{\lnot}$થી
    થાય છે અને મુખ્ય સૂત્ર~$\lnot !A$ છે.
    $\lnot !A$ પહેલા કે બીજા ભાગમાં છે
    તે પ્રમાણે~$\pi$નો અંત બેમાંથી
    એક રીતે થાય છે:
    \[
    \AxiomC{}
    \RightLabel{$\pi_1$}
    \Deduce$\maeh{!A, \Gamma_1}{\Gamma_2} \fCenter
      \maeh{\Delta_1}{\Delta_2}$
    \RightLabel{\RightR{\lnot}}
    \UnaryInf$\maeh{\Gamma_1}{\Gamma_2} \fCenter
      \maeh{\Delta_1, \lnot !A}{\Delta_2}$
    \DisplayProof
    \]
    અથવા
    \[
    \AxiomC{}
    \RightLabel{$\pi_1$}
    \Deduce$\maeh{\Gamma_1}{!A, \Gamma_2} \fCenter
      \maeh{\Delta_1}{\Delta_2}$
    \RightLabel{\RightR{\lnot}}
    \UnaryInf$\maeh{\Gamma_1}{\Gamma_2} \fCenter
      \maeh{\Delta_1}{\Delta_2, \lnot !A}$
    \DisplayProof
    \]
    ધ્યાન આપો: જો મુખ્ય સૂત્ર
    ~$\lnot !A$ પહેલા ભાગમાં હોય,
    તો આધારવિધાનમાં સક્રિય
    સૂત્ર~$!A$ પણ પહેલા
    ભાગમાં મૂક્યું છે. આ પસંદગી
    આપણી છે. સક્રિય સૂત્ર બીજા
    ભાગમાં મૂક્યું હોય તોય
    આગમન-પૂર્વધારણા લાગુ પડે,
    પરંતુ પછી અંતર્વેશક મળી
    શકતો નથી (કસરત: અજમાવી જુઓ).

    પહેલાં~$\lnot !A$ પહેલા ભાગમાં
    હોય તે સ્થિતિ લો.
    આગમન-પૂર્વધારણાથી~$\pi_1$ના અંતિમ
    સિક્વન્ટનો અંતર્વેશક, એટલે
    સૂત્ર~$!C_1$, એવો છે કે નીચેના
    બંને સિક્વન્ટ
    \begin{align*}
    !A, \Gamma_1 & \Sequent \Delta_1, !C_1 &&\text{અને} &
    !C_1, \Gamma_2 & \Sequent \Delta_2
    \intertext{સાબિત થઈ શકે છે. આપણને એવી સૂત્ર~$!C$ જોઈએ કે નીચેના બંને સિક્વન્ટ સાબિતીક્ષમ હોય:}
    \Gamma_1 & \Sequent \Delta_1, \lnot !A, !C &&\text{અને} &
    !C, \Gamma_2 & \Sequent \Delta_2
    \end{align*}
    સ્પષ્ટ રીતે~$!C \ident !C_1$ લઈ શકાય:
    જમણો સિક્વન્ટ સાબિતીક્ષમ છે
    એ આગમન-પૂર્વધારણાથી મળે છે,
    અને ડાબો સિક્વન્ટ
    આગમન-પૂર્વધારણાથી મળેલા સિક્વન્ટ
    પર~$\RightR{\lnot}$ લગાવતાં મળે છે.
    આગમન-પૂર્વધારણા એમ પણ કહે છે કે
    $\Lang L(!C_1) \subseteq \Lang L(!A, \Gamma_1,
    \Delta_1) \cap \Lang L(\Gamma_2, \Delta_2)$.
    કારણ કે $\Lang L(\lnot !A) = \Lang L(!A)$
    અને $\Lang L(!C) = \Lang L(!C_1)$,
    આપણને
    $\Lang L(!C_1) \subseteq
    \Lang L(\Gamma_1, \Delta_1, \lnot !A) \cap
    \Lang L(\Gamma_2, \Delta_2)$ મળે છે.

    બીજા કિસ્સામાં~$\lnot !A$ બીજા
    ભાગમાં હોવાથી સ્થિતિ થોડી જુદી છે.
    આધારવિધાનનું સક્રિય સૂત્ર
    પણ બીજા ભાગમાં મૂકીને
    આગમન-પૂર્વધારણા આના પર લગાવીએ:
    \begin{align*}
    \Gamma_1 & \Sequent \Delta_1, !C_1 &&\text{અને} &
    !C_1, !A, \Gamma_2 & \Sequent \Delta_2
    \intertext{આ બંને સાબિત થઈ શકે છે. બીજાં સિક્વન્ટ પર ફરી \RightR{\lnot} લગાવતાં આની સાબિતીઓ મળે છે:}
    \Gamma_1 & \Sequent \Delta_1, !C_1 &&\text{અને} &
    !C_1, \Gamma_2 & \Sequent \Delta_2, \lnot !A
    \end{align*}
    $!C_1$ અંતર્વેશક હોય ત્યારે
    આપણને સાબિતીક્ષમ જોઈએ એવા
    બરાબર આ જ સિક્વન્ટ છે;
    ફરી~$!C \ident !C_1$ લઈએ.
    અગાઉની જેમ
    $\Lang L(!C_1) \subseteq
    \Lang L(\Gamma_1, \Delta_1) \cap
    \Lang L(\Gamma_2, \Delta_2, \lnot !A)$.
    \item $\pi$નો અંત~$\RightR{\lif}$થી
    થાય છે અને મુખ્ય
    સૂત્ર~$!A \lif !B$ છે.
    $!A \lif !B$ પહેલા કે બીજા
    ભાગમાં છે તે પ્રમાણે બે કિસ્સા છે.
    સાબિતી~$\pi$નો અંત આ રીતે થાય છે:
    \[
    \AxiomC{}
    \RightLabel{$\pi_1$}
    \Deduce$\maeh{!A, \Gamma_1}{\Gamma_2} \fCenter
      \maeh{\Delta_1, !B}{\Delta_2}$
    \RightLabel{\RightR{\lif}}
    \UnaryInf$\maeh{\Gamma_1}{\Gamma_2} \fCenter
      \maeh{\Delta_1, !A \lif !B}{\Delta_2}$
    \DisplayProof
    \]
    અથવા
    \[
    \AxiomC{}
    \RightLabel{$\pi_1$}
    \Deduce$\maeh{\Gamma_1}{!A, \Gamma_2} \fCenter
      \maeh{\Delta_1}{\Delta_2, !B}$
    \RightLabel{\RightR{\lif}}
    \UnaryInf$\maeh{\Gamma_1}{\Gamma_2} \fCenter
      \maeh{\Delta_1}{\Delta_2, !A \lif !B}$
    \DisplayProof
    \]
    ફરી, આધારવિધાનનાં સક્રિય
    સૂત્રોને નિષ્કર્ષના મુખ્ય
    સૂત્રના જ ભાગમાં મૂક્યાં છે.
    પહેલા કિસ્સામાં આગમન-પૂર્વધારણાથી
    આની સાબિતીઓ મળે છે:
    \begin{align*}
    !A, \Gamma_1 & \Sequent \Delta_1, !B, !C_1 &&\text{અને} & !C_1, \Gamma_2 & \Sequent \Delta_2
    \intertext{ડાબો સિક્વન્ટ \RightR{\lif}નું યોગ્ય આધારવિધાન છે. તેથી નીચેના સિક્વન્ટ સાબિતીક્ષમ છે:}
    \Gamma_1 & \Sequent \Delta_1, !A \lif !B, !C_1 &&\text{અને} & !C_1, \Gamma_2 & \Sequent \Delta_2
    \end{align*}
    આથી~$!C_1$ એ~$\pi$ના અંતિમ
    સિક્વન્ટનો પણ અંતર્વેશક છે;
    ફરી~$!C \ident !C_1$ લઈ શકાય.
    
    બીજો કિસ્સો એ જ રીતે થાય છે.
    \item $\pi$નો અંત~$\LeftR{\lif}$થી
    થાય છે અને મુખ્ય સૂત્ર
    ~$!A \lif !B$ છે.
    જો $!A \lif !B$ પહેલા ભાગમાં હોય,
    તો સાબિતી~$\pi$નો અંત આ રીતે થાય છે:
    \[
    \AxiomC{}
    \RightLabel{$\pi_1$}
    \Deduce$\maeh{\Gamma_1}{\Gamma_2} \fCenter
      \maeh{\Delta_1, !A}{\Delta_2}$
    \AxiomC{}
    \RightLabel{$\pi_2$}
    \Deduce$\maeh{!B, \Gamma_1}{\Gamma_2} \fCenter
      \maeh{\Delta_1}{\Delta_2}$
    \BinaryInf$\maeh{!A \lif !B, \Gamma_1}{\Gamma_2} \fCenter
      \maeh{\Delta_1}{\Delta_2}$
    \DisplayProof
    \]
    આગમન-પૂર્વધારણા ચાર
    સાબિતીક્ષમ સિક્વન્ટ આપે છે:
    ડાબા આધારવિધાનના અંતર્વેશક~$!C_1$
    માટે બે અને બીજા આધારવિધાનના
    અંતર્વેશક~$!C_2$ માટે બે:
    \begin{align*}
    \Gamma_1 & \Sequent \Delta_1, !A, !C_1 &&\text{(1)} &
    !C_1, \Gamma_2 & \Sequent \Delta_2 && \text{(2)}\\
    !B, \Gamma_1 & \Sequent \Delta_1, !C_2 &&\text{(3)} &
    !C_2, \Gamma_2 & \Sequent \Delta_2 && \text{(4)}
    \end{align*}
    શિથિલીકરણ વડે (1)ના જમણે~$!C_2$
    અને (3)ના જમણે~$!C_1$ ઉમેરીએ,
    તો સિક્વન્ટ (1) અને (3)
    ~\LeftR{\lif}નાં આધારવિધાન બને:
    \[
    \AxiomC{}
    \Deduce$\Gamma_1 \fCenter \Delta_1, !A, !C_1, !C_2$
    \AxiomC{}
    \Deduce$!B, \Gamma_1 \fCenter \Delta_1, !C_1, !C_2$
    \RightLabel{\LeftR{\lif}}
    \BinaryInf$!A \lif !B, \Gamma_1 \fCenter \Delta_1, !C_1, !C_2$
    \DisplayProof
    \]
    ~$\RightR{\lor}$ લગાવતાં સિક્વન્ટ મળે:
    \[
    !A \lif !B, \Gamma_1 \fCenter \Delta_1, !C_1 \lor !C_2
    \]
    આથી આ કિસ્સામાં
    $!C_1 \lor !C_2$ અંતર્વેશક હોઈ શકે.
    ખરેખર, બાકી સિક્વન્ટ (2) અને~(4)
    વડે બીજું જરૂરી સિક્વન્ટ
    સાબિત કરી શકાય:
    \[
    \AxiomC{}
    \Deduce$!C_1, \Gamma_2 \fCenter \Delta_2$
    \AxiomC{}
    \Deduce$!C_2, \Gamma_2 \fCenter \Delta_2$
    \RightLabel{\LeftR{\lor}}
    \BinaryInf$!C_1 \lor !C_2, \Gamma_2 \fCenter \Delta_2$
    \DisplayProof
    \]
    $!C_1 \lor !C_2$ યોગ્ય ભાષામાં છે
    તે ચકાસવાનું રહે છે.
    આગમન-પૂર્વધારણાથી મળે છે કે
    \begin{align*}
    \Lang L(!C_1) & \subseteq
    \Lang L(\Gamma_1, \Delta_1, !A) \cap \Lang L(\Gamma_2, \Delta_2) &\text{અને}\\
    \Lang L(!C_2) & \subseteq
    \Lang L(!B, \Gamma_1, \Delta_1) \cap \Lang L(\Gamma_2, \Delta_2)
    \intertext{કારણ કે}
    \Lang L(!A) & \subseteq \Lang L(!A \lif !B) &\text{અને}\\
    \Lang L(!B) & \subseteq \Lang L(!A \lif !B)
    \intertext{તેથી બંને માટે}
    \Lang L(!C_1) & \subseteq
    \Lang L(\Gamma_1, \Delta_1, !A \lif !B) \cap \Lang L(\Gamma_2, \Delta_2) &\text{અને}\\
    \Lang L(!C_2) & \subseteq
    \Lang L(\Gamma_1, \Delta_1, !A \lif !B) \cap \Lang L(\Gamma_2, \Delta_2)
    \intertext{અને તેથી $\Lang L (!C_1 \lor !C_2) = {}$}
    \Lang L(!C_1) \cup \Lang L(!C_2) & \subseteq
    \Lang L(\Gamma_1, \Delta_1, !A \lif !B) \cap \Lang L(\Gamma_2, \Delta_2),
    \end{align*}
    જે જરૂરી હતું તે જ છે.

    જો~$!A \lif !B$ બીજા ભાગમાં હોય,
    તો સ્થિતિ જુદી છે, પણ દલીલ સમાન છે.
    આગમન-પૂર્વધારણા ફરી ચાર
    સાબિતીક્ષમ સિક્વન્ટ આપે છે:
    \begin{align*}
    \Gamma_1 & \Sequent \Delta_1, !C_1 &&\text{(1)} &
    !C_1, \Gamma_2 & \Sequent \Delta_2, !A && \text{(2)}\\
    \Gamma_1 & \Sequent \Delta_1, !C_2 &&\text{(3)} &
    !C_2, !B, \Gamma_2 & \Sequent \Delta_2 && \text{(4)}
    \end{align*}
    હવે સિક્વન્ટ (2) અને (4),
    શિથિલીકરણથી ઉમેરેલાં કેટલાંક
    સૂત્રો સાથે,
    ~\LeftR{\lif}નાં આધારવિધાન બને:
    \[
    \AxiomC{}
    \Deduce$!C_1, !C_2, \Gamma_2 \fCenter \Delta_2, !A$
    \AxiomC{}
    \Deduce$!B, !C_1, !C_2, \Gamma_2 \fCenter \Delta_2$
    \RightLabel{\LeftR{\lif}}
    \BinaryInf$!A \lif !B, !C_1, !C_2, \Gamma_2 \fCenter \Delta_2$
    \DisplayProof
    \]
    આ વખતે પણ~$!C_1$ અને~$!C_2$માંથી
    બનેલું એક જ સૂત્ર ધરાવતો
    સિક્વન્ટ જોઈએ, પણ તે પૂર્વાંગમાં
    જોઈએ (કારણ કે હવે બીજો ભાગ છે).
    ~$\LeftR{\land}$ લગાવતાં કામ થાય છે;
    $!C \ident (!C_1 \land !C_2)$
    અંતર્વેશક લેવો જોઈએ.
    બાકી સિક્વન્ટ (1) અને (3)
    વડે બીજા ભાગને અનુરૂપ સિક્વન્ટ
    સાબિત કરી શકાય:
    \[
    \AxiomC{}
    \Deduce$\Gamma_1 \fCenter \Delta_1, !C_1$
    \AxiomC{}
    \Deduce$\Gamma_1 \fCenter \Delta_1, !C_2$
    \RightLabel{\RightR{\land}}
    \BinaryInf$\Gamma_1 \fCenter \Delta_1, !C_1 \land !C_2$
    \DisplayProof
    \]
    અગાઉની જેમ
    $\Lang L(!C_1 \land !C_2) \subseteq \Lang
    L(\Gamma_1, \Delta_1) \cap
    \Lang L(\Gamma_2, \Delta_2, !A \lif !B)$.

    \item $\pi$નો અંત~$\RightR{\lforall}$થી
    થાય છે અને મુખ્ય સૂત્ર
    ~$\lforall[x][!A(x)]$ છે:
    \[
    \AxiomC{}
    \RightLabel{$\pi_1$}
    \Deduce$\maeh{\Gamma_1}{\Gamma_2} \fCenter
    \maeh{\Delta_1, !A(c)}{\Delta_2}$
    \RightLabel{\RightR{\lforall}}
    \UnaryInf$\maeh{\Gamma_1}{\Gamma_2} \fCenter
      \maeh{\Delta_1, \lforall[x][!A(x)]}{\Delta_2}$
    \DisplayProof
    \]
    અથવા
    \[
    \AxiomC{}
    \RightLabel{$\pi_1$}
    \Deduce$\maeh{\Gamma_1}{\Gamma_2} \fCenter
      \maeh{\Delta_1}{\Delta_2, !A(c)}$
    \RightLabel{\RightR{\lforall}}
    \UnaryInf$\maeh{\Gamma_1}{\Gamma_2} \fCenter
      \maeh{\Delta_1}{\Delta_2, \lforall[x][!A(x)]}$
    \DisplayProof
    \]
    પહેલા કિસ્સામાં આગમન-પૂર્વધારણાથી
    $\Gamma_1 \Sequent \Delta_1, !A(c), !C_1$
    અને $!C_1, \Gamma_2 \Sequent \Delta_2$ની
    સાબિતીઓ મળે છે. પહેલા સિક્વન્ટ પર
    $\RightR{\lforall}$ લગાવતાં
    $\Gamma_1 \Sequent \Delta_1,
    \lforall[x][!A(x)], !C_1$ મળે છે.
    (સ્વચલ-શરત પૂરી છે, કારણ કે~$\pi$ના
    અંતિમ \RightR{\lforall}માં તે
    પહેલેથી પૂરી હતી.)
    તેથી
    $\Lang L(!C_1) \subseteq
    \Lang L(\Gamma_1, \Delta_1, \lforall[x][!A(x)])
    \cap \Lang L(\Gamma_2, \Delta_2)$
    હોય તો~$!C_1$ અંતર્વેશક છે.
    આગમન-પૂર્વધારણા
    $\Lang L(!C_1) \subseteq
    \Lang L(\Gamma_1, \Delta_1, !A(c))
    \cap \Lang L(\Gamma_2, \Delta_2)$ આપે છે.
    એટલે~$c$ અંગે ચિંતા છે:
    શું~$c$ એ~$!C_1$માં આવી શકે?
    ના. જો~$c$ એ~$!C_1$માં આવે,
    તો $c \in \Lang L(\Gamma_1, \Delta_1, !A(c))$
    (જે સાચું છે) અને
    $c \in \Lang L(\Gamma_2, \Delta_2)$
    બંને થશે. પરંતુ ત્યારે~$c$
    એ~$\pi$ના $\RightR{\lforall}$
    અનુમાનના નિષ્કર્ષમાં આવશે,
    જે સ્વચલ-શરતનો ભંગ કરે છે.
    
    બીજો કિસ્સો સમાન છે.

    \item $\pi$નો અંત~$\RightR{\lexists}$થી
    થાય છે અને મુખ્ય
    સૂત્ર~$\lexists[x][!A(x)]$ છે.
    ફરી બે કિસ્સા છે.
    $\lexists[x][!A(x)]$ પહેલા ભાગમાં
    હોય તે કિસ્સો જોઈએ;
    બીજો કસરત તરીકે છોડીએ છીએ.

    પહેલા કિસ્સામાં~$\pi$નો અંત
    આ રીતે થાય છે:
    \[
    \AxiomC{}
    \RightLabel{$\pi_1$}
    \Deduce$\maeh{\Gamma_1}{\Gamma_2} \fCenter
      \maeh{\Delta_1, \lexists[x][!A(x)], !A(t)}{\Delta_2}$
    \RightLabel{\RightR{\lexists}}
    \UnaryInf$\maeh{\Gamma_1}{\Gamma_2} \fCenter
      \maeh{\Delta_1, \lexists[x][!A(x)]}{\Delta_2}$
    \DisplayProof
    \]
    આગમન-પૂર્વધારણાથી
    $\Gamma_1 \Sequent
    \Delta_1, \lexists[x][!A(x)], !A(t), !C_1$
    અને $!C_1, \Gamma_2 \Sequent \Delta_2$ની
    સાબિતીઓ મળે છે.
    પહેલા સિક્વન્ટ પર
    $\RightR{\lexists}$ લગાવતાં
    $\Gamma_1 \Sequent \Delta_1,
    \lexists[x][!A(x)], !C_1$ મળે છે.
    તેથી
    $\Lang L(!C_1) \subseteq
    \Lang L(\Gamma_1, \Delta_1,
    \lexists[x][!A(x)])
    \cap \Lang L(\Gamma_2, \Delta_2)$
    હોય તો~$!C_1$ ફરી અંતર્વેશક છે.
    અહીં આગમન-પૂર્વધારણાથી
    $\Lang L(!C_1) \subseteq
    \Lang L(\Gamma_1, \Delta_1,
    \lexists[x][!A(x)], !A(t))
    \cap \Lang L(\Gamma_2, \Delta_2)$ મળે છે.
    યાદ કરો કે ફલનો નથી.
    પદ~$t$ ચલ હોઈ શકે; એ
    ભાષામાં નવું અતાર્કિક ચિહ્ન ઉમેરતું નથી,
    એટલે એ કિસ્સામાં અંતર્વેશક સીધો ચાલે.
    અચલના કિસ્સામાં તે
    અચળ છે, એટલે
    $t \ident b$.
    \sourcecorrection{OLMD-206}{મૂળ લખાણમાં ફલનો ન હોવાથી $t$ માત્ર અચલ છે એવો દાવો છે; ચલ પણ પદ છે, પરંતુ તે અંતર્વેશકની અતાર્કિક શબ્દાવલી વધારતું નથી. નીચેની અચલ-વિભાગની દલીલ યથાવત્ છે.}
    જો $b \notin \Lang L(!C_1)$
    અથવા
    $b \in \Lang L(\Gamma_1, \Delta_1,
    \lexists[x][!A(x)]) \cap
    \Lang L(\Gamma_2, \Delta_2)$,
    તો ચિંતા નથી:~$!C_1$ને
    અંતર્વેશક લઈ શકાય.
    પરંતુ જો
    $b \in \Lang L(!C_1)$ અને
    $b \notin \Lang
    L(\Gamma_1, \Delta_1,
    \lexists[x][!A(x)])
    \cap \Lang L(\Gamma_2, \Delta_2)$ હોય તો?
    પહેલાં~$!C_1$માં~$b$ના
    દરેક દૃષ્ટાંતને ચલ~$y$થી
    બદલી
    $!C_1 \ident \Subst{!C_1'}{b}{y}$
    લખી શકાય, જ્યાં~$!C_1'$માં~$b$
    આવતો નથી. વધુમાં, કાં તો
    $b \notin
    \Lang L(\Gamma_1, \Delta_1,
    \lexists[x][!A(x)])$
    અથવા
    $b \notin \Lang L(\Gamma_2, \Delta_2)$.
    અનુક્રમે
    $!C \ident \lforall[y][!C_1'(y)]$
    અથવા
    $!C \ident \lexists[y][!C_1'(y)]$
    લઈ શકાય. પહેલા કિસ્સામાં:
    \[
    \AxiomC{}
    \Deduce$\Gamma_1 \fCenter
      \Delta_1, \lexists[x][!A(x)], !A(b), !C_1'(b)$
    \RightLabel{\RightR{\lexists}}
    \UnaryInf$\Gamma_1 \fCenter \Delta_1, \lexists[x][!A(x)], !C_1'(b)$
    \RightLabel{\RightR{\lforall}}
    \UnaryInf$\Gamma_1 \fCenter
      \Delta_1, \lexists[x][!A(x)], \lforall[y][!C_1'(y)]$
    \DisplayProof
    \]
    અને
    \[
    \AxiomC{}
    \Deduce$!C_1'(b), \lforall[y][!C_1'(y)], \Gamma_2 \fCenter \Delta_2$
    \RightLabel{\LeftR{\lforall}}
    \UnaryInf$\lforall[y][!C_1'(y)], \Gamma_2 \fCenter \Delta_2$
    \DisplayProof
    \]
    ઉપરના સિક્વન્ટની સાબિતીઓ
    આગમન-પૂર્વધારણાથી મળે છે
    (બીજામાં પૂર્વાંગમાં
    $\lforall[y][!C_1'(y)]$
    ઉમેરવા શિથિલીકરણની સ્વીકાર્યતા
    વાપરીએ છીએ).
    પહેલી સાબિતીમાં
    \RightR{\lforall}ની સ્વચલ-શરત
    પૂરી થાય છે: કારણ કે
    $b \notin \Lang L(\Gamma_1, \Delta_1,
    \lexists[x][!A(x)])$ અને
    $!C_1'$ની વ્યાખ્યા એવી છે કે
    તેમાં~$b$ નથી.
    
    બીજા કિસ્સામાં આગમન-પૂર્વધારણા
    અને~$\lexists[y][!C_1'(y)]$
    ઉમેરતા શિથિલીકરણની સ્વીકાર્યતા વડે:
    \[
    \AxiomC{}
    \Deduce$\Gamma_1 \fCenter
      \Delta_1, \lexists[x][!A(x)], !A(b), \lexists[y][!C_1'(y)], !C_1'(b)$
    \RightLabel{\RightR{\lexists}}
    \UnaryInf$\Gamma_1 \fCenter
      \Delta_1, \lexists[x][!A(x)], \lexists[y][!C_1'(y)], !C_1'(b)$
    \RightLabel{\RightR{\lexists}}
    \UnaryInf$\Gamma_1 \fCenter
      \Delta_1, \lexists[x][!A(x)], \lexists[y][!C_1'(y)]$
    \DisplayProof
    \]
    અને
    \[
    \AxiomC{}
    \Deduce$!C_1'(b), \Gamma_2 \fCenter \Delta_2$
    \RightLabel{\LeftR{\lexists}}
    \UnaryInf$\lexists[y][!C_1'(y)], \Gamma_2 \fCenter \Delta_2$
    \DisplayProof
    \]
    બીજી સાબિતીમાં
    \LeftR{\lexists}ની સ્વચલ-શરત
    પૂરી થાય છે, કારણ કે
    $b \notin \Lang L(\Gamma_2, \Delta_2)$
    અને~$!C_1'$ની વ્યાખ્યા એવી છે
    કે તેમાં~$b$ નથી.
  \end{enumerate}
  બીજા કિસ્સા સમાન છે;
  તેમને કસરત તરીકે છોડીએ છીએ.
\end{proof}

\begin{prob}
માનો કે~$\lnand$ સંયોજક માટે
આ નિયમો હોય:
\[
\Axiom$\Gamma \fCenter \Delta, !A$
\Axiom$\Gamma \fCenter \Delta, !B$
\RightLabel{\LeftR{\lnand}}
\BinaryInf$!A \lnand !B, \Gamma \fCenter \Delta$
\DisplayProof
\quad
\Axiom$!A, !B, \Gamma \fCenter \Delta$
\RightLabel{\RightR{\lnand}}
\UnaryInf$\Gamma \fCenter \Delta, !A \lnand !B$
\DisplayProof
\]
\olref{lem:maehara}ની સાબિતીમાં
જ્યાં~$\pi$નો અંત આમાંથી કોઈ
નિયમથી થાય છે તે કિસ્સા કરીને
માએહારાનું ઉપપ્રમેય હજી સાચું છે
એ બતાવો.
\end{prob}

માએહારાનું ઉપપ્રમેય મળ્યું છે;
હવે તે વડે ક્રેગનું અંતર્વેશન
પ્રમેય સાબિત કરીએ.

\begin{proof}[\olref{thm:interpolation}ની સાબિતી]
  માનો કે $!A \Sequent !B$
  સાબિતીક્ષમ છે.
  $\maeh{!A}{\ } \Sequent \maeh{\ }{!B}$
  પર \olref{lem:maehara} લગાવતાં
  અંતર્વેશક~$!C$ મળે છે.
  આપણને
  $\Proves !A \Sequent !C$ અને
  $\Proves !C \Sequent !B$ મળે છે.
\end{proof}

માનો કે~$\Gamma$ એ~$n$-સ્થાની
વિધેય~$R$ ધરાવતી
ભાષા~$\Lang L$નો
સિદ્ધાંત (વાક્યોનો ગણ) છે.
~$\Gamma$ એ~$R$ને
\emph{ગૂઢ રીતે વ્યાખ્યાયિત કરે છે}
ત્યારે અને માત્ર ત્યારે જ જ્યારે
\begin{align*}
\Gamma, \Gamma'
& \Entails \lforall[x_1][\dots\lforall[x_n][(R(x_1,\dots,x_n) \liff
R'(x_1, \dots, x_n))]],
\intertext{અહીં $\Gamma'$ એ $\Gamma$માં $R$ના બધા દૃષ્ટાંતોને $R' \notin \Lang L$થી બદલીને મળ્યું છે. $\Gamma$ એ $R$ને \emph{સ્પષ્ટ રીતે વ્યાખ્યાયિત કરે છે} ત્યારે અને માત્ર ત્યારે જ જ્યારે $R$ ન ધરાવતું એવું $!A(x_1, \dots, x_n)$ છે કે}
\Gamma &\Entails
\lforall[x_1][\dots\lforall[x_n][(R(x_1,\dots,x_n) \liff !A(x_1,
\dots, x_n))]].
\end{align*}
સ્પષ્ટ છે કે જો~$\Gamma$ એ~$R$ને
સ્પષ્ટ રીતે વ્યાખ્યાયિત કરે,
તો ગૂઢ રીતે પણ કરે છે.
~$R$ માટે વ્યસ્ત દિશા તરત સ્પષ્ટ નથી,
પણ તે પણ સાચી છે:

\begin{lem}[બેથનું વ્યાખ્યેયતા પ્રમેય]
  જો~$\Gamma$ એ~$R$ને ગૂઢ રીતે
  વ્યાખ્યાયિત કરે, તો સ્પષ્ટ
  રીતે પણ વ્યાખ્યાયિત કરે છે.
\end{lem}

\begin{proof}
  માનો કે~$\Gamma$ એ~$R$ને ગૂઢ રીતે
  વ્યાખ્યાયિત કરે છે.
  ઉપર મુજબ~$R'$ અને~$\Gamma'$
  લો તથા
  $\Lang L' = \Lang L(\Gamma') =
  \Lang L \setminus \{R\}
  \cup \{R'\}$.
  પૂર્વધારણાથી મળે છે:
  \begin{align*}
    \Gamma, \Gamma' & 
    \Sequent \lforall[x_1][\dots\lforall[x_n][(R(x_1,\dots,x_n) 
    \liff R'(x_1, \dots, x_n))]]
  \intertext{$c_1$, \dots,~$c_n$ એ $\Gamma$માં ન આવતાં અચળો હો. ઉલટાવી શકવાના ઉપપ્રમેયથી મળે છે:}
    \Gamma, \Gamma', R(c_1,\dots,c_n) & \Sequent  R'(c_1, \dots, c_n)
  \intertext{માએહારાનું ઉપપ્રમેય આના પર લગાવો:}
    \maeh{\Gamma, R(c_1,\dots,c_n)}{\Gamma'} & \Sequent
      \maeh{\ }{R'(c_1, \dots, c_n)}
  \intertext{ત્યારે એવું $!A \ident !A(x_1, \dots, x_n)$ મળે છે કે}
    \Gamma, R(c_1,\dots,c_n) & \Sequent !A(c_1, \dots, c_n) \\
    !A(c_1, \dots, c_n), \Gamma' & \Sequent R'(c_1, \dots, c_n)
  \intertext{આ બંનેની સાબિતીઓ છે. માએહારાના ઉપપ્રમેય મુજબ $\Lang L(!A) \subseteq \Lang L(\Gamma) \setminus \{R\}$; એટલે $!A$માં $R$ નથી. બીજા સિક્વન્ટની સાબિતીમાં $R'$ને સર્વત્ર $R$થી બદલો; ત્યારે આની સાબિતી મળે છે:}
    !A(c_1, \dots, c_n), \Gamma & \Sequent R(c_1, \dots, c_n)
  \end{align*}
  \sourcecorrection{OLMD-207}{મૂળમાં અંતર્વેશકની ભાષા સામાન્ય ભાષાના બરાબર છે અને અપરિભાષિત $L_1,L_2$ વાપરે છે; ઉપપ્રમેય માત્ર સમાવેશ આપે છે, જે અહીં પૂરતો છે.}
  બંને દિશાના અનુસરણો માટે
  \RightR{\lif}, પછી દ્વિમુખી
  સંબંધની રજૂઆત અને
  \RightR{\lforall} લગાવી
  આની સાબિતી મેળવો:
  \sourcecorrection{OLMD-208}{મૂળમાં માત્ર અનુસરણ-જમણો અને સર્વમાન્યક-જમણો નિયમ લખેલા છે; બે દિશા ભેગી કરીને દ્વિમુખી સંબંધ બનાવવાનું મધ્યવર્તી પગલું પણ જરૂરી છે.}
  \[
    \Gamma \Sequent \lforall[x_1][\dots\lforall[x_n][(R(x_1,\dots,x_n) \liff !A(x_1, \dots, x_n))]].
  \]
  આ બતાવે છે કે
  $!A(x_1, \dots, x_n)$
  એ~$\Gamma$માં~$R$ને
  સ્પષ્ટ રીતે વ્યાખ્યાયિત કરે છે.
  \sourcecorrection{OLMD-209}{મૂળ અંતિમ વાક્યમાં $!R$ લખાયું છે, જ્યારે $R$ અહીં વ્યાખ્યાયિત થતું વિધેયચિહ્ન છે.}
\end{proof}

અંતર્વેશનને સમકક્ષ ત્રીજું પરિણામ
(જેને પણ માએહારાના ઉપપ્રમેયથી
સાબિત કરી શકાય) આ છે:
જો બે સુસંગત સિદ્ધાંતો~$\Gamma_1$
અને~$\Gamma_2$માં તેમના સામાન્ય
ભાષાનો પૂર્ણ સિદ્ધાંત~$\Gamma$
સમાયેલો હોય, એટલે
$\Lang L(\Gamma) =
\Lang L(\Gamma_1) \cap \Lang L(\Gamma_2)$,
તો $\Gamma_1 \cup \Gamma_2$
સુસંગત છે.
\sourcecorrection{OLMD-210}{મૂળમાં $\Lang L(\Gamma)$ને માત્ર સામાન્ય ભાષાનો ઉપગણ કહે છે; આગળની પૂર્ણતાની દલીલ દરેક સામાન્ય-ભાષાના વાક્ય પર ચાલે છે, તેથી અહીં ભાષાઓની સમાનતા જરૂરી છે.}

યાદ કરો: પૂર્ણ સિદ્ધાંત એવો છે કે તેની
ભાષાના દરેક વાક્ય~$!A$ માટે
કાં તો $\Gamma \Proves !A$ અથવા
$\Gamma \Proves \lnot !A$.
કારણ કે~$\Gamma_1$ અને~$\Gamma_2$
એ~$\Gamma$ના સુસંગત વિસ્તારો છે,
$\Gamma$ પણ સુસંગત છે.
એથી જો~$!A$ એ
$\Lang L(\Gamma_1) \cap
\Lang L(\Gamma_2)$ની
વાક્ય હોય, તો
$\Gamma_1 \Entails !A$ અને
$\Gamma_2 \Entails \lnot !A$
બંને હોઈ શકે નહીં.
વિરોધ માટે માનો કે એવું~$!A$ છે.
$\Gamma_1 \Entails !A$ છે;
આથી પૂર્ણતા મુજબ
$\Gamma \Entails !A$ જ હોવું જોઈએ.
\sourcecorrection{OLMD-211}{મૂળમાં $\Gamma_1\Entails !A$ પરથી સીધું $\Gamma\Entails !A$ કાઢ્યું છે; નાનું સિદ્ધાંત મોટાના બધા પરિણામો આપતું નથી. પૂર્ણતા અને $\Gamma_1$ની સુસંગતતા વડે જ આ નિષ્કર્ષ મળે છે.}
નહીં તો~$\Gamma$ની પૂર્ણતાથી
$\Gamma \Entails \lnot !A$,
અને $\Gamma \subseteq \Gamma_1$
હોવાથી
$\Gamma_1 \Entails \lnot !A$ પણ,
જે~$\Gamma_1$ને અસુસંગત બનાવે.
તેથી $\Gamma \Entails !A$;
અને $\Gamma \subseteq \Gamma_2$
હોવાથી
$\Gamma_2 \Entails !A$.
પણ પૂર્વધારણા મુજબ
$\Gamma_2 \Entails \lnot !A$ પણ છે,
જે~$\Gamma_2$ની સુસંગતતા સામે છે.
\sourcecorrection{OLMD-212}{મૂળ અંતિમ સૂત્રમાં $\Gamma\Entails/\lnot !A$ લખ્યું છે; વિરોધ ખરેખર $\Gamma_2$માં $!A$ અને $\lnot !A$ બંને ફલિત થવાથી મળે છે.}

સાબિતી માટે એટલું જ જરૂરી છે કે
સામાન્ય
ભાષા~$\Lang L(\Gamma_1)
\cap \Lang L(\Gamma_2)$માં
એવું કોઈ વાક્ય~$!A$ ન હોય
જેના માટે
$\Gamma_1 \Entails !A$ અને
$\Gamma_2 \Entails \lnot !A$.
હવે માએહારાના ઉપપ્રમેય વડે
સાબિતી-સૈદ્ધાંતિક સાબિતી આપીએ.

\begin{thm}[રૉબિન્સનનું સંયુક્ત સુસંગતતા પ્રમેય]
માનો કે~$\Gamma_1$ અને~$\Gamma_2$
સુસંગત સિદ્ધાંતો છે, અને
ભાષા
$\Lang L(\Gamma_1) \cap
\Lang L(\Gamma_2)$માં એવું
કોઈ~$!A$ નથી કે
$\Gamma_1 \Entails !A$ અને
$\Gamma_2 \Entails \lnot !A$.
તો~$\Gamma_1 \cup \Gamma_2$
સુસંગત છે.
\end{thm}

\begin{proof}
  વિરોધ માટે માનો કે
  $\Gamma_1 \cup \Gamma_2$
  અસુસંગત છે. ત્યારે
  $\Gamma_1, \Gamma_2 \Sequent \ $
  સિક્વન્ટની સાબિતી છે.
  $\maeh{\Gamma_1}{\Gamma_2}
  \Sequent \maeh{\ }{\ }$
  પર માએહારાનું ઉપપ્રમેય
  લગાવતાં
  ભાષા~$\Lang L(\Gamma_1)
  \cap \Lang L(\Gamma_2)$માં
  એવી વાક્ય~$!A$ મળે છે કે
  $\Proves \Gamma_1 \Sequent !A$
  અને
  $\Proves !A, \Gamma_2 \Sequent \quad$.
  બીજા સિક્વન્ટથી
  $\Proves \Gamma_2 \Sequent \lnot !A$
  મળે છે. પરંતુ એવું કોઈ
  વાક્ય નથી એમ ધારેલું છે.
\end{proof}

ઉપર સિદ્ધાંતો~$\Gamma$,
$\Gamma_1$ અને~$\Gamma_2$
સાન્ત છે એવું ગર્ભિત રીતે
ધાર્યું છે. સઘનતાથી પરિણામો
અનંત સિદ્ધાંતો માટે પણ સાચાં છે.
% END OLP-0658
\endgroup


\OLEndChapterHook


\begin{editorial}
આગળનો સંબંધિત અંશ મૂળના સક્રિય આયાતક્રમમાં નથી. તેને અહીં સ્પષ્ટ વૈકલ્પિક સંદર્ભમાં જાળવ્યો છે.
\end{editorial}
\begingroup
\makeatletter\def\ol@sectioncs{section}\makeatother

% BEGIN OLP-0660
\phantomsection\label{guunit:OLP-0660}
\olfileid{pt}{cut}{alt0660}
જો $!A \ident (!B \lif !C)$ હોય,
તો~$\pi$નો અંત આ રીતે થાય છે:
\[
\AxiomC{}
\RightLabel{$\pi_1'$}
\Deduce$!B, \Gamma \fCenter \Delta, !C$
\RightLabel{\RightR{\lif}}
\UnaryInf$\Gamma \fCenter \Delta, !B \lif !C$
\LeftSubproofLabel{$\pi_1$}
\AxiomC{}
\RightLabel{$\pi_2'$}
\Deduce$!B \lif !C, \Pi \fCenter \Lambda, !B$
\AxiomC{}
\RightLabel{$\pi_2''$}
\Deduce$!C, \Pi \fCenter \Lambda$
\RightLabel{\LeftR{\lif}}
\BinaryInf$!B \lif !C, \Pi \fCenter \Lambda$
\RightSubproofLabel{$\pi_2$}
\RightLabel{\Cut}
\BinaryInf$\Gamma, \Pi \fCenter \Delta, \Lambda$
\DisplayProof
\]
\Cut{}માં પૂરી થતી આ સાબિતી લો:
\[
\AxiomC{}
\RightLabel{$\pi_1'$}
\Deduce$!B, \Gamma \fCenter \Delta, !C$
\RightLabel{\RightR{\lif}}
\UnaryInf$\Gamma \fCenter \Delta, !B \lif !C$
\LeftSubproofLabel{$\pi_1$}
\AxiomC{}
\RightLabel{$\pi_2'$}
\Deduce$!B \lif !C, \Pi \fCenter \Lambda, !B$
\RightLabel{\Cut}
\BinaryInf$\Gamma, \Pi \fCenter \Delta, \Lambda, !B$
\DisplayProof
\]
તેની કટ-ઊંચાઈ~$\pi$ કરતાં
ઓછી છે. આગમન-પૂર્વધારણા
$\Gamma, \Pi \fCenter \Delta, \Lambda, !B$ની
સાબિતી~$\pi_1''$ આપે છે.
હવે \Cut{}માં પૂરી થતી આ સાબિતી લો:
\[
\AxiomC{}
\RightLabel{$\pi_1'$}
\Deduce$!B, \Gamma \fCenter \Delta, !C$
\AxiomC{}
\RightLabel{$\pi_2''$}
\Deduce$!C, \Pi \fCenter \Lambda$
\RightLabel{\Cut}
\BinaryInf$!B, \Gamma, \Pi \fCenter \Delta, \Lambda$
\DisplayProof
\]
\sourcecorrection{OLMD-215}{મૂળ ત્રીજા વૃક્ષમાં સૂત્ર $B!$ લખાયું છે; સમગ્ર કિસ્સામાં તે $!B$ છે.}
તેનો કટ-ક્રમ અને કટની ઊંચાઈ
બંને~$\pi$ના કટ-ક્રમ અને
કટની ઊંચાઈ કરતાં
ઓછાં છે. તેથી આગમન-પૂર્વધારણાથી
$!B, \Gamma, \Pi \fCenter \Delta, \Lambda$ની
\Log{G3c}-સાબિતી~$\pi_2'''$ મળે છે.
\sourcecorrection{OLMD-216}{મૂળ આગમન-નિષ્કર્ષમાં $\Pi$ અને $\Lambda$ની વધારાની નકલો લખાઈ છે; ત્રીજા વૃક્ષનો નિષ્કર્ષ એક-એક નકલવાળો છે.}
હવે \Cut{}માં પૂરી થતી આ સાબિતી
પર આગમન-પૂર્વધારણા લગાવો:
\[
\AxiomC{}
\RightLabel{$\pi_1''$}
\Deduce$\Gamma, \Pi \fCenter \Delta, \Lambda, !B$
\AxiomC{}
\RightLabel{$\pi_2'''$}
\Deduce$!B, \Gamma, \Pi \fCenter \Delta, \Lambda$
\RightLabel{\Cut}
\BinaryInf$\Gamma, \Gamma, \Pi, \Pi \fCenter \Delta, \Delta, \Lambda, \Lambda$
\DisplayProof
\]
\sourcecorrection{OLMD-217}{મૂળ અંતિમ વૃક્ષનું જમણું આધારવિધાન અગાઉના આગમન-નિષ્કર્ષથી જુદું હતું; અહીં $\pi_2'''$નું એ જ સિક્વન્ટ મૂક્યું છે.}
\sourcecorrection{OLMD-218}{મૂળ અંતિમ વૃક્ષનો નિષ્કર્ષ કટ-નિયમના સંદર્ભ-સંયોજન સાથે બંધબેસતો નથી; બંને બાજુની સંદર્ભ-નકલો સ્પષ્ટ કરી છે.}
(તેનો કટ-ક્રમ $< \cutrank{\pi}$ છે.)
\sourcecorrection{OLMD-219}{મૂળમાં અપરિભાષિત $\cutr{\pi}$ છે; અગાઉના વિભાગની વ્યાખ્યા પ્રમાણે $\cutrank{\pi}$ છે.}
તેથી
$\Gamma, \Gamma, \Pi, \Pi
\fCenter \Delta, \Delta, \Lambda, \Lambda$ની
સાબિતી મળે છે.
\sourcecorrection{OLMD-220}{મૂળ વર્ણનમાં ત્રણ $\Pi$, ત્રણ $\Lambda$ અને એક $\Delta$ હતાં; સુધારેલા કટ-વૃક્ષના નિષ્કર્ષમાં દરેક સંદર્ભની બે નકલો છે.}
\olref[seq][inv]{prop:G3c-cont-adm}
લગાવતાં
$\Gamma, \Pi \fCenter \Delta, \Lambda$ની
સાબિતી~$\pi'$ મળે છે.
% END OLP-0660
\endgroup
% END OLP-0657
\endgroup

\begingroup
\makeatletter\def\ol@sectioncs{section}\makeatother

% BEGIN OLP-0664
\olchapter{pt}{nat}{સ્વાભાવિક નિગમન}
\olchapterid{pt}{nat}
\phantomsection\label{guunit:OLP-0664}


\begingroup
\makeatletter\def\ol@sectioncs{section}\makeatother

% BEGIN OLP-0663
\olfileid{pt}{nat}{int}

\olsection{પરિચય}
\phantomsection\label{guunit:OLP-0663}


સ્વાભાવિક નિગમન એવાં સાબિતી
તંત્રોનો પરિવાર છે જેમાં દરેક
તાર્કિક સંયોજક અથવા પરિમાણક
માટે અનુમાનોની એક જોડ હોય છે:
એક અનુમાન એ સંક્રિયકનો
``પરિચય'' કરાવે છે
(દાખલ નિયમ)
અને બીજું તેનો
``નિવારણ'' કરે છે
(લોપ નિયમ).
દાખલા તરીકે, \Intro{\lif}માં
$!A \lif !B$ સ્વરૂપનાં
સૂત્રો નિષ્કર્ષરૂપે હોય છે,
જ્યારે~\Elim{\lif}માં
$!A \lif !B$ સ્વરૂપનાં
સૂત્રો
આધારવિધાનરૂપે હોય છે.

સ્વાભાવિક નિગમનનાં પ્રથમ બે
રૂપો ૧૯૩૦ના દાયકામાં ગેરહાર્ડ
ગેન્ટ્ઝન અને સ્તાનિસ્વાવ
યાશ્કોવ્સ્કીએ રજૂ કર્યાં હતાં.
સાબિતી-સિદ્ધાંતમાં ગેન્ટ્ઝનનું
તંત્ર વધુ ચર્ચાય છે, જ્યારે
શિક્ષણમાં સૌથી વધુ વપરાતાં
તંત્રો યાશ્કોવ્સ્કીની રીતમાંથી
વિકસ્યાં છે. બંનેમાં
\emph{ધારણાઓ}થી શરૂઆત
કરી શકાય છે: એવાં
સૂત્રો જેમને
સાબિત કરવાં જરૂરી
નથી, પણ જે બીજા
સૂત્રોને સાબિત
કરવા વાપરી શકાય છે.
આખી સાબિતી આવી
ધારણાઓ પર આધારિત હોઈ
શકે છે. તે ધારણાઓ
અસાબિત હોવાથી સાબિતી
પોતાના નિષ્કર્ષને
(પોતાના \emph{અંતિમ-સૂત્ર}ને)
સ્થાપિત કરતી નથી;
માત્ર એટલું સ્થાપિત
કરે છે કે નિષ્કર્ષ
ધારણાઓમાંથી ફલિત થાય છે.

કેટલાક અનુમાન-નિયમોના
નિષ્કર્ષો પછી કેટલીક ધારણાઓ
પર આધારિત રહેતા નથી:
નિયમો એ ધારણાઓને
``મુક્ત'' કરે છે.
દાખલા તરીકે,
\Intro{\lif} નિયમ
ધારણા~$!A$માંથી
નિષ્કર્ષ~$!B$ની
સાબિતી લે છે અને
$!A \lif !B$ નિષ્કર્ષ આપે છે.
$!A \lif !B$માં પૂરી થતી
સાબિતી હવે~$!A$ પર
આધારિત નથી; ધારણા~$!A$
મુક્ત થાય છે.
ગેન્ટ્ઝનના તંત્રોમાં
સાબિતીઓ એ
સૂત્રોનાં વૃક્ષો છે;
ધારણાઓ સૌથી ઉપરનાં
સૂત્રો છે;
અને~\Intro{\lif} જેવા
નિયમો પરનાં ચિહ્નો કઈ
ધારણાઓ મુક્ત થાય છે
તે બતાવે છે.
યાશ્કોવ્સ્કીના તંત્રોમાં
ધારણા પર આધારિત
સાબિતીનો ભાગ અલગ
દેખાડાય છે, દાખલા તરીકે,
ઊભી રેખાની પાછળ
ગોઠવીને અથવા ચોકઠામાં
મૂકીને. એ ચોકઠાની
પહેલી પંક્તિમાં મુક્ત
થતી ધારણા~$!A$ અને
છેલ્લી પંક્તિમાં
અનુસરણનું પરિણામ~$!B$
હોય છે; \Intro{\lif}નો
નિષ્કર્ષ~$!A \lif !B$
ચોકઠાની બહાર હોય છે.

આ તંત્રો ``સ્વાભાવિક''
છે, કારણ કે તેમનાં
અનુમાન-નિયમો આપણે
(ઓછામાં ઓછા ગણિતશાસ્ત્રીઓ)
સ્વાભાવિક રીતે કેવી રીતે
તર્ક કરીએ છીએ તે દર્શાવે છે.
શરતી પરિણામ સ્થાપવા માટે
આપણે તેનું પૂર્વપદ
ધારીને પરિણામપદ સાબિત
કરીએ છીએ: એ~\Intro{\lif}.
$!A \lif !B$ અને તેનું
પૂર્વપદ~$!A$ જાણીએ
ત્યારે~$!B$ તારવીએ છીએ:
એ~\Elim{\lif}.
વિકલ્પ~$!A \lor !B$
ધારતાં કોઈ નિષ્કર્ષ
મળે છે તે કિસ્સાવાર
સાબિતીથી બતાવીએ છીએ:
એ~\Elim{\lor}.
સામાન્ય દાવો
$\lforall[x][!A(x)]$
સાબિત કરવા કોઈ મનસ્વી~$c$
માટે~$!A(c)$ બતાવીએ છીએ:
એ~\Intro{\lforall}.
અને આવી રીતે આગળ.

દાખલા તરીકે, ગેન્ટ્ઝનની
શૈલીના સ્વાભાવિક નિગમનમાં
$!A \lif (!A \lor !B)$ની
આ સરળ સાબિતી જુઓ:
\[
\AxiomC{$\Discharge{!A}{x}$}
\RightLabel{\Intro{\lor}}
\UnaryInfC{$!A \lor !B$}
\DischargeRule{\Intro{\lif}}{x}
\UnaryInfC{$!A \lif (!A \lor !B)$}
\DisplayProof
\]
તે ધારણા~$!A$થી શરૂ થાય છે,
\Intro{\lor}થી~$!A
\lor !B$ તારવે છે,
અને પછી~\Intro{\lif}થી
$!A \lif (!A \lor !B)$
તારવે છે. \Intro{\lif}
અનુમાન ધારણા~$!A$ને
મુક્ત કરે છે; ચિહ્ન~$x$
એ દર્શાવે છે.

સાબિતી-સિદ્ધાંતકારો
સૂત્રોનાં વૃક્ષો
વાપરતાં ગેન્ટ્ઝનનાં મૂળ
તંત્રોને પસંદ કરે છે,
છતાં એક અન્ય રૂપ પણ
સાબિતી-સિદ્ધાંતમાં
મહત્ત્વનું છે.
ધારણાઓ~$\Gamma$માંથી
$!A$ની સાબિતી
બતાવે છે કે~$!A$ એ
$\Gamma$માંથી ફલિત થાય છે,
અર્થાત્~$\Gamma \Entails !A$.
સ્વાભાવિક નિગમનના
અનુમાન-નિયમોને આ પ્રકારની
ફલિતતાઓ વિશેના નિયમો
તરીકે પણ વાંચી શકાય છે.
દાખલા તરીકે,
\Intro{\lif} કહે છે કે
જો $\Gamma + !A
\Entails !B$, તો
$\Gamma \Entails !A \lif !B$.
\sourcecorrection{OLMD-227}{મૂળમાં અનુસરણ-પરિચયનું નિષ્કર્ષ $\Gamma \Entails !B$ લખાયું છે; ધારણા $!A$ મુક્ત કર્યા પછી યોગ્ય નિષ્કર્ષ $\Gamma \Entails !A \lif !B$ છે.}
આથી નિયમોને સીધા
એ રીતે ઘડવું સ્વાભાવિક છે
કે તેઓ આધારવિધાન અને
નિષ્કર્ષ બંનેમાં
ધારણાઓ તથા તેમનામાંથી
મળતા પરિણામ પર કામ કરે;
અર્થાત્ તેઓ
$\Gamma \Sequent !A$
સિક્વન્ટો પર કામ કરે.
ઉપરની સરળ સાબિતી
આવા તંત્રમાં આ રીતે
દેખાય:
\[
\Axiom$x:!A \fCenter !A$
\RightLabel{\Intro{\lor}}
\UnaryInf$x:!A \fCenter !A \lor !B$
\DischargeRule{\Intro{\lif}}{x}
\UnaryInf$\fCenter !A \lif (!A \lor !B)$
\DisplayProof
\]
અહીં નિયમો એ જ છે:
તેઓ~$\Sequent$ની જમણી
બાજુના સૂત્ર પર કામ કરે છે;
ડાબી બાજુનાં સૂત્રો
દરેક પગલે આધારરૂપ
ધારણાઓની યાદી આપે છે.

દાખલ/લોપ
નિયમોની બધી જોડો
એકલી મળીને શાસ્ત્રીય
તર્ક માટે પૂર્ણ નથી.
$\lfalse$ને લગતા
બે વધારાના નિયમો
ઉમેરવા પડે છે.
પહેલો ``અંતઃપ્રજ્ઞાવાદી
અસંગતિ-નિયમ''~\FalseInt,
અથવા ``વિસ્ફોટ'' છે:
$\lfalse$માંથી
કંઈ પણ તારી શકાય.
બીજો પરોક્ષ સાબિતીનો
શાસ્ત્રીય સિદ્ધાંત~\FalseCl{}
(અથવા ``શાસ્ત્રીય
અસંગતિ-નિયમ'') છે:
જો~$\lnot !A$માંથી
$\lfalse$ સાબિત
કરી શકાય, તો~$!A$
સાબિત થયું છે.
જો~\FalseCl કાઢી નાખીએ,
તો અંતઃપ્રજ્ઞાવાદી
તર્ક માટે યોગ્ય તંત્ર મળે છે.
બંને કાઢી નાખીએ,
તો ``લઘુતમ તર્ક''
નામે ઓળખાતું તંત્ર મળે છે.

હવે શાસ્ત્રીય અને
અંતઃપ્રજ્ઞાવાદી તર્ક માટે
ગેન્ટ્ઝનની શૈલીનું
સ્વાભાવિક નિગમન જોઈશું.
આ તંત્રો~\Log{N1c}
અને~\Log{N1i} છે.
એકલાં સૂત્રો~$!A$
ને બદલે સિક્વન્ટો
$\Gamma \Sequent !A$
પર કામ કરતાં અનુરૂપ
તંત્રો~\Log{N2c}
અને~\Log{N2i} છે.
% END OLP-0663
\endgroup

\begingroup
\makeatletter\def\ol@sectioncs{section}\makeatother

% BEGIN OLP-0668
\olfileid{pt}{nat}{rp}

\olsection{નિયમો અને સાબિતીઓ}
\phantomsection\label{guunit:OLP-0668}


પહેલાં કેટલીક વ્યાખ્યાઓ
અને પરિભાષા નક્કી કરીએ.

\Log{N1c} અને~\Log{N1i}માંના
$!A \lif !B$ ધરાવતા
બે નિયમો જુઓ.
\[
\bottomAlignProof
\AxiomC{$\Discharge{!A}{x}$}
\shortDeduce
\DeduceC{$!B$}
\DischargeRule{\Intro{\lif}}{x}
\UnaryInfC{$!A \lif !B$}
\DisplayProof
\qquad
\bottomAlignProof
\AxiomC{$!A \lif !B$}
\AxiomC{$!A$}
\RightLabel{\Elim{\lif}}
\BinaryInfC{$!B$}
\DisplayProof
\]
\Elim{\lif} નિયમ સમજવો
સહેલો છે. રેખાની
ઉપરનાં સૂત્રો
\emph{આધારવાક્યો} છે.
ડાબી બાજુનું સૂત્ર
એ \emph{મુખ્ય}
સૂત્ર~$!A \lif !B$
છે. દરેક
લોપ નિયમમાં
આવું એક આધારવાક્ય હોય છે,
હંમેશાં સૌથી ડાબે.
તેને અનુમાનનું
\emph{મુખ્ય} આધારવાક્ય
કહે છે. આ કિસ્સાની
જેમ બીજાં આધારવાક્યો
હોય, તો તેમને
\emph{ગૌણ} આધારવાક્યો
કહે છે.
દાખલ નિયમોનાં
આધારવાક્યો પણ મુખ્ય
આધારવાક્યો કહેવાય છે.
રેખાની નીચેનું
સૂત્ર એ
\emph{ફલિતવાક્ય} છે.

\Intro{\lif} નિયમને
માત્ર એક આધારવાક્ય છે,
અને અહીં ફલિતવાક્ય
એ મુખ્ય સૂત્ર~$!A \lif !B$
છે. બધાં
દાખલ નિયમો
આવા છે: ફલિતવાક્ય
મુખ્ય સૂત્ર હોય છે
અને તેનો મુખ્ય સંક્રિયક
``દાખલ'' થાય છે.

સંજ્ઞા~$\Discharge{!A}{x}$નો
અર્થ આ છે. કોઈ
સાબિતીમાં
\Intro{\lif} અનુમાન
આધારવાક્યરૂપે
સૂત્ર~$!B$ને
લાગુ પડે છે.
$!B$માં પૂરી થતી
ઉપ-સાબિતીમાં
શરૂઆતનાં કેટલાંક
સૂત્રો હોય છે;
તેમને \emph{ધારણાઓ}
કહે છે.
$!A$ સ્વરૂપની એક
કે વધુ ધારણાઓમાંથી
$!B$ની સાબિતી
એ~$!A$માંથી~$!B$
ફલિત થાય છે તેની
સાબિતી છે.
\Intro{\lif} નિયમ
લગાવીએ ત્યારે
$!A \lif !B$
તારવીએ છીએ,
અને ફલિતવાક્ય
હવે ધારણા~$!A$
પર આધારિત રહેતું નથી.
તે બતાવવા~\Intro{\lif}
અનુમાન ધરાવતી
સાબિતીમાંનાં
$!A$ સ્વરૂપનાં
ધારણાઓને \emph{મુક્ત}
અથવા \emph{બંધ}
કહીએ છીએ.
કઈ ધારણાઓ મુક્ત થઈ
શકે તે નિયમ જણાવે છે.
નિષ્કર્ષ~$!A \lif !B$
ધરાવતા~\Intro{\lif}
માટે એ~$!A$ સ્વરૂપની,
એટલે અનુસરણના પૂર્વપદ
સાથે મળતી, ધારણાઓ છે.
કઈ ધારણા કયા નિયમથી
મુક્ત થાય છે તે
બતાવવું ક્યારેક જરૂરી
હોય છે; તે માટે
મુક્તિ-ચિહ્ન~$x$
વાપરીએ છીએ.
પરંતુ ધારણા મુક્ત
કરતો નિયમ જરૂરી
સ્વરૂપની \emph{બધી}
ધારણાઓ, કે એકેય
ધારણા, મુક્ત કરે
તે \emph{જરૂરી નથી}.
દાખલા તરીકે,
પહેલું~\Intro{\lif}
કોઈ ધારણા મુક્ત ન
કરે તોપણ નીચેની
સાબિતી યોગ્ય છે:
\[
\AxiomC{$\Discharge{!A}{y}$}
\RightLabel{\Intro{\lif}}
\UnaryInfC{$!B \lif !A$}
\DischargeRule{\Intro{\lif}}{y}
\UnaryInfC{$!A \lif (!B \lif !A)$}
\DisplayProof
\]
\sourcecorrection{OLMD-231}{મૂળ વૃક્ષના પ્રથમ અનુસરણ-પરિચય પર $x$ મુક્તિ-ચિહ્ન મૂકાયું છે, જ્યારે તે કોઈ ધારણા મુક્ત કરતું નથી અને તરતનું ગદ્ય ચિહ્ન છોડવાનું કહે છે; અહીં તે લેબલ કાઢ્યું છે.}
કોઈ નિયમ એકેય ધારણા
મુક્ત ન કરે ત્યારે
મુક્તિ-ચિહ્ન છોડીએ છીએ
(જેમ અહીં કર્યું છે).

આ સાબિતીમાં
\[
\AxiomC{$\Discharge{!A}{x}$}
\AxiomC{$\Discharge{!A}{y}$}
\RightLabel{\Intro{\land}}
\BinaryInfC{$!A \land !A$}
\RightLabel{\Elim{\land}}
\UnaryInfC{$!A$}
\DischargeRule{\Intro{\lif}}{x}
\UnaryInfC{$!A \lif !A$}
\DischargeRule{\Intro{\lif}}{y}
\UnaryInfC{$!A \lif (!A \lif !A)$}
\DisplayProof
\]
પહેલું~\Intro{\lif}
અનુમાન ડાબી
ધારણા~$!A^x$ને,
અને બીજું જમણી
ધારણા~$!A^y$ને
મુક્ત કરે છે.
એટલે~$x$ ચિહ્નવાળી
બધી ધારણાઓને
$x$ ચિહ્નવાળું
\Intro{\lif} મુક્ત કરે છે,
અને~$y$ ચિહ્નવાળીઓને
$y$ ચિહ્નવાળું અનુમાન.
આ રીતે કામ કરવા
એક જ ચિહ્નવાળી
બધી ધારણાઓનું
સૂત્ર પણ એક
જ હોવું જરૂરી છે.

પરિમાણકોના નિયમો
થોડા વધુ જટિલ છે.
દાખલા તરીકે,
\Log{N1i}માં
$\lexists$ના નિયમો આ છે:
\[
\bottomAlignProof
\AxiomC{$!A(t)$}
\RightLabel{\Intro{\lexists}}
\UnaryInfC{$\lexists[x][!A(x)]$}
\DisplayProof
\qquad
\bottomAlignProof
\AxiomC{$\lexists[x][!A(x)]$} 
\AxiomC{$\Discharge{!A(c)}{x}$}
\shortDeduce
\DeduceC{$!B$}
\DischargeRule{\Elim{\lexists}}{x}
\BinaryInfC{$!B$} \DisplayProof 
\bottomAlignProof
\]
\Intro{\lexists} નિયમમાં
આધારવાક્ય~$!A(t)$ છે,
જ્યાં~$t$ બંધ પદ છે,
અર્થાત્ ચલ વિનાનું પદ.
આવું અનુમાન ત્યારે યોગ્ય
છે, એટલે કે
\Intro{\lexists}ના
યોજનાત્મક નિયમને
ત્યારે અનુરૂપ છે,
જ્યારે કોઈ
સૂત્ર~$!A$,
કોઈ ચલ~$x$
અને કોઈ બંધ પદ~$t$
માટે નિષ્કર્ષ
$\lexists[x][!A]$
અને આધારવાક્ય
$\Subst{!A}{t}{x}$
હોય.
\Elim{\lexists}થી
$!A(c)$ સ્વરૂપની
ધારણાઓ મુક્ત કરી
શકાય છે:
દરેક અનુમાન માટે
કોઈ અચળ~$c$
હોય છે અને
$\Subst{!A}{c}{x}$
સ્વરૂપની ધારણાઓ
મુક્ત કરી શકાય છે.
એક અનુમાનમાં
મુક્ત થતી બધી
ધારણાઓમાં એ જ~$c$
વપરાવો જોઈએ.
આ અચળને
\Elim{\lexists}
અનુમાનનો
\emph{સ્વનિર્ધારક ચલ}
કહે છે.
આ અનુમાન માત્ર ત્યારે
યોગ્ય છે જ્યારે~$c$
મુક્ત થતી~$!A(c)$
સિવાયની કોઈ ખુલ્લી
ધારણામાં, $!B$માં
કે~$!A$માં ન આવતો હોય.
આને \emph{સ્વનિર્ધારક
ચલની શરત} કહે છે.

\Log{N1c} અને~\Log{N1i}માંની
સાબિતીઓ એ
સૂત્રોનાં સાન્ત વૃક્ષો છે, જે
સ્વયંસિદ્ધો અને
અનુમાન-નિયમોથી
આગમનાત્મક રીતે બને છે.
\sourcecorrection{OLMD-232}{મૂળના આ વાક્યમાં N1-સાબિતીઓને સિક્વન્ટોનાં વૃક્ષો કહેવાય છે, પણ આગળની વ્યાખ્યા અને N1 નિયમ-તકતી પ્રમાણે તે સૂત્રોનાં વૃક્ષો છે; સિક્વન્ટો N2 તંત્રમાં આવે છે.}
દરેક અંતિમ ડાળી
ધારણા છે; તેને
મુક્તિ-ચિહ્ન
લાગેલું હોઈ શકે.
બીજું દરેક
સૂત્ર
તેની તરત ઉપરનાં
સૂત્રોમાંથી
તંત્રના કોઈ
અનુમાન-નિયમથી
મળે છે.
વૃક્ષમાં કોઈ
અનુમાનની ઉપર
ધારણા તરીકે આવતા
સૂત્રને
તેની ઉપરના અનુમાનથી
મુક્ત ન કરાયું હોય,
તો તે અનુમાન સાપેક્ષે
\emph{ખુલ્લું} કહેવાય છે.

\begin{defn}\ollabel{defn:proof}
\Log{N1c} અને~\Log{N1i}માં
\emph{સાબિતીઓ} એ
સૂત્રોનાં
સાન્ત વૃક્ષો છે,
જે આગમનાત્મક રીતે
આ પ્રમાણે બને છે:
\begin{enumerate}
  \item\ollabel{defn:prf:assumption}
  $x$ અથવા~$y$ જેવા
  મુક્તિ-ચિહ્નથી ચિહ્નિત
  દરેક સૂત્ર પોતે
  સાબિતી છે.
  માત્ર એક સૂત્ર
  ધરાવતી સાબિતીમાં
  એ સૂત્ર ખુલ્લી
  ધારણા ગણાય છે.
  \item\ollabel{defn:prf:one-prem}
  જો~$R$ એ એક, બે
  અથવા ત્રણ આધારવાક્યો
  $!A_1$, $!A_2$, $!A_3$
  અને ફલિતવાક્ય~$!A$
  ધરાવતો નિયમ હોય,
  અને~$\delta_1$,
  $\delta_2$, $\delta_3$
  અનુક્રમે~$!A_1$,
  $!A_2$, $!A_3$માં
  પૂરી થતી સાબિતીઓ
  હોય, તો અનુક્રમે
  નીચેનાં વૃક્ષો
  \[
  \AxiomC{}
  \RightLabel{$\delta_1$}
  \DeduceC{$!A_1$}
  \RightLabel{$R$}
  \UnaryInfC{$!A$}
  \DisplayProof
\qquad
  \AxiomC{}
  \RightLabel{$\delta_1$}
  \DeduceC{$!A_1$}
  \AxiomC{}
  \RightLabel{$\delta_2$}
  \DeduceC{$!A_2$}
  \RightLabel{$R$}
  \BinaryInfC{$!A$}
  \DisplayProof
\qquad
  \AxiomC{}
  \RightLabel{$\delta_1$}
  \DeduceC{$!A_1$}
  \AxiomC{}
  \RightLabel{$\delta_2$}
  \DeduceC{$!A_2$}
  \AxiomC{}
  \RightLabel{$\delta_3$}
  \DeduceC{$!A_3$}
  \RightLabel{$R$}
  \TrinaryInfC{$!A$}
  \DisplayProof
  \]
  સાબિતી છે, જો
  તેમને રચતી સાબિતીઓની
  ધારણાઓનાં
  ચિહ્નો સુસંગત હોય
  (બે જુદાં સૂત્રોને
  એક જ મુક્તિ-ચિહ્ન
  ન અપાયું હોય)
  અને અનુમાનો~$R$નો
  યોગ્ય ઉપયોગ હોય.
  \sourcecorrection{OLMD-233}{મૂળના ત્રણ આધારવાક્યવાળા વૃક્ષમાં $\delta_3$નું નિષ્કર્ષ $!A_2$ લખાયું છે; ઉપર આપેલી અનુક્રમણીકા પ્રમાણે તે $!A_3$ છે.}
  \item નવી સાબિતીમાં
  સૌથી ઉપર આવતાં
  સૂત્રનાં સ્થાનો
  $\delta_1$, $\delta_2$
  અથવા~$\delta_3$માં
  ખુલ્લી ધારણા હોય તો
  નવી સાબિતીની
  \emph{ખુલ્લી ધારણાઓ}
  ગણાય છે, પરંતુ
  નિયમથી મુક્ત થયેલાં
  સ્થાનોને છોડીને.
  જો નિયમ પર~$x$
  (અથવા~$x\, y$)
  ચિહ્ન હોય,
  તો~\Intro{\lif}
  અને~\Intro{\lnot} માટે
  $\delta_1$માંનાં
  $x$ ચિહ્નવાળાં બધાં
  ખુલ્લાં ધારણાસ્થાનો
  મુક્ત થાય છે;
  \Elim{\lexists} માટે
  $\delta_2$માંનાં
  $x$ ચિહ્નવાળાં બધાં,
  અને~\Elim{\lor} માટે
  $\delta_2$માંનાં
  $x$ ચિહ્નવાળાં તથા
  $\delta_3$માંનાં
  $y$ ચિહ્નવાળાં બધાં
  ખુલ્લાં ધારણાસ્થાનો
  મુક્ત થાય છે.
\end{enumerate}
\end{defn}

સાબિતીની
આગમનાત્મક રચનામાં
વપરાતી સાબિતીઓને
તેની \emph{ઉપ-સાબિતીઓ}
કહે છે.
અનુમાનનાં આધારવાક્યોમાં
પૂરી થતી
ઉપ-સાબિતીઓને
તે અનુમાનની
\emph{તાત્કાલિક
ઉપ-સાબિતીઓ}
કહે છે.
\Log{N1c} અને~\Log{N1i}
ના નિયમો
\olref{tab:N1}માં છે.
\sourcecorrection{OLMD-234}{મૂળમાં N2c/N2iના નિયમો tab:N1માં છે એમ કહેવાયું છે; નીચે આયાત થતી N1 નિયમ-તકતી અને અહીંની N1 વ્યાખ્યા પ્રમાણે તંત્રોનાં નામ N1c/N1i છે.}

\begingroup
\makeatletter\def\ol@sectioncs{section}\makeatother

% BEGIN OLP-0666
\phantomsection\label{guunit:OLP-0666}


\begin{table}
\begin{center}
\begin{tabular}[t]{@{}cc@{}}
\bottomAlignProof
\AxiomC{$\Discharge{!A}{x}$}
\shortDeduce
\DeduceC{$!B$}
\DischargeRule{\Intro{\lif}}{x}
\UnaryInfC{$!A \lif !B$}
\DisplayProof
&
\bottomAlignProof
\AxiomC{$!A \lif !B$}
\AxiomC{$!A$}
\RightLabel{\Elim{\lif}}
\BinaryInfC{$!B$}
\DisplayProof
\\[4ex]
\bottomAlignProof
\AxiomC{$!A$}
\AxiomC{$!B$}
\RightLabel{\Intro{\land}}
\BinaryInfC{$!A \land !B$}
\DisplayProof
&
\begin{tabular}[b]{@{}l@{}}
\AxiomC{$!A \land !B$}
\RightLabel{\Elim{\land}[1]}
\UnaryInfC{$!A$}
\DisplayProof
\\[3ex]
\bottomAlignProof
\AxiomC{$!A \land !B$}
\RightLabel{\Elim{\land}[2]}
\UnaryInfC{$!B$}
\DisplayProof
\end{tabular}
\\[4ex]
\begin{tabular}{@{}l@{}}
\AxiomC{$!A$}
\RightLabel{\Intro{\lor}[1]}
\UnaryInfC{$!A \lor !B$}
\DisplayProof
\\[3ex]
\AxiomC{$!B$}
\RightLabel{\Intro{\lor}[2]}
\UnaryInfC{$!A \lor !B$}
\DisplayProof
\end{tabular}
&
\AxiomC{$!A \lor !B$}
\AxiomC{$\Discharge{!A}{x}$}
\shortDeduce
\DeduceC{$!C$}
\AxiomC{$\Discharge{!B}{y}$}
\shortDeduce
\DeduceC{$!C$}
\DischargeRule{\Elim{\lor}}{x\,y}
\TrinaryInfC{$!C$}
\DisplayProof
\bottomAlignProof
\\[6ex]
\bottomAlignProof
\AxiomC{$\Discharge{!A}{x}$}
\shortDeduce
\DeduceC{$\lfalse$}
\DischargeRule{\Intro{\lnot}}{x}
\UnaryInfC{$\lnot !A$}
\DisplayProof
&
\bottomAlignProof
\AxiomC{$\lnot !A$}
\AxiomC{$!A$}
\RightLabel{\Elim{\lnot}}
\BinaryInfC{$\lfalse$}
\DisplayProof
\\[4ex]
\bottomAlignProof
\AxiomC{$!A(a)$}
\RightLabel{\Intro{\lforall}}
\UnaryInfC{$\lforall[x][!A(x)]$}
\DisplayProof
&
\bottomAlignProof
\AxiomC{$\lforall[x][!A(x)]$}
\RightLabel{\Elim{\lforall}}
\UnaryInfC{$!A(t)$}
\DisplayProof
\\[4ex]
\bottomAlignProof
\AxiomC{$!A(t)$}
\RightLabel{\Intro{\lexists}}
\UnaryInfC{$\lexists[x][!A(x)]$}
\DisplayProof
&
\bottomAlignProof
\AxiomC{$\lexists[x][!A(x)]$}
\AxiomC{$\Discharge{!A(a)}{x}$}
\shortDeduce
\DeduceC{$!C$}
\DischargeRule{\Elim{\lexists}}{x}
\BinaryInfC{$!C$}
\DisplayProof
\\[4ex]
\bottomAlignProof
\AxiomC{$\lfalse$}
\RightLabel{\FalseInt}
\UnaryInfC{$!A$}
\DisplayProof
&
\bottomAlignProof
\AxiomC{$\Discharge{\lnot!A}{x}$}
\shortDeduce
\DeduceC{$\lfalse$}
\DischargeRule{\FalseCl}{x}
\UnaryInfC{$!A$}
\DisplayProof
\end{tabular}
\end{center}
\caption{\Log{N1i} અને~\Log{N1c}ના નિયમો.
$\Intro{\lforall}$માં
અચળ~$a$ નિષ્કર્ષમાં
અથવા કોઈ ખુલ્લી ધારણામાં
આવવો ન જોઈએ.
$\Elim{\lexists}$માં~$a$
આધારવિધાનોમાં અથવા કોઈ
ખુલ્લી ધારણામાં આવવો
ન જોઈએ. \Log{N1i}માં
\FalseCl સામેલ નથી.}
\ollabel{tab:N1}
\end{table}
% END OLP-0666
\endgroup


\begin{defn}
સાબિતી~$\delta$ની
\emph{ઊંચાઈ}~$\pheight{\delta}$
આગમનાત્મક રીતે આ
પ્રમાણે વ્યાખ્યાયિત થાય છે.
\begin{enumerate}
  \item જો~$\delta$માં માત્ર
  એક ધારણા હોય, તો
  $\pheight{\delta} = 0$.
  \item જો~$\delta$નો અંત
  એક આધારવાક્યવાળા અનુમાનમાં
  થાય અને તેની તાત્કાલિક
  ઉપ-સાબિતી~$\delta_1$ હોય,
  તો $\pheight{\delta} = \pheight{\delta_1} + 1$.
  \item જો~$\delta$નો અંત
  બે આધારવાક્યવાળા અનુમાનમાં
  થાય અને તેની તાત્કાલિક
  ઉપ-સાબિતીઓ~$\delta_1$
  અને~$\delta_2$ હોય, તો
  $\pheight{\delta} =
  \max(\pheight{\delta_1}, \pheight{\delta_2}) + 1$.
  \item જો~$\delta$નો અંત
  \Elim{\lor}માં થાય
  અને તેની તાત્કાલિક
  ઉપ-સાબિતીઓ~$\delta_1$,
  $\delta_2$ અને~$\delta_3$
  હોય, તો
  $\pheight{\delta} =
  \max(\pheight{\delta_1}, \pheight{\delta_2}, \pheight{\delta_3}) + 1$.
\end{enumerate}
જો સિક્વન્ટ~$S$ની
ઊંચાઈ~$\le n$ ધરાવતી
સાબિતી હોય, તો
$\Proves[n] S$ લખીએ છીએ.
\end{defn}

\begin{ex}
\Log{N1i}માં
કોઈ પણ સૂત્ર
પોતે ખુલ્લી ધારણા
હોઈ પોતાની
સાબિતી ગણાય છે.
એટલે
\[
\Discharge{!A \land !B}{x}
\]
એ~સાબિતી~$\delta_1$
છે. તેની ઊંચાઈ~$0$ છે.
$\delta_1$માંથી
\Elim{\land}નાં બે
રૂપોમાંથી એકેક લગાવી
બે સાબિતીઓ~$\delta_2$
અને~$\delta_3$ મળે છે:
\[
\AxiomC{$\Discharge{!A \land !B}{x}$}
\RightLabel{\Elim{\land}[2]}
\UnaryInfC{$!B$}
\DisplayProof
\qquad
\AxiomC{$\Discharge{!A \land !B}{x}$}
\RightLabel{\Elim{\land}[1]}
\UnaryInfC{$!A$}
\DisplayProof
\]
\sourcecorrection{OLMD-235}{મૂળના આરંભિક અને અનુગામી વૃક્ષોમાં $!A \land !B^x$ લખવાથી $x$ માત્ર $!B$ સાથે જોડાય છે; અહીં ધારણા આખા $!A \land !B$ સૂત્રની હોવાથી તેને એક જ મુક્તિ-ચિહ્નથી દર્શાવ્યું છે.}
$\delta_2$ અને~$\delta_3$ના
છેલ્લા અનુમાનની તાત્કાલિક
ઉપ-સાબિતીઓ સમાન છે:
$!A \land !B$.
બંને સાબિતીઓમાં
$!A\land !B$નું
સ્થાન ખુલ્લી ધારણા છે.
આપણે
$\pheight{\delta_2} = \pheight{\delta_3} = 1$
મેળવીએ છીએ.
\sourcecorrection{OLMD-236}{મૂળમાં $\delta_1$ની ઊંચાઈ એક જ વાક્યમાં $0$ અને પછી $1$ કહેવાઈ છે; અહીં એક પગલું ધરાવતી $\delta_2$ તથા $\delta_3$ની ઊંચાઈ $1$ બતાવી છે.}

\Intro{\land} નિયમથી
$!B \land !A$ સ્વરૂપનું
ફલિતવાક્ય મેળવવા
$!B$ અને~$!A$
સ્વરૂપનાં બે આધારવાક્યો
જોઈએ. તેથી નીચેનું
સાબિતી~$\delta_4$ છે:
\[
\AxiomC{$\Discharge{!A \land !B}{x}$}
\RightLabel{\Elim{\land}[2]}
\UnaryInfC{$!B$}
\LeftSubproofLabel{$\delta_2$}
\AxiomC{$\Discharge{!A \land !B}{x}$}
\RightLabel{\Elim{\land}[1]}
\UnaryInfC{$!A$}
\RightSubproofLabel{$\delta_3$}
\RightLabel{\Intro{\land}}
\BinaryInfC{$!B \land !A$}
\DisplayProof
\]
તેની ઊંચાઈ
$\pheight{\delta_4} = \max(\pheight{\delta_2},
\pheight{\delta_3})+1 = \max(1,1) + 1 = 2$ છે.
તેમાં~$!A \land !B$
ધારણા તરીકે બે વાર
આવે છે અને બંને
સ્થાનો ખુલ્લાં છે.

છેલ્લે~\Intro{\lif}
લગાવી~$\delta_5$
મેળવી શકાય છે:
\[
\AxiomC{$\Discharge{!A \land !B}{x}$}
\RightLabel{\Elim{\land}[2]}
\UnaryInfC{$!B$}
\LeftSubproofLabel{$\delta_2$}
\AxiomC{$\Discharge{!A \land !B}{x}$}
\RightLabel{\Elim{\land}[1]}
\UnaryInfC{$!A$}
\RightSubproofLabel{$\delta_3$}
\RightLabel{\Intro{\land}}
\BinaryInfC{$!B \land !A$}
\DischargeRule{\Intro{\lif}}{x}
\UnaryInfC{$(!A \land !B) \lif (!B \land !A)$}
\DisplayProof
\]
તેની ઊંચાઈ
$\pheight{\delta_4} + 1 = 3$ છે.
\Intro{\lif} અનુમાન
$x$ ચિહ્નવાળી
$!A \land !B$
સ્વરૂપની ધારણાઓનાં
બધાં સ્થાનો મુક્ત કરે છે;
અર્થાત્ તાત્કાલિક
ઉપસાબિતીમાં પહેલાં
ખુલ્લાં રહેલાં બંને
સ્થાનો મુક્ત થાય છે.
$\delta_5$ની એ જ
ધારણાઓ હોવાથી તેમાં
હવે કોઈ ખુલ્લી
ધારણા નથી.
આથી તે પોતાના
અંતિમ-સૂત્ર
$(!A \land !B) \lif (!B \land !A)$ની
સાબિતી છે.
\end{ex}
% END OLP-0668
\endgroup

\begingroup
\makeatletter\def\ol@sectioncs{section}\makeatother

% BEGIN OLP-0669
\olfileid{pt}{nat}{seq}

\olsection{સિક્વન્ટો સાથેનું સ્વાભાવિક નિગમન: \Log{N2c} અને~\Log{N2i}}
\phantomsection\label{guunit:OLP-0669}


\Log{N1c} અથવા~\Log{N1i}માંની
સાબિતી~$\delta$
બતાવે છે કે તેનું
અંતિમ-સૂત્ર~$!A$
તેની ખુલ્લી ધારણાઓમાંથી
ફલિત થાય છે.
\sourcecorrection{OLMD-237}{મૂળમાં અહીં G1c/G1i લખાયું છે, પણ ખુલ્લી ધારણાવાળી સૂત્ર-સાબિતીઓની ચર્ચા N1c/N1i અંગે છે; G1 તંત્ર સિક્વન્ટ-કલનનાં છે.}
જો~$\Gamma$ એવા
સૂત્રો~$!B$નો
ગણ હોય કે~$!B^x$
એ~$\delta$માં
ખુલ્લી ધારણા હોય,
તો આને
$\delta \Proves \Gamma
\Sequent !A$ એમ
લખી શકાય.
આથી સિક્વન્ટો માટે
પણ સ્વાભાવિક નિગમનનું
તંત્ર ઘડી શકાય એવું
સૂઝે છે.
આવા તંત્રમાં દરેક
સિક્વન્ટ જણાવે છે કે
જમણી બાજુનું સૂત્ર
કઈ ખુલ્લી ધારણાઓ
પર આધારિત છે,
એટલે એ સિક્વન્ટમાં
પૂરી થતી ઉપ-સાબિતીમાં
કઈ ધારણાઓ ખુલ્લી છે.
આ સિક્વન્ટો સાથેનું,
અથવા ``સિક્વન્ટ-શૈલીનું'',
સ્વાભાવિક નિગમન છે;
મળતાં તંત્રોને
\Log{N2c} અને~\Log{N2i}
કહીએ છીએ.

\Log{N1} અને~\Log{N2}
વચ્ચેનો અનુરૂપ સંબંધ
વધુ સીધો બનાવવા
ડાબી બાજુના ગણો~$\Gamma$
ચિહ્નિત
સૂત્રો~$x:!B$
ધરાવે એમ લઈએ.
\Log{N1}ના નિયમો
આધારવાક્યોનાં
સૂત્રો પર,
એટલે તાત્કાલિક
ઉપ-સાબિતીઓના
અંતિમ-સૂત્રો પર,
જ કામ કરે છે;
તેથી~\Log{N2}ના
નિયમો સિક્વન્ટના
ઉત્તરાંગ પર જ
કામ કરે છે.
આથી સિક્વન્ટની
ડાબી બાજુનાં
ચિહ્નિત સૂત્રોને
ફક્ત \emph{સંદર્ભ}
કહી શકાય.

\Log{N2}ની સાબિતીઓમાં
ધારણાઓને બદલે
સૌથી ઉપર
આ સ્વરૂપનાં
સ્વયંસિદ્ધ સિક્વન્ટો છે:
\[x:!A \Sequent !A.\]
આ નિયમો~\Log{N1}ના
નિયમોને ચોક્કસ અનુરૂપ
છે અને
\olref{tab:N2}માં
આપેલા છે.

\begingroup
\makeatletter\def\ol@sectioncs{section}\makeatother

% BEGIN OLP-0667
\phantomsection\label{guunit:OLP-0667}


\begin{table}
\begin{tabular}{@{}c@{}c@{}}
\Axiom$x: !A, \Gamma \fCenter !B$
\RightLabel{\Intro{\lif}}
\UnaryInf$\Gamma \fCenter !A \lif !B $
\DisplayProof
&
\Axiom$\Gamma_1 \fCenter !A \lif !B$
\Axiom$\Gamma_2 \fCenter !A$
\RightLabel{\Elim{\lif}}
\BinaryInf$\Gamma_1, \Gamma_2 \fCenter !B$
\DisplayProof
\\[4ex]
\Axiom$\Gamma_1 \fCenter !A$
\Axiom$\Gamma_2 \fCenter !B$
\RightLabel{\Intro{\land}}
\BinaryInf$\Gamma_1, \Gamma_2 \fCenter !A \land !B $
\DisplayProof
&
\begin{tabular}[b]{@{}l@{}}
\Axiom$\Gamma \fCenter !A \land !B$
\RightLabel{\Elim{\land}[1]}
\UnaryInf$\Gamma \fCenter !A$
\DisplayProof
\\[3ex]
\Axiom$\Gamma \fCenter !A \land !B$
\RightLabel{\Elim{\land}[2]}
\UnaryInf$\Gamma \fCenter !B$
\DisplayProof\\[3ex]
\end{tabular}
\\[4ex]
\Axiom$\Gamma \fCenter !A$
\RightLabel{\Intro{\lor}[1]}
\UnaryInf$\Gamma \fCenter !A \lor !B$
\DisplayProof
&
\Axiom$ \Gamma \fCenter !B$
\RightLabel{\Intro{\lor}[2]}
\UnaryInf$\Gamma \fCenter !A \lor !B$
\DisplayProof\\[3ex]
\multicolumn{2}{c}{
\begin{tabular}[b]{@{}l@{}}
\Axiom$\Gamma_1 \fCenter !A \lor !B$
\Axiom$x:!A, \Gamma_2 \fCenter !C$
\Axiom$y:!B, \Gamma_3 \fCenter !C$
\RightLabel{\Elim{\lor}}
\TrinaryInf$\Gamma_1, \Gamma_2, \Gamma_3 \fCenter !C$
\DisplayProof
\end{tabular}
}
\\[4ex]
\Axiom$x:!A, \Gamma \fCenter \lfalse$
\RightLabel{\Intro{\lnot}}
\UnaryInf$ \Gamma \fCenter \lnot !A $
\DisplayProof
&
\Axiom$\Gamma_1 \fCenter \lnot !A $
\Axiom$\Gamma_2 \fCenter !A $
\RightLabel{\Elim{\lnot}}
\BinaryInf$\Gamma_1, \Gamma_2 \fCenter \lfalse$
\DisplayProof
\\[4ex]
\Axiom$\Gamma \fCenter !A(a) $
\RightLabel{\Intro{\lforall}}
\UnaryInf$\Gamma \fCenter \lforall[x][!A(x)]$
\DisplayProof
&
\Axiom$\Gamma \fCenter \lforall[x][!A(x)]$
\RightLabel{\Elim{\lforall}}
\UnaryInf$\Gamma \fCenter !A(t)$
\DisplayProof
\\[4ex]
\Axiom$\Gamma \fCenter !A(t) $
\RightLabel{\Intro{\lexists}}
\UnaryInf$\Gamma \fCenter \lexists[x][!A(x)]$
\DisplayProof
&
\Axiom$\Gamma_1 \fCenter  \lexists[x][!A(x)]$
\Axiom$x: !A(a), \Gamma_2 \fCenter !C$
\RightLabel{\Elim{\lexists}}
\BinaryInf$\Gamma_1, \Gamma_2 \fCenter !C$
\DisplayProof
\\[4ex]
\Axiom$\Gamma \fCenter \lfalse$
\RightLabel{\FalseInt}
\UnaryInf$\Gamma \fCenter !A$
\DisplayProof
&
\Axiom$x: \lnot !A, \Gamma \fCenter \lfalse$
\RightLabel{\FalseCl}
\UnaryInf$\Gamma \fCenter !A$
\DisplayProof
\end{tabular}
\caption{\Log{N2c} અને~\Log{N2i}ના નિયમો.
$\Intro{\forall}$ અને
$\Elim{\exists}$માં
અચળ~$a$
નિષ્કર્ષ-સિક્વન્ટમાં
આવવો ન જોઈએ.
$\Gamma_i$ એ ચિહ્નિત
સૂત્રો~$x :
!A$ના ગણો છે.
\Intro{\lforall} અને
\Elim{\lexists} નિયમોએ
સ્વનિર્ધારક ચલની શરત
પૂરી કરવી જોઈએ:
$a$ નિષ્કર્ષ-સિક્વન્ટમાં
આવવો ન જોઈએ.
\Log{N2i}માં \FalseCl
સામેલ નથી.}
\ollabel{tab:N2}
\end{table}
% END OLP-0667
\endgroup


\Log{N1c} અને~\Log{N1i}ની
જેમ અહીં પણ
\Intro{\lif}, \Intro{\lnot},
\Elim{\lor} અને
\Elim{\lexists} નિયમો
ધારણાઓ મુક્ત કરી
શકે, પણ એમ કરવું
જરૂરી ન હોય,
એવી છૂટ જોઈએ.
તેથી~\Log{N2c}
અને~\Log{N2i}માં
સંબંધિત આધારવાક્યના
સંદર્ભમાં અનુરૂપ
સૂત્રો~$!A$,
$!B$ અને~$!A(a)$
હાજર ન હોય તોપણ
એ નિયમો યોગ્ય રીતે
લાગુ પાડી શકાય છે.

\begin{prop}
જો~$\delta$ એ~$!A$ની
\Log{N1c} અથવા~\Log{N1i}માંની
સાબિતી હોય,
તો~\Log{N2c}
(\Log{N2i})માં
$\Gamma \Sequent !A$ની
સાબિતી~$\delta'$
છે. અહીં~$\Gamma$
એ એવા બધા~$x:!B$નો
ગણ છે કે~$!B^x$
એ~$\delta$ની
ખુલ્લી ધારણા છે.
\end{prop}

\begin{proof}
  $n=\pheight{\delta}$
  પર આગમન કરીએ.
  જો~$n=0$, તો~$\delta$
  એક જ ખુલ્લી
  ધારણા~$!A^x$ છે.
  ત્યારે~$\Gamma = \{x:!A\}$
  અને~$\Gamma
  \Sequent !A$ એ
  \Log{N2c} અથવા~\Log{N2i}
  નું સ્વયંસિદ્ધ છે.

  જો~$n>0$, તો~$\delta$ના
  છેલ્લા અનુમાન પ્રમાણે
  કિસ્સા જુદા પાડીએ.
  અહીં માત્ર એક
  કિસ્સો જોઈએ;
  બાકીના કસરત
  તરીકે રાખીએ છીએ.

  માનો કે છેલ્લું
  અનુમાન~$\Intro{\lif}$
  છે. ત્યારે~$\delta$નો
  અંત આ રીતે થાય છે:
  \[
  \AxiomC{$\Discharge{!B}{x}$}
  \RightLabel{$\delta_1$}
  \DeduceC{$!C$}
  \DischargeRule{\Intro{\lif}}{x}
  \UnaryInfC{$!B \lif !C$}
  \DisplayProof
  \]
  આગમન-પરિકલ્પનાથી
  \Log{N2c} (\Log{N2i})માં
  $\Gamma' \Sequent !C$ની
  સાબિતી~$\delta_1'$
  છે, જ્યાં~$\Gamma'$
  એ~$\delta_1$માં
  ખુલ્લી ચિહ્નિત
  ધારણાઓનો ગણ છે.
  જો~$!B^x$,
  $\delta_1$ની ખુલ્લી
  ધારણા હોય, તો
  \Intro{\lif} અનુમાન
  તેને મુક્ત કરે છે,
  અને~$\delta$ની
  ખુલ્લી ધારણાઓ
  $\Gamma =
  \Gamma' \setminus \{x:!B\}$
  છે. એટલે~$\delta_1'$
  એ~$x:!B, \Gamma \Sequent !C$ની
  સાબિતી છે,
  અને~$\delta'$ને
  આ રીતે લઈ શકાય:
  \[
  \AxiomC{}
  \RightLabel{$\delta_1'$}
  \Deduce$x:!B, \Gamma \fCenter !C$
  \RightLabel{\Intro{\lif}}
  \UnaryInf$\Gamma \fCenter !B \lif !C$
  \DisplayProof
  \]
  \sourcecorrection{OLMD-238}{મૂળના આ N2 વૃક્ષમાં પૂર્વાંગમાં $x:B$ લખાયું છે; તરતની દલીલ અને ચિહ્નિત સૂત્રની વ્યાખ્યા પ્રમાણે $x:!B$ જોઈએ.}
  જો~$!B^x$,
  $\delta_1$ની ખુલ્લી
  ધારણા ન હોય,
  તો~\Intro{\lif}
  કોઈ ધારણા મુક્ત
  કરતું નથી, અને
  $\delta$ તથા~$\delta_1$ની
  ખુલ્લી ધારણાઓ સમાન છે.
  \Log{N2c} તથા~\Log{N2i}માં
  \Intro{\lif}ના
  આધારવાક્યના સંદર્ભમાં
  $x:!B$ હાજર હોવું
  જરૂરી ન હોવાથી,
  $\delta'$ આ રીતે
  લઈ શકાય:
  \[
  \AxiomC{}
  \RightLabel{$\delta_1'$}
  \Deduce$\Gamma \fCenter !C$
  \RightLabel{\Intro{\lif}}
  \UnaryInf$\Gamma \fCenter !B \lif !C$
  \DisplayProof
  \]

  $\delta$નો અંત બીજા
  કોઈ નિયમમાં થાય
  તે કિસ્સા સમાન છે.
\end{proof}

\begin{prop}
જો~$\delta$ એ
\Log{N2c} અથવા~\Log{N2i}માં
$\Gamma \Sequent !A$ની
સાબિતી હોય,
તો~\Log{N1c} (\Log{N1i})
માં એવી
સાબિતી~$\delta'$
છે જેનું
અંતિમ-સૂત્ર~$!A$
છે અને
દરેક
$x:!B \in \Gamma$
માટે~$!B^x$ તેની
ખુલ્લી ધારણા છે.
\end{prop}

\begin{proof}
  ફરી~$\delta$ની
  ઊંચાઈ~$n$ પર
  આગમન કરીએ.
  જો~$n=0$,
  તો~$\delta$ એ
  સ્વયંસિદ્ધ
  $x:!A \Sequent !A$ છે.
  સાબિતી~$\delta'$
  એ ધારણા~$!A^x$ છે.

  જો~$n>0$, તો~$\delta$ના
  છેલ્લા અનુમાન પ્રમાણે
  કિસ્સા જુદા પાડીએ.
  દાખલા તરીકે, અનુમાન
  $\Intro{\lif}$ હોય,
  તો~$\delta$નો અંત
  આ બેમાંથી એક રીતે થાય છે:
  \[
  \bottomAlignProof
  \AxiomC{}
  \RightLabel{$\delta_1$}
  \Deduce$x:!B, \Gamma \fCenter !C$
  \RightLabel{\Intro{\lif}}
  \UnaryInf$\Gamma \fCenter !B \lif !C$
  \DisplayProof
  \quad\text{ અથવા }\quad
  \bottomAlignProof
  \AxiomC{}
  \RightLabel{$\delta_1$}
  \Deduce$\Gamma \fCenter !C$
  \RightLabel{\Intro{\lif}}
  \UnaryInf$\Gamma \fCenter !B \lif !C$
  \DisplayProof
  \]
  આગમન-પરિકલ્પનાથી
  \Log{N1c} (\Log{N1i})
  માં~$!C$ની
  સાબિતી~$\delta_1'$
  છે. પહેલા કિસ્સામાં
  $!B^x$ એ તેની
  ખુલ્લી ધારણા છે;
  બીજામાં નથી.
  સાબિતી~$\delta'$ને
  આ રીતે લઈ શકાય:
  \[
  \bottomAlignProof
  \AxiomC{$\Discharge{!B}{x}$}
  \RightLabel{$\delta_1'$}
  \DeduceC{$!C$}
  \DischargeRule{\Intro{\lif}}{x}
  \UnaryInfC{$!B \lif !C$}
  \DisplayProof
  \quad\text{ અથવા }\quad
  \bottomAlignProof
  \AxiomC{}
  \RightLabel{$\delta_1'$}
  \DeduceC{$!C$}
  \RightLabel{\Intro{\lif}}
  \UnaryInfC{$!B \lif !C$}
  \DisplayProof
  \]
  અનુક્રમે.
  \sourcecorrection{OLMD-239}{મૂળના અંતિમ બે N1 વૃક્ષોમાં $\delta_1$ લખાયું છે, જે N2 ઉપસાબિતી છે; આગમનથી મળતી N1 સાબિતી $\delta_1'$ છે. પહેલીમાં $!B^x$ ખુલ્લું છે તે પણ એ જ ઉપસાબિતી માટે કહેવાય છે.}
\end{proof}
% END OLP-0669
\endgroup

\begingroup
\makeatletter\def\ol@sectioncs{section}\makeatother

% BEGIN OLP-0665
\olfileid{pt}{nat}{qua}

\olsection{નિયમિત સાબિતીઓ અને પ્રતિસ્થાપન}
\phantomsection\label{guunit:OLP-0665}


પરિમાણકોના નિયમોમાં
\Elim{\lexists} અને
\Intro{\lforall}ને લાગતી
``સ્વનિર્ધારક ચલની શરતો''
ને કારણે કેટલીક મુશ્કેલીઓ
ઊભી થાય છે. આ અનુમાનોમાં
આવતો અચળ~$c$
ફરી કદી ન વપરાય એ
સુનિશ્ચિત કરીએ તો
મોટાભાગની મુશ્કેલીઓ
ઉકેલી શકાય છે.
આ માટે નીચેની વ્યાખ્યા લો.

\begin{defn}
સ્વાભાવિક નિગમનની
સાબિતી~$\delta$
\emph{નિયમિત} છે, જો
\begin{enumerate}
  \item કોઈ સ્વનિર્ધારક ચલ~$a$
  એકથી વધુ \Elim{\lexists}
  અથવા~\Intro{\lforall}
  અનુમાનમાં સ્વનિર્ધારક ચલ
  તરીકે ન વપરાતો હોય, અને
  \item જો~$a$ કોઈ
  \Elim{\lexists} અથવા
  \Intro{\lforall} અનુમાનનો
  સ્વનિર્ધારક ચલ હોય,
  તો તે અનુક્રમે ગૌણ
  આધારવિધાન અથવા આધારવિધાન
  સુધી પહોંચતી તાત્કાલિક
  ઉપ-સાબિતીમાં જ આવતો હોય.
\end{enumerate}
\end{defn}

\begin{lem}\ollabel{lem:subst-regular}
જો~$\delta$ નિયમિત હોય,
તેનો અંત~$!C$માં થતો હોય,
$c$ એવો અચળ હોય
જે~$\delta$માં સ્વનિર્ધારક
ચલ તરીકે વપરાયો ન હોય,
અને~$t$ એવું પદ હોય
જેમાં~$\delta$નો કોઈ
સ્વનિર્ધારક ચલ ન હોય,
તો~$\delta$માં~$c$ના
દરેક સ્થાનને~$t$થી
બદલવાથી મળતી
$\Subst{\delta}{t}{c}$
નિયમિત સાબિતી છે.
$\Subst{\delta}{t}{c}$નું
અંતિમ-સૂત્ર
$\Subst{!C}{t}{c}$ છે.
જો~$!A^x$ એ~$\delta$ની
ખુલ્લી ધારણા હોય, તો
$\Subst{!A}{t}{c}^x$ એ
$\Subst{\delta}{t}{c}$ની
ખુલ્લી ધારણા છે.
\end{lem}

\begin{proof}
  $\pheight{\delta}$ પર
  આગમન કરીએ.
  જો~$\pheight{\delta} = 0$,
  તો તેમાં માત્ર
  ધારણા~$!A^x$ છે.
  ત્યારે~$\Subst{\delta}{t}{c}$
  એ~$\Subst{!A}{t}{c}^x$ છે.

જો~$\pheight{\delta} > 0$,
  તો~$\delta$ના છેલ્લા
  અનુમાન પ્રમાણે કિસ્સા
  જુદા પાડીએ.
  વિધાનાત્મક સંયોજકોના
  નિયમોના કિસ્સા સહેલા છે
  અને કસરત તરીકે રાખ્યા છે.
  અહીં~$\delta$નો અંત
  પરિમાણક-અનુમાનમાં થાય
  તે કિસ્સા જોઈએ.
  \begin{enumerate}
    \item છેલ્લું અનુમાન
    $\Intro{\lforall}$ હોય,
    તો~$\delta$નું સ્વરૂપ
    \[
    \AxiomC{}
    \RightLabel{$\delta'$}
    \DeduceC{$!A(a,c)$}
    \RightLabel{\Intro{\lforall}}
    \UnaryInfC{$\lforall[x][!A(x,c)]$}
    \DisplayProof
    \]
    આગમન-પરિકલ્પનાથી
    $\Subst{\delta'}{t}{c}$
    એ~$!A(a,t)$ની નિયમિત
    સાબિતી છે.
    $a$ એ~$\delta$નો
    સ્વનિર્ધારક ચલ હોવાથી
    $t$માં~$a$ આવતો નથી.
    તેથી~$\Subst{\delta}{t}{c}$નું
    છેલ્લું અનુમાન,
    \[
    \AxiomC{}
    \RightLabel{$\Subst{\delta'}{t}{c}$}
    \DeduceC{$!A(a,t)$}
    \RightLabel{\Intro{\lforall}}
    \UnaryInfC{$\lforall[x][!A(x,t)]$}
    \DisplayProof
    \]
    યોગ્ય છે: $t$માં~$a$
    ન હોવાથી નિષ્કર્ષમાં
    કે કોઈ ખુલ્લી ધારણામાં
    $a$ આવતો નથી.
    \item છેલ્લું અનુમાન
    $\Elim{\lforall}$ હોય,
    તો~$\delta$નું સ્વરૂપ
    \[
    \AxiomC{}
    \RightLabel{$\delta'$}
    \DeduceC{$\lforall[x][!A(x,c)]$}
    \RightLabel{\Elim{\lforall}}
    \UnaryInfC{$!A(s(c),c)$}
    \DisplayProof
    \]
    આગમન-પરિકલ્પનાથી
    $\Subst{\delta'}{t}{c}$
    એ~$\lforall[x][!A(x,t)]$ની
    નિયમિત સાબિતી છે.
    (નિષ્કર્ષમાંનું
    પદ~$s(c)$, $c$ ધરાવતું
    હોય પણ શકે અને ન પણ હોય.)
    $\Subst{\delta}{t}{c}$નું
    છેલ્લું અનુમાન,
    \[
    \AxiomC{}
    \RightLabel{$\Subst{\delta'}{t}{c}$}
    \DeduceC{$\lforall[x][!A(x,t)]$}
    \RightLabel{\Elim{\lforall}}
    \UnaryInfC{$!A(s(t),t)$}
    \DisplayProof
    \]
    યોગ્ય છે, અને
    $!A(s(t),t) \ident \Subst{!A(s(c),c)}{t}{c}$.
    \Intro{\lexists}નો
    કિસ્સો આને સમાન છે.
  \item છેલ્લું અનુમાન
    $\Elim{\lexists}$ હોય,
    તો~$\delta$નું સ્વરૂપ
    \[
    \AxiomC{}
    \RightLabel{$\delta_1'$}
    \DeduceC{$\lexists[x][!A(x,c)]$}
    \AxiomC{$\Discharge{!A(a,c)}{x}$}
    \RightLabel{$\delta_2'$}
    \DeduceC{$!C$}
    \DischargeRule{\Elim{\lexists}}{x}
    \BinaryInfC{$!C$}
    \DisplayProof
    \]
    આગમન-પરિકલ્પનાથી
    $\Subst{\delta_1'}{t}{c}$
    અને~$\Subst{\delta_2'}{t}{c}$
    અનુક્રમે
    $\lexists[x][!A(x,t)]$ની
    અને ખુલ્લી ધારણા
    $!A(a,t)$માંથી
    $\Subst{!C}{t}{c}$ની
    નિયમિત સાબિતીઓ છે.
    $a$ એ~$\delta$નો
    સ્વનિર્ધારક ચલ હોવાથી
    $t$માં~$a$ આવતો નથી.
    આથી~$\Subst{\delta}{t}{c}$નું
    છેલ્લું અનુમાન,
    \[
    \AxiomC{}
    \RightLabel{$\Subst{\delta_1'}{t}{c}$}
    \DeduceC{$\lexists[x][!A(x,t)]$}
    \AxiomC{$\Discharge{!A(a,t)}{x}$}
    \RightLabel{$\Subst{\delta_2'}{t}{c}$}
    \DeduceC{$\Subst{!C}{t}{c}$}
    \DischargeRule{\Elim{\lexists}}{x}
    \BinaryInfC{$\Subst{!C}{t}{c}$}
    \DisplayProof
    \]
    યોગ્ય છે:
    $\Subst{!A(x,t)}{a}{x} \ident \Subst{!A(a,c)}{t}{c}$,
    અને નિષ્કર્ષ
    $\Subst{!C}{t}{c}$માં
    કે કોઈ ખુલ્લી ધારણામાં
    $a$ આવતો નથી.
    \sourcecorrection{OLMD-228}{મૂળના અસ્તિત્વ-નિવારણ કિસ્સામાં પ્રતિસ્થાપિત પહેલા નિષ્કર્ષમાં $c$ જ રહી ગયું છે, બીજા નિષ્કર્ષમાં $C$નું પ્રતિસ્થાપન છૂટ્યું છે અને અંતિમ વૃક્ષ ભૂલથી સર્વમાન્યક-પરિચયનું છે; અહીં બંને પ્રતિસ્થાપિત ઉપસાબિતીથી અસ્તિત્વ-નિવારણનું યોગ્ય વૃક્ષ મૂક્યું છે.}
  \end{enumerate}
\end{proof}

\begin{prop}\ollabel{prop:regular}
જો~$\delta$ એ~\Log{N1c}
અથવા~\Log{N1i}માંની
સાબિતી હોય,
તો એ જેવી જ
રચના ધરાવતી
(ખાસ કરીને સરખી
ઊંચાઈ ધરાવતી)
નિયમિત સાબિતી~$\delta'$
છે, જે~$\delta$થી
માત્ર સ્વનિર્ધારક ચલો
બદલવામાં જુદી પડે છે.
\end{prop}

\begin{proof}
આ સાબિતી માટે~$\delta$માંના
સ્વનિર્ધારક ચલ ધરાવતા
અનુમાનને \emph{સ્વચ્છ}
કહીએ, જો તેનો સ્વનિર્ધારક
ચલ બીજા કોઈ અનુમાનનો
સ્વનિર્ધારક ચલ ન હોય
અને તે આ અનુમાનની
તાત્કાલિક ઉપ-સાબિતી
સિવાય~$\delta$માં ક્યાંય
ન આવતો હોય; નહિતર તેને
\emph{અસ્વચ્છ} કહીએ.
$n$ એ~$\delta$માંનાં
અસ્વચ્છ અનુમાનોની સંખ્યા
હોય. $n$ પર આગમનથી
બતાવીશું કે~$\delta$ને
જરૂરી નિયમિત
સાબિતી~$\delta'$માં
ફેરવી શકાય છે.
જો~$n=0$, તો~$\delta$નાં બધાં
સ્વનિર્ધારક ચલવાળાં અનુમાનો
સ્વચ્છ છે, એટલે~$\delta$
નિયમિત છે અને કંઈ સાબિત
કરવાનું રહેતું નથી.
નહિતર સૌથી ઉપરનું
અસ્વચ્છ સ્વનિર્ધારક ચલવાળું
અનુમાન પસંદ કરો;
માનો કે તે~\Intro{\lforall}
છે. તેનામાં પૂરી થતી
$\delta$ની
ઉપ-સાબિતી~$\delta_1$
આ સ્વરૂપની છે:
    \[
    \AxiomC{}
    \RightLabel{$\delta_2$}
    \DeduceC{$!A(c)$}
    \RightLabel{\Intro{\lforall}}
    \UnaryInfC{$\lforall[x][!A(x)]$}
    \DisplayProof
    \]
$c$ સ્વનિર્ધારક ચલ હોવાથી
તે~$\lforall[x][!A(x)]$માં
કે કોઈ ખુલ્લી ધારણામાં
આવતો નથી.
તેમાં અસ્વચ્છ
સ્વનિર્ધારક ચલવાળું અનુમાન
નથી, તેથી~$\delta_2$
નિયમિત છે. વળી પસંદ કરેલો
ચલ તેના કોઈ અનુમાનનો
સ્વનિર્ધારક ચલ નથી;
હોય તો તે અનુમાન પણ
અસ્વચ્છ બનત.
\sourcecorrection{OLMD-229}{મૂળમાં સૌથી ઉપરનું કોઈ પણ સ્વનિર્ધારક-ચલવાળું અનુમાન પસંદ થાય છે અને ઉપરની ઉપસાબિતીમાં આવાં અનુમાનો જ નથી એમ કહેવાયું છે; આગમન માટે સૌથી ઉપરનું \emph{અસ્વચ્છ} અનુમાન જોઈએ, જેથી તેની ઉપર સ્વચ્છ અનુમાનો રહી શકે પણ ઉપસાબિતી નિયમિત રહે.}
$\delta$માં ન આવતો
અચળ~$a$ લો.
\olref{lem:subst-regular}થી
$\Subst{\delta_2}{a}{c}$
એ~$!A(a)$ની નિયમિત
સાબિતી છે.
સ્વનિર્ધારક ચલની શરતથી
$\delta_2$ની કોઈ ખુલ્લી
ધારણામાં~$c$ નથી,
તેથી~$\Subst{\delta_2}{a}{c}$ની
કોઈ ખુલ્લી ધારણામાં~$a$
નથી. આથી~$\Intro{\lforall}$
લાગુ કરી શકાય છે.
જો~$\delta$માં
$\delta_1$ને બદલે
સાબિતી~$\delta_1'$ મૂકીએ,
    \[
    \AxiomC{}
    \RightLabel{$\Subst{\delta_2}{a}{c}$}
    \DeduceC{$!A(a)$}
    \RightLabel{\Intro{\lforall}}
    \UnaryInfC{$\lforall[x][!A(x)]$}
    \DisplayProof
    \]
તો અસ્વચ્છ સ્વનિર્ધારક
ચલવાળું અનુમાન સ્વચ્છ
અનુમાનથી બદલાયું છે,
અને ફરક માત્ર
સ્વનિર્ધારક ચલ~$c$ની
જગ્યાએ~$a$ મૂકવાનો છે.
મળેલી સાબિતીમાં
વધુમાં વધુ~$n-1$
અસ્વચ્છ સ્વનિર્ધારક
ચલવાળાં અનુમાનો છે;
આગમન-પરિકલ્પનાથી
પરિણામ મળે છે.
\sourcecorrection{OLMD-230}{મૂળમાં અસ્વચ્છ અનુમાનોની સંખ્યા ચોક્કસ $n-1$ થાય છે એમ કહેવાયું છે; તાજા ચલથી બીજાં અનુમાનો પણ સ્વચ્છ થઈ શકે, તેથી આગમન માટે જરૂરી વધુમાં વધુ $n-1$ દાવો લખ્યો છે.}
\end{proof}

\begin{prob}
\olref{prop:regular}ની
સાબિતીમાં સૌથી ઉપરનું
અસ્વચ્છ સ્વનિર્ધારક
ચલવાળું અનુમાન~\Elim{\lexists}
હોય તે કિસ્સો વર્ણવો.
\end{prob}

હમણાં સાબિત કરેલી
પ્રતિજ્ઞાપનનો સ્પષ્ટ
ઉલ્લેખ હવે નહીં કરીએ;
પણ ચર્ચામાં આવતી
બધી સાબિતીઓ નિયમિત
છે એમ માનીશું.
નિયમિત ન હોય તો
સ્વનિર્ધારક ચલો બદલીને
તેમને નિયમિત બનાવી
શકાય છે.

\Olref{lem:subst-regular}
અને~\olref{prop:regular}
\Log{N2c} તથા~\Log{N2i}
માટે પણ સાચાં છે.
તેમની સાબિતીઓ કસરત
તરીકે રાખીએ છીએ.
% END OLP-0665
\endgroup

\begingroup
\makeatletter\def\ol@sectioncs{section}\makeatother

% BEGIN OLP-0662
\olfileid{pt}{nat}{gra}

\olsection{સાબિતીઓનું કલમ-જોડાણ}
\phantomsection\label{guunit:OLP-0662}


\Log{N1c} અને~\Log{N1i}માંની
સાબિતીઓ એ ધારણાઓ~$!B$માંથી
સૂત્રો~$!A$ની સાબિતીઓ છે.
માનો કે~$!B$ની પણ સાબિતી છે.
તો~$!B$ની સાબિતી અને
$!B$માંથી~$!A$ની સાબિતી
જોડીને~$!A$ની એવી
સાબિતી કેવી રીતે
મેળવવી જેમાં~$!B$ની ધારણાનો
ઉપયોગ ન થતો હોય?

એક વિકલ્પ એ છે કે~\Intro{\lif}
લગાવી~$!A$ની સાબિતીમાંથી
ધારણા~$!B$ દૂર કરીએ.
પછી મળેલી~$!B \lif !A$ની
સાબિતીને~$!B$ની સાબિતી
સાથે જોડીને આ મળે:
\[
\AxiomC{$\Discharge{!B}{x}$}
\DeduceC{$!A$}
\DischargeRule{\Intro{\lif}}{x}
\UnaryInfC{$!B \lif !A$}
\AxiomC{}
\DeduceC{$!B$}
\RightLabel{\Elim{\lif}}
\BinaryInfC{$!A$}
\DisplayProof
\]
\sourcecorrection{OLMD-224}{મૂળ વૃક્ષમાં અનુસરણ-પ્રવેશનું નિષ્કર્ષ $!B \lif !C$ લખાયું છે; ઉપરની ઉપસાબિતીનું નિષ્કર્ષ $!A$ હોવાથી અહીં $!B \lif !A$ જોઈએ.}
આ રીત સીધી છે, પણ થોડી
અકુદરતી લાગે છે: તરત જ
દૂર કરવાના~$\lif$ને શા માટે
પહેલાં દાખલ કરવું?
$!B \lif !A$માંથી પસાર
થતો આવો ``ચકરાવો''
ટાળી શકાય એવો હોવો જોઈએ.
(આવા ચકરાવા હંમેશાં ટાળી
શકાય છે એ જ સામાન્યીકરણ
પ્રમેયનો આશય છે.)

\Log{N1c} અને~\Log{N1i}માં
આ કિસ્સામાં ચકરાવો સહેલાઈથી
ટાળી શકાય છે: ધારણાઓ~$!B^x$
જ્યાં આવે ત્યાં~$!B$ની
આખી સાબિતી મૂકી દઈએ.
બીજી સાબિતીની ધારણાઓ પર
સાબિતીઓનું આવું
``કલમ-જોડાણ'' ઉપયોગી છે,
એટલે તેને વ્યાખ્યા તરીકે
નોંધીએ છીએ.

\begin{defn}
જો~$\delta_1$ એ ખુલ્લી
ધારણા~$!B^x$માંથી~$!A$ની
\Log{N1c}-સાબિતી
(અથવા~\Log{N1i}-સાબિતી)
હોય, અને~$\delta_2$ એ~$!B$ની
\Log{N1c}-સાબિતી
(\Log{N1i}-સાબિતી) હોય,
તો~$\delta_1$માંનાં
મુક્ત ન કરાયેલાં~$!B^x$નાં
દરેક સ્થાનને~$\delta_2$થી
બદલવાથી મળતું પરિણામ
$\Subst{\delta_1}{\delta_2}{!B^x}$
છે.
\end{defn}

જો~$\delta_1$ અને~$\delta_2$માં
જુદાં ધારણા-સૂત્રોને
એક જ ચિહ્ન ન આપેલું હોય,
અને~$\delta_2$ની કોઈ ખુલ્લી
ધારણામાં~$\delta_1$નો
સ્વનિર્ધારક ચલ ન આવતો હોય,
તો મળેલી
$\Subst{\delta_1}{\delta_2}{!B^x}$
યોગ્ય સાબિતી છે.
તેની ખુલ્લી ધારણાઓ
$\delta_1$ તથા~$\delta_2$ની
ખુલ્લી ધારણાઓ છે.
સાબિતીઓનું કલમ-જોડાણ
કરતી વખતે આ શરતો સામાન્ય રીતે
પૂરી થાય છે. $\delta_1$ અને
$\delta_2$ કોઈ મોટી
સાબિતીની ઉપસાબિતીઓ હોય
તો પહેલી શરત પૂરી થાય છે.
બીજી શરત પૂરી ન થતી હોય
તો~$\delta_1$ના સ્વનિર્ધારક
ચલોનાં નામ બદલીને તેને
પૂરી કરી શકાય છે.

\Log{N2c} અને~\Log{N2i}
માટે પણ અનુરૂપ વ્યાખ્યા છે.

\begin{defn}
જો~$\delta_1$ એ
$x:!B, \Gamma_1 \Sequent !A$ની
\Log{N2c}-સાબિતી
(અથવા~\Log{N2i}-સાબિતી)
હોય, અને~$\delta_2$ એ
$\Gamma_2 \Sequent !B$ની
\Log{N2c}-સાબિતી
(\Log{N2i}-સાબિતી) હોય,
તો~$\delta_1$માંનાં
દરેક સ્વયંસિદ્ધ
$x:!B \Sequent !B$ને
$\delta_2$થી બદલીને,
અને એવાં સ્વયંસિદ્ધોની
નીચે આવતા દરેક સિક્વન્ટના
સંદર્ભમાં~$\Gamma_2$ ઉમેરવાથી
મળતું પરિણામ
$\Subst{\delta_1}{\delta_2}{x:!B}$
છે.
\sourcecorrection{OLMD-225}{મૂળની આ N2 વ્યાખ્યામાં $\delta_2$ને ભૂલથી N1c/N1i-સાબિતી કહેવાયું છે; સિક્વન્ટવાળી આધારસાબિતી માટે અનુરૂપ N2c/N2i જોઈએ.}
\end{defn}

\begin{prop}
જો $\delta_1 \Proves x:!B, \Gamma_1 \Sequent
!A$ અને $\delta_2 \Proves \Gamma_2 \Sequent !B$,
તો, જો~$\delta_1$નો કોઈ
સ્વનિર્ધારક ચલ~$\Gamma_2$માં
ન આવતો હોય,
$\Subst{\delta_1}{\delta_2}{x:!B} \Proves \Gamma_1, \Gamma_2 \Sequent
!A$.
\end{prop}


\begin{proof}
  $\pheight{\delta_1}$ પર આગમનથી.
  \sourcecorrection{OLMD-226}{મૂળની એક-વાક્યીય સાબિતીમાં અપરિભાષિત $\delta$ની ઊંચાઈ લખાઈ છે; પ્રતિજ્ઞાપનમાં આવેલી બદલાતી સાબિતી $\delta_1$ પર આગમનનો આશય છે.}
\end{proof}
% END OLP-0662
\endgroup

\begingroup
\makeatletter\def\ol@sectioncs{section}\makeatother

% BEGIN OLP-0670
\olfileid{pt}{nat}{tgi}

\olsection{\Log{G2i}માંથી \Log{N2i}માં રૂપાંતર}
\phantomsection\label{guunit:OLP-0670}


\Log{G2i}ના સિક્વન્ટોના
પૂર્વાંગો સૂત્રોનાં
બહુગણો~$\Gamma$ છે,
જ્યારે~\Log{N2i}ના
સિક્વન્ટોના પૂર્વાંગો
એવાં ચિહ્નિત સૂત્રોના
ગણો~$\Gamma'$ છે
જેમાં દરેક ચિહ્ન માત્ર
એક વાર આવે છે.
ચિહ્નિત સૂત્રોનો
ગણ~$\Gamma'$,
બહુગણ~$\Gamma$ને
\emph{અનુરૂપ} છે એમ
ત્યારે કહીએ જ્યારે
દરેક સૂત્ર~$!A$
માટે~$\Gamma'$માંના
$x : !A$ સ્વરૂપનાં
ઘટકોની સંખ્યા,
$\Gamma$માં~$!A$ની
બહુલતા (એટલે~$!A$ની
નકલોની સંખ્યા)
કરતાં વધુ ન હોય.

\begin{prop}\ollabel{prop:G2i-to-N2i}
જો $\Log{G2i} + \Cut \Proves \Gamma \Sequent \Delta$,
તો~$\Log{N2i}
\Proves \Gamma' \Sequent !C$,
જ્યાં~$\Delta$ ખાલી
હોય ત્યારે
$!C \ident \lfalse$
અને નહિતર
$\Delta = \{!A\}$ હોય,
તથા~$\Gamma'$ એ~$\Gamma$ને
અનુરૂપ હોય.
\end{prop}

\begin{proof}
  બતાવીએ કે જો~$\pi$
  એ~\Log{G2i}+\Cut{}માં
  $\Gamma \Sequent \Delta$ની
  સાબિતી હોય,
  તો~$\Gamma$ને અનુરૂપ
  એવો~$\Gamma'$ અને
  \Log{N2i}માં
  $\Gamma' \Sequent !C$ની
  સાબિતી~$\delta$
  છે.
  \sourcecorrection{OLMD-240}{મૂળ સાબિતીના આરંભમાં અપરિભાષિત G2ci તંત્ર લખાયું છે; પ્રતિજ્ઞાપનનું તંત્ર G2i સાથે Cut છે.}
  આ માટે
  $n = \pheight{\pi}$
  પર આગમન કરીએ.

  જો~$n = 0$, તો~$\pi$
  સ્વયંસિદ્ધ છે.
  તે~$\lfalse \Sequent
  \quad$ સ્વયંસિદ્ધ હોય,
  તો~$\delta$ એ
  $x:\lfalse \Sequent \lfalse$
  સ્વયંસિદ્ધ છે,
  અર્થાત્
  $\Gamma' = \{x:\lfalse\}$
  અને~$!C \ident \lfalse$.
  જો~$\pi$ એ
  $!A \Sequent !A$
  સ્વરૂપનું સ્વયંસિદ્ધ
  હોય, તો~$\delta$
  એ~$x: !A \Sequent !A$ છે.
  
  જો~$n>0$, તો~$\pi$નો
  અંત નિયમ~$R$માં થાય છે.
  આગમન-પરિકલ્પનાથી
  \Log{N2i}માં એ
  નિયમનાં આધારવિધાનોને
  અનુરૂપ સિક્વન્ટોની
  સાબિતીઓ મળે છે.
  \sourcecorrection{OLMD-241}{મૂળના આગમન-વાક્યમાં N2c લખાયું છે; આ દાવો અને આખી સાબિતી N2i માટે છે.}
  જરૂરી હોય તો ફરી
  ચિહ્નિત કરીને ધારી
  શકીએ કે એ સાબિતીઓનાં
  બધાં મુક્તિ-ચિહ્નો
  જુદાં છે, અને
  કોઈ ચિહ્ન~$x$
  મુક્ત થતું હોય
  તો તેને મુક્ત કરતા
  અનુમાનની ઉપર જ
  આવે છે.
  હવે એ સાબિતીઓને
  જોડીને યોગ્ય
  $\Gamma' \Sequent !C$ની
  સાબિતી કેવી રીતે
  મળે તે બતાવીએ.
  $R$ પ્રમાણે કિસ્સા
  જુદા પાડીએ.

  \begin{enumerate}
    \item જો~$R$ એ
    \RightR{\Weakening}
    હોય, તો~$\pi$નો
    અંત આ રીતે થાય છે:
    \[
    \AxiomC{}
    \RightLabel{$\pi_1$}
    \Deduce$\Gamma \fCenter$
    \RightLabel{\RightR{\Weakening}}
    \UnaryInf$\Gamma \fCenter !A$
    \DisplayProof
    \]
    $\pi_1$ પર
    આગમન-પરિકલ્પના
    લાગુ કરતાં
    $\Gamma' \Sequent \lfalse$ની
    સાબિતી મળે છે.
    $\FalseInt$ લાગુ કરી
    $\Gamma' \Sequent !A$ની
    સાબિતી મળે છે.
    \item જો~$R$ એ
    \LeftR{\Weakening}
    હોય, તો~$\pi$નો
    અંત આ રીતે થાય છે:
    \[
    \AxiomC{}
    \RightLabel{$\pi_1$}
    \Deduce$\Gamma \fCenter \Delta$
    \RightLabel{\LeftR{\Weakening}}
    \UnaryInf$!A, \Gamma \fCenter \Delta$
    \DisplayProof
    \]
    \sourcecorrection{OLMD-242}{મૂળના ડાબા શિથિલીકરણના વૃક્ષમાં નિયમ-લેબલ ભૂલથી જમણું શિથિલીકરણ છે; પૂર્વાંગમાં સૂત્ર ઉમેરાતું હોવાથી ડાબું શિથિલીકરણ લખ્યું છે.}
    $\pi_1$ પર
    આગમન-પરિકલ્પના
    લગાવતાં
    $\Gamma' \Sequent !C$ની
    સાબિતી મળે છે;
    તેને~$\delta$ લઈએ.
    નોંધો કે~$\Gamma'$
    એ~$\Gamma$ને અનુરૂપ
    હોય તો~$!A,
    \Gamma$ને પણ
    અનુરૂપ છે.
    કારણ કે~$\Gamma'$માં
    $x : !A$ની સંખ્યા
    $\Gamma$માં~$!A$ની
    નકલોની સંખ્યાથી
    વધુ નથી, તેથી
    $\Gamma \cup \{!A\}$માં
    $!A$ની નકલોની
    સંખ્યાથી પણ વધુ નથી.
    \item જો~$R$ એ
    \LeftR{\Contraction}
    હોય, તો~$\pi$નો
    અંત આ રીતે થાય છે:
    \[
    \AxiomC{}
    \RightLabel{$\pi_1$}
    \Deduce$!A, !A, \Gamma \fCenter \Delta$
    \RightLabel{\LeftR{\Contraction}}
    \UnaryInf$!A, \Gamma \fCenter \Delta$
    \DisplayProof
    \]
    $\pi_1$ પર
    આગમન-પરિકલ્પના
    લગાવતાં
    $\Pi, \Gamma' \Sequent !C$ની
    સાબિતી મળે છે,
    જ્યાં~$\Pi \subseteq \{x : !A,
    y:!A\}$.
    જો~$\Pi = \{x : !A, y:!A\}$,
    તો મળેલી~$\delta$માંના
    ચિહ્નિત
    સૂત્રો~$y:!A$
    ને~$x:!A$થી બદલો.
    $\delta$ના અંતિમ સિક્વન્ટના
    સંદર્ભમાં~$x$
    આવતો હોવાથી
    $\delta$માં કોઈ
    અનુમાન~$x:!B$
    સ્વરૂપનું
    સૂત્ર મુક્ત
    કરતું નથી એમ
    ધારી શકીએ;
    એટલે~$\delta$ના
    કોઈ સિક્વન્ટના
    ડાબે આવેલું~$x:!B$
    સક્રિય નથી.
    \sourcecorrection{OLMD-243}{મૂળમાં ચિહ્નિત N2 સૂત્રોનું નામબદલણ G2i ઉપસાબિતી $\pi_1$માં કરવાનું કહેવાયું છે; આગમનથી મળેલી N2 સાબિતી $\delta$માં તે કરવું જોઈએ.}
    તેથી પરિણામ
    \Log{N2i}માં
    હજી પણ સાબિતી છે.
    \item જો~$R$ એ
    $\RightR{\lif}$ હોય,
    તો~$\pi$નો અંત
    આ રીતે થાય છે:
    \[
    \AxiomC{}
    \RightLabel{$\pi_1$}
    \Deduce$!A, \Gamma \fCenter !B$
    \RightLabel{\RightR{\lif}}
    \UnaryInf$\Gamma \fCenter !A \lif !B$
    \DisplayProof
    \]
    આગમન-પરિકલ્પનાથી
    $x:
    !A, \Gamma' \Sequent !B$
    અથવા~$\Gamma' \Sequent !B$ની
    સાબિતી~$\delta_1$
    છે, જ્યાં~$\Gamma'$
    એ~$\Gamma$ને
    અનુરૂપ છે.
    $\delta_1$ પછી
    \Intro{\lif}
    લગાવી મળતી
    સાબિતીને~$\delta$
    લઈ શકાય.
    \sourcecorrection{OLMD-244}{મૂળમાં $\delta_1$ એ $\delta$ પછી અનુસરણ-પરિચય લગાવી મળે છે એમ ઉલટું લખાયું છે; આગમનથી $\delta_1$ મળે છે અને તેની નીચે નિયમ લગાવી $\delta$ બને છે.}
    \item જો~$R$ એ
    \LeftR{\lif} હોય,
    તો~$\pi$નો
    અંત આ રીતે થાય છે:
    \[
    \AxiomC{}
    \RightLabel{$\pi_1$}
    \Deduce$\Gamma_1 \fCenter !A$
    \AxiomC{}
    \RightLabel{$\pi_2$}
    \Deduce$!B, \Gamma_2 \fCenter \Delta$
    \RightLabel{\LeftR{\lif}}
    \BinaryInf$!A \lif !B, \Gamma_1, \Gamma_2 \fCenter \Delta$
    \DisplayProof
    \]
    આગમન-પરિકલ્પનાથી
    $\Gamma_1' \Sequent !A$ની
    સાબિતી~$\delta_1$
    અને~$x: !B, \Gamma_2' \Sequent !C$
    અથવા~$\Gamma_2' \Sequent !C$ની
    સાબિતી~$\delta_2$
    મળે છે.
    $\delta_2$ બીજા
    સ્વરૂપની હોય,
    તો~$\delta_2$ને જ~$\delta$
    લઈ શકાય.
    નહિતર~$\delta_1$માંનાં
    બધા સૂત્રો તથા
    મુક્તિ-ચિહ્નો ફરી
    ચિહ્નિત કરીને
    $\delta_1$ અને~$\delta_2$
    વચ્ચે કોઈ ચિહ્ન
    સામાન્ય ન રહે
    તે સુનિશ્ચિત કરીએ.
    \sourcecorrection{OLMD-245}{મૂળમાં આ N2 ચિહ્નબદલી G2i ઉપસાબિતી $\pi_1$માં કરવાનું કહેવાયું છે; ચિહ્નિત સૂત્રો અને મુક્તિ-ચિહ્નો ધરાવતી અનુવાદિત સાબિતી $\delta_1$ છે.}
    ત્યારે~$\Gamma_1' \cup
    \Gamma_2'$ એ
    $\Gamma_1 \cup \Gamma_2$
    ને અનુરૂપ છે.
    $x:!B$ એ~$\delta_2$ના
    અંતિમ સિક્વન્ટમાં
    હોવાથી~$\delta_2$માં
    $x$ મુક્તિ-ચિહ્નવાળા
    કોઈ અનુમાનમાં
    $x:!B$ સક્રિય
    હોઈ શકે નહીં.
    $\delta_2$માંનું
    દરેક સ્વયંસિદ્ધ
    $x:!B \Sequent !B$
    નીચેની
    સાબિતી~$\delta_3$થી
    બદલો:
    \[
    \Axiom$x: !A \lif !B \fCenter !A \lif !B$
    \AxiomC{}
    \RightLabel{$\delta_1$}
    \Deduce$\Gamma_1' \fCenter !A$
    \RightLabel{\Elim{\lif}}
    \BinaryInf$x: !A \lif !B, \Gamma_1' \fCenter !B$
    \DisplayProof
    \]
    \sourcecorrection{OLMD-246}{મૂળના $\delta_3$ અનુસરણ-નિવારણ વૃક્ષના નિષ્કર્ષમાં $\Gamma_1'$ છૂટે છે; બીજું આધારવિધાન એ સંદર્ભ ધરાવે છે અને નિયમ તેને નિષ્કર્ષમાં રાખે છે.}
    અને~$\delta_2$માં
    $x:!B$ના દરેક
    સ્થાનને
    $x:!A \lif !B$થી
    બદલો.
    પરિણામ
    $\Subst{\delta_2}{\delta_3}{x:!B}$
    એ
    $x:!A \lif !B, \Gamma_1', \Gamma_2' \Sequent !C$ની
    સાબિતી છે.
    \item જો~$R$ એ
    \RightR{\land} હોય,
    તો~$\pi$નો
    અંત આ રીતે થાય છે:
    \[
    \AxiomC{}
    \RightLabel{$\pi_1$}
    \Deduce$\Gamma_1 \fCenter !A$
    \AxiomC{}
    \RightLabel{$\pi_2$}
    \Deduce$\Gamma_2 \fCenter !B$
    \RightLabel{\RightR{\land}}
    \BinaryInf$\Gamma_1, \Gamma_2 \fCenter !A \land !B$
    \DisplayProof
    \]
    આગમન-પરિકલ્પનાથી
    $\Gamma_1' \Sequent !A$ની
    સાબિતી~$\delta_1$
    અને~$\Gamma_2' \Sequent !B$ની
    સાબિતી~$\delta_2$
    મળે છે.
    $\delta_2$માંનાં
    બધા સૂત્રો તથા
    મુક્તિ-ચિહ્નો ફરી
    ચિહ્નિત કરીને
    $\delta_1$ અને~$\delta_2$
    વચ્ચે કોઈ ચિહ્ન
    સામાન્ય ન રહે
    તે સુનિશ્ચિત કરીએ.
    \sourcecorrection{OLMD-247}{મૂળમાં બીજા આગમન-નિષ્કર્ષનું સૂત્ર $!A$ લખાયું છે, પણ બીજા આધારવિધાનમાં $!B$ છે; વળી N2 ચિહ્નબદલી $\pi_2$ને બદલે $\delta_2$માં થવી જોઈએ.}
    ત્યારે~$\Gamma_1' \cup \Gamma_2'$
    એ~$\Gamma_1 \cup
    \Gamma_2$ને અનુરૂપ છે.
    $\delta_1$ તથા~$\delta_2$ના
    અંતિમ સિક્વન્ટો પર
    $\Intro{\land}$
    લગાવી
    સાબિતી~$\delta$
    મળે છે.
    \item જો~$R$ એ
    $\LeftR{\land}$ હોય,
    તો, દાખલા તરીકે,
    $\pi$નો અંત આ રીતે થાય છે:
    \[
    \AxiomC{}
    \RightLabel{$\pi_1$}
    \Deduce$!A, \Gamma \fCenter \Delta$
    \RightLabel{\LeftR{\land}}
    \UnaryInf$!A \land !B, \Gamma \fCenter \Delta$
    \DisplayProof
    \]
    આગમન-પરિકલ્પનાથી
    $x:
    !A, \Gamma' \Sequent !C$
    અથવા~$\Gamma' \Sequent !C$ની
    સાબિતી~$\delta_1$
    છે, જ્યાં~$\Gamma'$
    એ~$\Gamma$ને
    અનુરૂપ છે.
    બીજા કિસ્સામાં
    $\delta$ને~$\delta_1$
    લઈ શકાય.
    \sourcecorrection{OLMD-248}{મૂળના બીજા કિસ્સામાં $\delta_1$ને $\delta$ કહેવાનું વાક્ય દિશા ઉલટી બતાવે છે; આગમનથી મળેલી $\delta_1$ જ માંગેલી $\delta$ તરીકે લેવાય છે.}
    પહેલા કિસ્સામાં
    $x:!A$ અંતિમ
    સિક્વન્ટમાં હોવાથી
    $\delta_1$માં~$x$
    મુક્તિ-ચિહ્નવાળા
    કોઈ અનુમાનમાં
    $x:!A$ સક્રિય
    નથી. દરેક
    સ્વયંસિદ્ધ
    $x:!A
    \Sequent !A$ને
    આ~સાબિતી~$\delta_2$થી
    બદલો:
    \[
    \Axiom$x: !A \land !B \fCenter !A \land !B$
    \RightLabel{\Elim{\land}}
    \UnaryInf$x: !A \land !B \fCenter !A$
    \DisplayProof
    \]
    અને~$x:!A$ના
    દરેક સ્થાનને
    $x:!A \land !B$થી
    બદલો; એટલે
    $\Subst{\delta_1}{\delta_2}{x:!A}$
    લો. આ
    $x:!A
    \land !B, \Gamma' \Sequent !C$ની
    સાબિતી છે.
    \item જો~$R$ એ
    \Cut હોય, તો~$\pi$નો
    અંત આ રીતે થાય છે:
    \[
    \AxiomC{}
    \RightLabel{$\pi_1$}
    \Deduce$\Gamma_1 \fCenter !A$
    \AxiomC{}
    \RightLabel{$\pi_2$}
    \Deduce$!A, \Gamma_2 \fCenter \Delta$
    \RightLabel{\Cut}
    \BinaryInf$\Gamma_1, \Gamma_2 \fCenter \Delta$
    \DisplayProof
    \]
    આગમન-પરિકલ્પનાથી
    $\Gamma_1' \Sequent !A$ની
    સાબિતી~$\delta_1$
    અને
    $x:!A, \Gamma_2' \Sequent !C$
    અથવા~$\Gamma_2' \Sequent !C$ની
    સાબિતી~$\delta_2$
    મળે છે.
    બીજા કિસ્સામાં
    $\delta$ને~$\delta_2$
    લો. નહિતર
    $\delta_2$માંનાં
    દરેક સ્વયંસિદ્ધ
    $x:!A \Sequent !A$
    ને~$\delta_1$થી
    બદલો. આમ
    $\Gamma_1',
    \Gamma_2' \Sequent !C$ની
    સાબિતી
    $\Subst{\delta_2}{\delta_1}{x:!A}$
    એ~$\delta$ બને છે.
  \end{enumerate}
\end{proof}
% END OLP-0670
\endgroup

\begingroup
\makeatletter\def\ol@sectioncs{section}\makeatother

% BEGIN OLP-0671
\olfileid{pt}{nat}{tni}

\olsection{\Log{N2}માંથી \Log{G2}માં રૂપાંતર}
\phantomsection\label{guunit:OLP-0671}


\begin{prop}\ollabel{prop:N2-to-G2}
જો $\Log{N2c}\ (\Log{N2i}) \Proves
\Gamma \Sequent !A$,
તો~$\Log{G2c}\ (\Log{G2i}) + \Cut \Proves \Gamma'
\Sequent \Delta$,
જ્યાં~$\Gamma$માંથી
ચિહ્નો દૂર કરતાં
મળતો સૂત્રોનો
બહુગણ~$\Gamma'$ છે,
અને~$!A$ એ~$\lfalse$
ન હોય તો~$\Delta = \{!A\}$,
જ્યારે હોય તો
$\Delta = \emptyset$.
\sourcecorrection{OLMD-249}{મૂળના પરિણામમાં G2c પછી કૌંસમાં N2i લખાયું છે; અનુરૂપ અંતઃપ્રજ્ઞાવાદી સિક્વન્ટ-કલન G2i છે.}
\end{prop}

\begin{proof}
ફરી~$\Gamma \Sequent !A$ની
સાબિતી~$\delta$ની
ઊંચાઈ પર આગમન કરીએ.
જો~$\delta$ એ
$x:!A \Sequent !A$
સ્વયંસિદ્ધ હોય,
તો~$\pi$ એ અનુરૂપ
$!A \Sequent !A$
સ્વયંસિદ્ધ છે;
$!A \ident \lfalse$
હોય ત્યારે
$\lfalse \Sequent \ $
લેવું.

જો~$\delta$નો અંત
અનુમાન-નિયમમાં થાય,
તો કિસ્સા જુદા પાડીએ:
\begin{enumerate}
  \item જો~$R$ એ
  મુખ્ય સૂત્ર~$!B \lif !C$
  ધરાવતું~\Intro{\lif}
  હોય અને~$x:!B$
  આધારવિધાનના સંદર્ભમાં
  હોય પણ નિષ્કર્ષના
  સંદર્ભમાં ન હોય,
  તો~$\delta$નો અંત
  આ રીતે થાય છે:
  \[
  \AxiomC{}
  \RightLabel{$\delta_1$}
  \Deduce$x:!B, \Gamma \fCenter !C$
  \RightLabel{\Intro{\lif}}
  \UnaryInf$\Gamma \fCenter !B \lif !C$
  \DisplayProof
  \]
  આગમન-પરિકલ્પનાથી
  $!B, \Gamma' \Sequent !C$ની
  (અથવા~$!C \ident \lfalse$
  હોય તો~$!B, \Gamma'
  \Sequent \ $ની)
  સાબિતી~$\pi_1$
  મળે છે. પછી
  \RightR{\lif} લગાવી
  $\pi$ મળે છે
  ($!C \ident \lfalse$
  હોય તો પહેલાં
  \RightR{\Weakening}
  લગાવીએ).
  \item જો~$R$ એ
  \Intro{\lif} હોય,
  પણ~$x:!B$
  આધારવિધાનના સંદર્ભમાં
  ન હોય, તો~$\delta$નો
  અંત આ રીતે થાય છે:
  \[
  \AxiomC{}
  \RightLabel{$\delta_1$}
  \Deduce$\Gamma \fCenter !C$
  \RightLabel{\Intro{\lif}}
  \UnaryInf$\Gamma \fCenter !B \lif !C$
  \DisplayProof
  \]
  આગમન-પરિકલ્પનાથી
  $\Gamma' \Sequent !C$ની
  સાબિતી~$\pi_1$
  મળે છે. તેમાં
  પહેલાં~$!B$ સાથે
  \LeftR{\Weakening}
  અને પછી~\RightR{\lif}
  લગાવી~$\pi$
  મેળવીએ છીએ.
  \item જો~$R$ એ
  \Elim{\lif} હોય,
  તો~$\delta$નો
  અંત આ રીતે થાય છે:
  \[
  \AxiomC{}
  \RightLabel{$\delta_1$}
  \Deduce$\Gamma_1 \fCenter !B \lif !A$
  \AxiomC{}
  \RightLabel{$\delta_2$}
  \Deduce$\Gamma_2 \fCenter !B$
  \RightLabel{\Elim{\lif}}
  \BinaryInf$\Gamma_1, \Gamma_2 \fCenter !A$
  \DisplayProof
  \]
  અહીં~$\Gamma = \Gamma_1 \cup \Gamma_2$.
  આગમન-પરિકલ્પનાથી
  $\Gamma_1' \Sequent !B \lif !A$ની
  સાબિતી~$\pi_1$
  અને~$\Gamma_2' \Sequent !B$ની
  સાબિતી~$\pi_2$
  મળે છે.
  સાબિતી~$\pi$
  આ રીતે મળે છે:
  \[
  \AxiomC{}
  \RightLabel{$\pi_1$}
  \Deduce$\Gamma_1' \fCenter !B \lif !A$
  \AxiomC{}
  \RightLabel{$\pi_2$}
  \Deduce$\Gamma_2' \fCenter !B$
  \Axiom$!A \fCenter !A$
  \RightLabel{\LeftR{\lif}}
  \BinaryInf$!B \lif !A, \Gamma_2' \fCenter !A$
  \RightLabel{\Cut}
  \BinaryInf$\Gamma_1', \Gamma_2' \fCenter !A$
  \DisplayProof
  \]
  બહુગણ~$\Gamma_1' \cup \Gamma_2'$
  કદાચ~$\Gamma' =
  (\Gamma_1 \cup \Gamma_2)'$
  જેટલો ન હોય:
  દાખલા તરીકે,
  $\Gamma_1$ અને~$\Gamma_2$
  બંનેમાં~$x:!C$ હોય,
  તો~$\Gamma_1' \cup \Gamma_2'$
  માં~$!C$ની બે નકલો
  છે, જ્યારે
  $(\Gamma_1 \cup \Gamma_2)'$
  માં એક જ છે.
  યોગ્ય
  \LeftR{\Contraction}
  અનુમાનો ઉમેરી
  આ ભેદ દૂર કરી શકાય.

  જો~$!A \ident \lfalse$
  હોય, તો~$\delta_2$ પર
  આગમન-પરિકલ્પના
  લગાવતાં
  $\Gamma_2' \Sequent !B$ની
  સાબિતી મળે છે.
  તેને પહેલા આધારવિધાનની
  સાબિતી અને અસત્યના
  ખાલી-ઉત્તરાંગ સ્વયંસિદ્ધ
  સાથે~\LeftR{\lif}
  અને પછી~\Cut વડે
  જોડીને
  સાબિતી~$\pi$
  મળે છે; અહીં~$!A$
  માટે ઉત્તરાંગ ખાલી છે.
  \sourcecorrection{OLMD-250}{મૂળમાં અસત્ય-નિષ્કર્ષના કિસ્સામાં $\delta_2$નું અનુવાદિત ઉત્તરાંગ ખાલી કહેવાયું છે; તે હજી $!B$ સાબિત કરે છે. ખાલી ઉત્તરાંગ અસત્ય-ડાબા સ્વયંસિદ્ધમાંથી આવે છે, જેને અનુસરણ-ડાબા નિયમ અને કટથી જોડવું પડે.}
\item જો~$R$ એ
\FalseCl હોય,
તો~$\delta$નો
અંત આ રીતે થાય છે:
  \[
  \AxiomC{}
  \RightLabel{$\delta_1$}
  \Deduce$x: \lnot !A, \Gamma \fCenter \lfalse$
  \RightLabel{\FalseCl}
  \UnaryInf$\Gamma \fCenter !A$
  \DisplayProof
  \]
  આગમન-પરિકલ્પનાથી
  $\lnot !A,
  \Gamma' \Sequent \ $ની
  સાબિતી~$\pi_1$
  છે. સાબિતી~$\pi$
  આ રીતે લઈએ:
  \[
  \Axiom$!A \fCenter !A$
  \RightLabel{\RightR{\lnot}}
  \UnaryInf$\fCenter !A, \lnot !A$
  \AxiomC{}
  \RightLabel{$\pi_1$}
  \Deduce$\lnot !A, \Gamma' \fCenter $
  \RightLabel{\Cut}
  \BinaryInf$\Gamma' \fCenter !A$
  \DisplayProof
  \]
  \sourcecorrection{OLMD-251}{મૂળના અંતિમ G2 વૃક્ષમાં $\pi_1$ના સિક્વન્ટમાં N2નું $x:$ ચિહ્ન રહી ગયું છે; G2 પૂર્વાંગ ચિહ્નવિહોણા સૂત્રોનું બહુગણ છે.}
\end{enumerate}
નોંધો કે ફક્ત
\FalseCl{}ના રૂપાંતરમાં
જમણી બાજુ એકથી વધુ
સૂત્ર ધરાવતી
સાબિતી~$\pi$
મળે છે. એટલે
\Log{N2i}ની
સાબિતીનું રૂપાંતર
\Log{G2i}-સાબિતી
છે.
\end{proof}
% END OLP-0671
\endgroup


\OLEndChapterHook
% END OLP-0664
\endgroup

\begingroup
\makeatletter\def\ol@sectioncs{section}\makeatother

% BEGIN OLP-0674
\olchapter{pt}{nor}{સામાન્યીકરણ}
\olchapterid{pt}{nor}
\phantomsection\label{guunit:OLP-0674}


\begingroup
\makeatletter\def\ol@sectioncs{section}\makeatother

% BEGIN OLP-0672
\olfileid{pt}{nor}{int}

\olsection{પરિચય}
\phantomsection\label{guunit:OLP-0672}


જો તમે અનુસરણ~$!A \lif !B$ને
સાબિત કર્યું હોય, તો
\Elim{\lif} વાપરીને
તેને~$!B$ની સાબિતીમાં
કામે લગાડી શકો.
અલબત્ત, આ માટે~$!A$ની
સાબિતી~$\delta_1$ પણ
જોઈએ. અનુસરણ~$!A \lif !B$ની
તમારી સાબિતી~$\delta$
મોટે ભાગે~\Intro{\lif}થી
પૂરી થતી હશે:
તે નિયમ ખુલ્લી
ધારણા~$!A$ હેઠળની
$!B$ની સાબિતી~$\delta_2$ને
$!A \lif !B$ની સાબિતીમાં
ફેરવે છે. એમ હોય તો
અંતિમ સાબિતીનું
સ્વરૂપ આ છે:
\[
\AxiomC{$\Discharge{!A}{x}$}
\RightLabel{$\delta_2$}
\DeduceC{$!B$}
\DischargeRule{\Intro{\lif}}{x}
\UnaryInfC{$!A \lif !B$}
\AxiomC{}
\RightLabel{$\delta_1$}
\DeduceC{$!A$}
\RightLabel{\Elim{\lif}}
\BinaryInfC{$!B$}
\DisplayProof
\]
\sourcecorrection{OLMD-252}{મૂળ ગદ્યમાં $\delta_1$ને પહેલે $!A$ની અને પછી $!A \lif !B$ની સાબિતી કહેવામાં આવી છે; વૃક્ષના બીજા આધાર ઉપર પણ $\delta_2$ ખોટી રીતે પુનરાવર્તિત છે. અહીં અનુસરણની સાબિતી $\delta$ અને $!A$નો આધાર $\delta_1$ છે.}
પહેલેથી મળેલી
$!A \lif !B$ની સાબિતી
ફરી વાપરવાથી તમારી
રીત કાર્યક્ષમ છે,
પરંતુ તેની સાબિતી
કાર્યક્ષમ નથી.
$!A \lif !B$ને રજૂ
કરીને તરત જ તેનો
લોપ કરવો ચકરાવો લાગે
છે. $!B$ને સાબિત
કરવાની વધુ સીધી
રીત હોવી જોઈએ.

ખરેખર, એવી રીત હંમેશાં
હોય છે. ઉપરના જેવા
``ચકરાવા''—દાખલ
નિયમ પછી તરત
લોપ નિયમ—સ્વાભાવિક
નિગમનની સાબિતીઓમાંથી
હંમેશાં દૂર કરી શકાય.
તેને \emph{સામાન્યીકરણ}
કહે છે અને મળતી
સાબિતીને \emph{સામાન્ય}
સાબિતી (અથવા
\emph{સામાન્ય સ્વરૂપમાં}
હોય એવી સાબિતી)
કહે છે.

આ પરિણામની સાબિતી,
સિક્વન્ટ-કલનના કટ-નિવારણ
પ્રમેયની જેમ, થોડી
જટિલ છે અને તેમાં
ઘણા કિસ્સા જોવા પડે છે.
અંતઃપ્રજ્ઞાવાદી~\Log{N1i}
કરતાં શાસ્ત્રીય~\Log{N1c}
માટે તે સાબિત કરવું
વધુ મુશ્કેલ છે.
સિક્વન્ટ-કલનમાં \Cut{}
નિયમ વાપરીને તેની
વિના કરતાં વધુ
પરિણામો સાબિત
કરી શકાતાં નથી એવું
કટ-નિવારણ બતાવે છે;
તેવી જ રીતે,
ચકરાવા વાપરીને તેમને
ટાળવા કરતાં વધુ
પરિણામો સાબિત કરી
શકાતાં નથી એવું
સામાન્યીકરણ બતાવે છે.
બીજી સામ્યતાઓ પણ છે:
એક જ પરિણામની ચકરાવાવાળી
સૌથી ટૂંકી સાબિતી
કરતાં સામાન્ય
સાબિતી ઘણી લાંબી
હોઈ શકે છે; તેવી જ
રીતે, સિક્વન્ટ-કલનમાં
કટવાળી સાબિતીઓ
સૌથી ટૂંકી કટ-રહિત
સાબિતી કરતાં ઘણી
ટૂંકી હોઈ શકે છે.
ઉપ-સૂત્ર
ગુણધર્મ કહે છે કે
પરિણામના ઉપ-સૂત્રો
વાપરીને જ તે પરિણામ
સાબિત કરતી
સાબિતી હંમેશાં
મેળવી શકાય છે.
સિક્વન્ટ-કલનમાં આ
ગુણધર્મ કટ-નિવારણમાંથી
અને સ્વાભાવિક નિગમનમાં
સામાન્યીકરણમાંથી ફલિત
થાય છે.
% END OLP-0672
\endgroup

\begingroup
\makeatletter\def\ol@sectioncs{section}\makeatother

% BEGIN OLP-0677
\olfileid{pt}{nor}{seg}

\olsection{ખંડો અને કટ}
\phantomsection\label{guunit:OLP-0677}


દાખલ નિયમ
પછી લોપ
નિયમ આવે તે
સાબિતીમાં ચકરાવો
છે; સાબિતીનું
સામાન્યીકરણ એટલે
એવા ચકરાવા
દૂર કરવા.
જોકે~\Elim{\lor}
અને~\Elim{\lexists}
હોવાને કારણે
``દૂર કરવાનો
ચકરાવો'' શું છે
તેની વ્યાખ્યા
પણ થોડી જટિલ
બને છે. સમસ્યા
એ છે કે
આ નિયમોને લીધે
ચકરાવામાં
લોપ
નિયમ દાખલ
નિયમની \emph{તરત}
પછી આવે જ
તે જરૂરી નથી.
વચ્ચે~\Elim{\lor}
કે~\Elim{\lexists}
આવી શકે છે,
જેમ કે:
\[
\AxiomC{}
\DeduceC{$\lexists[x][!C(x)]$}
\AxiomC{$\Discharge{!A}{x}$}
\DeduceC{$!B$}
\DischargeRule{\Intro{\lif}}{x}
\UnaryInfC{$!A \lif !B$}
\DischargeRule{\Elim{\lexists}}{y}
\BinaryInfC{$!A \lif !B$}
\AxiomC{}
\DeduceC{$!A$}
\RightLabel{\Elim{\lif}}
\BinaryInfC{$!B$}
\DisplayProof
\]
ખરેખર, ગમે તેટલા
\Elim{\lor} અથવા
\Elim{\lexists}
નિયમો~\Intro{\lif}
અને~\Elim{\lif}
વચ્ચે આવી શકે
છે, જેમ કે:
\[
\AxiomC{}
\DeduceC{$\lexists[x][!C(x)]$}
\AxiomC{}
\DeduceC{$!D \lor !E$}
\AxiomC{$\Discharge{!A}{x}$}
\DeduceC{$!B$}
\DischargeRule{\Intro{\lif}}{x}
\UnaryInfC{$!A \lif !B$}
\AxiomC{$\Discharge{!A}{x}$}
\DeduceC{$!B$}
\DischargeRule{\Intro{\lif}}{x}
\UnaryInfC{$!A \lif !B$}
\RightLabel{\Elim{\lor}}
\TrinaryInfC{$!A \lif !B$}
\DischargeRule{\Elim{\lexists}}{y}
\BinaryInfC{$!A \lif !B$}
\AxiomC{}
\DeduceC{$!A$}
\RightLabel{\Elim{\lif}}
\BinaryInfC{$!B$}
\DisplayProof
\]
આ દાખલો બતાવે
છે કે એક જ
\Elim{\lif} અનેક
ચકરાવામાં સામેલ
હોઈ શકે છે,
કારણ કે તેના
પહેલાં એકથી વધુ
\Intro{\lif} આવે
છે. આ બધી
શક્યતાઓ સમાવવા
\Log{N1i}-સાબિતીમાં
સૂત્ર આવૃત્તિઓના
અમુક પ્રકારના
ક્રમ, જેને
\emph{ખંડો} કહીએ,
તે વડે ચકરાવા
વ્યાખ્યાયિત કરીએ.
આપણા દાખલામાં
તે~$!A \lif !B$ની
આવૃત્તિઓનો એવો
ક્રમ છે જેમાં
$\Elim{\lor}$ કે
$\Elim{\lexists}$ના
ગૌણ આધારવિધાન
પછી એ નિયમનું
નિષ્કર્ષ આવે
છે. અહીં આવા
બે ખંડ છે,
અને દરેકમાં~$!A \lif !B$ની
ત્રણ આવૃત્તિઓ
છે. હવે ઔપચારિક
વ્યાખ્યા આપીએ.

\begin{defn}
\Log{N1i}ની સાબિતી~$\delta$માં
એક જ સૂત્ર~$!A$ની
આવૃત્તિઓનો ક્રમ
$!A_1$, \dots, $!A_n$
એ જ~$\delta$માં
નીચેની શરતો પૂરી
કરે તો તે
\emph{ખંડ} છે:
\begin{enumerate}
  \item $!A_1$ એ
  \Elim{\lor} કે
  \Elim{\lexists}નું
  નિષ્કર્ષ ન હોય;
  \item $!A_n$ એ
  \Elim{\lor} કે
  \Elim{\lexists}નું
  ગૌણ આધારવિધાન
  ન હોય;
  \item $i = 1$,
  \dots,~$n-1$
  માટે~$!A_i$ એ
  \Elim{\lor}
  અથવા~\Elim{\lexists}નું
  ગૌણ આધારવિધાન
  હોય અને
  $!A_{i+1}$
  તે જ અનુમાનનું
  નિષ્કર્ષ હોય.
\end{enumerate}
\sourcecorrection{OLMD-259}{મૂળની ત્રીજી શરતમાં પહેલાની આવૃત્તિ $!A_i$ પછીની આવૃત્તિ $A_{i+1}$ લખાઈ છે; બંને એક જ $!A$ સૂત્રની આવૃત્તિઓ હોવાથી બીજામાં પણ $!$ જરૂરી છે.}
તે ખંડની
\emph{લંબાઈ}~$n$ છે.
\end{defn}

દરેક ખંડ દૂર
કરવા યોગ્ય
ચકરાવો નથી.
જે એવા છે
તેમને \emph{કટ}
કહીએ છીએ.

\begin{defn}
\emph{કટ} એવો
ખંડ છે જેમાં
$!A_n$ એ
લોપ
અનુમાનનું મુખ્ય
આધારવિધાન હોય
અને કાં તો
$n=1$ હોય તથા
$!A_1 = !A_n$
એ દાખલ
અનુમાનનું અથવા
\FalseInt{}નું
નિષ્કર્ષ હોય,
અથવા~$n>1$
હોય.
\end{defn}

નોંધો કે કટ-ખંડની
શરૂઆત દાખલ
નિયમના નિષ્કર્ષથી
થવી જરૂરી નથી.
તેનો અંત
લોપ
નિયમના મુખ્ય
આધારવિધાનમાં
થવો પૂરતો છે.
જોકે લંબાઈ~$1$નો
ખંડ કટ ગણાય
તે માટે તે
દાખલ
નિયમનું અથવા
\FalseInt{}નું
નિષ્કર્ષ હોવું
જોઈએ.
\sourcecorrection{OLMD-260}{મૂળનો સારાંશ લંબાઈ એકના કિસ્સામાં ફક્ત પરિચય-નિયમ કહે છે, જ્યારે તરત પહેલાંની કટની વ્યાખ્યામાં \FalseInt{} પણ સ્પષ્ટ સામેલ છે.}

\begin{defn}
કટની \emph{કક્ષા}~$\depth{!A}$
છે. સાબિતી~$\delta$ની
\emph{કટ-કક્ષા}~$\cutrank{\delta}$
એ~$\delta$માંના
કટોની મહત્તમ
કક્ષા છે.
કટની કક્ષા
$\cutrank{\delta}$
હોય, એટલે
મહત્તમ કક્ષા
હોય, તો તે
$\delta$માં
\emph{મહત્તમ}
કટ છે.
$\delta$ની
\emph{કટ-લંબાઈ}
એ તેના બધા
મહત્તમ કટોની
લંબાઈનો સરવાળો
છે.
\end{defn}
% END OLP-0677
\endgroup

\begingroup
\makeatletter\def\ol@sectioncs{section}\makeatother

% BEGIN OLP-0675
\olfileid{pt}{nor}{per}

\olsection{ક્રમપરિવર્તન-રૂપાંતરો}
\phantomsection\label{guunit:OLP-0675}


\Log{N1i}માં કટ
એવા ખંડો છે
જેમનો અંત
લોપ નિયમના
મુખ્ય આધારવિધાનમાં
થાય છે. સામાન્યીકરણની
સાબિતીમાં એક કિસ્સો
મહત્તમ કટની લંબાઈ
અને તેની સાથે
આખી સાબિતીની
કટ-લંબાઈ ઘટાડવાનો છે.
$>1$ લંબાઈનો કટ
ઘણી રીતે પૂરો
થઈ શકે છે: કટમાં
$!A$ની છેલ્લી
સૂત્ર આવૃત્તિ
(\Elim{\lor} કે
\Elim{\lexists})નું
નિષ્કર્ષ છે કે નહીં,
અને તે કયા
લોપ અનુમાનનું
મુખ્ય આધારવિધાન
છે, તેના પર
તે આધાર રાખે છે.

દાખલા તરીકે, સંબંધિત
અનુમાનો~\Elim{\lor}
અને~\Elim{\land}
હોય, $!A \ident !B \land !C$
હોય, અને
ઉપ-સાબિતી~$\delta_1$નો
અંત આ રીતે
થતો હોય:
\sourcecorrection{OLMD-253}{મૂળ ગદ્યમાં $!A$ને $!B \lor !C$ કહે છે, પરંતુ નીચેના વૃક્ષમાં કટનું છેલ્લું મુખ્ય સૂત્ર $!B \land !C$ છે; દાખલાના અનુરૂપ અહીં સંયોજન રાખ્યું છે.}
\[
\AxiomC{}
\RightLabel{$\delta_2$}
\DeduceC{$!D \lor !E$}
\AxiomC{$\Discharge{!D}{x}$}
\RightLabel{$\delta_3$}
\DeduceC{$!B \land !C$}
\AxiomC{$\Discharge{!E}{y}$}
\RightLabel{$\delta_4$}
\DeduceC{$!B \land !C$}
\DischargeRule{\Elim{\lor}}{x\,y}
\TrinaryInfC{$!B \land !C$}
\RightLabel{\Elim{\land}[1]}
\UnaryInfC{$!B$}
\DisplayProof
\]
કટમાં છેલ્લી
સૂત્ર આવૃત્તિ~$!A_n$
એ~\Elim{\lor}નું નિષ્કર્ષ,
નીચેનું $!B \land !C$,
છે. છેલ્લાથી
આગળની
સૂત્ર આવૃત્તિ~$!A_{n-1}$
એ~\Elim{\lor}ના
ગૌણ આધારવિધાનોમાંથી
એક છે.

આ કટની લંબાઈ
ઘટાડવા~\Elim{\lor}
અને~\Elim{\land}
અનુમાનોનો ક્રમ
બદલી (``ક્રમપરિવર્તન''
કરી) શકાય છે.
એટલે~$\delta$માંની
ઉપસાબિતી~$\delta_1$ના
સ્થાને આ મૂકીએ:
\[
\AxiomC{}
\RightLabel{$\delta_2$}
\DeduceC{$!D \lor !E$}
\AxiomC{$\Discharge{!D}{x}$}
\RightLabel{$\delta_3$}
\DeduceC{$!B \land !C$}
\RightLabel{\Elim{\land}[1]}
\UnaryInfC{$!B$}
\AxiomC{$\Discharge{!E}{y}$}
\RightLabel{$\delta_4$}
\DeduceC{$!B \land !C$}
\RightLabel{\Elim{\land}[1]}
\UnaryInfC{$!B$}
\DischargeRule{\Elim{\lor}}{x\,y}
\TrinaryInfC{$!B$}
\DisplayProof
\]
અગાઉ~\Elim{\lor}ની
નીચેના~\Elim{\land}ના
આધારવિધાનમાં પૂરા
થતા કટો હવે
\Elim{\lor}ના ગૌણ
આધારવિધાનોની ઉપરના
\Elim{\land} અનુમાનોમાંથી
એકના આધારવિધાનમાં
પૂરા થાય છે.
એટલે દરેકની
લંબાઈ~$1$ ઘટે છે.

વળી, સંપૂર્ણપણે
$\delta_2$, $\delta_3$
અથવા~$\delta_4$ની
અંદરનો કટ, કે
$\delta$માં~$\delta_1$ની
સંપૂર્ણપણે બહારનો
કટ, યથાવત્ રહે
છે: તે જ
સૂત્રો અને
તે જ અનુમાનોમાંથી
પસાર થાય છે.
તેથી ખાસ કરીને,
$\delta$માંના અનુરૂપ
કટ જેટલી જ તેની
કક્ષા અને લંબાઈ
રહે છે. આમ નવી
સાબિતીની કટ-કક્ષા
વધતી નથી અને
નવી સાબિતીની
કટ-લંબાઈ
જૂની સાબિતીની
કટ-લંબાઈ કરતાં
ઓછી છે.

\begin{prob}
ક્રમપરિવર્તન-રૂપાંતર પછી
ઓછામાં ઓછા એક
કટની લંબાઈ~$1$
ઘટે છે.
\Elim{\land}ને~\Elim{\lor}
ની આરપાર ખસેડીએ
ત્યારે હકીકતમાં
બે કટ-ખંડોની
લંબાઈ ઘટે છે;
એથી આ કિસ્સામાં
સાબિતીની કટ-લંબાઈ
ઓછામાં ઓછી~$2$
ઘટે છે. એક જ
આવા ક્રમપરિવર્તનથી
સાબિતીની કટ-લંબાઈ
$2$થી વધુ ઘટી
શકે? કટની
\emph{કક્ષા} ઘટી શકે?
\end{prob}

લોપ નિયમને
\Elim{\lor} કે
\Elim{\lexists} નિયમની
ઉપર ખસેડવાની
પ્રક્રિયાને
\emph{ક્રમપરિવર્તન-રૂપાંતર}
કહે છે. કેટલાક
નિયમ-યુગ્મો માટેનાં
નમૂના
\olref{tab:permutation-conversions}માં
આપ્યાં છે.
(બાકીનાં કસરત
તરીકે મૂક્યાં છે.)


{\small
\begin{longtable}{@{}l@{}}
\AxiomC{$!D \lor !E$}
\AxiomC{$\Discharge{!D}{x}$}
\shortDeduce
\DeduceC{$!B \land !C$}
\AxiomC{$\Discharge{!E}{y}$}
\shortDeduce
\DeduceC{$!B \land !C$}
\DischargeRule{\Elim{\lor}}{x\,y}
\TrinaryInfC{$!B \land !C$}
\RightLabel{\Elim{\land}[1]}
\UnaryInfC{$!B$}
\DisplayProof
$\longrightarrow$\\[6ex]
\quad\AxiomC{$!D \lor !E$}
\AxiomC{$\Discharge{!D}{x}$}
\shortDeduce
\DeduceC{$!B \land !C$}
\RightLabel{\Elim{\land}[1]}
\UnaryInfC{$!B$}
\AxiomC{$\Discharge{!E}{y}$}
\shortDeduce
\DeduceC{$!B \land !C$}
\RightLabel{\Elim{\land}[1]}
\UnaryInfC{$!B$}
\DischargeRule{\Elim{\lor}}{x\,y}
\TrinaryInfC{$!B$}
\DisplayProof
\\[10ex]
\AxiomC{$!D \lor !E$}
\AxiomC{$\Discharge{!D}{x}$}
\shortDeduce
\DeduceC{$!B \lif !C$}
\AxiomC{$\Discharge{!E}{y}$}
\shortDeduce
\DeduceC{$!B \lif !C$}
\DischargeRule{\Elim{\lor}}{x\,y}
\TrinaryInfC{$!B \lif !C$}
\AxiomC{$!B$}
\RightLabel{\Elim{\lif}}
\BinaryInfC{$!C$}
\DisplayProof
$\longrightarrow$\\[6ex]
\AxiomC{$!D \lor !E$}
\AxiomC{$\Discharge{!D}{x}$}
\shortDeduce
\DeduceC{$!B \lif !C$}
\AxiomC{$!B$}
\RightLabel{\Elim{\lif}}
\BinaryInfC{$!C$}
\AxiomC{$\Discharge{!E}{y}$}
\shortDeduce
\DeduceC{$!B \lif !C$}
\AxiomC{$!B$}
\RightLabel{\Elim{\lif}}
\BinaryInfC{$!C$}
\DischargeRule{\Elim{\lor}}{x\,y}
\TrinaryInfC{$!C$}
\DisplayProof
\\[10ex]
\AxiomC{$!D \lor !E$}
\AxiomC{$\Discharge{!D}{x}$}
\shortDeduce
\DeduceC{$!B \lif !C$}
\AxiomC{$\Discharge{!E}{y}$}
\shortDeduce
\DeduceC{$!B \lif !C$}
\DischargeRule{\Elim{\lor}}{x\,y}
\TrinaryInfC{$!B \lif !C$}
\AxiomC{$!B$}
\RightLabel{\Elim{\lif}}
\BinaryInfC{$!C$}
\DisplayProof
$\longrightarrow$\\[6ex]
\AxiomC{$!D \lor !E$}
\AxiomC{$\Discharge{!D}{x}$}
\shortDeduce
\DeduceC{$!B \lif !C$}
\AxiomC{$!B$}
\RightLabel{\Elim{\lif}}
\BinaryInfC{$!C$}
\AxiomC{$\Discharge{!E}{y}$}
\shortDeduce
\DeduceC{$!B \lif !C$}
\AxiomC{$!B$}
\RightLabel{\Elim{\lif}}
\BinaryInfC{$!C$}
\DischargeRule{\Elim{\lor}}{x\,y}
\TrinaryInfC{$!C$}
\DisplayProof
\\[10ex]
\AxiomC{$!D \lor !E$}
\AxiomC{$\Discharge{!D}{x}$}
\shortDeduce
\DeduceC{$\lforall[x][!B(x)]$}
\AxiomC{$\Discharge{!E}{y}$}
\shortDeduce
\DeduceC{$\lforall[x][!B(x)]$}
\DischargeRule{\Elim{\lor}}{x\,y}
\TrinaryInfC{$\lforall[x][!B(x)]$}
\RightLabel{\Elim{\lforall}}
\UnaryInfC{$!B(t)$}
\DisplayProof
\quad$\longrightarrow$\\[8ex]
\quad\AxiomC{$!D \lor !E$}
\AxiomC{$\Discharge{!D}{x}$}
\shortDeduce
\DeduceC{$\lforall[x][!B(x)]$}
\RightLabel{\Elim{\lforall}}
\UnaryInfC{$!B(t)$}
\AxiomC{$\Discharge{!E}{y}$}
\shortDeduce
\DeduceC{$\lforall[x][!B(x)]$}
\RightLabel{\Elim{\lforall}}
\UnaryInfC{$!B(t)$}
\DischargeRule{\Elim{\lor}}{x\,y}
\TrinaryInfC{$!B(t)$}
\DisplayProof
\\
\caption{\Log{N1i} માટેનાં કેટલાંક ક્રમપરિવર્તન-રૂપાંતરો.}
\ollabel{tab:permutation-conversions}
\end{longtable}
}
\sourcecorrection{OLMD-254}{મૂળ કોષ્ટકમાં અનુસરણ-લોપના બે મધ્યવર્તી રૂપાંતર-નમૂના એકસરખા પુનરાવર્તિત થાય છે. મૂળના વૃક્ષો યથાવત્ રાખ્યાં છે; કોઈ ગાયબ નમૂનો અનુમાનથી ઉમેર્યો નથી.}

\begin{prob}
\Elim{\lor} પછી
\Elim{\lor} અથવા
\Elim{\lnot} આવે
ત્યારેનાં
ક્રમપરિવર્તન-રૂપાંતરો
આપો.
\end{prob}

\begin{prob}
\Elim{\lexists} પછી
\Elim{\exists} આવે
ત્યારેનાં
ક્રમપરિવર્તન-રૂપાંતરો
આપો. છેલ્લા
કિસ્સામાં કોઈ
સ્વનિર્ધારક ચલની
શરત કેમ ભંગાતી
નથી તે સમજાવો.
\end{prob}

ક્રમપરિવર્તન-રૂપાંતર
લગાવતાં કોઈ
સ્વનિર્ધારક ચલની
શરત ભંગાતી નથી,
તેની ખાતરી કરીએ.
આ જુઓ:
\[
\AxiomC{}
\RightLabel{$\delta_2$}
\DeduceC{$!D \lor !E$}
\AxiomC{$\Discharge{!D}{x}$}
\RightLabel{$\delta_3$}
\DeduceC{$\lexists[x][!B(x)]$}
\AxiomC{$\Discharge{!E}{y}$}
\RightLabel{$\delta_4$}
\DeduceC{$\lexists[x][!B(x)]$}
\DischargeRule{\Elim{\lor}}{x\,y}
\TrinaryInfC{$\lexists[x][!B(x)]$}
\AxiomC{$\Discharge{!B(c)}{z}$}
\RightLabel{$\delta_5$}
\DeduceC{$!C$}
\DischargeRule{\Elim{\lexists}}{z}
\BinaryInfC{$!C$}
\DisplayProof
\]
તેના સ્થાને
આ મૂકીએ:
\[
\AxiomC{}
\RightLabel{$\delta_2$}
\DeduceC{$!D \lor !E$}
\AxiomC{$\Discharge{!D}{x}$}
\RightLabel{$\delta_3$}
\DeduceC{$\lexists[x][!B(x)]$}
\AxiomC{$\Discharge{!B(c)}{z}$}
\RightLabel{$\delta_5$}
\DeduceC{$!C$}
\DischargeRule{\Elim{\lexists}}{z}
\BinaryInfC{$!C$}
\AxiomC{$\Discharge{!E}{y}$}
\RightLabel{$\delta_4$}
\DeduceC{$\lexists[x][!B(x)]$}
\AxiomC{$\Discharge{!B(c)}{z}$}
\RightLabel{$\delta_5$}
\DeduceC{$!C$}
\DischargeRule{\Elim{\lexists}}{z}
\BinaryInfC{$!C$}
\DischargeRule{\Elim{\lor}}{x\,y}
\TrinaryInfC{$!C$}
\DisplayProof
\]
તો કદાચ
સ્વનિર્ધારક ચલ~$c$
દાખલા તરીકે~$!D$માં
આવતું હોય એવી
ચિંતા થાય. મૂળ
વૃક્ષમાં ધારણા~$!D$
મૂળના~$\Elim{\lexists}$ની
ઉપર મુક્ત થાય છે.
તેથી મૂળ સાબિતીમાં
તે~\Elim{\exists}
અનુમાનના મુખ્ય
આધારવિધાનની ઉપર
ખુલ્લી રહેલી
ધારણામાં આવતું
નથી; ત્યાં
સ્વનિર્ધારક ચલની
શરત સંતોષાય છે.
પણ નવી સાબિતીમાં
$!D$~\Elim{\lexists}
અનુમાનો પછી જ
મુક્ત થાય છે,
એટલે શરત ભંગાતી
હોય એમ લાગે.
જોકે યાદ રાખો
કે આપણે
સાબિતીઓ નિયમિત
છે એમ માનીએ છીએ:
દરેક સ્વનિર્ધારક
ચલ માત્ર એક
અનુમાનનું સ્વનિર્ધારક
ચલ છે અને આખી
સાબિતીમાં તે
માત્ર~\Elim{\lexists}ના
ગૌણ આધારવિધાનની
ઉપર જ આવે છે.
ખાસ કરીને,
$c$~$\delta_3$માં
કે~$\delta_4$માં
આવી શકતું નથી.

આ દાખલો બીજી
એક વાત પણ
બતાવે છે: નવી
સાબિતીમાં~$\delta_5$ની
બે નકલો છે.
જો~$\delta_5$માં
મહત્તમ કટ હોય,
તો આપણે કટ-લંબાઈ
\emph{ઘટાડી નથી}!
તેથી~$\delta_5$માં
મહત્તમ કટ નથી
તેની ખાતરી હોય
ત્યારે જ
ક્રમપરિવર્તન-રૂપાંતર
લગાવવું જોઈએ.
આ માટે હંમેશાં
\emph{સૌથી ઉપરનો}
મહત્તમ કટ પસંદ
કરીશું: એટલે
એવો મહત્તમ કટ
જેની ઉપ-સાબિતી
(અહીં~$\delta_1$)માં
તે ઉપ-સાબિતીના
છેલ્લા અનુમાનની
ઉપર પૂરો થતો
બીજો કોઈ મહત્તમ
કટ ન હોય.
એથી~$\delta_5$માં
(અને ખરેખર
$\delta_2$, $\delta_3$
કે~$\delta_4$માં પણ)
બીજા મહત્તમ
કટ હોવાની શક્યતા
દૂર થાય છે.
% END OLP-0675
\endgroup

\begingroup
\makeatletter\def\ol@sectioncs{section}\makeatother

% BEGIN OLP-0676
\olfileid{pt}{nor}{red}

\olsection{ઘટાડા-રૂપાંતરો}
\phantomsection\label{guunit:OLP-0676}


\Log{N1i}ની
સાબિતીમાં કટની
લંબાઈ~$1$ હોય,
તો તેનો ખંડ
માત્ર એક
સૂત્ર આવૃત્તિનો
બને છે. એ
સૂત્ર એક
દાખલ
અનુમાનનું નિષ્કર્ષ
અને લોપ
અનુમાનનું મુખ્ય
આધારવિધાન બંને
છે. સામાન્યીકરણની
સાબિતીમાં, આવા
દાખલ/લોપ
યુગ્મમાં પૂરી થતી
ઉપ-સાબિતીના સ્થાને,
તે યુગ્મ કાઢી
નાખેલી નવી
સાબિતી મૂકીને
કટ ``ઘટાડીએ'' છીએ.
આ બદલાવથી નવા
કટ ઊભા થઈ
શકે છે. જોકે
નવા કટની કક્ષા
ઘટાડાતા કટની
કક્ષા કરતાં ઓછી
હોય છે. તેથી
નવી સાબિતીની
કટ-લંબાઈ ઓછી
હોય છે (કારણ
કે લંબાઈ~$1$નો
મહત્તમ કટ દૂર
થયો), અથવા
સાબિતીની કટ-કક્ષા
ઘટે છે (જો એ
જ એકમાત્ર
મહત્તમ કટ હોય).

દાખલા તરીકે,
સંબંધિત અનુમાનો
\Intro{\lif} અને
\Elim{\lif} હોય,
$!A \ident !B \lif !C$
હોય અને
ઉપ-સાબિતી~$\delta_1$નો
અંત આ રીતે
થાય:
\[
\AxiomC{$\Discharge{!B}{x}$}
\RightLabel{$\delta_2$}
\DeduceC{$!C$}
\DischargeRule{\Intro{\lif}}{x}
\UnaryInfC{$!B \lif !C$}
\AxiomC{}
\RightLabel{$\delta_3$}
\DeduceC{$!B$}
\RightLabel{\Elim{\lif}}
\BinaryInfC{$!C$}
\DisplayProof
\]
આ કટ દૂર કરવા
$\delta$માંની
ઉપસાબિતી~$\delta_1$ના
સ્થાને આ મૂકીએ:
\[
\AxiomC{}
\RightLabel{$\delta_3$}
\DeduceC{$!B$}
\RightLabel{$\Subst{\delta_2}{\delta_3}{!B^x}$}
\DeduceC{$!C$}
\DisplayProof
\]
$\Subst{\delta_2}{\delta_3}{!B^x}$
એટલે~$\delta_2$માં
$x$ ચિહ્નવાળી
દરેક ધારણા~$!B$ને
આખી સાબિતી~$\delta_3$થી
બદલતાં મળતી
સાબિતી. આ
સાબિતીમાં
$x$ ચિહ્નવાળી
$!B$ સ્વરૂપની
ખુલ્લી ધારણા
રહેતી નથી.
તેમાં~$x$
ચિહ્નવાળી બીજી
કોઈ ધારણા પણ
નથી, કારણ કે
હંમેશની જેમ
અલગ સૂત્રો
એક જ ચિહ્ન
સાથે ધારણા
તરીકે આવતાં
નથી એમ માનીએ
છીએ. તેથી નવી
સાબિતીની ખુલ્લી
ધારણાઓ~$\delta_1$ની
જેટલી જ છે.
$\delta_2$માં ધારણાઓ
મુક્ત કરતું
દરેક અનુમાન
$\Subst{\delta_2}{\delta_3}{!B^x}$
માં પણ એ જ
ધારણાઓ મુક્ત
કરે છે; અને
$\delta_3$માં ધારણાઓ
મુક્ત કરતું
દરેક અનુમાન
શક્ય હોય એવી
અનેક નકલોમાં
આવે છે, જે
$\Subst{\delta_2}{\delta_3}{!B^x}$
માં અનુરૂપ
ધારણાઓ મુક્ત
કરે છે.

સંપૂર્ણપણે
$\delta_2$ અથવા
$\delta_3$ની અંદરનો
કટ, કે
$\delta$માં~$\delta_1$ની
સંપૂર્ણપણે બહારનો
કટ, યથાવત્
રહે છે: તે
જ સૂત્રો
અને એ જ
અનુમાનોમાંથી
પસાર થાય છે.
તેથી ખાસ કરીને,
$\delta$ના અનુરૂપ
કટ જેટલી જ
તેની કક્ષા અને
લંબાઈ છે.
આપણે લંબાઈ~$1$નો
મહત્તમ કટ
દૂર કર્યો છે;
એટલે કટ-લંબાઈ
અથવા કટ-કક્ષા
ઘટાડી છે (જો
$\delta$ની કટ-લંબાઈ~$1$
હતી અને આ
એકમાત્ર મહત્તમ
કટ હતો).
\sourcecorrection{OLMD-255}{મૂળ ફકરો $\delta_4$ની અંદરના કટનો ઉલ્લેખ કરે છે, પરંતુ આ અનુસરણ-પરિચય/લોપના ઉદાહરણમાં ફક્ત $\delta_2$ અને $\delta_3$ ઉપસાબિતીઓ છે.}

જોકે બે બાબતો
બની શકે છે:
\begin{enumerate}
  \item બધી
  ધારણાઓ~$!B^x$ને
  $\delta_3$થી બદલતાં
  $\delta_3$ના બધા
  કટની નકલો વધી
  જાય છે. એમાંથી
  કેટલાક મહત્તમ
  કટ હોય તો
  નવી સાબિતીની
  કટ-લંબાઈ જૂની
  સાબિતીની કટ-લંબાઈ
  કરતાં વધુ થઈ
  શકે. આવું
  ન થાય તેની
  ખાતરી જોઈએ.
  તેથી ઘટાડા-રૂપાંતર
  ફક્ત \emph{સૌથી
  જમણા} મહત્તમ
  કટ પર લગાવીએ:
  એટલે~$\delta$ની
  શાખાઓમાં ઘટાડાતા
  કટની જમણી બાજુ
  બીજો કોઈ
  મહત્તમ કટ ન
  હોય. $\delta_3$માં
  કટ હોય તો
  તે અહીં ઘટાડાતા
  કટની જમણી
  બાજુએ હશે;
  તેથી પસંદ કરેલો
  કટ સૌથી જમણો
  ન ગણાય. હંમેશાં
  સૌથી જમણો
  કટ પસંદ કરતાં
  સાબિતીની
  કટ-લંબાઈ અથવા
  કટ-કક્ષા ઘટે
  છે તેની ખાતરી
  મળે છે.
  \item જો
  $\delta_2$માંની
  ધારણા~$!B^x$
  લોપ
  અનુમાનનું મુખ્ય
  આધારવિધાન હોય,
  અને~$\delta_3$નું
  છેલ્લું અનુમાન
  \Elim{\lor},
  \Elim{\lexists}
  અથવા દાખલ
  હોય, તો આ
  $\delta_2$માંની આ
  ધારણાની જગ્યાએ
  $\delta_3$નું
  કલમ-જોડાણ
  કરવાથી નવો
  કટ ઊભો થાય
  છે. $\delta_2$માં
  જે માત્ર ધારણા
  હતી તે હવે
  લંબાઈ~$1$નો
  કટ (એક
  દાખલ
  અનુમાનનું નિષ્કર્ષ
  અને લોપ
  અનુમાનનું મુખ્ય
  આધારવિધાન) અથવા
  કટ-ખંડની
  છેલ્લી સૂત્ર
  આવૃત્તિ
  (લોપ
  અનુમાનનું
  મુખ્ય આધારવિધાન
  અને~\Elim{\lor}
  કે~\Elim{\lexists}નું
  નિષ્કર્ષ) બને
  છે. જો~$!B^x$
  એ~\Elim{\lor}
  અથવા~\Elim{\lexists}નું
  ગૌણ આધારવિધાન
  હતું, તો તે
  પહેલેથી કોઈ
  ખંડની, અને
  કદાચ કટ-ખંડની,
  પ્રથમ સૂત્ર
  આવૃત્તિ હતું.
  નવી સાબિતીમાં
  એ ખંડ હવે
  લાંબો થઈ શકે.
  પણ નોંધો કે
  આ નવા કટ
  અથવા જૂના
  કટ-ખંડનું
  સૂત્ર~$!B$
  છે અને
  $\depth{!B} < \depth{!B \lif !C}$.
  તેથી કટ ઉમેરાયા
  હોય અથવા તેમની
  લંબાઈ વધી હોય
  તો પણ તેઓ
  મહત્તમ કટ નથી.
  આમ કટ-કક્ષા
  ઘટે છે; અથવા
  સમાન કક્ષાના
  કટ બાકી હોય
  તો કટ-લંબાઈ
  ઘટે છે.
  \sourcecorrection{OLMD-256}{મૂળમાં આ જ ચિહ્નિત ધારણા પહેલા $!B^x$ અને પછી $B^x$ લખાઈ છે; આખા ઉદાહરણમાં સૂત્રનું ચિહ્ન $!B^x$ જ છે.}
\end{enumerate}

લંબાઈ~$1$નો
કટ ઘટાડવાની
પ્રક્રિયાને
\emph{ઘટાડા-રૂપાંતર}
(અથવા
\emph{ચકરાવા-રૂપાંતર})
કહે છે.
કેટલાક નિયમ-યુગ્મો
માટેના નમૂના
\olref{tab:reduction-conversions}માં
આપ્યાં છે.
(બાકીનાં કસરત
તરીકે મૂક્યાં છે.)

\begin{table}
\begin{multline*}
\bottomAlignProof
\AxiomC{$\Discharge{!B}{x}$}
\RightLabel{$\delta_2$}
\shortDeduce
\DeduceC{$!C$}
\DischargeRule{\Intro{\lif}}{x}
\UnaryInfC{$!B \lif !C$}
\AxiomC{}
\RightLabel{$\delta_3$}
\shortDeduce
\DeduceC{$!B$}
\RightLabel{\Elim{\lif}}
\BinaryInfC{$!C$}
\DisplayProof
\quad \longrightarrow\quad
\shoveright{\bottomAlignProof
\AxiomC{}
\RightLabel{$\delta_3$}
\shortDeduce
\DeduceC{$!B$}
\RightLabel{$\Subst{\delta_2}{\delta_3}{!B^x}$}
\shortDeduce
\DeduceC{$!C$}
\DisplayProof}
\\[4ex]
\shoveleft{\bottomAlignProof
\AxiomC{}
\RightLabel{$\delta_2$}
\shortDeduce
\DeduceC{$!B$}
\AxiomC{}
\RightLabel{$\delta_3$}
\shortDeduce
\DeduceC{$!C$}
\RightLabel{\Intro{\land}}
\BinaryInfC{$!B \land !C$}
\RightLabel{\Elim{\land}[1]}
\UnaryInfC{$!B$}
\DisplayProof
\quad \longrightarrow\quad
\bottomAlignProof
\AxiomC{}
\RightLabel{$\delta_2$}
\shortDeduce
\DeduceC{$!B$}
\DisplayProof}\qquad
\shoveright{\bottomAlignProof
\AxiomC{}
\RightLabel{$\delta_2$}
\shortDeduce
\DeduceC{$!B$}
\AxiomC{}
\RightLabel{$\delta_3$}
\shortDeduce
\DeduceC{$!C$}
\RightLabel{\Intro{\land}}
\BinaryInfC{$!B \land !C$}
\RightLabel{\Elim{\land}[2]}
\UnaryInfC{$!C$}
\DisplayProof
\quad \longrightarrow\quad
\bottomAlignProof
\AxiomC{}
\RightLabel{$\delta_3$}
\shortDeduce
\DeduceC{$!C$}
\DisplayProof}
\\[4ex]
\bottomAlignProof
\AxiomC{}
\RightLabel{$\delta_2$}
\shortDeduce
\DeduceC{$!B$}
\RightLabel{\Intro{\lor}[1]}
\UnaryInfC{$!B \lor !C$}
\AxiomC{$\Discharge{!B}{x}$}
\RightLabel{$\delta_3$}
\shortDeduce
\DeduceC{$!D$}
\AxiomC{$\Discharge{!C}{y}$}
\RightLabel{$\delta_4$}
\shortDeduce
\DeduceC{$!D$}
\DischargeRule{\Elim{\lor}}{x\,y}
\TrinaryInfC{$!D$}
\DisplayProof
\quad \longrightarrow\quad
\shoveright{\bottomAlignProof
\AxiomC{}
\RightLabel{$\delta_2$}
\shortDeduce
\DeduceC{$!B$}
\RightLabel{$\Subst{\delta_3}{\delta_2}{!B^x}$}
\shortDeduce
\DeduceC{$!D$}
\DisplayProof}
\\[2ex]
\shoveleft{\bottomAlignProof
\AxiomC{}
\RightLabel{$\delta_2$}
\shortDeduce
\DeduceC{$!B(c)$}
\RightLabel{\Intro{\lforall}}
\UnaryInfC{$\lforall[x][!B(x)]$}
\RightLabel{\Elim{\lforall}}
\UnaryInfC{$!B(t)$}
\DisplayProof
\quad \longrightarrow\quad
\bottomAlignProof
\AxiomC{}
\RightLabel{$\Subst{\delta_2}{t}{c}$}
\shortDeduce
\DeduceC{$!B(t)$}
\DisplayProof}
\\[4ex]
\shoveright{\bottomAlignProof
\AxiomC{}
\RightLabel{$\delta_2$}
\shortDeduce
\DeduceC{$!B(t)$}
\RightLabel{\Intro{\lexists}}
\UnaryInfC{$\lexists[x][!B(x)]$}
\AxiomC{$\Discharge{!B(c)}{x}$}
\RightLabel{$\delta_3$}
\shortDeduce
\DeduceC{$!D$}
\DischargeRule{\Elim{\lexists}}{x}
\BinaryInfC{$!D$}
\DisplayProof
\quad \longrightarrow\quad
\bottomAlignProof
\AxiomC{}
\RightLabel{$\delta_2$}
\shortDeduce
\DeduceC{$!B(t)$}
\RightLabel{$\Subst{\Subst{\delta_3}{t}{c}}{\delta_2}{!B(t)^x}$}
\shortDeduce
\DeduceC{$!D$}
\DisplayProof}
\end{multline*}
\caption{\Log{N1i} માટેનાં કેટલાંક ઘટાડા-રૂપાંતરો.}
\ollabel{tab:reduction-conversions}
\end{table}

\begin{prob}
\Intro{\lnot} પછી
\Elim{\lnot} આવે
ત્યારેનું
ઘટાડા-રૂપાંતર
આપો.
\end{prob}

હજી એક કિસ્સો
જોવાનો બાકી છે.
કટની વ્યાખ્યામાં
લંબાઈ~$1$ હોય
અને~$!A$ એ
\FalseInt{}નું
નિષ્કર્ષ તથા
લોપ
અનુમાનનું મુખ્ય
આધારવિધાન હોય,
તે કિસ્સો પણ
આવે છે. $!A$
એ દાખલ
નિયમનું નિષ્કર્ષ
હોય તે કિસ્સાની
જેમ આને પણ
સંભાળીએ. દાખલા
તરીકે, આને
\[
\bottomAlignProof
\AxiomC{}
\RightLabel{$\delta_2$}
\DeduceC{$\lfalse$}
\RightLabel{\FalseInt}
\UnaryInfC{$!B \lif !C$}
\AxiomC{}
\RightLabel{$\delta_3$}
\DeduceC{$!B$}
\RightLabel{\Elim{\lif}}
\BinaryInfC{$!C$}
\DisplayProof
\quad\text{ના સ્થાને}\quad
\bottomAlignProof
\AxiomC{}
\RightLabel{$\delta_2$}
\DeduceC{$\lfalse$}
\RightLabel{\FalseInt}
\UnaryInfC{$!C$}
\DisplayProof
\]
\sourcecorrection{OLMD-257}{મૂળના પ્રથમ અસત્ય-પ્રવેશ રૂપાંતરના બે વૃક્ષોમાં અસત્યમાંથી નવું સૂત્ર કાઢતી એકપદી કડીઓ પર નિયમનું લેબલ છૂટ્યું છે; નીચેના સમાન દાખલાની જેમ \FalseInt{} દેખાડ્યું છે.}
$!A \ident (!B \lor
!C)$ અથવા
$!A \ident \lexists[x][!B(x)]$
હોય ત્યારે થોડી
સાવચેતી જોઈએ.
દાખલા તરીકે,
જો આને
\[
\bottomAlignProof
\AxiomC{}
\RightLabel{$\delta_2$}
\DeduceC{$\lfalse$}
\RightLabel{\FalseInt}
\UnaryInfC{$!B \lor !C$}
\AxiomC{$\Discharge{!B}{x}$}
\RightLabel{$\delta_3$}
\DeduceC{$!D$}
\AxiomC{$\Discharge{!C}{y}$}
\RightLabel{$\delta_4$}
\DeduceC{$!D$}
\DischargeRule{\Elim{\lor}}{x\,y}
\TrinaryInfC{$!D$}
\DisplayProof
\quad\text{ના સ્થાને}\quad
\bottomAlignProof
\AxiomC{}
\RightLabel{$\delta_2$}
\DeduceC{$\lfalse$}
\RightLabel{\FalseInt}
\UnaryInfC{$!D$}
\DisplayProof
\]
\sourcecorrection{OLMD-258}{મૂળના વિકલ્પ-લોપ વૃક્ષમાં $!B$ અને $!C$ની જુદી ધારણાઓ બંનેને $x$ ચિહ્ન આપ્યું છે અને લોપ-નિયમ પર મુક્તિ-ચિહ્નો નથી. અગાઉના N1 નિયમવૃક્ષ પ્રમાણે બીજી ધારણા $y$ અને નિયમ $x,y$ બંનેને મુક્ત કરે છે.}
તો કટનું
સૂત્ર~$!D$
ધરાવતો નવો
કટ ઊભો થઈ
શકે છે. પરંતુ
$\depth{!D} < \depth{!B \lor !C}$
હોવાની કોઈ
ખાતરી નથી!{}
તેથી અહીં પણ
\Intro{\lor}/\Elim{\lor}
ઘટાડા-રૂપાંતરની
જેમ આગળ વધવું
અને તેના
સ્થાને આ મૂકવું:
\[
\bottomAlignProof
\AxiomC{}
\RightLabel{$\delta_2$}
\DeduceC{$\lfalse$}
\RightLabel{\FalseInt}
\UnaryInfC{$!B$}
\RightLabel{$\delta_3$}
\DeduceC{$!D$}
\DisplayProof
\quad\text{અથવા}\quad
\bottomAlignProof
\AxiomC{}
\RightLabel{$\delta_2$}
\DeduceC{$\lfalse$}
\RightLabel{\FalseInt}
\UnaryInfC{$!C$}
\RightLabel{$\delta_4$}
\DeduceC{$!D$}
\DisplayProof
\]
% END OLP-0676
\endgroup

\begingroup
\makeatletter\def\ol@sectioncs{section}\makeatother

% BEGIN OLP-0673
\olfileid{pt}{nor}{thm}

\olsection{સામાન્યીકરણ પ્રમેય}
\phantomsection\label{guunit:OLP-0673}


સાબિતીમાં કટ-ખંડ
ન હોય તો તે
\emph{સામાન્ય} સ્વરૂપમાં
છે. સ્પષ્ટ છે કે
આ ત્યારે અને ત્યારે
જ બને જ્યારે તેના
પર ક્રમપરિવર્તન-રૂપાંતર
કે ઘટાડા-રૂપાંતર,
બેમાંથી એકેય લગાડી
ન શકાય.
સાબિતીનું
\emph{સામાન્યીકરણ} કરવું
એટલે કટ ન રહે
ત્યાં સુધી આ બંને
પ્રકારનાં રૂપાંતર
લગાવી તેને સામાન્ય
સ્વરૂપની સાબિતીમાં
ફેરવવી. આવું હંમેશાં
કરી શકાય છે, એ
\emph{સામાન્યીકરણ પ્રમેય}
છે.

\begin{thm}\ollabel{thm:normalization}
દરેક \Log{N1i}-સાબિતી
$\delta$નું સામાન્યીકરણ
થાય છે; એટલે તેને
કટ વિનાની
સાબિતી~$\delta'$માં
ફેરવી શકાય છે.
\end{thm}

\begin{proof}
  $\delta$ની કટ-કક્ષા
  અને કટ-લંબાઈ પર
  આગમન કરીએ.
  આગમન-પરિકલ્પના એવી
  છે કે ઓછી કટ-કક્ષા
  ધરાવતી, અથવા એ
  જ કટ-કક્ષા સાથે
  ઓછી કટ-લંબાઈ ધરાવતી,
  દરેક સાબિતીનું
  સામાન્યીકરણ થાય છે.

  કટ-લંબાઈ~$0$
  હોય તો કટ જ
  નથી અને~$\delta$
  પહેલેથી સામાન્ય
  સ્વરૂપમાં છે.
  એટલે માનો કે
  $\delta$ની કટ-કક્ષા
  અને કટ-લંબાઈ
  બંને~$> 0$ છે.
  સૌથી ઉપરનો અને
  સૌથી જમણો મહત્તમ
  કટ-ખંડ પસંદ કરો.
  તેની લંબાઈ~$>1$
  હોય તો
  ક્રમપરિવર્તન-રૂપાંતર
  લગાવો. લંબાઈ~$1$
  હોય તો અનુરૂપ
  ઘટાડા-રૂપાંતર
  લગાવો. આમ મળતી
  નવી સાબિતીની
  કટ-કક્ષા સમાન
  અને કટ-લંબાઈ ઓછી
  હોય છે, અથવા
  કટ-કક્ષા જ ઓછી
  હોય છે; તેથી
  આગમન-પરિકલ્પના
  લાગુ પડે છે.
\end{proof}

ઉપરની સાબિતીની રીતમાં
હંમેશાં સૌથી ઉપરના,
સૌથી જમણા મહત્તમ
કટને પહેલાં હાથ
ધરીએ છીએ. તેથી
સાબિતી પર રૂપાંતરો
ચોક્કસ ક્રમે જ
લગાવવાં પડે એવું
લાગે, જેથી સામાન્ય
સાબિતી મળે.
આશ્ચર્યની વાત એ
છે કે ક્રમ ખરેખર
મહત્ત્વનો નથી.
ક્રમપરિવર્તન અને
ઘટાડા-રૂપાંતરો કોઈપણ
ક્રમે લગાવતાં છેવટે
સામાન્ય સ્વરૂપની
સાબિતી મળે છે.

ખંડ અને કટની
વ્યાખ્યાઓ~\Log{N2i}
માટે સહેલાઈથી આપી
શકાય છે; અને
\Log{N1i} માટેનાં
ક્રમપરિવર્તન અને
ઘટાડા-રૂપાંતરો સીધાં
\Log{N2i}માં રૂપાંતરિત
થાય છે. તેથી
સામાન્યીકરણ પ્રમેય
\Log{N2i} માટે પણ
લાગુ પડે છે.

\begin{thm}\ollabel{thm:normalization-N2i}
દરેક \Log{N2i}-સાબિતી
$\delta$નું સામાન્યીકરણ
થાય છે; એટલે તેને
કટ વિનાની
સાબિતી~$\delta'$માં
ફેરવી શકાય છે.
\end{thm}
% END OLP-0673
\endgroup

\begingroup
\makeatletter\def\ol@sectioncs{section}\makeatother

% BEGIN OLP-0678
\olfileid{pt}{nor}{tr}

\olsection{સામાન્ય \Log{N2i}-સાબિતીઓ અને \Log{G2i} વચ્ચેનું રૂપાંતર}
\phantomsection\label{guunit:OLP-0678}


સાબિતીઓને
\Log{N2i}માંથી
\Log{G2i}ની
સાબિતીઓમાં
રૂપાંતરિત કરી
શકાય છે અને
ઊલટું પણ.
પહેલાં, કટ-રહિત
\Log{G2i}ની
સાબિતીઓનું
રૂપાંતર સામાન્ય
\Log{N2i}-સાબિતીઓ
આપે છે.
\sourcecorrection{OLMD-261}{મૂળમાં કટ-રહિત સાબિતીના તંત્રનું નામ \Log{2i} છે, જે અપરિભાષિત છે; અહીં સંદર્ભ અને પૂર્વવર્તી રૂપાંતર મુજબ \Log{G2i} છે.}

પરિણામ કહેવા
\Log{N2i} અને
\Log{G2i}ના
સિક્વન્ટોના
પૂર્વાંગો વચ્ચેનો
સંબંધ સરળતાથી
દર્શાવતું થોડું
પરિભાષિક ગદ્ય
જોઈએ: ચિહ્નિત
સૂત્રોનો
ગણ~$\Gamma'$
ચિહ્નવિહોણાં
સૂત્રોના
બહુગણ~$\Gamma$ને
\emph{અનુરૂપ છે},
જો દરેક
સૂત્ર~$!A$
માટે~$\Gamma'$માં
$x : !A$
સ્વરૂપના
ઘટકોની
સંખ્યા~$\Gamma$માં
$!A$ની બહુલતા
કરતાં ઓછી કે
સમાન હોય.

\begin{cor}
જો
$\Log{G2i} \Proves \Gamma \Sequent \Delta$
હોય, તો
$\Gamma' \Sequent !C$ની
સામાન્ય
\Log{N2i}-સાબિતી~$\delta$
હોય છે. પહેલાના
સિક્વન્ટના જમણા
ભાગમાં સૂત્ર
હોય તો તે
જ પરિણામનું
સૂત્ર છે; જમણો
ભાગ ખાલી હોય
તો પરિણામ
અસત્ય છે.
અહીં
ચિહ્નિત
સૂત્રોનો
ગણ~$\Gamma'$
બહુગણ~$\Gamma$ને
અનુરૂપ છે:
એટલે દરેક
સૂત્ર~$!A$
માટે~$\Gamma'$માં
$x : !A$
સ્વરૂપના
ઘટકોની
સંખ્યા~$\Gamma$માં
$!A$ની બહુલતા
(એટલે~$!A$ની
નકલોની સંખ્યા)
કરતાં ઓછી કે
સમાન છે.
\sourcecorrection{OLMD-262}{મૂળ ઉપપત્તિ કટ-રહિત G2iમાંથી N2cની સામાન્ય સાબિતી કહે છે; અગાઉના રૂપાંતર મુજબ લક્ષ્ય \Log{N2i} છે. મૂળમાં જમણો ભાગ ખાલી હોય ત્યારે પરિણામ અસત્ય લેવાય અને એક સૂત્ર હોય ત્યારે તે જ લેવાય તે શરત પણ છૂટી છે.}
\end{cor}

\begin{proof}
\olref[nat][tgi]{prop:G2i-to-N2i}ની
સાબિતી તપાસો:
દાખલ નિયમો
સાબિતીઓના
અંતે ઉમેરાય છે,
અને લોપ
નિયમો ફક્ત
સાબિતીઓની
ટોચે ઉમેરાય
છે. તેથી
લોપ
નિયમની ઉપર
દાખલ
નિયમ કદી
આવતો નથી.
\end{proof}

જોકે
\olref[nat][tgi]{prop:G2i-to-N2i}માં
આપેલું~\Cut{}
નિયમનું રૂપાંતર
\Log{N2i}-સાબિતી~$\delta$માં
કટ ઊભા કરી
શકે છે.
\Log{G2i}ની
સાબિતીનો
અંત~\Cut{}માં
થતો હોય અને
તેની તત્કાલીન
ઉપ-સાબિતીઓ
$\pi_1$ અને
$\pi_2$ અનુક્રમે
$\Gamma_1
\Sequent !A$
અને
$!A, \Gamma_2 \Sequent \Delta_2$
ની હોય, તો
$\pi_1$નું
રૂપાંતર~$\delta_1$,
$\pi_2$ના
રૂપાંતર~$\delta_2$માં
આવતા
$x:!A
\Sequent !A$
સ્વયંસિદ્ધો પર
કલમ-જોડીએ છીએ.
$\delta_1$નો
અંત મુખ્ય
સૂત્ર~$!A$
ધરાવતા
દાખલ
નિયમમાં થાય
અને~$\delta_2$માંનો
$x:!A \Sequent !A$
સ્વયંસિદ્ધ
લોપ
નિયમનું મુખ્ય
આધારવિધાન હોય,
તો~$\delta$માં
કટ ઊભો થાય.

\Log{N2i}ની
સાબિતીઓને
\Log{G2i}ની
સાબિતીઓમાં
પણ રૂપાંતરિત
કરી શકાય છે.
\olref[nat][tni]{prop:N2-to-G2}ની
સાબિતીમાં
\Elim{\land},
\Elim{\lif},
\Elim{\lnot} અને
\Elim{\lforall}નું
રૂપાંતર પરિણામની
\Log{G2i}
સાબિતીમાં
\Cut{} અનુમાનો
આપે છે.
જોકે શરૂઆતની
\Log{N2i}-સાબિતી
સામાન્ય સ્વરૂપમાં
હોય, તો આ
ટાળી શકાય છે.
\sourcecorrection{OLMD-263}{મૂળનું છેલ્લું વાક્ય ``If the N2i-proof we start from, ...'' અધૂરું છે; આગળનું પ્રતિજ્ઞાપન સ્પષ્ટ કરે છે કે કટ ટાળવા શરૂઆતની સાબિતી સામાન્ય સ્વરૂપમાં હોવી જોઈએ.}

સિક્વન્ટોની
આવૃત્તિઓની
યાદી
$S_1$, \dots,~$S_n$
ને~\Log{N2i}-સાબિતી~$\delta$માં
\emph{શાખા} કહીએ,
જો~$S_1$
સ્વયંસિદ્ધ હોય
અને~$S_i$
કોઈ અનુમાનનું
મુખ્ય આધારવિધાન
હોય ત્યારે
$S_{i+1}$
તેનું નિષ્કર્ષ
હોય. બીજા
શબ્દોમાં, શાખા
સ્વયંસિદ્ધોથી
નીચે તરફ
$\delta$માં
સિક્વન્ટો અનુસરે
છે, પણ
લોપ
નિયમના ગૌણ
આધારવિધાન પર
અટકે છે.
$S_n$ અંતિમ
સિક્વન્ટ હોય
એવી~$\delta$ની
શાખાને
\emph{મુખ્ય શાખા}
કહીએ. દરેક
સાબિતી~$\delta$માં
ઓછામાં ઓછી
એક મુખ્ય શાખા
હોય છે, કારણ
કે દરેક નિયમને
ઓછામાં ઓછું
એક મુખ્ય
આધારવિધાન છે
(ફક્ત~\Intro{\land}ને
બે છે).


\begin{prop}\ollabel{prop:normal-N2i-to-G2i}
જો~$\delta$ એ
$\Gamma \Sequent !A$ની
સામાન્ય
$\Log{N2c}$ અથવા
$\Log{N2i}$-સાબિતી
હોય, તો
$\Gamma' \Sequent \Delta$ની
(કટ-રહિત)
$\Log{G2c}$-સાબિતી
(અથવા
$\Log{G2i}$-સાબિતી)~$\pi$
હોય છે. અહીં
$\Gamma$માંથી
ચિહ્નો દૂર કરતાં
મળતો સૂત્રોનો
બહુગણ~$\Gamma'$
છે; અને~$!A$
એ~$\lfalse$
ન હોય તો
$\Delta = \{!A\}$,
જ્યારે હોય તો
$\Delta = \emptyset$.
\end{prop}

\begin{proof}
આ વખતે
$\Gamma \Sequent !A$ની
સાબિતી~$\delta$માં
અનુમાનોની સંખ્યા
પર આગમન કરીએ.
$\delta$માં
એકેય અનુમાન
ન હોય, તો
તે સ્વયંસિદ્ધ
$x:!A \Sequent !A$
છે અને~$\pi$
અનુરૂપ સ્વયંસિદ્ધ
$!A \Sequent !A$
છે; અથવા
$!A
\ident \lfalse$
હોય તો
$\lfalse \Sequent \ $
છે.

જો~$\delta$નો
અંત અનુમાન-નિયમમાં
થાય, તો કિસ્સા
જુદા પાડીએ:
\begin{enumerate}
  \item $R$ એ
  દાખલ
  નિયમ અથવા~\FalseInt
  હોય, તો
  રૂપાંતર
  \olref[nat][tni]{prop:N2-to-G2}ની
  સાબિતીની જેમ
  જ થાય છે
  અને~\Cut
  જરૂરી નથી.

  $\delta$નો
  છેલ્લો નિયમ~$R$
  લોપ
  નિયમ હોય,
  તો નીચેની
  વાતો નોંધો:
  \begin{enumerate}
    \item $\delta$ની
    મુખ્ય શાખા
    કોઈ દાખલ
    નિયમમાંથી
    પસાર થતી
    નથી. નહિતર
    એ શાખા પર
    દાખલ
    નિયમ અથવા
    \FalseInt{}
    પછી
    લોપ
    નિયમ આવત
    અને~$\delta$
    સામાન્ય ન
    રહેત.
    \item $\delta$ની
    મુખ્ય શાખાઓ
    પર દાખલ
    નિયમો નથી,
    તેથી મુખ્ય
    શાખા એક
    જ છે.
    (માત્ર
    \Intro{\land}ને
    બે મુખ્ય
    આધારવિધાન છે;
    એટલે એકથી
    વધુ મુખ્ય
    શાખાઓ હોય
    તો તે
    \Intro{\land}ના
    આધારવિધાનોમાંથી
    પસાર થવી
    પડે.)
    \item મુખ્ય
    શાખામાં
    \Elim{\lor}
    અથવા
    \Elim{\lexists}નું
    મુખ્ય આધારવિધાન
    આવે, તો
    એવો ફક્ત
    એક જ
    નિયમ હોય
    અને તે
    છેલ્લો
    નિયમ~$R$
    હોય.
    (નહિતર
    મુખ્ય શાખામાં
    \Elim{\lor}
    કે
    \Elim{\lexists}
    એવો હોય
    કે તેનું
    નિષ્કર્ષ
    લોપ
    નિયમનું મુખ્ય
    આધારવિધાન
    બને; પછી
    ક્રમપરિવર્તન-રૂપાંતર
    લગાડી શકાય,
    એટલે સાબિતી
    સામાન્ય ન
    હોય.)
  \item \Elim{\lor}
  અને
  \Elim{\lexists}
  સિવાયના
  લોપ
  નિયમો ધારણાઓ
  મુક્ત કરતા
  નથી, એટલે
  સિક્વન્ટના
  ડાબા સંદર્ભને
  બદલતા નથી.
  આ બંને નિયમો
  ફક્ત ગૌણ
  આધારવિધાનની
  ઉપરની ધારણાઓ
  મુક્ત કરે
  છે; એટલે
  ફક્ત પોતાના
  ગૌણ આધારવિધાનોનો
  ડાબો સંદર્ભ
  બદલે છે.
  બીજા શબ્દોમાં,
  $\delta$ની
  મુખ્ય શાખા
  પરના
  સિક્વન્ટ~$S_i$નો
  સંદર્ભ~$\Gamma_1$
  હોય, તો
  $S_i$ જે
  અનુમાનનું
  મુખ્ય આધારવિધાન
  છે તેના
  નિષ્કર્ષ
  $S_{i+1}$નો
  સંદર્ભ
  $\Gamma_1' \supseteq \Gamma_1$
  હોય છે.
  \item તેથી
  મુખ્ય શાખાનો
  સૌથી ઉપરનો
  સિક્વન્ટ~$S_1$
  જો
  $x: !C \Sequent !C$
  હોય, તો
  $x:!C \in \Gamma$.
  \end{enumerate}
  હવે~$!C$ના
  સ્વરૂપ પ્રમાણે
  કિસ્સા જુદા
  પાડીએ. મુખ્ય
  શાખાનું અંતિમ
  સિક્વન્ટ સિવાયનું
  દરેક~$S_i$
  લોપ
  નિયમનું મુખ્ય
  આધારવિધાન છે;
  તેથી આ વાત
  $x:!C \Sequent !C$
  માટે પણ
  સાચી છે.
  \sourcecorrection{OLMD-264}{મૂળ કહે છે કે મુખ્ય શાખાનો દરેક $S_i$ લોપ-નિયમનું મુખ્ય આધારવિધાન છે; અંતિમ સિક્વન્ટ પછી કોઈ અનુમાન જ નથી. દાવો અંતિમ સિવાયના $S_i$ માટે છે.}

  \item $!C \ident !D \land !E$
  હોય. ત્યારે~$\delta$નું
  સ્વરૂપ આ છે:
  \[
  \Axiom$x:!D \land !E \fCenter !D \land !E$
  \RightLabel{\Elim{\land}[1]}
  \UnaryInf$x:!D \land !E \fCenter !D$
  \Deduce$x:!D \land !E, \Gamma_1 \fCenter !A$
  \DisplayProof
  \]
  મુખ્ય શાખામાં
  $x:!D \land !E \Sequent !D$
  આવે છે. તેમાં
  \Intro{\lif}
  કે~\Intro{\lnot}નું
  મુખ્ય આધારવિધાન
  નથી, તેમજ
  \Elim{\lor}
  કે~\Elim{\lexists}નું
  ગૌણ આધારવિધાન
  પણ નથી; તેથી
  $x:!D
  \land !E$માંનાં
  સંદર્ભ-સૂત્રો
  આખી મુખ્ય
  શાખામાં યથાવત્
  રહે છે. તેથી
  અંતિમ સિક્વન્ટના
  સંદર્ભમાં
  $x:!D \land !E$
  હોવું જ પડે.
  વળી,
  સાબિતી~$\delta$માં
  \Elim{\land}માં
  પૂરી થતી
  ઉપ-સાબિતીના
  સ્થાને
  $x:!D \Sequent !D$
  મૂકીને અને
  મુખ્ય શાખાના
  દરેક સિક્વન્ટના
  સંદર્ભમાં
  $x:!D \land !E$ના
  સ્થાને
  $x:!D$
  મૂકીને
  સાબિતી~$\delta_1$
  મેળવીએ; એટલે
  આને
  \[
  \bottomAlignProof
  \Axiom$x:!D \land !E \fCenter !D \land !E$
  \RightLabel{\Elim{\land}[1]}
  \UnaryInf$x:!D \land !E \fCenter !D$
  \Deduce$x:!D \land !E, \Gamma_1 \fCenter !A$
  \DisplayProof
  \quad\text{માં ફેરવો}\quad
  \bottomAlignProof
  \Axiom$x:!D \fCenter !D$
  \Deduce$x:!D, \Gamma_1 \fCenter !A$
  \DisplayProof
  \]
  આગમન-પરિકલ્પના
  $\delta_1$ને
  લાગુ પડે છે,
  કારણ કે~$\delta$માં
  અનુમાનોની સંખ્યા
  $\delta_1$માંની
  સંખ્યા કરતાં~$1$
  વધુ છે.
  \sourcecorrection{OLMD-265}{મૂળમાં આગમનથી બનેલી સાબિતીને $\delta_1$ કહીને પછી અનુમાનોની ગણતરીમાં અપરિભાષિત $\delta'$ વપરાયું છે; અહીં તે જ $\delta_1$ છે.}

  આગમન-પરિકલ્પનાથી
  $!D, \Gamma_1' \Sequent !A$ની
  \Log{G2i}-સાબિતી~$\pi_1$
  મળે છે.
  \LeftR{\land} લગાવી
  સાબિતી~$\pi$
  મેળવીએ:
  \sourcecorrection{OLMD-266}{મૂળ અહીં G2iની ઉપસાબિતીનું પૂર્વાંગ $x:!D, \Gamma_1$ લખે છે; G2iમાં $x:$ ચિહ્ન નથી અને સંદર્ભ $\Gamma_1'$ છે, જેમ નીચેના વૃક્ષમાં છે.}
  \[
  \AxiomC{}
  \RightLabel{$\pi_1$}
  \Deduce$!D, \Gamma_1' \fCenter !A$
  \RightLabel{\LeftR{\land}}
  \UnaryInf$!D \land !E, \Gamma_1' \fCenter !A$
  \DisplayProof
  \]
  $!A \ident \lfalse$
  હોય તો
  \RightR{\Weakening}
  ઉમેરીએ, જેમ કે:
  \[
  \AxiomC{}
  \RightLabel{$\pi_1$}
  \Deduce$!D, \Gamma_1' \fCenter $
  \RightLabel{\LeftR{\land}}
  \UnaryInf$!D \land !E, \Gamma_1' \fCenter $
  \RightLabel{\RightR{\Weakening}}
  \UnaryInf$!D \land !E, \Gamma_1' \fCenter !A$
  \DisplayProof
  \]
  \item જો
  $!C \ident !D \lor !E$
  હોય, તો તે
  \Elim{\lor}નું
  મુખ્ય આધારવિધાન
  હોવું જોઈએ
  અને આ અનુમાન
  મુખ્ય શાખાનું
  છેલ્લું અનુમાન~$R$
  હોવું જોઈએ.
  એટલે~$\delta$નું
  સ્વરૂપ આ છે:
  \[
  \Axiom$x:!D \lor !E \fCenter !D \lor !E$
  \AxiomC{}
  \RightLabel{$\delta_1$}
  \Deduce$y:!D, \Gamma_1 \fCenter!A$
  \AxiomC{}
  \RightLabel{$\delta_2$}
  \Deduce$z:!E, \Gamma_2 \fCenter !A$
  \RightLabel{\Elim{\lor}}
  \TrinaryInf$x:!D \lor !E, \Gamma_1, \Gamma_2 \fCenter !A$
  \DisplayProof
  \]
  આગમન-પરિકલ્પના
  સાબિતીઓ~$\delta_1$
  અને~$\delta_2$ને
  લાગુ પડે છે;
  તેથી
  $!D, \Gamma_1' \Sequent !A$
  અને
  $!E, \Gamma_2' \Sequent !A$ની
  \Log{G2i}-સાબિતીઓ
  $\pi_1$ અને
  $\pi_2$ મળે છે.
  $\LeftR{\lor}$
  લગાવી~$\pi$
  મેળવીએ:
  \[
  \AxiomC{}
  \RightLabel{$\pi_1$}
  \Deduce$!D, \Gamma_1' \fCenter !A$
  \AxiomC{}
  \RightLabel{$\pi_2$}
  \Deduce$!E, \Gamma_2' \fCenter !A$
  \RightLabel{\LeftR{\lor}}
  \BinaryInf$!D \lor !E, \Gamma_1', \Gamma_2' \fCenter !A$
  \DisplayProof
  \]
  $\Elim{\lor}$
  $y: !D$ અથવા
  $z: !E$ને
  મુક્ત ન કરતું
  હોય (એટલે
  તે ગૌણ
  આધારવિધાનોના
  સંદર્ભમાં
  ન હોય), અથવા
  $!A$ એ
  $\lfalse$
  હોય, તો
  કેટલાક
  \LeftR{\Weakening}
  અને
  \RightR{\Weakening}
  ઉમેરવા પડે.
  દાખલા તરીકે:
  \[
  \AxiomC{}
  \RightLabel{$\pi_1$}
  \Deduce$\Gamma_1' \fCenter $
  \RightLabel{\LeftR{\Weakening}}
  \UnaryInf$!D, \Gamma_1' \fCenter $
  \AxiomC{}
  \RightLabel{$\pi_2$}
  \Deduce$\Gamma_2' \fCenter $
  \RightLabel{\LeftR{\Weakening}}
  \UnaryInf$!E, \Gamma_2' \fCenter $
  \RightLabel{\LeftR{\lor}}
  \BinaryInf$!D \lor !E, \Gamma_1', \Gamma_2' \fCenter $
  \RightLabel{\RightR{\Weakening}}
  \UnaryInf$!D \lor !E, \Gamma_1', \Gamma_2' \fCenter !A$
  \DisplayProof
  \]
  
  \item $!C \ident !D \lif !E$
  હોય. ત્યારે~$\delta$નું
  સ્વરૂપ આ છે:
  \[
  \Axiom$x:!D \lif !E \fCenter !D \lif !E$
  \AxiomC{}
  \RightLabel{$\delta_1$}
  \Deduce$\Gamma_1 \fCenter!D$
  \RightLabel{\Elim{\lif}}
  \BinaryInf$x:!D \lif !E, \Gamma_1 \fCenter !E$
  \RightLabel{$\delta_2$}
  \Deduce$x:!D \lif !E, \Gamma_1, \Gamma_2 \fCenter !A$
  \DisplayProof
  \]
  અહીં નોંધો કે
  $x:!D \lif !E,
  \Gamma_1 \Sequent !E$
  ધરાવતી મુખ્ય
  શાખામાં
  \Intro{\lif}
  કે~\Intro{\lnot}નું
  મુખ્ય આધારવિધાન
  નથી, તેમજ
  \Elim{\lor}
  કે~\Elim{\lexists}નું
  ગૌણ આધારવિધાન
  પણ નથી.
  તેથી
  $x:!D \lif !E, \Gamma_1$ના
  સંદર્ભ-સૂત્રો
  મુખ્ય શાખામાં
  યથાવત્ રહે
  છે. આથી
  અંતિમ સિક્વન્ટના
  સંદર્ભમાં
  $x:!D \lif !E, \Gamma_1$
  આવવું જ જોઈએ.
  વળી,
  સાબિતી-ખંડ~$\delta_2$માં
  \Elim{\lif}માં
  પૂરી થતી
  ઉપ-સાબિતીના
  સ્થાને
  $x:!E \Sequent !E$
  મૂકીએ અને
  મુખ્ય શાખાના
  દરેક સિક્વન્ટના
  સંદર્ભમાં
  $x:!D \lif !E,
  \Gamma_1$ના
  સ્થાને~$x:!E$
  મૂકીએ, તો
  યોગ્ય
  સાબિતી~$\delta_2'$
  મળે છે;
  એટલે આને
  \[
  \bottomAlignProof
  \Axiom$x:!D \lif !E, \Gamma_1 \fCenter !E$
  \RightLabel{$\delta_2$}
  \Deduce$x:!D \lif !E, \Gamma_1, \Gamma_2 \fCenter !A$
  \DisplayProof
  \quad\text{માં ફેરવો}\quad
  \bottomAlignProof
  \Axiom$x:!E \fCenter !E$
  \RightLabel{$\delta_2'$}
  \Deduce$x:!E, \Gamma_2 \fCenter !A$
  \DisplayProof
  \]
  આગમન-પરિકલ્પના
  $\delta_1$ અને
  $\delta_2'$ને
  લાગુ પડે છે,
  કારણ કે~$\delta$માં
  અનુમાનોની સંખ્યા
  $\delta_1$ અને
  $\delta_2$માંની
  અનુમાનોની
  સંખ્યાના સરવાળાથી~$1$
  વધુ છે; આ
  એક અનુમાન
  મુખ્ય શાખાની
  શરૂઆતનું
  \Elim{\lif}
  છે. (જો
  આ જ
  $\delta$નું
  છેલ્લું
  અનુમાન~$R$
  હોય, તો
  $\delta_1$માં
  $0$ અનુમાનો
  છે.)

  આગમન-પરિકલ્પનાથી
  $\Gamma_1' \Sequent !D$ની
  \Log{G2i}-સાબિતી~$\pi_1$
  અને
  $!E, \Gamma_2' \Sequent !A$ની
  સાબિતી~$\pi_2$
  મળે છે.
  \LeftR{\lif} લગાવી
  સાબિતી~$\pi$
  મેળવીએ:
  \sourcecorrection{OLMD-267}{મૂળના આગમન-પરિણામમાં $\pi_1$ માટે $\Gamma_1$નું ચિહ્નવિહોણું સ્વરૂપ દર્શાવતું પ્રાઇમ નથી અને $\pi_2$ માટે $\Gamma_2'$ને બદલે $\Gamma_1'$ લખાયું છે; નીચેનું ડાબા અનુસરણનું વૃક્ષ યોગ્ય સંદર્ભો બતાવે છે.}
  \[
  \AxiomC{}
  \RightLabel{$\pi_1$}
  \Deduce$\Gamma_1' \fCenter !D$
  \AxiomC{}
  \RightLabel{$\pi_2$}
  \Deduce$!E, \Gamma_2' \fCenter !A$
  \RightLabel{\LeftR{\lif}}
  \BinaryInf$!D \lif !E, \Gamma_1', \Gamma_2' \fCenter !A$
  \DisplayProof
  \]
  $!D \ident \lfalse$
  અને/અથવા
  $!A \ident \lfalse$
  હોય, તો
  એક કે બે
  \RightR{\Weakening}
  ઉમેરીએ,
  જેમ કે:
  \sourcecorrection{OLMD-268}{મૂળમાં શરતનો પહેલો ભાગ માત્ર $!D$ લખાયો છે; ખાલી ઉત્તરાંગનો કિસ્સો $!D \ident \lfalse$ છે, જેમ તરત નીચેની જમણી શિથિલીકરણની કડી બતાવે છે.}
  \[
  \AxiomC{}
  \RightLabel{$\pi_1$}
  \Deduce$\Gamma_1' \fCenter $
  \RightLabel{\RightR{\Weakening}}
  \UnaryInf$\Gamma_1' \fCenter !D$
  \AxiomC{}
  \RightLabel{$\pi_2$}
  \Deduce$!E, \Gamma_2' \fCenter $
  \RightLabel{\RightR{\Weakening}}
  \UnaryInf$!E, \Gamma_2' \fCenter !A$
  \RightLabel{\LeftR{\lif}}
  \BinaryInf$!D \lif !E, \Gamma_1', \Gamma_2' \fCenter !A$
  \DisplayProof
  \]

\end{enumerate}
બાકીના કિસ્સા
આવી જ રીતે
સંભાળી શકાય
છે અને તેમને
કસરત તરીકે
મૂક્યા છે.
\end{proof}

\begin{cor}
સામાન્ય
\Log{N1i} અથવા
\Log{N2i}-સાબિતીમાં
આવતું દરેક
સૂત્ર
અંતિમ
સૂત્ર અથવા
અંતિમ સિક્વન્ટનું
ઉપ-સૂત્ર
છે.
\end{cor}

\begin{proof}
  \olref{prop:normal-N2i-to-G2i}
  પ્રમાણે દરેક
  સામાન્ય
  \Log{N2i}-સાબિતીને
  \Log{G2i}-સાબિતીમાં
  રૂપાંતરિત કરી
  શકાય છે.
  સાબિતી જોવાથી
  દેખાય છે કે
  \Log{N2i}-સાબિતીમાં
  આવતું દરેક
  સૂત્ર
  \Log{G2i}-સાબિતીમાં
  પણ આવે છે.
  \Log{G2i}માં
  \Cut{} નિયમ
  નથી, તેથી
  તેમાં
  ઉપ-સૂત્ર
  ગુણધર્મ છે.
\end{proof}
% END OLP-0678
\endgroup


\OLEndChapterHook
% END OLP-0674
\endgroup

\begingroup
\makeatletter\def\ol@sectioncs{section}\makeatother

% BEGIN OLP-0690
\olchapter{int}{pty}{પ્રકારો તરીકે પ્રતિજ્ઞાપનો}
\olchapterid{int}{pty}
\phantomsection\label{guunit:OLP-0690}


\begin{editorial}
  કરી--હાવર્ડ અનુરૂપતા અંગેના અધ્યાયનો આ
  \emph{અત્યંત પ્રાયોગિક} મુસદ્દો છે.
  વધુ ખુલાસા અને પ્રેરણા જોઈએ, અને તેમાં
  સંભવતઃ ભૂલો તથા ખૂટતી વિગતો છે.
  સામાન્યીકરણની સાબિતીની સમીક્ષા કરીને
  તેને વિસ્તૃત કરવી જોઈએ. ગુણાકાર-પ્રકારનાં
  ઉદાહરણો નથી. ક્રમચય અને સરળીકરણનાં
  રૂપાંતરો આવરી લીધાં નથી. (પ્રકારિત)
  લૅમ્બડા-કલન વિશેની સામગ્રી પણ મળે
  ત્યારે આ વધુ સમજાશે, કારણ કે અહીં
  મૂળભૂત રીતે તે જાણીતું માની લીધું છે.
  અત્યંત સાવચેતીથી વાપરો.
  \sourcecorrection{OLMD-290}{આ મૂળની સંપાદકીય નોંધ યથાવત્ અનુવાદી છે. જોકે સાબિતીથી પદોમાં રૂપાંતરના વિભાગમાં (\olref[pt]{sec}) ગુણાકાર-પ્રકારની અદલાબદલીનું એક ઉદાહરણ છે અને ઘટાડાના વિભાગમાં (\olref[red]{sec}) ક્રમચય-રૂપાંતરનું એક ઉદાહરણ છે. એટલે આ વિષયોનું સંપૂર્ણ આવરણ નથી, પરંતુ એકેય ઉદાહરણ નથી એવો મૂળનો દાવો હાલની અધ્યાય-સામગ્રી સાથે મેળ ખાતો નથી.}
\end{editorial}

\begingroup
\makeatletter\def\ol@sectioncs{section}\makeatother

% BEGIN OLP-0686
\olfileid{int}{pty}{int}

\olsection{પરિચય}
\phantomsection\label{guunit:OLP-0686}


લૅમ્બડા-કલન ગણતરીનું એક સરળ નિદર્શ છે. તે
ચલોથી બનેલાં પદો પર કામ કરે છે. $N$ અને $M$
પદો હોય, તો $(NM)$ પણ એક પદ છે
(``$N$ને~$M$ પર લાગુ કર્યું''). $N$ પદ
હોય, તો $\lambd[x][N]$ પણ પદ છે. $N$માં
ચલ~$x$ આવતો હોય, તો $\lambd[x][N]$માં
$x$ની સંબંધિત આવૃત્તિઓ $\lambd[x]$ ક્રિયાકારકથી
બદ્ધ થાય છે અને હવે મુક્ત રહેતી નથી.
$(\lambd[x][N]M)$ સ્વરૂપનું પદ
$\Subst{N}{M}{x}$માં $\beta$-સંકોચાય છે
(એટલે~$N$માં $x$ની દરેક મુક્ત આવૃત્તિને
$M$થી બદલી તેનું પરિણામ મળે છે).
$N$ના ઉપપદને $\beta$-સંકોચવાથી $N'$
મળે, તો $N$ એ~$N'$માં $\beta$-રૂપાંતરિત
થાય છે; તેને $N \redone N'$ લખીએ.
$\lambd$-કલનમાં ગણતરી એ
$N \redone N_1 \redone N_2 \redone \dots$
અનુક્રમ છે. કોઈ ઉપપદ $\beta$-સંકોચી ન
શકાય એવા પદ~$N_k$ પર આ અનુક્રમ સમાપ્ત
થાય, તો તે પદ \emph{સામાન્ય} કહેવાય,
અથવા~$N$નું સામાન્ય સ્વરૂપ કહેવાય. $\lambd$-કલનની
ગણતરીને આવા અનુક્રમ તરીકે વિચારો. સામાન્ય
સ્વરૂપ મળે તો ગણતરી થોભે છે અને એ સ્વરૂપ
ગણતરીનું પરિણામ છે. ગણતરીનું આ સરળ
નિદર્શ ટ્યુરિંગ-પૂર્ણ છે: દરેક આંશિક પુનરાવર્તી
વિધેય $\lambd$-કલનના કોઈ પદથી ગણી શકાય છે.

હાસ્કેલ કરી અને વિલિયમ હાવર્ડે સ્વતંત્ર રીતે
સ્વાભાવિક નિગમન અને $\lambd$-કલન વચ્ચેનું
નજીકનું સામ્ય શોધ્યું. સહજ રીતે,
$\lambd[x][N]$ પદ એક વિધેય છે: $x$ તેનો
આર્ગ્યુમેન્ટ છે અને~$x$ આપવામાં આવે ત્યારે
મૂલ્ય કેવી રીતે ગણવું તે~$N$ કહે છે.
$(\lambd[x][N]M)$ પદ એટલે
$\lambd[x][N]$ વિધેયને આર્ગ્યુમેન્ટ~$M$
પર લાગુ કરવું; તેનું મૂલ્યાંકન કેવી રીતે કરવું
તે $\beta$-સંકોચન કહે છે: $N$માં $x$ને~$M$થી
બદલો (અને મળેલા પદનું મૂલ્યાંકન ચાલુ રાખો).
હવે આ વિધેયોનો પ્રદિશ અને મૂલ્ય-પ્રદેશ
નિશ્ચિત છે એમ વિચારો. $\lambd[x][N]$નો
પ્રદિશ~$A$ અને મૂલ્ય-પ્રદેશ~$B$ હોય, તો
ચલ~$x$ ફક્ત~$A$માંથી આર્ગ્યુમેન્ટ લે છે
અને~$N$ ફક્ત~$B$માંથી મૂલ્યો આપે છે.
એટલે $(\lambd[x][N]M)$ને અર્થ ત્યારે જ
છે જ્યારે $\lambd[x][N]$ એ $A \to B$
વિધેય હોય અને $M \in A$ હોય. આ માહિતી
$\lambd$-કલનમાં ઉમેરીએ, તો
\emph{પ્રકારિત} $\lambd$-કલન મળે છે.
પ્રકારિત $\lambd$-કલનમાં દરેક
$(\lambd[x][N]M)$ અભિવ્યક્તિ માન્ય નથી;
ફક્ત $\lambd[x][N]$ અને $M$ના
``પ્રકારો'' પરસ્પર બંધબેસે ત્યારે, એટલે
$\lambd[x][N]$નો પ્રકાર $A \lif B$
અને $M$નો પ્રકાર~$A$ હોય ત્યારે, માન્ય છે.
$x$નો પ્રકાર~$A$ અને $N$નો પ્રકાર~$B$
હોય તો જ $\lambd[x][N]$નો પ્રકાર
$A \to B$ છે. સામાન્ય રીતે દરેક
$(M_1M_2)$ અભિવ્યક્તિ માન્ય નથી.
$M_1$નો પ્રકાર $A \to B$ અને
$M_2$નો પ્રકાર~$A$ હોય એવા
$(M_1M_2)$ પદો જ માન્ય છે;
એમ હોય તો $(M_1M_2)$નો પ્રકાર~$B$ છે.

આ પ્રકાર-નિયમો સ્વાભાવિક નિગમનની
નિષ્પત્તિઓને સ્પષ્ટ રીતે અનુરૂપ છે:
પ્રકારોને સૂત્રો દર્શાવે છે અને
નિષ્પત્તિઓ પોતે $\lambd$-પદોને
અનુરૂપ છે. દાખલા તરીકે,
$!A \lif !B$ની નિષ્પત્તિ
$\lambd[\typeof{x}{!A}][N]$ પદને
અનુરૂપ છે; આ પદનો વિધેય-પ્રકાર
$!A \lif !B$ છે. સ્વાભાવિક નિગમનના
અનુમાન-નિયમો પ્રકારિત લૅમ્બડા-કલનના
પ્રકાર-નિયમોને અનુરૂપ છે. દાખલા તરીકે,
\begin{prooftree}
  \AxiomC{$\Discharge{!A}{x}$}
  \DeduceC{$!B$}
  \DischargeRule{\Intro\lif}{x}
  \UnaryInfC{$!A \lif !B$}
  \DisplayProof\bottomAlignProof
  \quad\text{ને અનુરૂપ}\quad
  \Axiom$x: !A \fCenter N: !B$
  \UnaryInf$ \fCenter \lambd[\typeof{x}{!A}][N] : !A \lif !B$
\end{prooftree}
અહીં જમણી બાજુનો નિયમ કહે છે કે $x$નો
પ્રકાર~$!A$ અને $N$નો પ્રકાર~$!B$
હોય, તો $\lambd[\typeof{x}{!A}][N]$નો
પ્રકાર~$!A \lif !B$ છે.

$\Elim\lif$ નિયમ સંયુક્ત પદો માટેના
પ્રકાર-નિયમને અનુરૂપ છે, એટલે
\[
  \AxiomC{$!A \lif !B$}
  \AxiomC{$!A$}
  \RightLabel{\Elim\lif}
  \BinaryInfC{$!B$}
  \DisplayProof
  \quad\text{ને અનુરૂપ}\quad
  \Axiom$\fCenter M_1: !A \lif !B$
  \Axiom$\fCenter M_2: !A$
  \BinaryInf$ \fCenter (M_1M_2) : !B$
  \DisplayProof
\]
\Intro\lif{} નિયમ પછી તરત \Elim\lif{}
નિયમ આવે, તો નિષ્પત્તિને સરળ બનાવી શકાય:
\begin{gather*}
  \AxiomC{$\Discharge{!A}{x}$}
  \DeduceC{$!B$}
  \DischargeRule{\Intro\lif}{x}
  \UnaryInfC{$!A \lif !B$}
  \AxiomC{}
  \DeduceC{$!A$}
  \RightLabel{\Elim\lif}
  \BinaryInfC{$!B$}
  \DisplayProof
  \quad\redone\quad
  \AxiomC{}
  \DeduceC{$!A$}
  \DeduceC{$!B$}
  \DisplayProof
\end{gather*}
આ લૅમ્બડા પદોના $\beta$-સંકોચનને અનુરૂપ છે:
\[
(\lambd[\typeof{x}{!A}][N]M) \quad\redone\quad
\Subst{N}{M}{x}.
\]

આવું જ અનુરૂપતા $\land$ના નિયમો અને
``ગુણાકાર'' પ્રકારો વચ્ચે, તથા
$\lor$ના નિયમો અને ``સરવાળા''
પ્રકારો વચ્ચે છે.

સરળ પ્રકારિત લૅમ્બડા-કલનના પદો અને
સ્વાભાવિક નિગમનની નિષ્પત્તિઓ વચ્ચેની
આ અનુરૂપતાને ``કરી--હાવર્ડ'',
``પ્રોગ્રામ તરીકે સાબિતીઓ'' અથવા
``પ્રકારો તરીકે પ્રતિજ્ઞાપનો''ની
અનુરૂપતા કહે છે. સૂત્રો, એટલે
પ્રતિજ્ઞાપનો, પ્રકારોને અને સાબિતીઓ
પદોને અનુરૂપ છે; બાકીની અનુરૂપતાઓનો
સાર આ છે:
\begin{center}
  \begin{tabular}{c c c}
    તર્ક & પ્રોગ્રામ \\
    \hline
    પ્રતિજ્ઞાપન/સૂત્ર & પ્રકાર \\
    સાબિતી & પદ \\
    ધારણા & ચલ \\
    મુક્ત કરેલી ધારણા & બદ્ધ ચલ \\
    મુક્ત ન કરેલી ધારણા & મુક્ત ચલ \\
    $!A \lif !B$ & વિધેય-પ્રકાર \\
    $!A \land !B$ & ગુણાકાર-પ્રકાર \\
    $!A \lor !B$ & સરવાળો-પ્રકાર \\
    $\lfalse$ & ખાલી પ્રકાર \\
    \hline
  \end{tabular}
\end{center}

કરી--હાવર્ડ અનુરૂપતા સ્વચાલિત સાબિતી-સહાયકો
અને પ્રોગ્રામના પ્રકાર-ચકાસકોનો એક આધારસ્તંભ
છે. કારણ કે, પ્રતિજ્ઞાપનની સાબિતી તપાસવી
(જેમ ઉપર કર્યું) એ પ્રોગ્રામ, એટલે પદ,
જાહેર કરેલા પ્રકારનો છે કે નહીં તે તપાસવા
બરાબર છે.
% END OLP-0686
\endgroup

\begingroup
\makeatletter\def\ol@sectioncs{section}\makeatother

% BEGIN OLP-0688
\olfileid{int}{pty}{ter}

\olsection{સાબિતી-પદો}
\phantomsection\label{guunit:OLP-0688}


\Log{N1i}માં ધારણાઓને મુક્તિનાં ચિહ્ન તરીકે
ચલો આપીએ છીએ (દાખલા તરીકે $!A^x$).
\Log{N2i}માં દરેક સિક્વન્ટ $\Gamma \Sequent !A$ની
ડાબી બાજુના સંદર્ભ~$\Gamma$માં ચલ--સૂત્ર
જોડીઓની યાદી છે. તેથી \Log{N2i}માં
સાબિતીનો દરેક સિક્વન્ટ એટલું જ નથી
જણાવતો કે અત્યાર સુધીની સાબિતીએ શું
સાબિત કર્યું છે (એટલે $!A$), પણ દરેક
સ્થાને કઈ ધારણાઓ ખુલ્લી છે (એટલે $\Gamma$)
તે પણ જણાવે છે. દરેક સિક્વન્ટમાં~$!A$
\emph{કેવી રીતે} સાબિત થયું છે તે
માહિતી પણ ઉમેરીએ તો? આ માટે
\emph{પદથી ચિહ્નિત} તંત્ર \Log{tN2i}
વ્યાખ્યાયિત કરી શકીએ છીએ, જેમાં સિક્વન્ટનું
સ્વરૂપ $\Gamma \Sequent N : !A$ છે.
$\Gamma$ અને $!A$ પહેલાં જેવા જ છે:
$x: !B$ સ્વરૂપની જોડીઓનો ગણ અને
સિક્વન્ટ પર પૂરી થતી ઉપ-સાબિતીએ
સંદર્ભ~$\Gamma$ની ધારણાઓમાંથી અનુસરતું
બતાવેલું સૂત્ર. $N$ એ આ સિક્વન્ટ
પર પૂરી થતી સાબિતીને સંકેતરૂપે
રજૂ કરતું ચોક્કસ પદ હશે.

દરેક સિક્વન્ટને સોંપેલા પદો $x$, $y$,
\dots, ચલોથી બને છે. તેમાં અત્યાર સુધીનાં
અનુમાનો સંકેતરૂપે રજૂ કરતાં વિધેય-ચિહ્નો
વાપરીએ છીએ. દરેક અનુમાન-નિયમ માટે અલગ
વિધેય-ચિહ્ન જોઈએ. દાખલ
નિયમોને અનુરૂપ ચિહ્નો ``રચકો'' અને
લોપ નિયમોને અનુરૂપ ચિહ્નો
``વિઘટકો'' કહેવાય છે. નિયમનાં જેટલાં
આધારવિધાનો હોય, દરેક ચિહ્ન એટલાં
આર્ગ્યુમેન્ટ લે છે. દાખલા તરીકે,
\Intro{\land} અને \Elim{\land}[1]નાં
ચિહ્નો $c^\land$ (બે-સ્થાની) અને
$d^\land_1$ (એક-સ્થાની) હોઈ શકે.
\Log{tN2i}માં આ નિયમો આમ થાય છે:
\[
\Axiom$\Gamma_1 \fCenter N: !A$
\Axiom$\Gamma_2 \fCenter M: !B$
\RightLabel{\Intro{\land}}
\BinaryInf$\Gamma_1, \Gamma_2 \fCenter c^\land(N, M) : !A \land !B $
\DisplayProof
\quad
\Axiom$\Gamma \fCenter N: !A \land !B$
\RightLabel{\Elim{\land}[1]}
\UnaryInf$\Gamma \fCenter d^\land_1(N) : !A$
\DisplayProof
\]
\sourcecorrection{OLMD-287}{મૂળ સંયોજનના રચકમાં નાનું $m$ છે; બીજા આધારવિધાનનું પદ $M$ હોવાથી અહીં મોટું $M$ કર્યું છે.}
અનુમાન-નિયમ ધારણા મુક્ત કરે, એટલે તેના
આધારવિધાનમાં ચિહ્નિત સૂત્ર~$x:!A$
હોય પણ નિષ્કર્ષના સંદર્ભમાં ન હોય, તો
સાબિતી-પદમાં પણ આ બતાવવું જોઈએ.
આવા નિયમોનાં વિધેય-ચિહ્નો સંબંધિત
ચલોને \emph{બદ્ધ કરે} છે. દાખલા તરીકે,
\Intro{\lif} નિયમ આધારવિધાન
$x:!A, \Gamma \Sequent !B$માંથી
$\Gamma \Sequent !A \lif !B$ આપે છે.
રચક $c^{\lif}$ એક આર્ગ્યુમેન્ટ લે છે,
પણ ચલ~$x$ બદ્ધ થાય છે તે પણ
બતાવવું પડે. સાબિતી-પદ આ રીતે
ચિહ્ન~$x$વાળી ધારણા મુક્ત થઈ છે
તે વ્યક્ત કરે છે. દાખલા તરીકે,
આધારવિધાનનું સાબિતી-પદ~$N$ હોય,
તો નિષ્કર્ષનું પદ $c^{\lif}([x]N)$
અને નિયમ આમ લખી શકીએ:
\[
\Axiom$x:!A, \Gamma \fCenter N: !B$
\RightLabel{\Intro{\lif}}
\UnaryInf$\Gamma \fCenter c^{\lif}([x]N) : !A \lif !B$
\DisplayProof
\]

સાબિતીને સંપૂર્ણ રીતે ફરી રચવા માટે
વધારાની માહિતી જોઈએ. નિયમના આધારવિધાનનાં
સૂત્રોથી નિષ્કર્ષનું સૂત્ર
સંપૂર્ણ નક્કી ન થતું હોય, તો તે માહિતી
પદમાં ઉમેરવી પડે. \Intro{\lif}માં આવું
થાય છે, એટલે $c^{\lif}([x:!A]N)$
લખીશું. \Intro{\lor} અને~\FalseInt{}માં
પણ આવું થાય છે, એટલે એ માહિતી ધરાવતાં
વિધેય-ચિહ્નો વાપરીએ છીએ. દાખલા તરીકે,
\[
\Axiom$\Gamma \fCenter N:!A$
\RightLabel{\Intro{\lor}[1]}
\UnaryInf$\Gamma \fCenter c^{\lor{!B}}_1(N):!A \lor !B$
\DisplayProof
\quad
\Axiom$\Gamma \fCenter N:\lfalse$
\RightLabel{\FalseInt}
\UnaryInf$\Gamma \fCenter d^\lfalse_{!A}(N):!A$
\DisplayProof
\]
\sourcecorrection{OLMD-288}{મૂળ વિકલ્પ-પરિચયના નિષ્કર્ષમાં રચકનો આર્ગ્યુમેન્ટ છૂટી ગયો છે; આધારવિધાનનું પદ $N$ હોવાથી તેને ઉમેર્યું છે.}

આ વિચારો પરથી સિક્વન્ટ-શૈલીનું સ્વાભાવિક
નિગમન તંત્ર રચી શકાય છે, જેમાં દરેક
સિક્વન્ટનું સ્વરૂપ $\Gamma \Sequent N : !A$
છે અને~$N$ એક \emph{સાબિતી-પદ} છે.

આવાં સાબિતી-પદો અને \emph{પ્રકારિત લૅમ્બડા-કલન}નાં
પદો વચ્ચે નજીકનો સંબંધ છે. આ સંબંધને
અનુસરવા ઉપરનાં સામાન્ય રચક અને વિઘટકનાં
ચિહ્નોની જગ્યાએ લૅમ્બડા-કલનના સાહિત્યમાં
વપરાતાં ચિહ્નો લઈશું. દાખલા તરીકે,
$c^{\lif}([x:!A]N)$ની જગ્યાએ
$\lambd[\typeof{x}{!A}][N]$,
$d^{\lif}(N,M)$ની જગ્યાએ $(NM)$,
$c^\land(N,M)$ની જગ્યાએ $\pair{N}{M}$,
અને~$d^\land_i(N)$ની જગ્યાએ $\proj{i}{N}$
લખીએ.

\begin{defn}
  સાબિતી-પદો આ નિયમો વડે આગમનાત્મક રીતે બને છે:
  \begin{enumerate}
  \item દરેક ચલ~$x$ સાબિતી-પદ છે (સ્વયંસિદ્ધો).
  \item $x$ ચલ હોય, $!A$ સૂત્ર હોય અને
    $N$ સાબિતી-પદ હોય, તો $\lambd[\typeof{x}{!A}][N]$
    પણ સાબિતી-પદ છે~(\Intro\lif).
  \item $N$ અને $M$ સાબિતી-પદો હોય, તો
    $(NM)$ પણ સાબિતી-પદ છે~(\Elim\lif).
  \item $N$ અને $M$ સાબિતી-પદો હોય, તો
    $\pair{N}{M}$ સાબિતી-પદ છે~(\Intro\land).
  \item $N$ સાબિતી-પદ હોય, તો $\proj{i}{N}$
    પણ સાબિતી-પદ છે; અહીં $i$ એ $1$ અથવા~$2$
    છે~(\Elim\land[i]).
  \item $N$ સાબિતી-પદ હોય અને $!A$ સૂત્ર હોય,
    તો $\inj[!A]{i}{N}$ સાબિતી-પદ છે;
    અહીં $i$ એ $1$ અથવા~$2$ છે~(\Intro\lor[i]).
    \sourcecorrection{OLMD-289}{મૂળ આ કિસ્સામાં સાબિતી-પદ $N$થી શરૂ કરે છે, પણ નિષ્કર્ષમાં અસંબંધિત $!M$ આવે છે. અહીં નિષ્કર્ષનું પદ $N$ કર્યું છે.}
  \item $M$, $N_1$, $N_2$ સાબિતી-પદો હોય
    અને $x_1$, $x_2$ ચલો હોય, તો
    $\dcase{M}{x_1}{N_1}{x_2}{N_2}$
    સાબિતી-પદ છે~(\Elim\lor).
  \item $N$ સાબિતી-પદ હોય અને $!A$
    સૂત્ર હોય, તો $\abort{!A}{N}$
    સાબિતી-પદ છે~(\FalseInt).
  \end{enumerate}
\end{defn}

દરેક સાબિતી-પદ કોઈ સાબિતીનું
સાબિતી-પદ હોતું નથી. દાખલા તરીકે,
$\pair{N}{M}$ પદ \Intro{\land}
અનુમાનથી પૂરી થતી સાબિતીને સંકેતરૂપે
રજૂ કરે છે, એટલે તેનું અંતિમ સૂત્ર
સંયોજન~$!A \land !B$ છે. તે
\Elim{\lif} નિયમનું મુખ્ય આધારવિધાન
ન હોઈ શકે, એટલે $(\pair{N}{M}P)$
યોગ્ય સાબિતીનું પદ ન હોઈ શકે.
\Log{tN2ip}ના નિયમો અનુસરતાં
સાબિતી-પદો જ \emph{યોગ્ય} છે.
આ નિયમો \olref{tab:tN2ip}માં આપ્યા છે.

\begingroup
\makeatletter\def\ol@sectioncs{section}\makeatother

% BEGIN OLP-0692
\phantomsection\label{guunit:OLP-0692}


\begin{table}
\begin{tabular}{@{}c@{}c@{}}
\multicolumn{2}{l}{\textbf{સ્વયંસિદ્ધો:}}\\
\multicolumn{2}{c}{$x:!A \Sequent x:!A$}\\[3ex]
\multicolumn{2}{l}{\textbf{અનુમાન-નિયમો:}}\\[3ex]
\multicolumn{2}{c}{\Axiom$x: !A, \Gamma \fCenter N:!B$
\RightLabel{\Intro{\lif}}
\UnaryInf$\Gamma \fCenter \lambd[\typeof{x}{!A}][N] : !A \lif !B $
\DisplayProof}
\\[4ex]
\multicolumn{2}{c}{\Axiom$\Gamma_1 \fCenter N:!A \lif !B$
\Axiom$\Gamma_2 \fCenter M:!A$
\RightLabel{\Elim{\lif}}
\BinaryInf$\Gamma_1, \Gamma_2 \fCenter (NM) : !B$
\DisplayProof}
\\[4ex]
\multicolumn{2}{c}{\Axiom$\Gamma_1 \fCenter N_1 : !A$
\Axiom$\Gamma_2 \fCenter N_2 : !B$
\RightLabel{\Intro{\land}}
\BinaryInf$\Gamma_1, \Gamma_2 \fCenter \pair{N_1}{N_2} : !A \land !B $
\DisplayProof}
\\[4ex]
\Axiom$\Gamma \fCenter N: !A \land !B$
\RightLabel{\Elim{\land}[1]}
\UnaryInf$\Gamma \fCenter \proj{1}{N} : !A$
\DisplayProof
&
\Axiom$\Gamma \fCenter N:!A \land !B$
\RightLabel{\Elim{\land}[2]}
\UnaryInf$\Gamma \fCenter \proj{2}{N} : !B$
\DisplayProof\\[3ex]
\\[4ex]
\Axiom$\Gamma \fCenter N: !A$
\RightLabel{\Intro{\lor}[1]}
\UnaryInf$\Gamma \fCenter \inj[!B]{1}{N} : !A \lor !B$
\DisplayProof
&
\Axiom$ \Gamma \fCenter N : !B$
\RightLabel{\Intro{\lor}[2]}
\UnaryInf$\Gamma \fCenter \inj[!A]{2}{N} : !A \lor !B$
\DisplayProof\\[3ex]
\multicolumn{2}{c}{
\begin{tabular}[b]{@{}l@{}}
\Axiom$\Gamma_1 \fCenter M : !A \lor !B$
\Axiom$x:!A, \Gamma_2 \fCenter N_1 : !C$
\Axiom$y:!B, \Gamma_3 \fCenter N_2: !C$
\RightLabel{\Elim{\lor}}
\TrinaryInf$\Gamma_1, \Gamma_2, \Gamma_3 \fCenter
\dcase{M}{x}{N_1}{y}{N_2} : !C$
\DisplayProof
\end{tabular}
}
\\[4ex]
\multicolumn{2}{c}{\Axiom$\Gamma \fCenter N: \lfalse$
\RightLabel{\FalseInt}
\UnaryInf$\Gamma \fCenter \abort{!A}{N} : !A$
\DisplayProof}
\end{tabular}
\caption{પદથી ચિહ્નિત અંતઃપ્રજ્ઞાવાદી વિધાનતર્કીય
સ્વાભાવિક નિગમન તંત્ર \Log{tN2ip}ના નિયમો.}
\ollabel{tab:tN2ip}
\end{table}
% END OLP-0692
\endgroup
% END OLP-0688
\endgroup

\begingroup
\makeatletter\def\ol@sectioncs{section}\makeatother

% BEGIN OLP-0689
\olfileid{int}{pty}{pt}

\olsection{નિષ્પત્તિઓને સાબિતી-પદોમાં રૂપાંતરિત કરવી}
\phantomsection\label{guunit:OLP-0689}


સ્વાભાવિક નિગમનની સાબિતીઓને સાબિતી-પદોમાં
રૂપાંતરિત કરવી, એટલે~\Log{N2i}ની સાબિતીઓને
\Log{tN2i}ની સાબિતીઓમાં ફેરવવી, ખૂબ
સરળ છે. નિષ્પત્તિના દરેક સિક્વન્ટની
જમણી બાજુના સૂત્રની પહેલાં
સાબિતી-પદ લખવાનું જ છે; આમ
$\Gamma \Sequent M : !A$ સ્વરૂપના
સિક્વન્ટો મળે છે. પછી કહીએ કે
સંદર્ભ~$\Gamma$માં $M$ એ~$!A$નું
\emph{સાક્ષ્ય આપે છે}. કયું સાબિતી-પદ
લખવું તે \Log{tN2i}ના નિયમોથી
સંપૂર્ણ નક્કી થાય છે.

\begin{prop}
\Log{N2i}માં $\Gamma \Sequent !A$ની
દરેક સાબિતી~$\delta$ને અનુરૂપ
\Log{tN2i}માં $\Gamma \Sequent N : !A$ની
સાબિતી~$\delta'$ છે. સાબિતી-પદ~$N$
એ~$\delta$થી અનન્ય રીતે નક્કી થાય છે.
\end{prop}

\begin{proof}
$\delta$ની ઊંચાઈ પર આગમન કરીએ.
$\delta$ એ $x:!A \Sequent !A$ સ્વયંસિદ્ધ હોય,
તો $\delta'$ એ~$x:!A \Sequent x:!A$
સ્વયંસિદ્ધ છે. $\delta$ અનુમાન-નિયમથી
પૂરી થતી હોય, તો આગમન-ધારણા દરેક
આધારવિધાનનું સાબિતી-પદ અને તેને ધરાવતા
અનુરૂપ સિક્વન્ટની \Log{tN2i}-સાબિતીઓ
આપે છે. અનુરૂપ \Log{tN2i} નિયમ
લાગુ કરીએ; આ \Log{tN2i} અનુમાન-નિયમ સંબંધિત
સાબિતી-પદ નક્કી કરે છે.
\end{proof}

\begin{ex}
  \Log{N1i}માં $(!A \land !B) \lif (!B \land !A)$ની
  આ ખૂબ સરળ સાબિતી લો:
  \[
  \AxiomC{$\Discharge{!A\land !B}{x}$}
  \RightLabel{\Elim\land[2]}
  \UnaryInfC{$!B$}
  \AxiomC{$\Discharge{!A\land !B}{x}$}
  \RightLabel{\Elim\land[1]}
  \UnaryInfC{$!A$}
  \RightLabel{\Intro\land}
  \BinaryInfC{$!B \land !A$}
  \DischargeRule{\Intro\lif}{x}
  \UnaryInfC{$(!A \land !B) \lif (!B \land !A)$}
  \DisplayProof
  \]
  \Log{N2i}માં તે આવી છે:
  \[
  \Axiom$x : !A\land !B \fCenter !A\land !B$
  \RightLabel{\Elim\land[2]}
  \UnaryInf$x : !A\land !B \fCenter !B$
  \Axiom$x : !A\land !B \fCenter !A\land !B$
  \RightLabel{\Elim\land[1]}
  \UnaryInf$x : !A\land !B \fCenter !A$
  \RightLabel{\Intro\land}
  \BinaryInf$x : !A\land !B \fCenter !B \land !A$
  \RightLabel{\Intro\lif}
  \UnaryInf$\fCenter (!A \land !B) \lif (!B \land !A)$
  \DisplayProof
  \]
  ઉપરથી નીચે સાબિતી-પદો સોંપીએ.
  સ્વયંસિદ્ધોનાં પદો સંદર્ભના ચલ~$x$
  જ છે. પછી $\Elim\land[i]$ને અનુરૂપ
  પદો $\proj{2}{x}$ અને~$\proj{1}{x}$
  તરીકે નક્કી થાય છે. \Intro{\land}
  લાગુ કરીએ ત્યારે આ બંને પદો જોડીએ:
  $\pair{\proj{2}{x}}{\proj{1}{x}}$.
  છેલ્લે \Intro{\lif} નિયમ
  $\lambd$-અમૂર્તીકરણ લાગુ કરે છે:
  તે~$x$ને બદ્ધ કરે છે અને~$x$ એ
  ધારણા~$!A \land !B$નું ચિહ્ન હતું
  તે નોંધે છે:
  $\lambd[\typeof{x}{!A \land
  !B}][\pair{\proj{2}{x}}{\proj{1}{x}}]$.
  આમ \Log{tN2i}ની સાબિતી છે:
  \[
  \Axiom$x : !A\land !B \fCenter x: !A\land !B$
  \RightLabel{\Elim\land[2]}
  \UnaryInf$x : !A\land !B \fCenter \proj{2}{x} : !B$
  \Axiom$x : !A\land !B \fCenter x: !A\land !B$
  \RightLabel{\Elim\land[1]}
  \UnaryInf$x : !A\land !B \fCenter \proj{1}{x} : !A$
  \RightLabel{\Intro\land}
  \BinaryInf$x : !A\land !B \fCenter \pair{\proj{2}{x}}{\proj{1}{x}} : 
    !B \land !A$
  \RightLabel{\Intro\lif}
  \UnaryInf$\fCenter \lambd[\typeof{x}{!A \land
  !B}][\pair{\proj{2}{x}}{\proj{1}{x}}] : (!A \land !B) \lif (!B \land !A)$
  \DisplayProof
  \]
\end{ex}
% END OLP-0689
\endgroup

\begingroup
\makeatletter\def\ol@sectioncs{section}\makeatother

% BEGIN OLP-0695
\olfileid{int}{pty}{tp}

\olsection{સાબિતી-પદોમાંથી નિષ્પત્તિઓ પુનઃપ્રાપ્ત કરવી}
\phantomsection\label{guunit:OLP-0695}


હવે બીજી દિશા વિચારીએ: સાબિતી-પદ~$N$માંથી
\Log{N2i}ની સાબિતી પાછી મેળવવી.
સામાન્ય રીતે આ ફક્ત એવા સંદર્ભની સાપેક્ષે શક્ય છે,
જેમાં આપેલા સાબિતી-પદમાં મુક્ત રીતે આવતાં દરેક ચલ~$x$
માટે $x : !A$ની જોડી હોય.
પરંતુ મુક્ત ચલ વિનાના ઉદાહરણથી શરૂ કરીએ, એટલે આ પદથી:
\[
\lambd[\typeof{x}{!A \land
  !B}][\pair{\proj{2}{x}}{\proj{1}{x}}]
\]
તેનું સ્વરૂપ $\lambd[\typeof{x}{!A \land !B}][N_1]$ છે.
\Log{tN2i}માં આ સ્વરૂપનું પદ આપતો એકમાત્ર નિયમ
\Intro{\lif} છે; એટલે $N$માં સંકેતિત સાબિતીનો
અંત~\Intro{\lif}થી જ થાય છે, એટલે તે આમ પૂરો થાય છે:
  \[
  \Axiom$x : !A\land !B \fCenter \pair{\proj{2}{x}}{\proj{1}{x}} : 
    !C$
  \RightLabel{\Intro\lif}
  \UnaryInf$\fCenter \lambd[\typeof{x}{!A \land
  !B}][\pair{\proj{2}{x}}{\proj{1}{x}}] : (!A \land !B) \lif !C$
  \DisplayProof
  \]
હજુ $!C$ શું છે તે જાણતા નથી, પરંતુ પૂર્વાંગ
$!A \land !B$ છે અને આધારના સંદર્ભમાં
$x:!A \land !B$ આવવું જ જોઈએ તે જાણીએ છીએ.

હવે આધાર સાથે સંકળાયેલું સાબિતી-પદ નક્કી થઈ ચૂક્યું છે:
$\pair{\proj{2}{x}}{\proj{1}{x}}$.
\emph{તેમાં} સંકેતિત સાબિતીનો અંત \Intro{\land}થી જ થાય.
હજુ $!C$ શું છે તે જાણતા નથી, પરંતુ
તે~\Intro{\land}નું નિષ્કર્ષ હોવાથી સંયોજન જ હોય,
એટલે $!C \ident (!C_1 \land !C_2)$.
પુનઃરચના આગળ વધારીએ:
  \[
  \Axiom$x : !A\land !B \fCenter \proj{2}{x} : !C_1$
  \Axiom$x : !A\land !B \fCenter \proj{1}{x} : !C_2$
  \RightLabel{\Intro\land}
  \BinaryInf$x : !A\land !B \fCenter \pair{\proj{2}{x}}{\proj{1}{x}} : 
    !C_1 \land !C_2$
  \RightLabel{\Intro\lif}
  \UnaryInf$\fCenter \lambd[\typeof{x}{!A \land
  !B}][\pair{\proj{2}{x}}{\proj{1}{x}}] : (!A \land !B) \lif (!C_1 \land !C_2)$
  \DisplayProof
  \]
સાબિતી-પદ $\pi_2(x)$ સૂચવે છે કે ડાબી તરફનો સૌથી ઉપરનો
સિક્વન્ટ \Elim{\land}[2]થી મળ્યો હતો, એટલે આધારનું
સ્વરૂપ $!D \land !C_1$ છે.
જમણી તરફ $!C_2$ આપતું અનુમાન \Elim{\land}[1] જ હોય
અને આધારનું સ્વરૂપ $!C_2 \land !E$ હોય તે નક્કી કરી શકીએ.
  \[
  \Axiom$x : !A\land !B \fCenter x: !D\land !C_1$
  \RightLabel{\Elim\land[2]}
  \UnaryInf$x : !A\land !B \fCenter \proj{2}{x} : !C_1$
  \Axiom$x : !A\land !B \fCenter x: !C_2\land !E$
  \RightLabel{\Elim\land[1]}
  \UnaryInf$x : !A\land !B \fCenter \proj{1}{x} : !C_2$
  \RightLabel{\Intro\land}
  \BinaryInf$x : !A\land !B \fCenter \pair{\proj{2}{x}}{\proj{1}{x}} : 
    !C_1 \land !C_2$
  \RightLabel{\Intro\lif}
  \UnaryInf$\fCenter \lambd[\typeof{x}{!A \land
  !B}][\pair{\proj{2}{x}}{\proj{1}{x}}] : (!A \land !B) \lif (!C_1 \land !C_2)$
  \DisplayProof
  \]
સૌથી ઉપરના સિક્વન્ટોનાં સાબિતી-પદો હવે ચલ છે,
એટલે તેઓ સ્વયંસિદ્ધો જ હોવા જોઈએ.
આથી $!C_1$, $!C_2$, $!D$ અને $!E$ નક્કી થાય છે:
$!D \ident !A$ અને $!C_1 \ident !B$, નહીં તો
ડાબો સિક્વન્ટ સ્વયંસિદ્ધ નથી;
અને $!C_2 \ident !A$ તથા $!E \ident !B$,
નહીં તો જમણો સિક્વન્ટ સ્વયંસિદ્ધ નથી.
આપણે સાબિતી સંપૂર્ણપણે પુનઃપ્રાપ્ત કરી છે:
  \[
  \Axiom$x : !A\land !B \fCenter x: !A\land !B$
  \RightLabel{\Elim\land[2]}
  \UnaryInf$x : !A\land !B \fCenter \proj{2}{x} : !B$
  \Axiom$x : !A\land !B \fCenter x: !A\land !B$
  \RightLabel{\Elim\land[1]}
  \UnaryInf$x : !A\land !B \fCenter \proj{1}{x} : !A$
  \RightLabel{\Intro\land}
  \BinaryInf$x : !A\land !B \fCenter \pair{\proj{2}{x}}{\proj{1}{x}} : 
    !B \land !A$
  \RightLabel{\Intro\lif}
  \UnaryInf$\fCenter \lambd[\typeof{x}{!A \land
  !B}][\pair{\proj{2}{x}}{\proj{1}{x}}] : (!A \land !B) \lif (!B \land !A)$
  \DisplayProof
  \]
% END OLP-0695
\endgroup

\begingroup
\makeatletter\def\ol@sectioncs{section}\makeatother

% BEGIN OLP-0697
\olfileid{int}{pty}{typ}

\olsection{પ્રકારો}
\phantomsection\label{guunit:OLP-0697}


$(MN)$ જેવા સાબિતી-પદથી એકલેથી તે જે સાબિતીનું પદ છે
તે સાબિતી અનન્ય રીતે નક્કી થતી નથી.
કારણ કે મુક્ત ચલો, એટલે સ્વયંસિદ્ધોનાં સાબિતી-પદો,
અનુરૂપ સ્વયંસિદ્ધોમાંનાં સૂત્રો જણાવતા નથી.
પરંતુ તેમને સંદર્ભ~$\Gamma$માં જણાવીએ અને પદ યોગ્ય રીતે
પ્રકારિત હોય, તો સાબિતી અનન્ય રીતે નક્કી થાય છે.

\begin{defn}
સાબિતી-પદ~$N$ માટેનો \emph{સંદર્ભ} એવો સમૂહ $\Gamma$ છે
કે~$N$ના દરેક મુક્ત ચલ માટે
$\Gamma$માં બરાબર એક જ $x : !A$ની જોડી હોય.
\end{defn}

નોંધો કે વ્યાખ્યા મુજબ $\Gamma$માં
$x:!A$ની જોડી \emph{ફક્ત}~$N$ના મુક્ત ચલો માટે જ
હોવી જરૂરી નથી.
દાખલા તરીકે $\Gamma = \{x:!A, y:!B, z:!C\}$
એ $\pair{x}{y}$ માટેનો સંદર્ભ છે.

\begin{defn}\ollabel{defn:type}
  $\Gamma$ એ~$N$ માટેનો સંદર્ભ છે એમ ધારો.
  $N$નો \emph{પ્રકાર} ($\Gamma$ની સાપેક્ષે)
  આ રીતે આગમનથી વ્યાખ્યાયિત થાય છે:
  \begin{enumerate}
  \item $x$નો પ્રકાર $!A$ છે જો અને ફક્ત જો $x:!A \in \Gamma$.
  \item $N$નો પ્રકાર $x:!A, \Gamma$ની સાપેક્ષે~$!B$ હોય,
    તો $\lambd[\typeof{x}{!A}][N]$નો પ્રકાર $!A \lif !B$ છે.
  \item $N$નો પ્રકાર $!A \lif !B$ અને $M$નો પ્રકાર~$!A$ હોય,
    તો $(NM)$નો પ્રકાર~$!B$ છે.
  \item $N$નો પ્રકાર $!A$ અને $M$નો પ્રકાર~$!B$ હોય,
    તો $\pair{N}{M}$નો પ્રકાર $!A \land !B$ છે.
  \item $N$નો પ્રકાર~$!A_1 \land !A_2$ હોય,
    તો $\proj{i}{N}$નો પ્રકાર~$!A_i$ છે.
  \item $N$નો પ્રકાર $!A$ હોય, તો
    $\inj[!B]{i}{N}$નો પ્રકાર $i=1$ હોય ત્યારે $!A \lor !B$
    અને $i=2$ હોય ત્યારે $!B \lor !A$ છે.
    \sourcecorrection{OLMD-297}{મૂળમાં આ નિયમના આધારનું પદ $N$ છે પરંતુ વિકલ્પ-દાખલમાં $!M$ લખ્યું છે. દાખલનું પદ એ જ $N$ કર્યું છે, જેમ બંને અનુરૂપ અનુમાન-નિયમોમાં છે.}
  \item $M$નો પ્રકાર $!A \lor !B$ હોય,
    $N_1$નો પ્રકાર $x:!A, \Gamma$ની સાપેક્ષે $!C$
    અને $N_2$નો પ્રકાર $y:!B, \Gamma$ની સાપેક્ષે~$!C$ હોય,
    તો $\dcase{M}{x}{N_1}{y}{N_2}$નો પ્રકાર~$!C$ છે.
    \sourcecorrection{OLMD-298}{મૂળમાં શાખાઓના સંદર્ભમાં $x$ અને $y$ છે પરંતુ કેસ-પદમાં $x_1$ અને $x_2$ છે. કેસ-પદમાં શાખાઓનાં એ જ $x$ અને $y$ બદ્ધ ચલ રાખ્યાં છે.}
  \item $N$નો પ્રકાર~$\lfalse$ હોય,
    તો $\abort{!A}{N}$નો પ્રકાર~$!A$ છે.
  \end{enumerate}
\end{defn}

\begin{prop}
દરેક યોગ્ય રીતે પ્રકારિત સાબિતી-પદ~$N$ને
$N$ માટેના સંદર્ભ~$\Gamma$ની સાપેક્ષે બરાબર એક પ્રકાર હોય છે.
\sourcecorrection{OLMD-299}{મૂળ દરેક સાબિતી-પદ માટે પ્રકારનું અસ્તિત્વ તથા અનન્યતા કહે છે. પરંતુ આ જ વિભાગનું સ્વ-પ્રયોજનનું ઉદાહરણ બતાવે છે કે પદને પ્રકાર મળી ન શકે. તેથી અહીં યોગ્ય રીતે પ્રકારિત પદની શરત સ્પષ્ટ કરી છે; આગમન પ્રકારની અનન્યતા આપે છે, મનસ્વી પદ માટે પ્રકારનું અસ્તિત્વ નહીં. શરૂઆતના સાબિતીની અનન્યતા અંગેના વાક્યમાં પણ આ શરત સ્પષ્ટ કરી છે.}
\end{prop}

\begin{proof}
  $N$ પર આગમનથી.
\end{proof}

$N$નો પ્રકાર $\Gamma$ની સાપેક્ષે~$!A$ હોય,
તો $\Gamma \Sequent N : !A$ લખીએ.
\olref{defn:type}ની કલમો પરિચિત લાગવી જોઈએ:
તેઓ એક નાના તફાવત સાથે બરાબર \Log{tN2i}ના નિયમો જ છે;
સંદર્ભ~$\Gamma$ બધે એ જ છે.
આ વ્યાખ્યાને \olref{tab:tN2ip}માં આપેલા નિયમો ધરાવતા
\Log{tN3i} કલન તરીકે રજૂ કરી શકીએ.

\begingroup
\makeatletter\def\ol@sectioncs{section}\makeatother

% BEGIN OLP-0693
\phantomsection\label{guunit:OLP-0693}


\begin{table}
\begin{tabular}{@{}c@{}c@{}}
\multicolumn{2}{l}{\textbf{સ્વયંસિદ્ધો:}}\\
\multicolumn{2}{c}{$x:!A \in \Gamma$ હોય તો $\Gamma \Sequent x:!A$}\\[3ex]
\multicolumn{2}{l}{\textbf{અનુમાન-નિયમો:}}\\[3ex]
\multicolumn{2}{c}{\Axiom$x: !A, \Gamma \fCenter N:!B$
\RightLabel{\Intro{\lif}}
\UnaryInf$\Gamma \fCenter \lambd[\typeof{x}{!A}][N] : !A \lif !B $
\DisplayProof}
\\[4ex]
\multicolumn{2}{c}{\Axiom$\Gamma \fCenter N:!A \lif !B$
\Axiom$\Gamma \fCenter M:!A$
\RightLabel{\Elim{\lif}}
\BinaryInf$\Gamma \fCenter (NM) : !B$
\DisplayProof}
\\[4ex]
\multicolumn{2}{c}{\Axiom$\Gamma \fCenter N_1 : !A$
\Axiom$\Gamma \fCenter N_2 : !B$
\RightLabel{\Intro{\land}}
\BinaryInf$\Gamma \fCenter \pair{N_1}{N_2} : !A \land !B $
\DisplayProof}
\\[4ex]
\Axiom$\Gamma \fCenter N: !A \land !B$
\RightLabel{\Elim{\land}[1]}
\UnaryInf$\Gamma \fCenter \proj{1}{N} : !A$
\DisplayProof
&
\Axiom$\Gamma \fCenter N:!A \land !B$
\RightLabel{\Elim{\land}[2]}
\UnaryInf$\Gamma \fCenter \proj{2}{N} : !B$
\DisplayProof\\[3ex]
\\[4ex]
\Axiom$\Gamma \fCenter N: !A$
\RightLabel{\Intro{\lor}[1]}
\UnaryInf$\Gamma \fCenter \inj[!B]{1}{N} : !A \lor !B$
\DisplayProof
&
\Axiom$ \Gamma \fCenter N : !A$
\RightLabel{\Intro{\lor}[2]}
\UnaryInf$\Gamma \fCenter \inj[!B]{2}{N} : !B \lor !A$
\DisplayProof\\[3ex]
\multicolumn{2}{c}{
\begin{tabular}[b]{@{}l@{}}
\Axiom$\Gamma \fCenter M : !A \lor !B$
\Axiom$x:!A, \Gamma \fCenter N_1 : !C$
\Axiom$y:!B, \Gamma \fCenter N_2: !C$
\RightLabel{\Elim{\lor}}
\TrinaryInf$\Gamma \fCenter
\dcase{M}{x}{N_1}{y}{N_2} : !C$
\DisplayProof
\end{tabular}
}
\\[4ex]
\multicolumn{2}{c}{\Axiom$\Gamma \fCenter N: \lfalse$
\RightLabel{\FalseInt}
\UnaryInf$\Gamma \fCenter \abort{!A}{N} : !A$
\DisplayProof}
\end{tabular}
\caption{પ્રકાર-સોંપણી તંત્ર તરીકે કામ કરતા પદથી ચિહ્નિત
અંતઃપ્રજ્ઞાવાદી વિધાનતર્કીય સ્વાભાવિક નિગમન તંત્ર
\Log{tN3ip}ના નિયમો.}
\ollabel{tab:tN2ip}
\end{table}
% END OLP-0693
\endgroup


સાબિતી-પદો $\lambd$-કલનના એક સ્વરૂપનાં પદો છે,
જેને \emph{ગુણાકારો, સરવાળાઓ અને ખાલી પ્રકાર ધરાવતું
પ્રકારિત $\lambd$-કલન} કહે છે.
શાસ્ત્રીય $\lambd$-કલનમાં ફક્ત ચલોથી,
\emph{$\lambd$-અમૂર્ત પદો}~$\lambd[x][N]$થી
અને પ્રયોજન $(MN)$થી બનેલાં પદો હોય છે.
તેમાં પ્રકારિત $\lambd$-અમૂર્ત પદ
$\lambd[\typeof{x}{!A}][N]$ની~$!A$ની નોંધ હોતી નથી.
$\lambd$-કલન સંગણનાનું એવું નિદર્શ છે જે ટ્યુરિંગ-સંપૂર્ણ છે
(એટલે દરેક આંશિક સંગણનીય ફલનને તેમાં કોઈ પદથી રજૂ કરી શકાય).
સાબિતી-પદો અને પ્રકાર-નિયમોથી આપેલું $\lambd$-કલનનું સ્વરૂપ
સંગણનાનું મર્યાદિત નિદર્શ છે, પરંતુ મૂળભૂત વિચારો એ જ છે.
પ્રકાર-સોંપણી તેને મર્યાદિત કરે છે:
ફક્ત \emph{યોગ્ય રીતે પ્રકારિત} પદો,
એટલે સંદર્ભ~$\Gamma$ની સાપેક્ષે પ્રકાર ધરાવતાં સાબિતી-પદો,
માન્ય પદો છે.
દાખલા તરીકે શાસ્ત્રીય અપ્રકારિત $\lambd$-કલનમાં
$\lambd[x][(xx)]$ માન્ય પદ છે, પરંતુ
કોઈ પણ $!A$ માટે $\lambd[\typeof{x}{!A}][(xx)]$
યોગ્ય રીતે પ્રકારિત નથી.
કારણ કે $(xx)$ને પ્રકાર મળવા માટે પ્રકાર-નિયમો મુજબ
પહેલા~$x$નો પ્રકાર $!A \lif !B$
અને બીજા $x$નો પ્રકાર~$!A$ હોવો જોઈએ.
આ અશક્ય છે, કારણ કે $!A$ એ~$!A \lif !B$નું
યથાર્થ ઉપ-સૂત્ર છે.

સહજ રીતે પ્રકારને તે પ્રકારના પદથી નામિત થતી વસ્તુનો
પ્રકાર ગણી શકીએ.
એટલે પ્રકાર $!A \land !B$નું પદ એવી જોડી પસંદ કરે છે
જેનો પહેલો ઘટક પ્રકાર~$!A$નો અને બીજો ઘટક પ્રકાર~$!B$નો હોય,
એટલે ગુણાકાર~$!A \times !B$નું ઘટક.
પ્રકાર $!A \lif !B$નું પદ પ્રકાર $!A$ની વસ્તુઓમાંથી
પ્રકાર~$!B$ની વસ્તુઓમાં જતું ફલન છે.
પ્રકાર $!A \lor !B$નું પદ પ્રકાર $!A$ની કે
પ્રકાર~$!B$ની વસ્તુ છે, સાથે તે બેમાંથી કઈ છે તેની માહિતી હોય છે.
આ બે સમૂહો $!A$ અને $!B$ના વિયુક્ત સંઘ જેવું છે,
જેને $\{ \tuple{i,x}
: i = 1 \text{ અથવા } 2, x \in !A \text{ જો } i = 1, x \in !B \text{ જો
} i=2\}$ તરીકે વ્યાખ્યાયિત કરીએ.
પ્રકાર $\lfalse$ ખાલી સમૂહ છે, એટલે પદનો એ પ્રકાર હોય,
તો તે કોઈ વસ્તુને નામિત કરી શકતું નથી.
સ્વાભાવિક નિગમન (\Log{tN3i} સહિત) સુસંગત હોવાથી,
ખુલ્લી ધારણાઓ ન હોય ત્યાં સુધી~$\lfalse$ની સાબિતી નથી.
એટલે પ્રકાર~$\lfalse$નું કોઈ બંધ સાબિતી-પદ
(મુક્ત ચલ વિનાનું પદ) નથી.

પ્રકારિત $\lambd$-કલનના સંકારકો સહજ રીતે
યોગ્ય પ્રકારની વસ્તુને નામિત કરતાં પદો આપે છે,
જો તેમના પર લાગુ પાડવામાં આવતાં પદો અમુક વસ્તુઓને નામિત કરતાં હોય.
દાખલા તરીકે ``$N_1 : !A$ અને $N_2 : !B$ હોય,
તો $\pair{N_1}{N_2} : !A \land !B$'' એવું કહે છે
કે $N_1$ પ્રકાર~$!A$ની વસ્તુને અને
$N_2$ પ્રકાર~$!B$ની વસ્તુને નામિત કરે,
તો $\pair{N_1}{N_2}$ પ્રકાર $!A \land !B$ની
વસ્તુને નામિત કરે છે.
% END OLP-0697
\endgroup

\begingroup
\makeatletter\def\ol@sectioncs{section}\makeatother

% BEGIN OLP-0691
\olfileid{int}{pty}{red}

\olsection{ઘટાડો}
\phantomsection\label{guunit:OLP-0691}


સ્વાભાવિક નિગમનની નિષ્પત્તિઓમાં
દાખલ નિયમ પછી
લોપ નિયમ આવે, તો તે
આડમાર્ગ છે જેને ``ઘટાડી'' શકાય છે.
દાખલા તરીકે, $\Intro{\land}$ પછી
$\Elim{\land}$ આવે ત્યારે તે
``એકબીજાને રદ કરે છે'': \Intro{\land}થી
બનતા સંયોજનના બે સંયોજ્યો પૈકી
એક $\Elim{\land}$નો નિષ્કર્ષ હોય છે,
અને $\Intro{\land}$નાં આધારવિધાનો એ
બંને સંયોજ્યો છે. તેથી કોઈ પણ
સંયોજ્યની નિષ્પત્તિ જોઈએ, તો
સંયોજન દાખલ કરીને પછી તેને દૂર
કરવાને બદલે સીધી એ જ નિષ્પત્તિ
વાપરી શકીએ. આમ:
\begin{gather*}
  \bottomAlignProof
  \AxiomC{}
  \RightLabel{$\delta_1$}
  \DeduceC{$!A_1$}
  \AxiomC{}
  \RightLabel{$\delta_2$}
  \DeduceC{$!A_2$}
  \RightLabel{$\Intro{\land}$}
  \BinaryInfC{$!A_1 \land !A_2$}
  \RightLabel{$\Elim{\land}_i$}
  \UnaryInfC{$!A_i$}
  \DisplayProof\bottomAlignProof
  \quad
  \redone
  \quad
  \AxiomC{}
  \RightLabel{$\delta_i$}
  \DeduceC{$!A_i$}
  \DisplayProof
\end{gather*}

આ સરળીકરણ પરથી સંબંધિત સાબિતી-પદો પર
ઘટાડાનો સમાન સંબંધ લઈ શકીએ.
$\delta_1$નું સાબિતી-પદ~$N_1$ અને
$\delta_2$નું પદ~$N_2$ હોય, તો
ડાબી સાબિતીનું પદ એક પગલામાં
જમણી સાબિતીના પદમાં ઘટે છે:
\[
  \proj{i}{\pair{N_1}{N_2}} \redone N_i
\]
પ્રકારિત લૅમ્બડા-કલનમાં આ
ગુણાકાર-પ્રકારનો \emph{બીટા-ઘટાડાનો}
નિયમ છે.
\sourcecorrection{OLMD-291}{મૂળ જોડીના બંને પદ એક જ આર્ગ્યુમેન્ટમાં લખે છે. સાબિતી-પદની વ્યાખ્યા અને નિયમ-કોષ્ટક મુજબ જોડીના બંને આર્ગ્યુમેન્ટ અલગ કર્યા છે.}

સ્વાભાવિક નિગમનમાં \Intro{\land} પછી
\Elim{\land} જેવા બે અનુમાનો, એટલે
ઘટાડી શકાય એવાં અનુમાનોની જોડી,
\emph{કટ} કહેવાય છે. પ્રકારિત
લૅમ્બડા-કલનમાં $\redone$ની ડાબી
બાજુનું પદ \emph{સંકોચ્ય પદ} અને
જમણી બાજુનું પદ \emph{સંકોચનનું પરિણામ}
કહેવાય છે. અપ્રકારિત લૅમ્બડા-કલનમાં
ફક્ત $(\lambd[x][N])Q$ સંકોચ્ય પદ
ગણાય છે; પ્રકારિત લૅમ્બડા-કલનમાં
વાક્યરચનામાં $\pair{N}{M}$,
$\proj{i}{N}$, $\inj[!A]{i}{N}$,
$\dcase{N}{x_1}{M_1}{x_2}{M_2}$ અને
$\abort{!A}{N}$ ધરાવતાં પદો પણ
ઉમેરાય છે, જેમનાં અનુરૂપ સંકોચ્ય પદો છે.
\sourcecorrection{OLMD-292}{મૂળ વિકલ્પ-દાખલના ચિહ્ન માટે ત્રણ ફરજિયાત આર્ગ્યુમેન્ટ લખે છે; વ્યાખ્યા અને નિયમ-કોષ્ટક મુજબ પ્રકારની નોંધ વૈકલ્પિક આર્ગ્યુમેન્ટમાં મૂકી છે.}

એ જ રીતે વિકલ્પ માટે ઘટાડો છે:
\begin{gather*}
  \bottomAlignProof
  \AxiomC{}
  \RightLabel{$\delta$}
  \DeduceC{$!A_i$}
  \RightLabel{$\Intro{\lor}$}
  \UnaryInfC{$!A_1 \lor !A_2$}
  \AxiomC{$\Discharge{!A_1}{x_1}$}
  \RightLabel{$\delta_1$}
  \DeduceC{$!C$}
  \AxiomC{$\Discharge{!A_2}{x_2}$}
  \RightLabel{$\delta_2$}
  \DeduceC{$!C$}
  \DischargeRule{$\Elim{\lor}$}{x_1,x_2}
  \TrinaryInfC{$!C$}
  \DisplayProof\bottomAlignProof
  \quad
  \redone
  \quad
  \AxiomC{}
  \RightLabel{$\delta$}
  \DeduceC{$!A_i$}
  \RightLabel{$\delta_i$}
  \DeduceC{$!C$}
  \DisplayProof
\end{gather*}
તે સાબિતી-પદો પરના આ ઘટાડાને અનુરૂપ છે:
\begin{gather*}
  \dcase{\inj[!A_{3-i}]{i}{M}}{x_1}{N_1}{x_2}{N_2} \redone \Subst{N_i}{M}{x_i}
\end{gather*}
અહીં $M$ એ~$\delta$નું સાબિતી-પદ અને
$N_i$ એ~$\delta_i$નું પદ છે.
આ સરવાળા-પ્રકારનો બીટા-ઘટાડાનો નિયમ છે.
અહીં $\Subst{M}{N_i}{x}$ એટલે~$M$માં
ચલ~$x$ની દરેક મુક્ત આવૃત્તિને~$N_i$થી
બદલવી.
\sourcecorrection{OLMD-293}{મૂળ વિકલ્પ-દાખલની વૈકલ્પિક નોંધમાં દાખલ થતા પદનો પોતાનો પ્રકાર આપે છે. નિયમ-કોષ્ટક પ્રમાણે આ નોંધ બીજો, હજી નિર્દિષ્ટ ન થયેલો વિકલ્પ આપે છે; તેથી અહીં $!A_{3-i}$ કર્યું છે.}

વિધેય-પ્રકારનો ઘટાડો $\Intro\lif$
પછી $\Elim\lif$ના આડમાર્ગને
દૂર કરવાને અનુરૂપ છે.
\begin{gather*}
  \bottomAlignProof
  \AxiomC{$\Discharge{!A}{x}$}
  \RightLabel{$\delta$}
  \DeduceC{$!B$}
  \DischargeRule{$\Intro{\lif}$}{x}
  \UnaryInfC{$!A \lif !B$}
  \AxiomC{}
  \RightLabel{$\delta'$}
  \DeduceC{$!A$}
  \RightLabel{$\Elim{\lif}$}
  \BinaryInfC{$!B$}
  \DisplayProof\bottomAlignProof
  \quad
  \redone
  \quad
  \AxiomC{}
  \RightLabel{$\delta'$}
  \DeduceC{$!A$}
  \RightLabel{$\delta$}
  \DeduceC{$!B$}
  \DisplayProof
\end{gather*}
સાબિતી-પદો માટે આ સામાન્ય બીટા-ઘટાડો છે:
\begin{gather*}
  (\lambd[\typeof{x}{!A}][N]M)
  \redone \Subst{N}{M}{x}
\end{gather*}

સાબિતીઓનાં ઘટાડા-રૂપાંતરો ઉપરાંત
ક્રમચય-રૂપાંતરો પણ લઈ શકીએ.
તે એવી (ઉપ-)સાબિતીઓ પર લાગુ
થાય છે જેના અંતે \Elim{\lor}
પછી બીજો લોપ નિયમ
આવે છે. દાખલા તરીકે:
\begin{multline*}
  \bottomAlignProof
  \AxiomC{}
  \RightLabel{$\delta$}
  \DeduceC{$!A_i$}
  \RightLabel{$\Intro{\lor}$}
  \UnaryInfC{$!A_1 \lor !A_2$}
  \AxiomC{$\Discharge{!A_1}{x_1}$}
  \RightLabel{$\delta_1$}
  \DeduceC{$!C_1 \land !C_2$}
  \AxiomC{$\Discharge{!A_2}{x_2}$}
  \RightLabel{$\delta_2$}
  \DeduceC{$!C_1 \land !C_2$}
  \DischargeRule{$\Elim{\lor}$}{x_1,x_2}
  \TrinaryInfC{$!C_1 \land !C_2$}
  \RightLabel{\Elim{\land}[i]}
  \UnaryInfC{$!C_i$}
  \DisplayProof\bottomAlignProof
  \quad
  \redone
  \\
  \AxiomC{}
  \RightLabel{$\delta$}
  \DeduceC{$!A_i$}
  \RightLabel{$\Intro{\lor}$}
  \UnaryInfC{$!A_1 \lor !A_2$}
  \AxiomC{$\Discharge{!A_1}{x_1}$}
  \RightLabel{$\delta_1$}
  \DeduceC{$!C_1 \land !C_2$}
  \RightLabel{\Elim{\land}[i]}
  \UnaryInfC{$!C_i$}
  \AxiomC{$\Discharge{!A_2}{x_2}$}
  \RightLabel{$\delta_2$}
  \DeduceC{$!C_1 \land !C_2$}
  \RightLabel{\Elim{\land}[i]}
  \UnaryInfC{$!C_i$}
  \DischargeRule{$\Elim{\lor}$}{x_1,x_2}
  \TrinaryInfC{$!C_i$}
  \DisplayProof
\end{multline*}
$\delta$, $\delta_1$ અને $\delta_2$નાં
સાબિતી-પદો અનુક્રમે $M$, $N_1$ અને
$N_2$ હોય, તો સાબિતી-પદો, એટલે
$\lambd$-પદો, માટેનો અનુરૂપ
ઘટાડા-નિયમ છે:
\[
  \proj{i}{\dcase{M}{x_1}{N_1}{x_2}{N_2}} \redone 
\dcase{M}{x_1}{\proj{i}{N_1}}{x_2}{\proj{i}{N_2}}
\]

આડમાર્ગો ફક્ત સાબિતીઓના અંતે જ
નહીં, પણ સાબિતીઓની અંદર પણ
દૂર કરી શકીએ: આડમાર્ગથી પૂરી થતી
ઉપ-સાબિતી પર ઘટાડા-રૂપાંતર લગાવો.
એ જ રીતે $\beta$-ઘટાડો
$\lambd$-પદોના ઉપપદો પર લાગુ
થાય છે. $N$ એ~$M$નું ઉપપદ હોય
અને $N \redone N'$ થાય, તો
$M$માંનું ઉપપદ~$N$ એ~$N'$થી
બદલી $M'$ મેળવી શકીએ.
ત્યારે પણ $M \redone M'$ લખીએ.
$M = M_1 \redone M_2 \redone M_3 \redone \dots \redone M_k = M'$
અનુક્રમ હોય ($k=1$, એટલે $M = M'$
હોય તે કિસ્સો સહિત), તો
$M \red M'$ લખીએ.
\sourcecorrection{OLMD-294}{મૂળ અનુક્રમમાં બીજું પદ બે વખત લખાયું છે; ક્રમશઃ ત્રીજું પદ મૂક્યું છે.}

જેના કોઈ પણ ઉપ-સાબિતી પર રૂપાંતર
લાગુ ન થઈ શકે તેવી સાબિતી
\emph{સામાન્ય} કહેવાય છે.
એ જ રીતે જેના કોઈ પણ ઉપપદ પર
રૂપાંતર લાગુ ન થઈ શકે એવું
$\lambd$-પદ \emph{સામાન્ય} કહેવાય છે.
સામાન્ય $\lambd$-પદને
\emph{સામાન્ય સ્વરૂપ} પણ કહે છે.
% END OLP-0691
\endgroup

\begingroup
\makeatletter\def\ol@sectioncs{section}\makeatother

% BEGIN OLP-0696
\olfileid{int}{pty}{tpr}

\olsection{પ્રકારનું સંરક્ષણ}
\phantomsection\label{guunit:OLP-0696}


વિચારેલા ઘટાડા અને ક્રમચય-રૂપાંતરો
\Log{tN2i} તથા \Log{tN3i} તંત્રો માટે સીધા પણ વ્યાખ્યાયિત કરી શકાય.
\Log{tN3i} તંત્ર સંદર્ભ~$\Gamma$ની સાપેક્ષે
$\lambd$-પદ~$N$નો પ્રકાર એવું સૂત્ર~$!A$ ગણે છે,
જે માટે $\Gamma \Sequent N : !A$ની
સાબિતી \Log{tN3i}માં હોય.
$\lambd$-પદોના $\beta$-ઘટાડાના નિયમો
સાબિતીઓના ઘટાડા અને ક્રમચય-રૂપાંતરોને બરાબર અનુરૂપ છે.
આ હકીકતનું એક મહત્ત્વનું પરિણામ છે:
સંદર્ભ~$\Gamma$ની સાપેક્ષે $\lambd$-પદનો પ્રકાર
$\beta$-રૂપાંતર પછી એ જ રહે છે.
આ ગુણધર્મને \emph{પ્રકારનું સંરક્ષણ} કહે છે.

\begin{prop}
$N$નો પ્રકાર $\Gamma$ની સાપેક્ષે~$!A$ હોય
અને $N \redone N'$ હોય,
તો $N'$નો પ્રકાર~$\Gamma$ની સાપેક્ષે~$!A$ છે.
\end{prop}

\begin{proof}
  $N$ અને $\Gamma \Sequent N : !A$ની
  સાબિતીઓ~$\delta$ની ઊંચાઈ પર આગમનથી
  આ સરળતાથી સ્થાપી શકાય:
  $M$ એ~$N$નું ઉપપદ હોય, તો $\delta$માં
  $\Gamma' \Sequent M : !B$થી પૂરી થતી
  ઉપ-સાબિતી~$\delta_1$ છે.
  $M$ સંકોચ્ય પદ હોય, તો~$\delta_1$ પર રૂપાંતર લાગુ પડે છે.
  $\delta_1$ પર ઘટાડો લગાવીને
  $\Gamma' \Sequent M' : !B$ની સાબિતી~$\delta_1'$ મળે છે.
  \Log{tN3i}નાં રૂપાંતરો તપાસવાથી
  $M'$ એ~$M$ના સંકોચનનું પરિણામ છે તે દેખાય છે.
  $\delta$માં $\delta_1$ને બદલે $\delta_1'$ મૂકવાથી
  $\Gamma \Sequent N' : !A$ની સાબિતી~$\delta'$ મળે છે;
  અહીં $N$માં ઉપપદ~$M$ને બદલે~$M'$ મૂકતાં $N'$ મળે છે.

  દાખલા તરીકે~\Intro{\land} પછી
  \Elim{\land} આવે તે માટેનું ઘટાડા-રૂપાંતર વિચારીએ:
  \[
  \bottomAlignProof
  \AxiomC{}
  \Deduce$\Gamma \fCenter N_1 : !B_1$
  \AxiomC{}
  \Deduce$\Gamma \fCenter N_2 : !B_2$
  \RightLabel{\Intro{\land}}
  \BinaryInf$\Gamma \fCenter \pair{N_1}{N_2} : !B_1 \land !B_2$
  \RightLabel{\Elim{\land}[i]}
  \UnaryInf$\Gamma \fCenter \proj{i}{\pair{N_1}{N_2}} : !B_i$
  \DisplayProof
  \quad\redone\quad
  \bottomAlignProof
  \AxiomC{}
  \Deduce$\Gamma \fCenter N_i : !B_i$
  \DisplayProof
  \]
  જમણી તરફની સાબિતી, એટલે $\delta_1'$નું સાબિતી-પદ
  $N_i$ છે તે દેખાય છે.
  તે બરાબર ડાબી તરફની સાબિતી~($\delta_1$)ના
  સાબિતી-પદ $\proj{i}{\pair{N_1}{N_2}}$ પર
  $\beta$-રૂપાંતર લગાવવાનું પરિણામ છે.

  અથવા \Intro{\lif} પછી \Elim{\lif} અનુમાન આવે
  અને તેના પર ઘટાડા-રૂપાંતર લગાવીએ તે વિચારીએ:
  \begin{multline*}
  \bottomAlignProof
  \AxiomC{$\delta_2$}
  \Deduce$x: !B_1, \Gamma \fCenter N : !B_2$
  \RightLabel{\Intro{\lif}}
  \UnaryInf$\Gamma \fCenter \lambd[\typeof{x}{!B_1}][N] : !B_1 \lif !B_2$
  \AxiomC{}
  \RightLabel{$\delta_3$}
  \Deduce$\Gamma \fCenter M : !B_1$
  \RightLabel{\Elim{\lif}}
  \BinaryInf$\Gamma \fCenter (\lambd[\typeof{x}{!B_1}][N] M) : !B_2$
  \DisplayProof
  \quad\redone\\
  \bottomAlignProof
  \AxiomC{}
  \RightLabel{$\delta_3$}
  \Deduce$\Gamma \fCenter M : !B_1$
  \RightLabel{$\Subst{\delta_2}{\delta_3}{x:!B_1}$}
  \Deduce$\Gamma \fCenter \Subst{N}{M}{x} : !B_2$
  \DisplayProof
  \end{multline*}
  અહીં પણ જમણી તરફની સાબિતીનું સાબિતી-પદ
  બરાબર ડાબી તરફની સાબિતીના સાબિતી-પદ પર
  $\beta$-ઘટાડો લગાવવાનું પરિણામ છે.
  અલબત્ત,~$\delta_2$ની ઊંચાઈ પર આગમનથી
  સાવધાનીપૂર્વક આ તપાસવું પડે:
  $\delta_2$માં $\Gamma' \Sequent x:!B_1$ સ્વરૂપના બધા
  સ્વયંસિદ્ધોને બદલે $\Gamma \Sequent M:!B_1$ની
  સાબિતી~$\delta_3$ મૂકતાં ખરેખર
  $\Gamma \Sequent \Subst{N}{M}{x} : !B_2$ની
  સાબિતી $\Subst{\delta_2}{\delta_3}{x:!B_1}$ મળે છે.
\end{proof}

સાબિતીઓ અને $\lambd$-પદો વચ્ચેની નિકટ અનુરૂપતાનું
બીજું પણ પરિણામ છે.
$N$ સંદર્ભ~$\Gamma$ની સાપેક્ષે પ્રકાર $!A$નું
$\lambd$-પદ હોય અને તેમાં સંકોચ્ય પદ~$M$ હોય
(એટલે તે સામાન્ય સ્વરૂપ ન હોય), તો તેને અનુરૂપ
$\Gamma \Sequent N : !A$ની \Log{tN3i} સાબિતી~$\delta$માં
$\Gamma' \Sequent M: !B$થી પૂરી થતી એવી
ઉપ-સાબિતી હોય છે જેના પર રૂપાંતર લગાવી શકાય.
$M$ને બદલે તેના સંકોચનનું પરિણામ મૂકવાથી
$N \redone N'$ હોય, તો~$\delta$ પર
અનુરૂપ ઘટાડો લગાવવાનું પરિણામ
$\Gamma \Sequent N' : !A$ની સાબિતી છે.
એટલે સામાન્ય સ્વરૂપમાં ન હોય એવું દરેક $\lambd$-પદ
$\beta$-ઘટાડાથી બીજા પદમાં જાય છે.
આ ગુણધર્મને \emph{પ્રગતિ} કહે છે.

પ્રગતિ અને પ્રકારનું સંરક્ષણ મળીને ખાતરી આપે છે
કે $\lambd$-પદોનો $\beta$-ઘટાડો ક્યારેય ``અટકી જતો'' નથી:
પદ યોગ્ય રીતે પ્રકારિત હોય અને સામાન્ય સ્વરૂપમાં ન હોય,
તો તે બીજા યોગ્ય રીતે પ્રકારિત પદમાં ઘટે છે
(અને એ પદનો પ્રકાર એ જ હોય છે).
% END OLP-0696
\endgroup

\begingroup
\makeatletter\def\ol@sectioncs{section}\makeatother

% BEGIN OLP-0687
\olfileid{int}{pty}{nor}

\olsection{સામાન્યીકરણ}
\phantomsection\label{guunit:OLP-0687}


આ વિભાગમાં આપણે સાબિત કરીશું કે ઘટાડાના
કોઈ ક્રમથી દરેક નિગમનને સામાન્ય નિગમનમાં
ઘટાડી શકાય છે. તેને \emph{સામાન્યીકરણનો ગુણધર્મ}
કહે છે. આપણે પ્રતિજ્ઞાપનો અને પ્રકારો વચ્ચેની
અનુરૂપતાનો ઉપયોગ કરીશું: દરેક સાબિતી-પદને
સામાન્ય સ્વરૂપમાં ઘટાડી શકાય છે તે બતાવીશું;
પછી સ્વાભાવિક નિગમનની નિષ્પત્તિઓનું
સામાન્યીકરણ અનુસરે છે.

પહેલા પદોની જટિલતા માપતાં કેટલાંક વિધેયો
વ્યાખ્યાયિત કરીએ. સૂત્રની
\emph{લંબાઈ}~$\len{!A}$ આ પ્રમાણે છે:
\begin{align*}
  \len{p} &= 0\\
  \len{!A \land !B} &= \len{!A} + \len{!B} + 1\\
  \len{!A \lor !B} &= \len{!A} + \len{!B} + 1 \\
  \len{!A \lif !B} &= \len{!A} + \len{!B} + 1.
\end{align*}
સંકોચ્ય પદ~$M$ની જટિલતા તેના
\emph{કટ-ક્રમ}~$\cutrank{M}$થી માપીએ:
\begin{align*}
  \cutrank{(\lambd[\typeof{x}{!A}][\typeof{N}{!B}])Q} &=
  \len{!A} + \len{!B} + 1 \\
  \cutrank{\ande{i}{\andi{\typeof{M}{!A}}{\typeof{N}{!B}}}} &=
  \len{!A} + \len{!B} + 1 \\
  \cutrank{\ore{\ori{i}{!A_{i}}{\typeof{M}{!A_i}}}
    {\typeof{x_1}{!A_1}}{\typeof{N_1}{!C}}
    {\typeof{x_2}{!A_2}}{\typeof{N_2}{!C}}} &=
  \len{!A_1} + \len{!A_2} + 1
\end{align*}
\sourcecorrection{OLMD-280}{મૂળની સરવાળા-પ્રકારની છેલ્લી વ્યાખ્યામાં અસંબંધિત $!A$ અને $!B$ આવે છે; ઉપરના પદના પ્રકારો મુજબ $!A_1$ અને $!A_2$ કર્યા છે.}
સાબિતી-પદની જટિલતા તેમાંના સૌથી જટિલ
સંકોચ્ય ઉપપદથી માપીએ; પદ સામાન્ય હોય,
તો માપ~$0$ છે:
\[
  \maxrank{M} = \max \bigl(\{0\} \cup \{\cutrank{N} | N \text{ એ } M \text{
    નું સંકોચ્ય ઉપપદ છે}\}\bigr)
\]
\sourcecorrection{OLMD-286}{મૂળ ગદ્ય સામાન્ય પદનો મહત્તમ કટ-ક્રમ શૂન્ય કહે છે, પરંતુ દર્શાવેલું સૂત્ર ખાલી ગણનું મહત્તમ લે છે. ગદ્યના આ કિસ્સાને સૂત્રમાં પણ આવરી લેવા શૂન્ય ઉમેર્યું છે.}

\begin{lem}\ollabel{lem:subst}
  $\Subst{M}{\typeof{N}{!A}}{\typeof{x}{!A}}$ સંકોચ્ય પદ
  હોય અને $M \not\ident x$ હોય, તો આ પૈકીનો
  એક કિસ્સો બને:
  \begin{enumerate}
  \item $M$ પોતે સંકોચ્ય પદ છે, અથવા
  \item $M$નું સ્વરૂપ $\ande{i}{x}$ છે અને
    $N$નું સ્વરૂપ $\andi{P_1}{P_2}$ છે;
  \item $M$નું સ્વરૂપ $\ore{x}{x_1}{P_1}{x_2}{P_2}$ છે અને
    $N$નું સ્વરૂપ $\ori{i}{}{Q}$ છે;
    \sourcecorrection{OLMD-281}{મૂળમાં વિકલ્પ-લોપનો પ્રથમ આર્ગ્યુમેન્ટ $i$ છે; અહીં જે ચલનું પ્રતિસ્થાપન થાય છે તે $x$ હોવો જોઈએ.}
  \item $M$નું સ્વરૂપ $x Q$ છે અને
    $N$નું સ્વરૂપ $\lambd[x][P]$ છે.
  \end{enumerate}
  પહેલા કિસ્સામાં $\cutrank{\Subst{M}{N}{x}} = \cutrank{M}$;
  બીજા કિસ્સાઓમાં $\cutrank{\Subst{M}{N}{x}} = \len{!A}$ છે.
  \sourcecorrection{OLMD-282}{મૂળની છેલ્લી સમતામાં વધારાનો બંધ કૌંસ કાઢ્યો છે.}
\end{lem}
\begin{proof}
  $M$ પર આગમનથી સાબિતી કરીએ.
  \begin{enumerate}
  \item $M$ એક જ ચલ $y$ હોય અને $y \not\ident x$ હોય,
    તો $\Subst{y}{N}{x}$ એ $y$ જ છે, એટલે સંકોચ્ય નથી.
  \item $M$નું સ્વરૂપ $\andi{N_1}{N_2}$, $\lambd[x][N]$,
    અથવા $\ori{i}{!A}{N}$ હોય, તો
    $\Subst{M}{\typeof{N}{!A}}{\typeof{x}{!A}}$નું પણ
    એ જ સ્વરૂપ છે, એટલે તે સંકોચ્ય નથી.
  \item $M$નું સ્વરૂપ $\ande{i}{P}$ હોય, તો બે કિસ્સા લઈએ:
    \begin{enumerate}
      \item $P$નું સ્વરૂપ $\andi{P_1}{P_2}$ હોય, તો
        $M \ident \ande{i}{\andi{P_1}{P_2}}$ સંકોચ્ય પદ છે;
        સ્પષ્ટ રીતે
        \[
        \Subst{M}{N}{x} \ident
        \ande{i}{\andi{\Subst{P_1}{N}{x}}{\Subst{P_2}{N}{x}}}
        \]
        પણ સંકોચ્ય પદ છે. બંનેના કટ-ક્રમ સમાન છે.
      \item $P$ એક જ ચલ હોય, તો પ્રતિસ્થાપનથી સંકોચ્ય
        પદ બનાવવા તે $x$ જ હોવો જોઈએ, અને $N$નું
        સ્વરૂપ $\andi{P_1}{P_2}$ હોવું જોઈએ. હવે
        \[
        \Subst{M}{N}{x} \ident \Subst{\ande{i}{x}}{\andi{P_1}{P_2}}{x},
        \]
        લો; આ $\ande{i}{\andi{P_1}{P_2}}$ છે.
        તેનો કટ-ક્રમ પ્રકારની લંબાઈ $\len{!A}$ છે.
        \sourcecorrection{OLMD-283}{મૂળ અહીં ચલ $x$નો કટ-ક્રમ લે છે, પરંતુ કટ-ક્રમ ફક્ત સંકોચ્ય પદ માટે વ્યાખ્યાયિત છે. પ્રતિસ્થાપનથી બનેલા કટનો ક્રમ $x$ના પ્રકારની લંબાઈ છે.}
    \end{enumerate}
  \end{enumerate}
  $\ore{N}{x_1}{N_1}{x_2}{N_2}$ અને $P Q$ના કિસ્સા સમાન છે.
  \sourcecorrection{OLPG-0687-SUBST}{મૂળ સાબિતી કેટલાક કિસ્સા સમાન કહીને છોડે છે અને અધ્યાયમાં રજૂ થયેલ અસત્ય-લોપ તથા ક્રમચય-રૂપાંતરો માટે આ કટ-ક્રમની દલીલ આપતી નથી. અનુવાદ એ ખાલી જગ્યા છુપાવતો નથી; અહીંના પૂર્ણ સામાન્યીકરણ-દાવાની સ્રોત-સાબિતી અધૂરી છે.}
\end{proof}

\begin{lem}
  $M$ સંકોચાઈને $M'$ થાય અને $M$ના દરેક પોતાના કરતાં
  નાના સંકોચ્ય ઉપપદ~$N$ માટે $\cutrank{M} > \cutrank{N}$
  હોય, તો $\cutrank{M} > \maxrank{M'}$ છે.
\end{lem}

\begin{proof}
  કિસ્સાવાર સાબિતી કરીએ.
  \begin{enumerate}
  \item $M$નું સ્વરૂપ $\ande{i}{\andi{M_1}{M_2}}$ હોય,
    તો $M'$ એ $M_i$ છે. $M_i$નું દરેક ઉપપદ
    $M$નું પણ પોતાના કરતાં નાનું ઉપપદ છે,
    એટલે દાવો સાચો છે.
  \item $M$નું સ્વરૂપ
    $(\lambd[\typeof{x}{!A}][N])\typeof{Q}{!A}$ હોય, તો
    $M'$ એ $\Subst{N}{\typeof{Q}{!A}}{\typeof{x}{!A}}$ છે.
    $M'$માંનું એક સંકોચ્ય પદ લો. કાં તો~$N$માં
    તેને અનુરૂપ સંકોચ્ય પદનો કટ-ક્રમ સમાન છે,
    જે ધારણા મુજબ $\cutrank{M}$ કરતાં નાનો છે;
    અથવા તેનો કટ-ક્રમ $\len{!A}$ છે, જે વ્યાખ્યા મુજબ
    $\cutrank{(\lambd[x^{!A}][N])Q}$ કરતાં નાનો છે.
  \item $M$નું સ્વરૂપ
    \[
    \ore{\ori{i}{}{\typeof{N}{!A_i}}}
        {\typeof{x_1}{!A_1}}{\typeof{N_1}{!C}}
        {\typeof{x_2}{!A_2}}{\typeof{N_2}{!C}},
    \]
    હોય, તો $M' \ident \Subst{N_i}{N}{\typeof{x_i}{!A_i}}$ છે.
    $M'$માંનું એક સંકોચ્ય પદ લો. કાં તો~$N_i$માં
    તેને અનુરૂપ સંકોચ્ય પદનો કટ-ક્રમ સમાન છે,
    જે ધારણા મુજબ $\cutrank{M}$ કરતાં નાનો છે;
    અથવા તેનો કટ-ક્રમ $\len{!A_i}$ છે, જે વ્યાખ્યા મુજબ
    $\cutrank{\ore{\ori{i}{}{\typeof{N}{!A_i}}}
      {\typeof{x_1}{!A_1}}{\typeof{N_1}{!C}}
      {\typeof{x_2}{!A_2}}{\typeof{N_2}{!C}}}$ કરતાં નાનો છે.
  \end{enumerate}
  \sourcecorrection{OLPG-0687-COPIED}{મૂળના બીજા અને ત્રીજા કિસ્સાની સંકોચ્ય પદોની યાદીમાં પ્રતિસ્થાપિત પદમાંથી નકલ થતા આંતરિક સંકોચ્ય ઉપપદોનો કિસ્સો સ્પષ્ટ કરેલો નથી. દાવો ચકાસવા આ વધારાનો કિસ્સો પણ જરૂરી છે; સ્રોતમાં તે અધૂરો છે.}
\end{proof}

\begin{thm}
  બધાં સાબિતી-પદો સામાન્ય સ્વરૂપમાં ઘટે છે;
  બધી નિષ્પત્તિઓ સામાન્ય નિષ્પત્તિઓમાં ઘટે છે.
\end{thm}

\begin{proof}
  બીજો દાવો પહેલામાંથી અનુસરે છે. પહેલા દાવા માટે
  $m=\maxrank{M}$ પર પૂર્ણ આગમન કરીએ; અહીં $M$
  સાબિતી-પદ છે.
  \begin{enumerate}

  \item $m=0$ હોય, તો $M$ પહેલેથી સામાન્ય છે.
  \item
    નહિતર~$n$ પર આગમન કરીએ; તે~$M$માં
    કટ-ક્રમ~$m$ હોય એવાં સંકોચ્ય પદોની સંખ્યા છે.
    \begin{enumerate}
    \item $n=1$ હોય, તો એવું સંકોચ્ય પદ~$N$ પસંદ કરો
      કે સંકોચ્ય હોય એવા તેના દરેક પોતાના કરતાં
      નાના ઉપપદ~$P$ માટે $m = \cutrank{N} > \cutrank{P}$
      હોય. દરેક પદનાં ફક્ત સીમિત ઉપપદો હોય છે,
      એટલે આવું સંકોચ્ય પદ હોવું જ જોઈએ.

      $N$ના સંકોચનનું પરિણામ~$N'$ હોય.
      ઉપપાદ્યથી $\maxrank{N'} < \maxrank{N}$ છે.
      તેથી $\cutrank=m$ હોય એવાં સંકોચ્ય પદોની
      સંખ્યા~$n$ ઘટે છે. એટલે $m$ પણ ઘટે છે
      ($1$ કે વધુ જેટલું), અને~$m$ માટેની
      આગમન-ધારણા લાગુ કરી શકાય છે.
    \item આગમનના પગલા માટે $n>1$ ધારો.
      પ્રક્રિયા સમાન છે; ફક્ત~$n$ ઘટીને પણ ધન
      રહે છે, એટલે~$m$ બદલાતો નથી.
      હવે~$n$ માટેની આગમન-ધારણા લાગુ કરીએ.
    \end{enumerate}
\end{enumerate}
\sourcecorrection{OLPG-0687-CONTEXT}{મૂળની આ આગમન-દલીલ ફક્ત ઘટાડેલા ઉપપદની અંદરનો કટ-ક્રમ નિયંત્રિત કરે છે. તે ઉપપદને તેના સંકોચનના પરિણામથી બદલતાં બહારના આવરણમાં નવો સંકોચ્ય પદ ખુલ્લો પડી શકે છે, ખાસ કરીને વિકલ્પ-લોપનું પરિણામ વિધેય હોય ત્યારે. આવા નવા કટના ક્રમ અને મહત્તમ-ક્રમવાળા કટની કુલ સંખ્યા માટે મૂળ દલીલ આપતું નથી.}
\end{proof}

પ્રકારિત $\lambd$-પદોનું સામાન્યીકરણ થાય છે તેનો
અર્થ એ કે દરેક આવા પદનું \emph{સામાન્ય સ્વરૂપ છે}.
$\lambd$-પદને $\beta$-ઘટાડાથી સામાન્ય સ્વરૂપમાં
લાવવું ગણતરી કરવી છે, અને સામાન્ય સ્વરૂપ
ગણતરીનું મૂલ્ય છે એમ માનીએ, તો તેનો અર્થ
એ થાય કે પ્રકારિત $\lambd$-કલનની દરેક ગણતરી
આખરે થોભે છે અને મૂલ્ય આપે છે. વળી,
પ્રકારનું સંરક્ષણ એ મૂલ્યનો પ્રકાર યોગ્ય
હોવાની ખાતરી આપે છે. દાખલા તરીકે, $N$નો
પ્રકાર $!A \lif !B$ હોય (એટલે તે~$!A$
પ્રકારના આર્ગ્યુમેન્ટ પર કામ કરતું અને~$!B$
પ્રકારનાં મૂલ્યો આપતું વિધેય હોય) અને તેને~$!A$
પ્રકારના પદ~$M$ (ઇનપુટ) પર લાગુ કરીએ,
તો $(NM)$ પદ~$N'$ (આઉટપુટ)માં ઘટે છે,
અને તે આઉટપુટનો પ્રકાર~$!B$ હોવાની ખાતરી છે.
\sourcecorrection{OLPG-0687-STRONG}{આગળનો સબળ સામાન્યીકરણનો દાવો દરેક ઘટાડાક્રમ માટે છે, પરંતુ ઉપરની સાબિતી ફક્ત એક નિર્ધારિત પસંદગીથી સામાન્ય સ્વરૂપ મળે તે બતાવવાનો પ્રયાસ કરે છે. સબળ સામાન્યીકરણની અલગ સાબિતી મૂળમાં નથી; તેથી દરેક ગણતરી થોભવાના દાવાને ઉપરના પ્રમેયથી સાબિત ગણ્યો નથી.}

હકીકતમાં પ્રકારિત $\lambd$-કલનનાં પદો માટે
ઉપપદો પર $\beta$-ઘટાડા લગાવવાની કોઈ એક
રીતથી જ સામાન્ય સ્વરૂપ મળતું નથી; ઉપપદો
પર $\beta$-ઘટાડા લગાવવાની \emph{દરેક}
રીત આખરે સામાન્ય સ્વરૂપમાં સમાપ્ત થાય છે.
આ ગુણધર્મ \emph{સબળ સામાન્યીકરણ} કહેવાય.
તેનો અર્થ એ કે ગણતરી આખરે થોભીને મૂલ્ય
આપે તેની ખાતરી આપતી કોઈ એક રીત જ નથી,
પણ ગણતરી કરવાની \emph{દરેક} રીત
આખરે મૂલ્ય આપે છે.

ગણતરી કરવાની જુદી રીતો \emph{એક જ} મૂલ્ય
આપે છે તે કેવી રીતે જાણીએ? અત્યાર સુધી
પ્રકારના સંરક્ષણથી ફક્ત એટલું જાણીએ છીએ
કે મળતાં બધાં મૂલ્યોનો પ્રકાર સમાન છે.
પ્રકારિત $\lambd$-કલનનો બીજો અગત્યનો
ગુણધર્મ છે: તે \emph{સંગમક્ષમ}, એટલે
\emph{ચર્ચ--રોસર}, છે.
$N \red N_1$ અને $N \red N_2$ હોય,
તો એવું પદ~$N'$ છે કે $N_1 \red N'$
અને $N_2 \red N'$ થાય. યાદ કરો કે
$N \red N'$ એટલે શૂન્ય કે વધુ
$\redone$-ઘટાડા પછી $N'$ મળે છે.
$N_1$ સામાન્ય સ્વરૂપમાં હોય, તો
$N_1 \red N'$ થાય એવું એકમાત્ર પદ~$N'$
એ~$N_1$ પોતે જ છે (શૂન્ય
$\redone$-ઘટાડા વડે). તેથી $N_1$ અને
$N_2$ બંને સામાન્ય સ્વરૂપમાં હોય, તો
$N_1 = N' = N_2$ છે. એટલે~$\red$
સબળ સામાન્યીકરણ ધરાવતો હોવાથી પદને
ઘટાડવાની દરેક રીત આખરે સામાન્ય સ્વરૂપ આપે છે.
વળી~$\red$ ચર્ચ--રોસર હોવાથી કોઈ પણ
બે સામાન્ય સ્વરૂપ સમાન છે. આમ, પ્રકારિત
$\lambd$-કલનમાં ગણતરી હંમેશાં થોભે છે
અને તે કેવી રીતે આગળ વધી હોય તોપણ
મૂલ્ય, એટલે સામાન્ય સ્વરૂપ, એક જ હોય છે.

સંબંધ સંગમક્ષમ છે તે સાબિત કરવું સામાન્ય રીતે
અઘરું છે. જોકે આ નબળી શરત સહેલાઈથી
સ્થાપી શકીએ:

\begin{prop}[નબળો ચર્ચ--રોસર ગુણધર્મ]
$N \redone N_1$ અને $N \redone N_2$ હોય,
તો એવું પદ~$N'$ છે કે $N_1 \red N'$
અને $N_2 \red N'$ થાય.
\end{prop}

\begin{proof}
  $N$માં સંકોચ્ય પદ~$M_1$ અથવા~$M_2$ને
  તેના સંકોચનના પરિણામથી બદલી અનુક્રમે
  $N_1$ અને $N_2$ મળે છે. સંકોચ્ય પદો
  $M_1$ અને $M_2$ એ~$N$નાં ઉપપદો છે,
  એટલે તેઓ આંશિક રીતે પરસ્પર આવરી લેતાં નથી.
  કાં તો એક બીજાનું ઉપપદ હોય છે, અથવા
  તેઓ~$N$માં જુદાં, એકબીજાને ન આવરી લેતાં
  સ્થાનો પર આવે છે. બીજો કિસ્સો લો:
  $N$નું સ્વરૂપ $N(M_1,M_2)$ છે.
  $M_1$ને તેના સંકોચનના પરિણામ~$M_1'$થી
  બદલો, તો $N_1$ એટલે $N(M_1',M_2)$
  મળે; અને $M_2$ને તેના પરિણામ~$M_2'$થી
  બદલો, તો $N_2$ એટલે $N(M_1,M_2')$
  મળે. બંને $N(M_1',M_2')$માં ઘટે છે.
  પહેલા કિસ્સા માટે $M_2$ એ~$M_1$નું
  ઉપપદ છે એમ ધારો (બીજો કિસ્સો સમમિત છે),
  અને $M_1$ કયા પ્રકારનું સંકોચ્ય પદ છે
  તે મુજબ કિસ્સા પાડો. દાખલા તરીકે,
  $M_1 \ident \proj{1}{\pair{O_1}{O_2}}$
  હોય, તો તે $O_1$માં ઘટે છે. $M_2$
  એ~$O_1$નું ઉપપદ હોય, તો $M_2$ને
  રૂપાંતરિત કરવાથી $O_1'$ મળે છે.
  એટલે પહેલા $M_1$ અને પછી $M_2$ને
  રૂપાંતરિત કરતાં $O_1'$ મળે છે.
  પણ પહેલા $M_2$ને રૂપાંતરિત કરતાં
  $\proj{1}{\pair{O_1'}{O_2}}$ મળે છે,
  જે પણ~$O_1'$માં ઘટે છે.
  \sourcecorrection{OLMD-284}{મૂળ છેલ્લે $O_2$ કહે છે, પરંતુ પ્રથમ પ્રક્ષેપનું પરિણામ $O_1'$ જ છે, જે બીજી ઘટાડા-દિશામાં પણ મળ્યું છે.}
  \sourcecorrection{OLPG-0687-LOCAL}{મૂળની આ સાબિતી આંતરિક સંકોચન માટે ફક્ત પ્રથમ પ્રક્ષેપના રાખેલા ઘટકનો કિસ્સો આપે છે. કાઢી નાખાતા ઘટક, વિધેય-અરજી અને વિકલ્પ-લોપના અંદર-બહારના ઘટાડા તથા આગળના વિભાગના ક્રમચય-રૂપાંતરોના કિસ્સા અહીં સાબિત નથી.}
\end{proof}

\begin{lem}[ન્યૂમનનું ઉપપાદ્ય]
$\red$ સબળ સામાન્યીકરણ ધરાવતો હોય અને નબળો
ચર્ચ--રોસર ગુણધર્મ ધરાવતો હોય, તો તે સંગમક્ષમ છે.
\end{lem}

\begin{proof}
  $\red$ સબળ સામાન્યીકરણ ધરાવે છે, એટલે
  દરેક ઘટાડા-અનુક્રમ સામાન્ય સ્વરૂપ આપે છે.
  સામાન્ય સ્વરૂપ અનન્ય છે જો અને તો જ~$\red$
  સંગમક્ષમ છે. તેથી ધારો કે કોઈ પદ~$N$ને
  બે જુદાં સામાન્ય સ્વરૂપો $N'$ અને $N''$
  છે, જે આ ઘટાડા-અનુક્રમોથી મળે છે:
  \begin{align*}
  N & = N_0' \redone N_1' \redone \dots \redone N_k' = N' \text{ અને}\\
  N & = N_0'' \redone N_1'' \redone \dots \redone N_l'' = N''
  \end{align*}
  (ઘટાડા-અનુક્રમની લંબાઈ $>0$ હોવી જ જોઈએ;
  નહિતર $N$ પહેલેથી સામાન્ય સ્વરૂપમાં હોત.)
  \sourcecorrection{OLMD-285}{મૂળ બંને અનુક્રમમાં શરૂઆતને અનુક્રમણિકાંક 1 આપે છે, પરંતુ આગળની દલીલ 1વાળાં પદોને પ્રથમ ઘટાડા પછીનાં પદો માને છે. શરૂઆતને 0 અને પ્રથમ ઉત્તરગામીને 1 આપ્યાં છે.}

  $N_1' = N_1''$ હોય, તો તેને બે જુદાં
  સામાન્ય સ્વરૂપો ($N'$ અને $N''$) છે.
  $N_1' \neq N_1''$ હોય, તો નબળા
  ચર્ચ--રોસર ગુણધર્મથી એવું $N'''$ મળે છે
  કે $N_1' \red N'''$ અને $N_1'' \red N'''$
  થાય. $N' \neq N''$ હોવાથી,
  કાં તો $N''' \neq N'$ અથવા
  $N''' \neq N''$ છે. પરિણામે,
  $N_1'$ અથવા $N_1''$ને બે જુદાં
  સામાન્ય સ્વરૂપો છે. આપણે બતાવ્યું કે
  $N$ને બે જુદાં સામાન્ય સ્વરૂપો હોય,
  તો એવું $N_1$ છે કે $N \redone N_1$
  થાય અને $N_1$ને પણ બે જુદાં સામાન્ય
  સ્વરૂપો હોય. આ રીતે આગળ વધીને
  $N \redone N_1 \redone N_2 \redone \dots$
  મેળવી શકાય. પરંતુ~$\redone$ સબળ
  સામાન્યીકરણ ધરાવે છે, એટલે આ અશક્ય છે.
  \sourcecorrection{OLPG-0687-NEWMAN}{મૂળની આ દલીલમાં સામાન્ય સંયુક્ત ઉત્તરગામી $N'''$ પોતે સામાન્ય સ્વરૂપ હોવાની ખાતરી નથી. બે જુદાં સામાન્ય સ્વરૂપોનો દાવો આગળ વધારવા તેના સામાન્ય સ્વરૂપ સુધી ઘટાડાનો પગલો જરૂરી છે; મૂળ તે પગલો આપતું નથી.}
\end{proof}
% END OLP-0687
\endgroup


\OLEndChapterHook


\begin{editorial}
આગળનો સંબંધિત અંશ મૂળના સક્રિય આયાતક્રમમાં નથી. તેને અહીં સ્પષ્ટ વૈકલ્પિક સંદર્ભમાં જાળવ્યો છે.
\end{editorial}
\begingroup
\makeatletter\def\ol@sectioncs{section}\makeatother

% BEGIN OLP-0694
\olfileid{int}{pty}{snd}

\olsection{સિક્વન્ટ-શૈલીનું સ્વાભાવિક નિગમન}
\phantomsection\label{guunit:OLP-0694}


સ્વાભાવિક નિગમનની એવી નિષ્પત્તિ હોય
કે તેની અનિવૃત્ત ધારણાઓ~$\Gamma$
અને નિષ્કર્ષ~$!A$ હોય, તો
$\Gamma \Sequent !A$ લખીએ;
$\Gamma$ ખાલી હોય, તો $\Sequent !A$ લખીએ.

$\Gamma \cup \{!A_1, \dots, !A_n\}$ માટે
$\Gamma, !A_1, \dots, !A_n$ અને
$\Gamma \cup \Delta$ માટે $\Gamma, \Delta$ લખીએ.

નોંધો કે $\Gamma \Sequent !A \land !B$ હોય,
એટલે અનિવૃત્ત ધારણાઓ~$\Gamma$ અને
અંતિમ સૂત્ર~$!A \land !B$ ધરાવતી
નિષ્પત્તિ હોય, તો સૌથી નીચે
$\Elim{\land}$ લાગુ કરીને એ જ
અનિવૃત્ત ધારણાઓ અને નિષ્કર્ષ~$!A$
ધરાવતી નિષ્પત્તિ મળે છે.
એટલે $\Gamma \Sequent !A \land !B$
હોય, તો $\Gamma \Sequent !A$ છે.
\sourcecorrection{OLMD-295}{મૂળ આ ખુલાસાની શરૂઆતમાં બંને સંયોજ્યો $!A$ લખે છે, પરંતુ સારવાક્ય અને નિયમવૃક્ષમાં $!A$ તથા $!B$ છે. ખુલાસાના બંને સ્થાનોને એ જ સામાન્ય સંયોજન સાથે સુસંગત કર્યા છે.}
\begin{prooftree}
  \Axiom$\Gamma \fCenter !A \land !B$
  \RightLabel{$\Elim{\land}$}
  \UnaryInf$\Gamma \fCenter !A$
  \DisplayProof\qquad\bottomAlignProof
  \Axiom$\Gamma \fCenter !A \land !B$
  \RightLabel{$\Elim{\land}$}
  \UnaryInf$\Gamma \fCenter !B$
\end{prooftree}
$\Elim{\land}$ ચિહ્ન સ્વાભાવિક નિગમનના
એ જ નામના નિયમ સાથેનો સંબંધ સૂચવે છે.

એ જ રીતે $\Gamma, !A \Sequent !B$ હોય,
એટલે અનિવૃત્ત ધારણાઓ $\Gamma, !A$
અને અંતિમ સૂત્ર~$!B$ ધરાવતી
નિષ્પત્તિ હોય એમ ધારો.
$\Intro{\lif}$ નિયમ લગાવીએ, તો
અનિવૃત્ત ધારણાઓ~$\Gamma$
અને અંતિમ સૂત્ર~$!A \lif !B$
ધરાવતી નિષ્પત્તિ મળે છે,
એટલે $\Gamma \Sequent !A \lif !B$.
આ રીતે ધારણાઓની નિવૃત્તિ
વધુ સ્પષ્ટ થાય છે તે નોંધો.
\begin{prooftree}
  \Axiom $\Gamma, !A \fCenter !B$
  \RightLabel{$\Intro{\lif}$}
  \UnaryInf $\Gamma \fCenter !A \lif !B$
\end{prooftree}

બીજા નિયમોમાંથી પણ આ રીતે નિષ્કર્ષો
મેળવી શકીએ. તે આમ સ્પષ્ટ થાય છે:
\begin{gather*}
  \Axiom $\Gamma \fCenter !A$
  \Axiom $\Delta \fCenter !B$
  \RightLabel{$\Intro{\land}$}
  \BinaryInf $\Gamma,\Delta \fCenter !A \land !B$
  \DisplayProof\\
  \Axiom $\Gamma \fCenter !A \land !B$
  \RightLabel{$\Elim{\land}_1$}
  \UnaryInf $\Gamma \fCenter !A$
  \DisplayProof
  \qquad
  \Axiom $\Gamma \fCenter !A \land !B$
  \RightLabel{$\Elim{\land}_2$}
  \UnaryInf $\Gamma \fCenter !B$
  \DisplayProof
  \\
  \Axiom $\Gamma \fCenter !A$
  \RightLabel{$\Intro{\lor}_1$}
  \UnaryInf $\Gamma \fCenter !A \lor !B$
  \DisplayProof\qquad
  \Axiom $\Gamma \fCenter !B$
  \RightLabel{$\Intro{\lor}_2$}
  \UnaryInf $\Gamma \fCenter !A \lor !B$
  \DisplayProof\\
  \Axiom $\Gamma \fCenter !A \lor !B$
  \Axiom $\Delta, !A \fCenter !C$
  \Axiom $\Delta', !B \fCenter !C$
  \RightLabel{$\Elim{\lor}$}
  \TrinaryInf $\Gamma, \Delta, \Delta' \fCenter !C$
  \DisplayProof
  \\
  \Axiom $\Gamma, !A \fCenter !B$
  \RightLabel{$\Intro{\lif}$}
  \UnaryInf $\Gamma \fCenter !A \lif !B$
  \DisplayProof
  \qquad
  \Axiom $\Delta \fCenter !A \lif !B$
  \Axiom $\Gamma \fCenter !A$
  \RightLabel{$\Elim{\lif}$}
  \BinaryInf $\Gamma, \Delta \fCenter !B$
  \DisplayProof
  \\
  \Axiom $\Gamma \fCenter \lfalse$
  \RightLabel{$\FalseInt$}
  \UnaryInf $\Gamma \fCenter !A$
  \DisplayProof
\end{gather*}

કોઈ પણ ધારણા પોતે~$!A$માંથી $!A$ની
નિષ્પત્તિ છે; એટલે હંમેશાં
$!A \Sequent !A$ છે.

\begin{prooftree}
  \AxiomC{$ $}
  \UnaryInfC{$!A \fCenter !A$}
\end{prooftree}

આ બધા નિયમોને સ્વાભાવિક નિગમનની કઈ
નિષ્પત્તિઓ અસ્તિત્વ ધરાવે છે તે
વિશેનું કલન ગણી શકાય. તેને સ્વાભાવિક
નિગમનની નોંધની વૈકલ્પિક રીત પણ
ગણી શકાય, જેમાં દરેક પગલે ફક્ત
નિષ્પન્ન થયેલું સૂત્ર જ
નહીં, પણ જે અનિવૃત્ત ધારણાઓમાંથી
તે નિષ્પન્ન થયું હોય તે પણ નોંધાય છે.

\begin{prooftree}
  \Axiom$!A \fCenter !A$
  \UnaryInf $!A \fCenter !A \lor (!A \lif \lfalse)$
  \Axiom $!B \fCenter !B$
  \BinaryInf$!A, !B \fCenter \lfalse$
  \UnaryInf$!B \fCenter !A \lif \lfalse$
  \UnaryInf$!B \fCenter !A \lor (!A \lif \lfalse)$
  \Axiom$!B \fCenter !B$
  \BinaryInf$!B \fCenter \lfalse$
  \UnaryInf$\fCenter !B \lif \lfalse$
\end{prooftree}
અહીં $!B$ એ
$(!A \lor (!A \lif \lfalse))\lif \lfalse$નું
ટૂંકું રૂપ છે.
\sourcecorrection{OLMD-296}{મૂળના અંતિમ વૃક્ષમાં પૂર્વાંગના સૂત્ર $!B$ પછી એક વધારાનો અનુસરણ-સંયોજક તથા ચાર વધારાના ખુલ્લા કૌંસ છે. તેમને કાઢ્યાં છે; તેથી દરેક પગલાનો સંદર્ભ આપેલી $!B$ ધારણા છે અને અંતિમ નિષ્કર્ષ એ ધારણાથી અસત્ય મળે છે તે જ કહે છે.}
% END OLP-0694
\endgroup
% END OLP-0690
\endgroup


\OLEndPartHook


\begin{editorial}
આગળનું સંબંધિત પ્રકરણ મૂળ ભાગના સક્રિય આયાતક્રમમાં નથી. સંપૂર્ણ સ્રોત-આવરણ માટે તેને તેના વિષયના ભાગમાં જાળવ્યું છે.
\end{editorial}
\begingroup
\makeatletter\def\ol@sectioncs{section}\makeatother

% BEGIN OLP-0681
\olchapter{pt}{ps}{સાબિતી-શોધ}
\olchapterid{pt}{ps}
\phantomsection\label{guunit:OLP-0681}


\begingroup
\makeatletter\def\ol@sectioncs{section}\makeatother

% BEGIN OLP-0680
\olfileid{pt}{ps}{int}

\olsection{પરિચય}
\phantomsection\label{guunit:OLP-0680}


\Log{G3c} અથવા
\Log{N1i} જેવા
સાબિતી-તંત્રોના
સ્વયંસિદ્ધો અને
અનુમાન-નિયમો
નક્કી કરે છે
કે સિક્વન્ટો
અથવા સૂત્રોનાં
કયાં વૃક્ષો
યોગ્ય સાબિતીઓ
ગણાય. આવું
સિક્વન્ટો અથવા
સૂત્રોનું
વૃક્ષ અપાયું
હોય (દરેક
પગલે કયો નિયમ
લાગ્યો છે
અથવા ધારણાઓ
અને તેમને
મુક્ત કરતા
નિયમોનાં
ચિહ્નો જેવી
વધારાની માહિતી
સાથે પણ),
તો સ્વયંસિદ્ધો
અને નિયમોની
વ્યાખ્યાઓથી
તે યોગ્ય
સાબિતી છે
કે નહીં તે
તપાસી શકીએ.
પણ અપાયેલા
સિક્વન્ટની
અથવા અપાયેલી
ધારણાઓમાંથી
અપાયેલા
સૂત્રની
સાબિતી
\emph{શોધવી}
અલગ અને
સામાન્ય રીતે
વધુ મુશ્કેલ
સમસ્યા છે.
આ \emph{સાબિતી-શોધ}
ની સમસ્યા છે.

સાબિતી-શોધ માટે
ખૂબ સાદી
રીત છે.
સૂત્રોને
સીમિત વર્ણમાળાના
ચિહ્નોના ક્રમ
તરીકે વિચારી
શકાય. સિક્વન્ટો
અને સિક્વન્ટો
અથવા
સૂત્રોનાં
વૃક્ષોને પણ
સૂત્રોના
ક્રમ તરીકે
વિચારી શકાય
(વૃક્ષની શાખાઓ
બતાવવા ખાસ
કૌંસ વાપરીને).
સ્વયંસિદ્ધો
અને નિયમો
વડે ચિહ્નોનો
ક્રમ યોગ્ય
સાબિતી
દર્શાવે છે
કે નહીં તે
યાંત્રિક રીતે
તપાસી શકાય.
તેથી \emph{બધી
શક્ય} યોગ્ય
સાબિતીઓ
યાંત્રિક રીતે
બનાવી શકાય:
પ્રથમ સાબિતીઓ
દર્શાવતી
વર્ણમાળાના
ચિહ્નોના બધા
ક્રમ બનાવીએ
અને દરેક
ઉમેદવાર યોગ્ય
સાબિતી છે
કે નહીં તપાસીએ.
આ સાદી
સાબિતી-શોધ
એવી છે કે
અપાયેલા સિક્વન્ટની
સાબિતી મળે
ત્યાં સુધી
બધી યોગ્ય
સાબિતીઓ
બનાવતાં રહીએ
(અથવા અપાયેલી
ધારણાઓમાંથી
ખુલ્લી રહે
તેવી ધારણાઓ
સાથે અપાયેલા
સૂત્રની
સાબિતી મળે
ત્યાં સુધી).

વધુ સારી
રીત એ છે
કે હાથથી
સાબિતી
શોધવા વપરાતી
સહજ રીત
અપનાવીએ:
અંતિમ સિક્વન્ટથી
શરૂ કરી
સૂત્ર
પસંદ કરીએ
અને અનુમાન-નિયમ
``ઊલટી દિશામાં''
લગાવી એક
અથવા વધુ
આધારવિધાન
બનાવીએ; તેમની
સાબિતીઓ
હોય, તો
છેલ્લા અનુમાન
સાથે જોડીને
અપાયેલા સિક્વન્ટની
સાબિતી
મળે. સ્વાભાવિક
નિગમન-તંત્રમાં
સાબિતી
શોધવા આવી
જ, જોકે
વધુ જટિલ,
રીત વપરાય.

પણ આ રીત
અચૂક નથી:
ખોટો નિયમ
અથવા
સૂત્રોનો
ખોટો ક્રમ
પસંદ કરો,
તો આગળ
માર્ગ ન
રહે. તે
સંપૂર્ણ નિર્દિષ્ટ
પણ નથી:
દરેક તબક્કે
પસંદ કરવા
એકથી વધુ
સૂત્ર
અથવા ઊલટી
દિશામાં
લગાવવા એકથી
વધુ નિયમ
હોઈ શકે.
\emph{સાબિતી-શોધ
અલ્ગોરિધમ}
એવી સંપૂર્ણ
નિર્દિષ્ટ
પ્રક્રિયા છે
જે સાબિતી
હોય તો
તે શોધે,
એટલે અચૂક
હોય.

કેટલાંક
સાબિતી-તંત્રોમાં
સાબિતી-શોધના
અલ્ગોરિધમ
બીજાં કરતાં
સરળ બને છે.
સિક્વન્ટ-તંત્રોમાં
સૂત્રનો
મુખ્ય કારક
અને તે ડાબે
છે કે જમણે
એ ઊલટી
દિશામાં લગાવવાનો
નિયમ નક્કી
કરે છે.
દાખલા તરીકે,
$!A \land
!B$ જમણી
બાજુ હોય એવા
$\Gamma \Sequent \Delta, !A \land !B$ની
સાબિતી
શોધવી હોય,
તો
\RightR{\land}
ઊલટી દિશામાં
લગાવી
$\Gamma \Sequent \Delta, !A$
અને
$\Gamma \Sequent \Delta, !B$
એ બે અનુરૂપ
આધારવિધાન
મેળવવાં.
પણ તંત્ર
\Log{G1c}
હોય તો
સંકોચનને લીધે
વિકલ્પો વધે
છે: તમે
\RightR{\Contraction}
ઊલટી દિશામાં
લગાવી તેના
સ્થાને
$\Gamma \Sequent \Delta, !A \land
!B, !A \land !B$
આધારવિધાન
પણ બનાવી
શકો. પછી
$\Gamma \Sequent
\Delta, !A \land !B, !A \land !B, !A \land !B$
અને એમ
આગળ.
\Log{G2c}માં
તો વધુ
મુશ્કેલ છે.
\RightR{\land}
ઊલટી દિશામાં
લગાવવા
$\Gamma_1$,
$\Gamma_2$
એવા અનુમાનવા
પડે કે
$\Gamma
= \Gamma_1 \cup \Gamma_2$;
એ જ રીતે
$\Delta_1$,
$\Delta_2$
એવા કે
$\Delta
= \Delta_1 \cup \Delta_2$.
પછી
$\Gamma_1 \Sequent \Delta_1, !A$
અને
$\Gamma_2 \Sequent \Delta_2, !B$
આધારવિધાનો
અજમાવવાં
પડે. ખોટું
અનુમાન પાછળથી
બંધ માર્ગે
લઈ જાય,
તો શરૂઆતમાં
પાછા ફરી
અલગ
$\Gamma_1$,
$\Gamma_2$,
$\Delta_1$,
$\Delta_2$
પસંદ કરવા
પડે.
સૂત્ર
$!A \lor !B$
હોય તો
\RightR{\lor}
માટેના બે
વિકલ્પોમાંથી
એક પસંદ
કરી
$\Gamma \Sequent \Delta,
!A$
અથવા
$\Gamma \Sequent \Delta, !B$
આધારવિધાન
બનાવવું પડે.
ખોટી પસંદગી
ફરી પાછળ
ફરવા મજબૂર
કરે છે.

\Log{G3c}માં
આ સમસ્યાઓ
નથી. તેમાં
બંધારણીય
નિયમો નથી,
અને મુખ્ય કારક
તથા સૂત્ર
કઈ બાજુ છે
તે પરથી
અનુમાન-નિયમ
નક્કી થાય
છે. તેથી
\Log{G1c} કે
\Log{G2c} કરતાં
\Log{G3c}
સાબિતી-શોધ
માટે વધુ
અનુકૂળ છે.
સફળ શોધ
\Log{G3c}માં
સાબિતી
આપે છે,
અને પછી
તે સાબિતીને
\Log{G1c}
અથવા
\Log{G2c}
(અથવા
\Log{N1c})
માં રૂપાંતરિત
કરી શકાય.
આમ \Log{G3c}
માટે સાબિતી-શોધ
અલ્ગોરિધમ
અને
\Log{G3c}ની
સાબિતીઓને
બીજાં તંત્રોની
સાબિતીઓમાં
રૂપાંતરિત
કરતા અલ્ગોરિધમ
હોય, તો
તે તંત્રો
માટે પણ
સાબિતી-શોધ
ઉકેલાય છે.

સાબિતી-શોધ
અલ્ગોરિધમ
અચૂક છે
તે કેવી
રીતે જાણીએ?
જો અલ્ગોરિધમ
સાબિતી
શોધવામાં
નિષ્ફળ જાય,
તો કોઈ
સાબિતી
હોઈ જ
ન શકે
તે બતાવવું
પડે. ખરાબ
અલ્ગોરિધમ
પસંદ કરીએ
તો સાબિતી
હોવા છતાં
તે શોધવામાં
નિષ્ફળ જઈ
શકે. સાદા
અલ્ગોરિધમને
આ મુશ્કેલી
નથી, કારણ
કે તે
બધી શક્ય
સાબિતીઓ
ચકાસે છે.
પણ સાબિતી-શોધ
અલ્ગોરિધમ
કાર્યક્ષમ
અને ઓછામાં
ઓછો પ્રયત્ન
કરતો હોવો
જોઈએ.

સાબિતી-શોધ
અલ્ગોરિધમ
અચૂક છે
તે બતાવવાનો
મુખ્ય વિચાર
આ છે: અલ્ગોરિધમ
નિષ્ફળ જાય,
તો તેની
નિષ્ફળતાની
રીત પરથી
સિક્વન્ટની
સાબિતી
શક્ય નથી
તે બતાવીએ.
તે માટે
સિક્વન્ટ
પ્રામાણ્ય નથી
એવું બતાવીએ.
શાસ્ત્રીય
પ્રથમ-ક્રમ
તર્કમાં
$\Gamma \Sequent \Delta$
પ્રામાણ્ય ત્યારે
છે જ્યારે
દરેક
સંરચના
$\Gamma$ના
ઓછામાં ઓછા
એક સૂત્રને
અસત્ય અથવા
$\Delta$ના
ઓછામાં ઓછા
એક સૂત્રને
સત્ય બનાવે.
\Log{G3c}
\emph{યથાર્થ}
છે: તે ફક્ત
પ્રામાણ્ય
સિક્વન્ટો
સાબિત
કરે છે.
આ સાબિતીઓ
પર આગમનથી
સ્થાપિત થાય
છે: સ્વયંસિદ્ધો
સ્પષ્ટ રીતે
પ્રામાણ્ય છે,
અને નિયમનાં
આધારવિધાનો
પ્રામાણ્ય હોય
તો તેનું
નિષ્કર્ષ પણ
પ્રામાણ્ય છે.
સાબિતી-શોધ
અલ્ગોરિધમ
અચૂક છે
તે માટે
બતાવીએ કે
$\Gamma \Sequent \Delta$ની
સાબિતી
શોધ નિષ્ફળ
જાય, તો
સંરચના
એવી વ્યાખ્યાયિત
કરી શકાય
જેમાં
$\Gamma$નાં
બધાં
સૂત્રો
સત્ય અને
$\Delta$નાં
બધાં
સૂત્રો
અસત્ય હોય.
આથી
\Log{G3c}
\emph{પૂર્ણ}
પણ છે:
$\Gamma \Sequent \Delta$
પ્રામાણ્ય
હોય, તો
તેની
\Log{G3c}-સાબિતી
હોય છે.
દરેક નિષ્ફળ
શોધ
$\Gamma \Sequent \Delta$
અપ્રામાણ્ય
છે તે બતાવે
છે, એટલે
પ્રામાણ્ય
સિક્વન્ટની
દરેક શોધ
સફળ થાય.
% END OLP-0680
\endgroup

\begingroup
\makeatletter\def\ol@sectioncs{section}\makeatother

% BEGIN OLP-0683
\olfileid{pt}{ps}{sal}

\olsection{\Log{G3c} માટે સાબિતી-શોધ અલ્ગોરિધમ}
\phantomsection\label{guunit:OLP-0683}


\Log{G3c} માટે
સાબિતી-શોધ
અલ્ગોરિધમ વર્ણવીએ.
શક્ય એવો
આ એકમાત્ર
અલ્ગોરિધમ નથી,
અને સૌથી
કાર્યક્ષમ પણ
નથી. પણ તે
સાચો છે:
સિક્વન્ટ
$\Gamma \Sequent \Delta$ની
સાબિતી
હોય, તો
તે આખરે
થોભીને
સાબિતી
શોધશે. તે
અચૂક પણ
છે: સાબિતી
મળ્યા વિના
થોભે અથવા
કદી થોભે
જ નહીં,
તો
$\Gamma
\Sequent \Delta$
શરૂઆતથી જ
અપ્રામાણ્ય
છે.

અલ્ગોરિધમનું
આવશ્યક લક્ષણ
એ છે કે
$\Gamma \Sequent \Delta$
અથવા શોધ
દરમિયાન મળતા
સિક્વન્ટનું
દરેક
સૂત્ર
જરૂર પડે
તેટલી વાર
``ઘટાડાય'':
એટલે તેને
ઊલટી દિશામાં
લાગુ કરેલા
અનુમાન-નિયમનું
મુખ્ય સૂત્ર
ગણાય. આ
લક્ષણને
\emph{ન્યાયીપણું}
કહીએ છીએ.

અલ્ગોરિધમ
તબક્કાવાર
કામ કરે છે.
દરેક તબક્કાનું
પરિણામ
સિક્વન્ટોનું
વૃક્ષ છે;
આગલા તબક્કામાં
પહેલેથી સ્વયંસિદ્ધ
ન હોય એવા
દરેક સૌથી
ઉપરના સિક્વન્ટમાં
સૂત્ર
ઘટાડીને
નવું વૃક્ષ
બનાવાય છે.

તબક્કા~$0$માં
$\Gamma \Sequent \Delta$થી
શરૂ કરીએ.
ન્યાયીપણું
સુનિશ્ચિત કરવા
$\Gamma$
અને~$\Delta$નાં
સૂત્રોને
અનુક્રમણિકાંક
તરીકે સંખ્યાઓ
આપીએ, જેથી
અલગ
સૂત્રોને
અલગ સંખ્યા
મળે. દાખલા
તરીકે ડાબેથી
જમણે
સૂત્રોની
ગણતરી કરો.
સંખ્યાઓ
ઉપરિચિહ્ન
તરીકે
લખીએ.
$!A, !B \Sequent !C$ની
સાબિતી
શોધતા હોઈએ,
તો તબક્કા~$0$નું
પરિણામ
$!A^1, !B^2 \Sequent !C^3$
છે. આવા
અનુક્રમણિકાંકવાળા
સિક્વન્ટમાં
ઉપરિચિહ્ન
તરીકે આવતો
સૌથી મોટો~$n$
તેનો
\emph{મહત્તમ
અનુક્રમણિકાંક}
છે (દાખલામાં~$3$).

સિક્વન્ટમાં
પરિમાણકાર
હોય, તો જમણે
$\lexists[x][!A(x)]$
અથવા ડાબે
$\lforall[x][!A(x)]$
ઘટાડતાં કયાં
પદો વાપરવાં
તે પણ જાણવું
પડે. માટે
ફક્ત
અચળો
$\Obj c_0$,
$\Obj c_1$,
$\Obj c_2$,
\dots, અને
$\Gamma$
તથા~$\Delta$માં
આવતાં ફલન-ચિહ્નો
વડે બનેલાં
પદો જરૂરી
છે. આ
બધાં પદો
કોઈ સ્થિર
ક્રમમાં ગોઠવેલાં
છે એમ
માનીએ, જેથી
દાખલા તરીકે
``$\Gamma \Sequent \Delta$માં
ન આવતો
પ્રથમ
અચળ''
કહી શકાય.
અલ્ગોરિધમના
વૃક્ષમાં દરેક
અનુક્રમણિકાંકવાળા
સિક્વન્ટ સાથે
અચળોનો
ગણ સંકળાયેલો
છે. તબક્કા~$0$ના
સિક્વન્ટને
તેમાં આવતા
બધા
અચળોનો
ગણ આપીએ
(એવા અચળો
હોય તો);
એકેય અચળ
ન
હોય તો
$\{\Obj c_0\}$
આપીએ.

માનો કે
તબક્કા~$n$એ
અલ્ગોરિધમને
અનુક્રમણિકાંકવાળા
સિક્વન્ટોનું
વૃક્ષ અને
તેમની સાથે
અચળોના
ગણો મળ્યા
છે. તબક્કા~$n+1$એ
દરેક સૌથી
ઉપરના સિક્વન્ટ
માટે આ
કરીએ:

સૂત્ર~$!A$
સિક્વન્ટની
ડાબી અને
જમણી બંને
બાજુ હોય,
અથવા ડાબે
$\lfalse$
હોય, તો
કંઈ કરતા
નથી: આવા
સિક્વન્ટોની
\Log{G3c}માં
સાબિતીઓ
છે, એટલે
આ સૌથી
ઉપરના સિક્વન્ટ
માટે કામ
પૂરું થયું.

સિક્વન્ટમાં
બિનઆણ્વિક
સૂત્ર
ન હોય,
તો અલ્ગોરિધમ
અટકાવો.
$\Gamma \Sequent \Delta$ની
સાબિતી નથી.

નહિતર સૌથી
નાના
અનુક્રમણિકાંક~$i$
વાળું
બિનઆણ્વિક
સૂત્ર~$!A$
પસંદ કરો.
ત્યારે સંબંધિત
સૌથી ઉપરનો
સિક્વન્ટ
$!A^i, \Pi
\Sequent \Lambda$
અથવા
$\Pi \Sequent \Lambda, !A^i$
છે.
$k$ એ
આ સિક્વન્ટના
મહત્તમ
અનુક્રમણિકાંક
કરતાં એક
વધુ, એટલે
આગલો ન
વપરાયેલો
અનુક્રમણિકાંક,
અને~$C$
તેને સોંપેલો
અચળોનો
ગણ હોય.
\sourcecorrection{OLMD-275}{મૂળ $k$ને હાલનો મહત્તમ અનુક્રમણિકાંક કહે છે, પરંતુ નીચેના નિયમો નવા સૂત્રોને $k$ અને $k+1$ આપે છે. તેમની તાજગી તથા ન્યાયી પસંદગી માટે $k$ હાલના મહત્તમથી એક વધુ હોવો જોઈએ.}

$!A^i$ને
``ઘટાડીએ'':
એટલે~$!A$ના
મુખ્ય કારક
અને~$!A^i$
ડાબે છે
કે જમણે
તે મુજબ
નક્કી થતા
અનુમાન-નિયમથી
નવા સૌથી
ઉપરના સિક્વન્ટો
બનાવીએ.

\begin{enumerate}
  \item $!A \ident !B \land !C$ હોય અને ડાબે આવતું હોય:
  $(!B \land !C)^i, \Pi \Sequent \Lambda$ની ઉપર એક નવો
  સૌથી ઉપરનો સિક્વન્ટ ઉમેરો:
  \[
  \Axiom$!B^k, !C^{k+1}, \Pi \fCenter \Lambda$
  \RightLabel{\LeftR{\land}}
  \UnaryInf$(!B \land !C)^i, \Pi \fCenter \Lambda$
  \DisplayProof
  \]
  નવા સિક્વન્ટને સોંપેલો અચળોનો ગણ~$C$ છે.
  \item $!A \ident !B \land !C$ હોય અને જમણે આવતું હોય:
  $\Pi \Sequent \Lambda, (!B \land !C)^i$ની ઉપર બે નવા
  સૌથી ઉપરના સિક્વન્ટો ઉમેરો:
  \[
  \Axiom$\Pi \fCenter \Lambda, !B^k$
  \Axiom$\Pi \fCenter \Lambda, !C^k$
  \RightLabel{\RightR{\land}}
  \BinaryInf$\Pi \fCenter \Lambda, (!B \land !C)^i$
  \DisplayProof
  \]
  નવા સિક્વન્ટોને સોંપેલો અચળોનો ગણ~$C$ છે.

  \item $!A \ident \lnot !B$, $!A \ident !B \lor !C$,
  અને $!A \ident !B \lif !C$ના કિસ્સા સમાન છે અને
  અભ્યાસ માટે છોડ્યા છે.

  \item $!A \ident \lexists[x][!B(x)]$ હોય અને ડાબે આવતું હોય:
  $\lexists[x][!B(x)]^i, \Pi \Sequent \Lambda$ની ઉપર એક નવો
  સૌથી ઉપરનો સિક્વન્ટ ઉમેરો:
  \[
  \Axiom$!B(c)^k, \Pi \fCenter \Lambda$
  \RightLabel{\LeftR{\lexists}}
  \UnaryInf$\lexists[x][!B(x)]^i, \Pi \fCenter \Lambda$
  \DisplayProof
  \]
  અહીં $c$ એ~$C$માં ન હોય એવો પ્રથમ અચળ છે.
  નવા સિક્વન્ટને સોંપેલો અચળોનો ગણ $C \cup \{c\}$ છે.
  \item $!A \ident \lexists[x][!B(x)]$ હોય અને જમણે આવતું હોય:
  $\Pi \Sequent \Lambda, \lexists[x][!B(x)]^i$ની ઉપર એક નવો
  સૌથી ઉપરનો સિક્વન્ટ ઉમેરો:
  \sourcecorrection{OLMD-277}{મૂળ ગદ્યમાં આ નિષ્કર્ષના અસ્તિત્વ-સૂત્રનો અનુક્રમણિકાંક છૂટી ગયો છે; નીચેના નિયમવૃક્ષ મુજબ અહીં $i$ ઉમેર્યો છે.}
  \[
  \Axiom$\Pi \fCenter \Lambda, \lexists[x][!B(x)]^k, !B(t)^{k+1}$
  \RightLabel{\RightR{\lexists}}
  \UnaryInf$\Pi \fCenter \Lambda, \lexists[x][!B(x)]^i$
  \DisplayProof
  \]
  અહીં $t$ એ ફક્ત~$C$માંના અચળો ધરાવતું એવું
  પ્રથમ પદ છે જે આ સિક્વન્ટની નીચે જમણી બાજુના
  $\lexists[x][!B(x)]$ના ઘટાડામાં અગાઉ વપરાયું નથી
  (અથવા આવાં બધાં પદો વપરાઈ ગયાં હોય તો $\Obj c_0$).
  નવા સિક્વન્ટને સોંપેલો અચળોનો ગણ $C$ છે.
  \sourcecorrection{OLMD-276}{મૂળ નિયમવૃક્ષમાં જમણી બાજુના અસ્તિત્વ-સૂત્ર માટે ભૂલથી ડાબો નિયમ \LeftR{\lexists} લખાયો છે; તેને યોગ્ય \RightR{\lexists} કર્યો છે.}
  \item $!A \ident \lforall[x][!B(x)]$ના કિસ્સા સમાન છે અને
  અભ્યાસ માટે છોડ્યા છે.
\end{enumerate}

તબક્કા~$n+1$એ કોઈ પણ સૌથી ઉપરના સિક્વન્ટમાં નવો
સિક્વન્ટ ઉમેરાયો ન હોય, તો અલ્ગોરિધમ સફળતાપૂર્વક
થોભે છે. આ કિસ્સામાં તબક્કા~$n+1$ના વૃક્ષમાંથી
બધા અનુક્રમણિકાંક ભૂંસીને $\Gamma \Sequent \Delta$ની
સાબિતી મળે છે. સૌથી ઉપરના બધા સિક્વન્ટો
$!A, \Pi \Sequent \Lambda, !A$ અથવા
$\lfalse, \Pi \Sequent \Lambda$ સ્વરૂપના છે.
બધાં અનુમાનો \Log{G3c}નાં યોગ્ય અનુમાનો છે.
વિધાનતર્કીય અનુમાનો માટે આ સ્પષ્ટ છે.
નોંધો કે દરેક પગલે સિક્વન્ટને સોંપેલા
અચળોના ગણ~$C$માં તે સિક્વન્ટમાં આવતા
બધા અચળો હોય છે. તેથી
\LeftR{\lexists} અને \RightR{\lforall} વાપરીને
કરેલા ઘટાડામાં સ્વનિર્ધારક ચલની શરત પૂરી થાય છે:
સ્વનિર્ધારક ચલ~$c$ નિષ્કર્ષને સોંપેલા ગણ~$C$માં
ન હોય એવો અચળ હતો, એટલે નિષ્કર્ષમાં
એટલે $c$ નિષ્કર્ષમાં
પણ આવતો નહોતો. આમ:

\begin{prop}
\Log{G3c} માટેનો સાબિતી-શોધ અલ્ગોરિધમ સાચો છે:
તે સફળતાપૂર્વક થોભે, તો શરૂઆતના
સિક્વન્ટ~$\Gamma \Sequent \Delta$ની સાબિતી છે.
\end{prop}
% END OLP-0683
\endgroup

\begingroup
\makeatletter\def\ol@sectioncs{section}\makeatother

% BEGIN OLP-0679
\olfileid{pt}{ps}{cpl}

\olsection{પૂર્ણતા}
\phantomsection\label{guunit:OLP-0679}


હવે સાબિતી-શોધ
અલ્ગોરિધમ અચૂક
છે તે બતાવીએ:
જો તે અધવચ્ચે
અટકી જાય અથવા
કદી પૂર્ણ ન
થાય, તો શરૂઆતના
સિક્વન્ટ
$\Gamma \Sequent \Delta$ની
કોઈ સાબિતી
નથી. તે માટે
સંરચના~$\Struct M$
એવી બનાવીએ કે
દરેક~$!A \in \Gamma$
માટે~$\Sat{M}{!A}$
અને દરેક
$!A \in \Delta$
માટે~$\Sat/{M}{!A}$.
બીજા શબ્દોમાં,
$\Gamma$નાં બધાં
સૂત્રો
અને~$\Delta$નાં
બધાં સૂત્રો
અનુક્રમે
$\Struct M$માં
સત્ય અને અસત્ય
છે. એટલે
$\Gamma \Sequent \Delta$
પ્રામાણ્ય નથી.
\Log{G3c} યથાર્થ
હોવાથી
$\Gamma \Sequent \Delta$ની
સાબિતી નથી.
\sourcecorrection{OLMD-269}{મૂળમાં $\Delta$ના દરેક સભ્ય માટેની શરતમાં $! \in \Delta$ લખાયું છે; અહીં સૂત્રચલ $!A \in \Delta$ હોવો જોઈએ.}

$\Struct{M}$ને
સૂત્રોની
જોડી~$\Theta, \Xi$
અને
અચળોના
ગણ~$C$માંથી
વ્યાખ્યાયિત કરીએ.
આ~$\Theta, \Xi$
અને~$C$
સાબિતી-શોધનો
અસફળ (અટકેલી
અથવા અપૂર્ણ
રહેલી) દોડની
\emph{નિષ્ફળતા-શાખા}
માંથી મળે છે.
નિષ્ફળતા-શાખા
સિક્વન્ટોનો
ક્રમ~$\Pi_n \Sequent
\Lambda_n$
અને અનુરૂપ
અચળોના
ગણ~$C_n$નો
બને છે, જ્યાં
$\Pi_0 \Sequent \Lambda_0$
એ અંતિમ
સિક્વન્ટ
$\Gamma \Sequent
\Delta$ છે,
અને
$\Pi_{n+1} \Sequent \Lambda_{n+1}$
એ~$\Pi_n
\Sequent \Lambda_n$નું
તબક્કા~$n+1$એ
બનેલું એવું
આધારવિધાન છે
જેની સાબિતી
અલ્ગોરિધમ શોધતું
નથી. શાખા
ચાલુ હોય તેવા
દરેક~$n$
માટે આવું
આધારવિધાન હોય
જ છે, નહીં તો
અલ્ગોરિધમને
સાબિતી
મળી ગઈ હોત.
$C_n$ એ
$\Pi_n
\Sequent \Lambda_n$
સાથે સંકળાયેલો
અચળોનો
ગણ છે.
\sourcecorrection{OLMD-270}{મૂળનું છેલ્લું સિક્વન્ટ $\Pi_n \Sequent \Delta_n$ લખાયું છે; નિષ્ફળતા-શાખાનો જમણો ભાગ સમગ્ર વ્યાખ્યામાં $\Lambda_n$ છે.}
\sourcecorrection{OLMD-271}{મૂળ કહે છે કે દરેક $n$ માટે આગળનું આધારવિધાન હોવું જ જોઈએ; સીમિત નિષ્ફળતા-શાખા આણ્વિક સિક્વન્ટે પૂરી થાય છે. આગળનું આધારવિધાન ફક્ત શાખા ચાલુ હોય એવા તબક્કે જરૂરી છે.}

અલ્ગોરિધમ માત્ર
આણ્વિક
સૂત્રોવાળા
સૌથી ઉપરના
સિક્વન્ટ પર
અટકી જાય,
તો ત્યાં પૂરી
થતી શાખા
નિષ્ફળતા-શાખા
છે. જો
અલ્ગોરિધમ
કદી પૂર્ણ ન
થાય, તો તે
વધુ ને વધુ
મોટાં વૃક્ષો
બનાવતું રહે.
આને અનંત
વૃક્ષ વધતું
હોય તેમ
વિચારી શકાય
(અલ્ગોરિધમનાં
અનંત તબક્કા
પછીનું વૃક્ષ).
આ અનંત વૃક્ષ
સીમિત શાખાવાળું
છે: દરેક ગાંઠને
માત્ર એક કે
બે સંતાન
સિક્વન્ટ છે,
જે તેની તરત
ઉપર છે.
K\H{o}nigના
ઉપપાદ્ય પ્રમાણે
દરેક અનંત,
સીમિત શાખાવાળા
વૃક્ષમાં અનંત
શાખા હોય છે.
શોધ અલ્ગોરિધમ
કાયમ ચાલે
ત્યારે આવી
અનંત શાખા
નિષ્ફળતા-શાખા
છે.

$\Pi_n \Sequent \Lambda_n$
અને~$C_n$નો
ક્રમ નિષ્ફળતા-શાખા
હોય, તો
$\Theta = \bigcup_n \Pi_n$
અને
$\Xi = \bigcup_n \Lambda_n$
લઈએ
(સૂત્રોના
અનુક્રમણિકા અને
પુનરાવર્તિત
આવૃત્તિઓ દૂર
કરીને), અને
$C =
\bigcup_n C_n$
લઈએ.

હવે~$\Struct{M}$
વ્યાખ્યાયિત કરીએ.
\begin{enumerate}
\item પ્રદિશ
$\Domain{M}$ એ
$C$ અને
$\Gamma \Sequent \Delta$માં
આવતાં
ફલનો
હોય તો તેમની
મદદથી, પણ
ચલો વિના,
બનતા પદોનો
ગણ છે.
\item $c \in \Domain{M}$
હોય તો
$\Assign{c}{M} = c$.
\item $f$ એ
$\Gamma \Sequent \Delta$માંનું
$m$-સ્થાની
ફલન
હોય, અને
$t_1$,
\dots,
$t_m \in \Domain{M}$
હોય, તો
$\Assign{f}{M}(t_1, \dots, t_m) =
f(t_1, \dots, t_m)$.
\sourcecorrection{OLMD-272}{મૂળમાં $m$-સ્થાની વિધેયના મૂલ્ય માટે છેલ્લે $f(t_1,\dots,t_n)$ લખાયું છે; ઇનપુટ અને વિધેયની સ્થાનસંખ્યા મુજબ અંતિમ પદ $t_m$ છે.}
\item $R$ એ
$\Gamma \Sequent
\Delta$માંનું
$m$-સ્થાની
વિધેય
હોય, તો
$\Assign{R}{M} = \Setabs{\tuple{t_1, \dots, t_m} \in
\Domain{M}^m}{R(t_1, \dots, t_m) \in \Theta}$.
\sourcecorrection{OLMD-273}{મૂળમાં $m$-સ્થાની સંબંધ માટે પ્રદિશની ઘાત અને $R$ના અંતિમ પદમાં $n$ છે; બંને જગ્યાએ સ્થાનસંખ્યા $m$ જ હોવી જોઈએ.}
\end{enumerate}
$\Struct{M}$ એક
\emph{પદ-નિદર્શ}
છે: તેનો પ્રદિશ
પદોનો ગણ છે,
અને
અચળો તથા
ફલનોનું
અર્થઘટન એવી
રીતે થાય છે
કે
$\Value{t}{M} = t$.
$\Theta$નાં
આણ્વિક
સૂત્રો
વડે
$\Assign{R}{M}$
એવી રીતે
વ્યાખ્યાયિત
થાય છે કે
$\Sat{M}{R(t_1, \dots, t_n)}$
ત્યારે અને
ત્યારે જ,
જ્યારે
$R(t_1, \dots, t_n)
\in \Theta$.

\begin{lem}\ollabel{lem:truth}
દરેક
$!A \in \Theta \cup \Xi$
માટે:
\begin{enumerate}
  \item $!A \in \Theta$
  હોય તો
  $\Sat{M}{!A}$.
  \item $!A \in \Xi$
  હોય તો
  $\Sat/{M}{!A}$.
\end{enumerate}
\end{lem}

\begin{proof}
  $\depth{!A}$
  પર આગમન
  કરીએ.

પહેલાં માનો કે
$!A$ આણ્વિક છે,
એટલે કાં તો
$!A \ident \lfalse$
અથવા
$!A
\ident R(t_1, \dots, t_m)$.

$\Pi_n \Sequent \Lambda_n$
નિષ્ફળતા-શાખા
હોવાથી કોઈ
$\Pi_n$માં
$\lfalse$ નથી
(નહિતર
$\Pi_n \Sequent \Lambda_n$
સ્વયંસિદ્ધ
થાય).
$\Struct{M}$ની
વ્યાખ્યા પ્રમાણે
$\Sat{M}{R(t_1, \dots, t_m)}$
ત્યારે અને
ત્યારે જ,
જ્યારે
$R(t_1, \dots, t_m) \in
\Theta$.
બંને વાતો
મળીને આપે છે:
$!A \in
\Theta$ હોય,
તો
$\Sat{M}{!A}$.

$!A$ આણ્વિક
હોય અને
$!A \in \Pi_n$
હોય, તો
$!A \in \Pi_{n+1}$;
તેવી જ રીતે
$!A \in \Lambda_n$
હોય તો
$!A \in \Lambda_{n+1}$.
(સાબિતી-શોધ
અલ્ગોરિધમ
ઉત્તરાધિકારી
સિક્વન્ટો જે
રીતે બનાવે
છે તેમાંથી
આ ફલિત થાય
છે.) તેથી
$!A \in \Theta \cap \Xi$
હોય, તો કોઈ
$n$ માટે
$!A \in \Pi_n \cap
\Lambda_n$.
એટલે
$\Pi_n \Sequent \Lambda_n$
સ્વયંસિદ્ધ
થાય, પરંતુ
નિષ્ફળતા-શાખામાં
સ્વયંસિદ્ધ ન
હોઈ શકે.
આથી આણ્વિક
$!A$ માટે
$!A \notin \Theta
\cap \Xi$,
અને પરિણામે
$!A \in \Xi$
હોય તો
$!A \notin \Theta$.
$\Struct{M}$ની
વ્યાખ્યાથી
આનો અર્થ
એ કે
$!A \in \Xi$
હોય, તો
$\Sat/{M}{!A}$.

આગમન-પગલામાં
માનો કે~$!A$
કરતાં ઓછી
ઊંડાઈ ધરાવતા
બધા સૂત્રો
માટે (1)
અને (2)
સાચાં છે.
$!A$ માટે
(1) અને
(2) કિસ્સા
જુદા પાડીને
સાબિત કરીએ.

પહેલાં નોંધો કે
$!A^i$
સિક્વન્ટ
$\Pi_n \Sequent \Lambda_n$માં
હોય, તો~$i$
તે જ
સિક્વન્ટ~$\Pi_n \Sequent \Lambda_n$નો
સૌથી
નાનો અનુક્રમણિકાંક
ન હોય ત્યાં
સુધી~$!A^i$
પછીના
$\Pi_{n+1} \Sequent \Lambda_{n+1}$
માં પણ રહે
છે. દરેક તબક્કે
ઉમેરાતા સૌથી
ઉપરના સિક્વન્ટોમાં
સૌથી નાનો
અનુક્રમણિકાંક
વધે છે.
તેથી આખરે
એવું
$\Pi_n \Sequent \Lambda_n$
આવે છે જેમાં
$!A^i$ હોય
અને~$i$
સૌથી નાનો
હોય. તબક્કા
$n+1$એ
અલ્ગોરિધમ~$!A^i$ને
ઘટાડે છે.
એટલે
$!A \in \Theta$
હોય, તો કોઈ
$n$ માટે
$\Pi_n = \Pi_n',
!A^i$
અને~$i$
સૌથી નાનો
અનુક્રમણિકાંક
હોય. તેવી જ
રીતે
$!A \in \Xi$
હોય, તો કોઈ
$n$ માટે
$\Lambda_n = \Lambda_n', !A^i$
અને~$i$
સૌથી નાનો
અનુક્રમણિકાંક
હોય.

\begin{enumerate}
  \item $!A \in \Theta$
  અને
  $!A \ident !B \land !C$
  હોય. ત્યારે
  કોઈ~$n$ માટે
  $\Pi_n = \Pi_n', (!B \land !C)^i$.
  $\Pi_{n+1} \Sequent
  \Lambda_{n+1}$
  એ~\LeftR{\land}નું
  અનુરૂપ
  આધારવિધાન
  છે, એટલે
  $\Pi_{n+1} = \Pi_n', !B^k, !C^{k+1}$.
  તેથી
  $!B, !C \in
  \Theta$.
  આગમન-પરિકલ્પનાથી
  $\Sat{M}{!B}$
  અને
  $\Sat{M}{!C}$.
  $\Sat{M}{}$ની
  વ્યાખ્યા પ્રમાણે
  $\Sat{M}{!B \land !C}$.

  \item $!A \in \Xi$
  અને
  $!A \ident !B \land !C$
  હોય. ત્યારે
  કોઈ~$n$ માટે
  $\Lambda_n = \Lambda_n', (!B \land !C)^i$.
  $\Pi_{n+1} \Sequent
  \Lambda_{n+1}$
  એ નિષ્કર્ષ
  $\Pi_n \Sequent \Lambda'_n, (!B \land !C)^i$
  ધરાવતા
  \RightR{\land}ના
  બે આધારવિધાનોમાંથી
  એક છે, એટલે
  $\Lambda_{n+1} = \Lambda_n', !B^k$
  અથવા
  $\Lambda_{n+1} = \Lambda_n',
  !C^k$.
  તેથી
  $!B \in \Xi$
  અથવા
  $!C \in \Xi$.
  આગમન-પરિકલ્પનાથી
  $\Sat/{M}{!B}$
  અથવા
  $\Sat/{M}{!C}$.
  $\Sat{M}{}$ની
  વ્યાખ્યાથી
  $\Sat/{M}{!B \land !C}$.
  \sourcecorrection{OLMD-274}{મૂળના જમણા સંયોજન-કિસ્સામાં ઘટાડાતા મુખ્ય સૂત્રના બે ઉલ્લેખમાં અનુક્રમણિકાંક $i$ છૂટ્યો છે; અલ્ગોરિધમ અનુક્રમણિકાંકવાળા સિક્વન્ટો પર ચાલે છે.}
  
  \item
  $!A \ident \lnot !B$,
  $!A \ident !B \lor !C$
  અને
  $!A \ident !B \lif !C$
  કિસ્સા સમાન
  છે અને
  કસરત તરીકે
  મૂક્યાં છે.

  \item $!A \in \Theta$
  અને
  $!A \ident \lexists[x][!B(x)]$
  હોય. ત્યારે
  કોઈ~$n$ માટે
  $\Pi_n = \Pi_n', \lexists[x][!B(x)]^i$.
  $\Pi_{n+1}
  \Sequent \Lambda_{n+1}$
  એ~\LeftR{\lexists}નું
  અનુરૂપ
  આધારવિધાન
  છે, એટલે
  કોઈ~$c$ માટે
  $\Pi_{n+1} = \Pi_n', !B(c)^k$.
  તેથી
  $!B(c) \in \Theta$
  અને
  $c \in C$.
  આગમન-પરિકલ્પનાથી
  $\Sat{M}{!B(c)}$.
  $\Sat{M}{}$ની
  વ્યાખ્યાથી
  $\Sat{M}{\lexists[x][!B(x)]}$.

  \item $!A \in \Xi$
  અને
  $!A \ident \lexists[x][!B(x)]$
  હોય. ત્યારે
  કોઈ~$n$ માટે
  $\Lambda_n = \Lambda_n', \lexists[x][!B(x)]^i$.
  જ્યારે જ્યારે
  $\lexists[x][!B(x)]$
  ઘટાડાય છે,
  ત્યારે
  $\lexists[x][!B(x)]$
  નવા અનુક્રમણિકાંક
  સાથે આધાર
  સિક્વન્ટની
  જમણી બાજુએ
  રહે છે.
  તેથી નિષ્ફળતા-શાખામાં
  તે જમણી બાજુ
  આવે તો
  $\lexists[x][!B(x)]$
  અનંત વાર
  ઘટાડાય છે.
  દરેક
  $t \in
  \Domain{M}$
  ક્યારેક
  નિષ્ફળતા-શાખામાં
  ``ફક્ત~$C_n$નાં
  અચળો ધરાવતું,
  જમણી બાજુના
  $\lexists[x][!B(x)]$ના
  ઘટાડામાં હજી
  ન વપરાયેલું
  પ્રથમ પદ''
  બને છે.
  તેથી દરેક
  $t \in \Domain{M}$
  માટે કોઈ~$n$
  પર
  $!B(t) \in \Lambda_n$,
  અને પરિણામે
  $!B(t) \in \Xi$.
  આગમન-પરિકલ્પનાથી
  દરેક
  $t \in
  \Domain{M}$
  માટે
  $\Sat/{M}{!B(t)}$.
  $\Sat{M}{}$ની
  વ્યાખ્યાથી
  $\Sat/{M}{\lexists[x][!B(x)]}$.

  \item
  $!A \ident \lforall[x][!B(x)]$
  કિસ્સા સમાન
  છે અને
  કસરત તરીકે
  મૂક્યાં છે.
\end{enumerate}
\end{proof}

\begin{cor}\ollabel{cor:complete}
\Log{G3c}નો
સાબિતી-શોધ
અલ્ગોરિધમ પૂર્ણ
છે: તે અધવચ્ચે
અટકે અથવા
કદી પૂર્ણ ન
થાય, તો શરૂઆતનો
સિક્વન્ટ
$\Gamma \Sequent \Delta$
પ્રામાણ્ય નથી
અને તેથી
તેની સાબિતી
નથી.
\end{cor}

\begin{proof}
  અલ્ગોરિધમ
  અટકે અથવા
  કદી પૂર્ણ ન
  થાય, તો
  નિષ્ફળતા-શાખા
  મળે છે;
  તેમાંથી
  $\Struct{M}$
  વ્યાખ્યાયિત
  કરી શકાય.
  \olref{lem:truth}
  મુજબ દરેક
  $!A \in \Gamma$
  માટે
  $\Sat{M}{!A}$
  અને દરેક
  $!A \in \Delta$
  માટે
  $\Sat/{M}{!A}$.
  (યાદ રાખો:
  $\Pi_0 = \Gamma$
  અને
  $\Lambda_0 =
  \Delta$,
  તેથી
  $\Gamma \subseteq \Theta$
  અને
  $\Delta \subseteq \Xi$.)
  એટલે
  $\Gamma \Sequent \Delta$
  પ્રામાણ્ય નથી.
  \Log{G3c}ની
  યથાર્થતા મુજબ
  $\Gamma \Sequent \Delta$ની
  સાબિતી
  નથી.
\end{proof}

\begin{cor}
\Log{G3c}
પૂર્ણ છે: જો
$\Gamma \Sequent \Delta$
પ્રામાણ્ય હોય, તો
$\Log{G3c} \Proves \Gamma \Sequent \Delta$.
\end{cor}

\begin{proof}
  પ્રતિવિધાન
  સાબિત કરીએ.
  માનો કે
  $\Log{G3c} \Proves/ \Gamma
  \Sequent \Delta$.
  શોધવાની
  સાબિતી
  નથી, તેથી
  સાબિતી-શોધ
  અલ્ગોરિધમ
  અટકે અથવા
  કદી પૂર્ણ ન
  થાય. પછી
  \olref{cor:complete}
  મુજબ
  $\Gamma \Sequent \Delta$
  પ્રામાણ્ય નથી.
\end{proof}
% END OLP-0679
\endgroup

\begingroup
\makeatletter\def\ol@sectioncs{section}\makeatother

% BEGIN OLP-0684
\olfileid{pt}{ps}{tab}

\olsection{કોષ્ટક-વૃક્ષો}
\phantomsection\label{guunit:OLP-0684}


કોષ્ટક-વૃક્ષ તંત્રો સિક્વન્ટ-કલન સાથે નજીકનો સંબંધ ધરાવતાં
સાબિતી તંત્રો છે. હાથથી અને સૉફ્ટવેરમાં બંને રીતે
સાબિતી-શોધ માટે તે ખાસ અનુકૂળ છે. કોષ્ટક-વૃક્ષમાં
સિક્વન્ટોનું વૃક્ષ બનાવવાને બદલે સાબિતી એ
સૂત્રોનું વૃક્ષ છે. જોકે સ્વાભાવિક નિગમનથી
વિપરીત, તેના નિયમો સિક્વન્ટ-કલનના નિયમોને અનુસરે છે.
કોષ્ટક-વૃક્ષની સાબિતી હકીકતમાં નિદર્શની સ્પષ્ટ શોધ
છે. તેથી તેને ક્યારેક \emph{અર્થલક્ષી} કોષ્ટક-વૃક્ષ
અથવા \emph{સત્ય-વૃક્ષ} પણ કહે છે.

\emph{ચિહ્નિત} કોષ્ટક-વૃક્ષ ચિહ્નિત સૂત્રોનું
વૃક્ષ છે; તે $\sFmla{\True}{!A}$ અથવા
$\sFmla{\False}{!A}$ સ્વરૂપનાં હોઈ શકે.
વૃક્ષો નીચેની દિશામાં વધે છે. શરૂઆતમાં જે
સૂત્રો હોય છે તે આપણે શું સાબિત કરવું છે
તે નક્કી કરે છે. દાખલા તરીકે,
$!A, !B \Sequent !C, !D$ સિક્વન્ટને સાબિત
કરવા માટે તે
$\sFmla{\True}{!A}, \sFmla{\True}{!B},
\sFmla{\False}{!C}, \sFmla{\False}{!D}$થી
શરૂ થાય.
\sourcecorrection{OLMD-278}{મૂળ ઉદાહરણમાં $!D$ને ભૂલથી સત્ય-ચિહ્ન અપાયું છે; તરત પછીનું ગદ્ય $!C$ અને $!D$ બંને અસત્ય રાખે છે, તેથી અહીં અસત્ય-ચિહ્ન કર્યું છે.}
$!A$ અને $!B$ને સત્ય તથા $!C$ અને $!D$ને
અસત્ય કરતી સંરચના એ
$!A, !B \Sequent !C, !D$ અપ્રામાણ્ય છે તે
દર્શાવતી સંરચના છે.

કોષ્ટક-વૃક્ષના પાન, એટલે સૌથી નીચેના શિરોબિંદુ,
શાખા-વિસ્તારના નિયમોથી આગળ વધારી શકાય છે. આ
નિયમો એ જ શાખામાં નવા શિરોબિંદુઓ ઉમેરે છે અથવા
નીચે બે સમકક્ષ શિરોબિંદુઓ ઉમેરી શાખાને બે ભાગે
વહેંચે છે. એક જ શાખામાં પાનની ઉપરના કોઈ પણ
ચિહ્નિત સૂત્ર પર નિયમ લાગુ થઈ શકે છે.
શાસ્ત્રીય કોષ્ટક-વૃક્ષ તંત્ર \Log{Tc}ના
શાખા-વિસ્તારના નિયમો \olref{tab:Tc}માં આપ્યા છે.

\begingroup
\makeatletter\def\ol@sectioncs{section}\makeatother

% BEGIN OLP-0682
\phantomsection\label{guunit:OLP-0682}


\begin{table}
\begin{center}
\begin{tabular}{@{}cc@{}}
\AxiomC{\sFmla{\True}{\lnot !A}}
\RightLabel{\TRule{\True}{\lnot}}
\UnaryInfC{\sFmla{\False}{!A}}
\DisplayProof
&
\AxiomC{\sFmla{\False}{\lnot !A}}
\RightLabel{\TRule{\False}{\lnot}}
\UnaryInfC{\sFmla{\True}{!A}}
\DisplayProof
\\[6ex]
\AxiomC{\sFmla{\True}{!A \land !B}}
\RightLabel{\TRule{\True}{\land}}
\UnaryInfC{\sFmla{\True}{!A}}
\noLine
\UnaryInfC{\sFmla{\True}{!B}}
\DisplayProof
&
\AxiomC{\sFmla{\False}{!A \land !B}}
\RightLabel{\TRule{\False}{\land}}
\UnaryInfC{$\sFmla{\False}{!A} \quad \mid \quad \sFmla{\False}{!B}$}
\DisplayProof
\\[6ex]
\AxiomC{\sFmla{\True}{!A \lor !B}}
\RightLabel{\TRule{\True}{\lor}}
\UnaryInfC{$\sFmla{\True}{!A} \quad \mid \quad \sFmla{\True}{!B}$}
\DisplayProof
&
\AxiomC{\sFmla{\False}{!A \lor !B}}
\RightLabel{\TRule{\False}{\lor}}
\UnaryInfC{\sFmla{\False}{!A}}
\noLine
\UnaryInfC{\sFmla{\False}{!B}}
\DisplayProof
\\[6ex]
\AxiomC{\sFmla{\True}{!A \lif !B}}
\RightLabel{\TRule{\True}{\lif}}
\UnaryInfC{$\sFmla{\False}{!A} \quad \mid \quad \sFmla{\True}{!B}$}
\DisplayProof
&
\AxiomC{\sFmla{\False}{!A \lif !B}}
\RightLabel{\TRule{\False}{\lif}}
\UnaryInfC{\sFmla{\True}{!A}}
\noLine
\UnaryInfC{\sFmla{\False}{!B}}
\DisplayProof
\\[6ex]
\AxiomC{\sFmla{\True}{\lforall[x][!A(x)]}}
\RightLabel{\TRule{\True}{\forall}}
\UnaryInfC{\sFmla{\True}{!A(t)}}
\DisplayProof
&
\AxiomC{\sFmla{\False}{\lforall[x][!A(x)]}}
\RightLabel{\TRule{\False}{\lforall}}
\UnaryInfC{\sFmla{\False}{!A(a)}}
\DisplayProof
\\[6ex]
\AxiomC{\sFmla{\True}{\lexists[x][!A(x)]}}
\RightLabel{\TRule{\True}{\lexists}}
\UnaryInfC{\sFmla{\True}{!A(a)}}
\DisplayProof
&
\AxiomC{\sFmla{\False}{\lexists[x][!A(x)]}}
\RightLabel{\TRule{\False}{\lexists}}
\UnaryInfC{\sFmla{\False}{!A(t)}}
\DisplayProof
\end{tabular}
\end{center}
\caption{\Log{Tc}ના નિયમો. $\TRule{\False}{\lforall}$ અને
$\TRule{\True}{\lexists}$માં અચળ~$a$ ઉપરની
શાખામાં આવતો ન હોવો જોઈએ; એટલે તે શાખા માટે નવો હોવો જોઈએ.}
\ollabel{tab:Tc}
\end{table}
% END OLP-0682
\endgroup


જોઈ શકાય છે કે આ નિયમો \Log{G3c}ના નિયમો ઊલટા
મૂકેલા છે. $\True$ અને $\False$ ચિહ્નો સૂચવે છે કે
સૂત્ર સિક્વન્ટની અનુક્રમે ડાબી અને જમણી
બાજુએ મુખ્ય સૂત્ર છે.

કોષ્ટક-વૃક્ષની શાખામાં
$\sFmla{\True}{!A}$ અને $\sFmla{\False}{!A}$
બંને હોય, તો તે \emph{બંધ} છે. દરેક શાખા બંધ
હોય, તો આખું કોષ્ટક-વૃક્ષ બંધ કહેવાય.
દાખલા તરીકે,
$!A \lor (!B \land !C) \Sequent (!A \lor !B) \land
(!A \lor !C)$ માટેનું આ બંધ કોષ્ટક-વૃક્ષ છે;
તે $\sFmla{\True}{!A \lor (!B \land !C)}$
અને $\sFmla{\False}{(!A \lor !B) \land (!A \lor !C)}$થી
શરૂ થાય છે:
\sourcecorrection{OLMD-279}{મૂળમાં clause tableau લખાયું છે, પરંતુ દર્શાવેલા વૃક્ષની બધી શાખાઓ બંધ છે; અહીં બંધ કોષ્ટક-વૃક્ષ કર્યું છે.}
\begin{oltableau}
  [\sFmla{\True}{\formula{A} \lor (\formula{B} \land \formula{C})},
    just=\TAss, checked
    [\sFmla{\False}{(\formula{A} \lor \formula{B}) \land
        (\formula{A} \lor \formula{C})}, just=\TAss, checked
      [\sFmla{\True}{\formula{A}}, just={\TRule{\True}{\lor}[1]}
        [\sFmla{\False}{\formula{A} \lor \formula{B}},
          just={\TRule{\False}{\land}[2]}, checked
          [\sFmla{\False}{\formula{A}},
            just={\TRule{\False}{\lor}[4]}
            [\sFmla{\False}{\formula{B}},
              just={\TRule{\False}{\lor}[4]}, close]
          ]
        ]
        [\sFmla{\False}{\formula{A} \lor \formula{C}},
          just={\TRule{\False}{\land}[2]}, checked
          [\sFmla{\False}{\formula{A}},
            just={\TRule{\False}{\lor}[4]}
            [\sFmla{\False}{\formula{C}},
              just={\TRule{\False}{\lor}[4]}, close]
          ]
        ]
      ]
      [\sFmla{\True}{\formula{B} \land \formula{C}},
        just={\TRule{\True}{\lor}[1]}, checked
        [\sFmla{\False}{\formula{A} \lor \formula{B}},
          just={\TRule{\False}{\land}[2]}, checked
          [\sFmla{\True}{\formula{B}},
            just={\TRule{\True}{\land}[3]}, move by=2
            [\sFmla{\True}{\formula{C}},
              just={\TRule{\True}{\land}[3]}
              [\sFmla{\False}{\formula{A}},
                just={\TRule{\False}{\lor}[4]}
                [\sFmla{\False}{\formula{B}},
                  just={\TRule{\False}{\lor}[4]},close
                ]
              ]
            ]
          ]
        ]
        [\sFmla{\False}{\formula{A} \lor \formula{C}},
          just={\TRule{\False}{\land}[2]}, checked
          [\sFmla{\True}{\formula{B}},
            just={\TRule{\True}{\land}[3]}, move by=2
              [\sFmla{\True}{\formula{C}},
                just={\TRule{\True}{\land}[3]}
                [\sFmla{\False}{\formula{A}},
                  just={\TRule{\False}{\lor}[4]}
                  [\sFmla{\False}{\formula{C}},
                  just={\TRule{\False}{\lor}[4]},close
                ]
              ]
            ]
          ]
        ]
      ]
    ]
  ]
\end{oltableau}

ખરાની નિશાનીઓ સૂચવે છે કે તેમની નીચેની બધી
શાખાઓમાં સંબંધિત સૂત્ર પર નિયમ લાગુ
કરી દેવાયો છે. $\otimes$ ચિહ્ન બંધ શાખા સૂચવે છે.
નિયમોની પાછળની પંક્તિ-સંખ્યાઓ કયા ચિહ્નિત
સૂત્રો પર નિયમ લાગુ કર્યો છે તે બતાવે છે,
અર્થાત્ એ જ શાખામાં દર્શાવેલી પંક્તિ પરના
સૂત્ર પર.

કોષ્ટક-વૃક્ષમાં સાબિતી-શોધ સરળ છે. કોષ્ટક-વૃક્ષ
તબક્કાવાર બનાવીએ. તબક્કા~$0$એ શરૂઆતનાં
સૂત્રો સૌથી ઉપર લખો. તબક્કા~$n+1$એ
દરેક પાન માટે આ કરો: જે ચિહ્નિત સૂત્ર પર
હજી ખરાની નિશાની નથી તેમાંથી સૌથી ઉપરનું
પ્રથમ સૂત્ર શોધો. તેના પર શાખા-વિસ્તારનો નિયમ
લાગુ કરો. એ ચિહ્નિત સૂત્ર ધરાવતી દરેક
શાખામાં નિયમ લાગુ થઈ જાય, પછી તેના પર ખરાની
નિશાની કરો. (ઉપરનું ઉદાહરણ આ અલ્ગોરિધમથી
બનાવ્યું છે.)

$\sFmla{\True}{\lforall}$ અને
$\sFmla{\False}{\lexists}$ સ્વરૂપનાં ચિહ્નિત
સૂત્રો પર કદી ખરાની નિશાની થતી નથી,
એટલે તેમની વિશેષ સંભાળ લેવી પડે. તે હાજર હોય,
તો $\TRule{\True}{\forall}$ અને
$\TRule{\False}{\lexists}$ નિયમો દરેક લાગુ પડતા
પદ માટે આખરે લાગુ થાય તેની ખાતરી કરવી પડે.
દાખલા તરીકે, દરેક તબક્કે દરેક શાખાના આવાં બધાં
સૂત્રો પર આ નિયમો લાગુ કરી શકાય છે
(આ સ્વરૂપનું ન હોય એવા સૌથી ઉપરના ચિહ્નિત
સૂત્ર પર કામ કર્યા પછી). વાપરવાનું પદ
$t$ એ શાખામાં પહેલેથી આવતાં અચળો
(અથવા કોઈ અચળ ન આવતો હોય તો
$\Obj c_0$) અને શરૂઆતનાં સૂત્રોમાં
આવતાં ફલનો વડે બની શકે એવું પ્રથમ
પદ~$t$ છે, જેના માટે સંબંધિત ચિહ્નિત
સૂત્ર $\sFmla{\True}{!B(t)}$ અથવા
$\sFmla{\False}{!B(t)}$ શાખામાં અગાઉ આવતું નથી.

આ સાબિતી-શોધથી બંધ કોષ્ટક-વૃક્ષ ન બને, તો
કાં તો તેમાં એવી સીમિત ખુલ્લી શાખા હોય છે
જેના બધા બિનઆણ્વિક સૂત્રો પર ખરાની
નિશાની છે, અથવા શોધ અનંત પણ સીમિત-શાખાવાળું
વૃક્ષ બનાવે છે, જેમાં અનંત શાખા હોવી જ પડે.
\olref[cpl]{sec}ની જેમ, આપણે એવી
સંરચના~$\Struct{M}$ વ્યાખ્યાયિત કરી
શકીએ કે શાખામાં $\sFmla{\True}{!A}$ આવે તો
$\Sat{M}{!A}$, અને $\sFmla{\False}{!A}$
આવે તો $\Sat/{M}{!A}$.
($\Theta =
\Setabs{!A}{\sFmla{\True}{!A} \text{ શાખા પર}}$
અને $\Xi =
\Setabs{!A}{\sFmla{\False}{!A} \text{ શાખા પર}}$ લો.)

આ રીતે શોધીને બનાવેલાં કોષ્ટક-વૃક્ષના દરેક
શિરોબિંદુને સિક્વન્ટ સોંપવાથી તેને સરળતાથી
\Log{G3c}ની સાબિતીમાં રૂપાંતરિત કરી શકાય છે.
શરૂઆતનાં સૂત્રો $\Gamma \Sequent \Delta$
સિક્વન્ટને અનુરૂપ હોય, તો દરેક શરૂઆતના
સૂત્રને $\Gamma \Sequent \Delta$ ચિહ્ન આપો.
$\sFmla{\True}{!A}$ ચિહ્નિત સૂત્ર
પર શાખા-વિસ્તારનો નિયમ લાગુ કરવાથી શાખા
વધે, તો પાનને $!A, \Pi \Sequent \Lambda$
ચિહ્ન આપો; ચિહ્નિત સૂત્ર
$\sFmla{\False}{!A}$ પર તે
લાગુ થાય, તો પાનને $\Pi \Sequent \Lambda, !A$
ચિહ્ન આપો. કોષ્ટક-વૃક્ષમાં
\TRule{\True}{\star} નિયમ લાગુ થાય, ત્યાં
સિક્વન્ટ પર \LeftR{\star} નિયમ લગાવો
($!A$ને મુખ્ય સૂત્ર રાખીને); અને
\TRule{\False}{\star} લાગુ થાય, ત્યાં
\RightR{\star} નિયમ લગાવો. નિયમના
આધારવિધાનો કોષ્ટક-વૃક્ષમાં ઉમેરાયેલા
સંબંધિત શિરોબિંદુઓને ચિહ્નિત કરતા સિક્વન્ટો
છે. $\sFmla{\True}{!A}$ અને
$\sFmla{\False}{!A}$ વડે કોષ્ટક-વૃક્ષ બંધ
થાય, ત્યારે અંતિમ પાનનું ચિહ્ન
$!A, \Pi \Sequent \Lambda, !A$ સ્વરૂપનો
સિક્વન્ટ હશે. કોષ્ટક-વૃક્ષમાં ક્રમશઃ આવતા
કેટલાંક શિરોબિંદુઓને એક જ સિક્વન્ટ સોંપાશે;
પુનરાવર્તિત નકલો કાઢી નાખો. દાખલા તરીકે,
ઉપરનું કોષ્ટક-વૃક્ષ આ સાબિતીમાં રૂપાંતરિત થાય છે:
\[\small
\Axiom$!A \fCenter !A, !B$
\RightLabel{\RightR{\lor}}
\UnaryInf$!A \fCenter !A \lor !B$
\Axiom$!A \fCenter !A, !C$
\RightLabel{\RightR{\lor}}
\UnaryInf$!A \fCenter !A \lor !C$
\RightLabel{\RightR{\land}}
\BinaryInf$!A \fCenter (!A \lor !B) \land (!A \lor !C)$
\Axiom$!B, !C \fCenter !A, !B$
\RightLabel{\RightR{\lor}}
\UnaryInf$!B, !C \fCenter !A \lor !B$
\RightLabel{\LeftR{\land}}
\UnaryInf$!B \land !C \fCenter !A \lor !B$
\Axiom$!B, !C \fCenter !A, !C$
\RightLabel{\RightR{\lor}}
\UnaryInf$!B, !C \fCenter !A \lor !C$
\RightLabel{\LeftR{\land}}
\UnaryInf$!B \land !C \fCenter !A \lor !C$
\RightLabel{\RightR{\land}}
\BinaryInf$!B \land !C \fCenter (!A \lor !B) \land
(!A \lor !C)$
\RightLabel{\LeftR{\lor}}
\BinaryInf$!A \lor (!B \land !C) \fCenter (!A \lor !B) \land
(!A \lor !C)$
\DisplayProof
\]
% END OLP-0684
\endgroup


\OLEndChapterHook
% END OLP-0681
\endgroup
% END OLP-0685
